% !TeX root = ../Book3.tex
% 纲要标题：# 第一编　素谱、有限性与整性
% 纲要位置：工作记录/计划纲要.md，第 1812 行。


\BookPart{素谱、有限性与整性}{b3:part:01}
\BookPartGuide
一个多项式环既是可以加、乘、取理想的环，也是多项式函数的集合。方程的零点集与函数环分别保留什么信息，怎样在各点的局部环中检验整体结论，哪些结论必须依靠有限生成条件？本编把第一卷的环与方程、第二卷的模及其构造结合起来，围绕这些问题研究交换环。除明确说明外，环均指含单位元的交换环，环同态保持单位元。

\ProseParagraphBreak
\chapref{b3:ch:01}从素理想开始建立素谱。这里的点不只包括熟悉的坐标点：一个非极大素理想还可以记录整条不可约闭子集的一般情形。读者应先熟悉闭集 \(V(I)\)、基本开集 \(D(f)\) 与理想运算之间的对应，再阅读谱的函子性。\chapref{b3:ch:02}把局部化推广到模，解释如何在一点处检查零元、正合性和有限生成性，并特别区分局部化所得的茎与模去极大理想所得的纤维。

\ProseParagraphBreak
\chapref{b3:ch:03}整理 Noether 性、Artin 性\footnote{阿廷（Emil Artin，1898-1962），奥地利裔数学家，研究代数数论、域论与环论。}及 Nakayama 引理\footnote{中山正（假名：なかやま ただし，罗马音：Tadashi Nakayama，1912-1964），日本数学家，研究环论与表示论。}。有限生成假设常在清除公共分母、选取有限关系或控制幂零指数时真正发挥作用，阅读证明时应找到它被使用的那一步。\chapref{b3:ch:04}以元素的零化子定义关联素理想，证明 Noether 环中准素分解的存在性，并说明极小分解的关联素理想集合与孤立准素分量唯一，而嵌入准素分量一般不唯一。底拓扑空间不能辨认嵌入关联结构，这正是仅看零点集会遗漏的代数信息之一。

\ProseParagraphBreak
\chapref{b3:ch:05}讨论整扩张中的素理想及整闭性，并在任意底域上证明 Noether 正规化引理。\chapref{b3:ch:06}由此证明 Hilbert 零点定理，再处理仿射代数的正规化与导子理想（conductor）。两种“正规化”务必区分：Noether 正规化选取一个多项式子环，使原代数成为其上的有限生成模；环的正规化则在分式域或全分式环中取整闭包。后者的有限性需要证明，本编会处理任意底域及约化多分支情形。

\ProseParagraphBreak
\chapref{b3:ch:07}集中解释这些结论的几何意义。前半从代数闭域上的经典仿射代数集出发，证明正则函数、坐标环和态射的对应，不预设概形或层论。后半为保留纤维与基域扩张中的幂零信息，转向任意域 \(k\) 上的有限型 \(k\) 代数，允许非约化结构，并对域扩张 \(K/k\) 使用 \(K\) 值点和几何点描述扩域后的点。尖点、节点、交叉直线与加厚点检验了这一对象变化的必要性：相同的点数、双射的映射乃至同胚的空间，都可能伴随不同的函数环。

\begin{center}
\begin{tikzcd}[column sep=large,row sep=large]
\begin{gathered}\text{第1章}\\\text{素谱}\end{gathered}\arrow[r]
&\begin{gathered}\text{第2章}\\\text{局部化}\end{gathered}\arrow[r]
&\begin{gathered}\text{第3章}\\\text{有限性}\end{gathered}\arrow[r]\arrow[d]
&\begin{gathered}\text{第4章}\\\text{准素分解}\end{gathered}\arrow[d]\\
&&\begin{gathered}\text{第5章}\\\text{整扩张}\end{gathered}\arrow[r]
&\begin{gathered}\text{第6章}\\\text{零点与正规化}\end{gathered}\arrow[d]\\
&&&\begin{gathered}\text{第7章}\\\text{仿射代数集}\end{gathered}
\end{tikzcd}
\end{center}

图中列出主要依赖。第6章处理约化环的多分支正规化时使用关联素理想；第7章使用第1章的拓扑结论和第2章的局部化判据。其中\secref{b3:sec:0704}的交叉直线计算还用到\subsecref{b3:subsec:0605:04}的全分式环分解。

\ProseParagraphBreak
阅读本编需要熟悉第一卷的理想、商环、多项式、整性和域扩张，以及第二卷的模、正合列、张量积与有限展示。一般点和无限维反例有时需要选择原理，正规化有限性还会调用第一卷有限可分扩张的迹配对；使用处均给出确切引用。每章习题附有解答，建议先独立写出理想、同态或分式的具体表达式，再比较答案中对有限性及良定义的检查。

\ProseParagraphBreak
读过本编，应能判断一条局部结论怎样回到原环，辨认有限生成条件在论证中的作用，并说明为什么零点集不足以恢复函数环。比较两个代数对象时，除了点与不可约分量，还应检查局部环、幂零元和嵌入关联结构；这些信息决定了外观相同的方程是否具有相同的代数性质。

% !TeX root = ../../Book3.tex
% 纲要标题：## 第1章　素理想、根理想与素谱
% 纲要位置：工作记录/计划纲要.md，第 1912 行。

\chapter{素理想、根理想与素谱}
\label{b3:ch:01}
\BookChapterNumber{b3:ch:01}

\begin{chapterabstract}
模去一个素理想得到整环，模去一个极大理想得到域；两种商环保留的代数信息不同。本章把全部素理想组成一个拓扑空间，以理想的根描述闭集，以素理想的包含关系辨认一般点、闭点和不可约分量。商环与局部化的谱对应进一步说明环的运算怎样改变这个空间，而幂零元的例子说明底拓扑空间不能恢复全部环结构。
\end{chapterabstract}

\begin{learninggoals}
\begin{itemize}
\item 用商环辨认素理想与极大理想，说明选择原理在一般存在性证明中的作用。
\item 证明根理想的素理想交表示，区分幂零元、幂零理想与约化性。
\item 在闭集、基本开集和理想之间往返计算，辨认一般点、闭点与不可约分量。
\item 写出环同态诱导的素谱映射，并区分主开集与一般局部化的像。
\end{itemize}
\end{learninggoals}

设 \(k\) 为域。在 \(k[x,y]\) 中，模去 \((x,y)\) 就把两个变量都置为零，得到域 \(k\)；模去 \((x)\) 则只把 \(x\) 置为零，留下多项式环 \(k[y]\)：
\[
k[x,y]/(x,y)\cong k,\qquad k[x,y]/(x)\cong k[y].
\]
前者对应原点，后者保留整条直线上的多项式。包含 \((x)\) 的极大理想仍给出这条直线的闭点，例如每个 \(a\in k\) 都给出 \((x,y-a)\)。只收集极大理想时，遗漏的是 \((x)\) 本身作为一个点；它记录整条直线共同满足的关系，而没有再指定 \(y\) 的取值。素谱把这个点也纳入其中，\cref{b3:prop:generic-and-closed-points}将证明它的闭包恰为整条 \(V(x)\)，因而称它为这条不可约直线的一般点。这里的“点”不只意味着给变量代入一组数。

\ProseParagraphBreak
整数环也有这种区别。\((2)\) 的商为域，而 \((0)\) 的商仍是整环 \(\mathbb Z\)。在 \(\operatorname{Spec}\mathbb Z\) 中，\((0)\) 的闭包是全谱，各个 \((p)\)（\(p\) 为素数）却是闭点；理想的包含关系说明这些点之间怎样联系。

\ProseParagraphBreak

同一空间还会忘记一些方程中的幂次。例如 \((6)\) 与 \((12)\) 在 \(\operatorname{Spec}\mathbb Z\) 中给出相同的闭集 \(\{(2),(3)\}\)。两者的根都是 \((6)\)，因此包含它们的素理想相同，尽管原来的两个理想不同。根理想将精确描述闭集所记录的信息；商环中的幂零元则说明取根时失去了什么。

\ProseParagraphBreak

本章使用第一卷的理想、商环、多项式和环局部化，也沿用其中的 Noether 升链条件。一般极大对象的存在性明确调用第一卷附录A的 Zorn 引理\footnote{佐恩（Max August Zorn，1906-1993），德国裔美国数学家，提出与选择公理等价的极大原则。}。此时不需要概形、层或同调代数；“点”暂时就是一个素理想，几何的直观只用来帮助解释已经给出的代数定义。

\ProseParagraphBreak
可参阅松村英之的《Commutative Ring Theory》\cite[\S 1, pp. 1--3; \S 4, pp. 24--25]{Matsumura1989CommutativeRingTheory}。该书 \S 4 的 Exercises 4.3--4.5、4.9--4.12 分别讨论局部化谱、商环谱、拟紧性及 Noether 素谱的不可约分解\cite[p. 29]{Matsumura1989CommutativeRingTheory}。该书的根记号与本卷略有不同，本卷以 \(\operatorname{Nil}(R)\) 表示幂零根，并将其与 Jacobson 根\footnote{雅各布森（Nathan Jacobson，1910-1999），美国数学家，研究环论和 Lie 代数，证明了稠密性定理。}区分。后者的模论作用在第三章讨论。

% !TeX root = ../../Book3.tex
% 纲要标题：### 1.1 素理想与极大理想
% 纲要位置：工作记录/计划纲要.md，第 1816 行。
% 纲要细目（写作范围）：
% * 素理想及商整环
% * 极大理想及剩余域
% * 素理想的存在与回避

\section{素理想与极大理想}
\label{b3:sec:0101}

在交换环中，模去一个理想以后，可能得到整环，也可能直接得到域。这两种情形分别由素理想和极大理想刻画。它们将成为本卷所研究空间的点，因此先把代数条件、商环及存在性说清楚。除另有说明外，本卷的环都交换且含幺，同态保持单位；允许零环，但素理想和极大理想必须是真理想。

\subsection{素理想与商整环}
\label{b3:subsec:0101:01}

\begin{definition}\label{b3:def:prime-maximal}
设 \(R\) 为交换含幺环。若 \(R\) 的真理想 \(\mathfrak p\) 满足
\[
ab\in\mathfrak p\ \Longrightarrow\ a\in\mathfrak p\text{ 或 }b\in\mathfrak p
\qquad(a,b\in R),
\]
则称 \(\mathfrak p\) 为\term[sulixiang]{素理想}[prime ideal]。若 \(R\) 的真理想 \(\mathfrak m\) 与 \(R\) 之间不存在其他理想，则称 \(\mathfrak m\) 为\term[jidalixiang]{极大理想}[maximal ideal]。
\end{definition}

\begin{proposition}\label{b3:prop:prime-maximal-quotient}
设 \(R\) 为交换含幺环，\(\mathfrak p,\mathfrak m\) 为 \(R\) 的真理想。\(\mathfrak p\) 为素理想当且仅当 \(R/\mathfrak p\) 为整环；\(\mathfrak m\) 为极大理想当且仅当 \(R/\mathfrak m\) 为域。特别地，极大理想必为素理想。
\end{proposition}
\begin{proof}
真理想保证商中的单位元不为零。商中 \((a+\mathfrak p)(b+\mathfrak p)=0\) 等价于 \(ab\in\mathfrak p\)，故第一个断言正是无零因子的条件。

若 \(\mathfrak m\) 极大，对 \(a\notin\mathfrak m\) 有 \(\mathfrak m+(a)=R\)，可写 \(1=u+ra\)、\(u\in\mathfrak m\)，所以 \(r+\mathfrak m\) 是 \(a+\mathfrak m\) 的逆。反之，若商为域，设 \(J\) 严格包含 \(\mathfrak m\)，取 \(a\in J\setminus\mathfrak m\)。将商中 \(a+\mathfrak m\) 的逆提升为 \(b\in R\)，则 \(ab-1\in\mathfrak m\subseteq J\)，而 \(ab\in J\)，相减得 \(1\in J\)。所以 \(J=R\)，证明极大性。域是整环，再由第一个断言即得最后结论。
\end{proof}

这重述了第一卷命题10.2.6的商环判据。实际检验时，寻找一个方便辨认的商环，常比逐一检查所有乘积更直接。例如 \(k[x,y]/(x)\cong k[y]\)，所以 \((x)\) 素而不极大；\(k[x,y]/(x,y)\cong k\)，所以 \((x,y)\) 极大。\(\mathbb Z\) 的素理想为 \((0)\) 和 \((p)\)、\(p\) 为素数，其中只有后者极大。

\ProseParagraphBreak
还可用理想的乘积表述素性：若 \(I J\subseteq\mathfrak p\)，则 \(I\subseteq\mathfrak p\) 或 \(J\subseteq\mathfrak p\)。否则分别取 \(a\in I\setminus\mathfrak p\)、\(b\in J\setminus\mathfrak p\)，有 \(ab\in\mathfrak p\)，违反素性。这个形式将在闭集并的计算中使用。

\subsection{极大理想与剩余域}
\label{b3:subsec:0101:02}

设 \(R\) 为交换含幺环，\(\mathfrak m\) 为 \(R\) 的极大理想。域 \(\kappa(\mathfrak m)=R/\mathfrak m\) 称为 \(\mathfrak m\) 处的\term[shengyu-yu]{剩余域}[residue field]。若 \(R\) 是域 \(k\) 上的非零代数，这个域含有 \(k\) 的像，却未必等于 \(k\)。例如 \(\mathbb R[t]\) 中 \((t^2+1)\) 极大，剩余域为 \(\mathbb C\)，而不存在实数 \(a\) 使这个理想等于 \((t-a)\)。第六章将说明有限型代数在代数闭域上的情形为何不同。
\symidx{\(\kappa(\mathfrak m)=R/\mathfrak m\)：极大理想处的剩余域}

\begin{theorem}\label{b3:thm:maximal-existence}
设 \(R\) 为交换含幺环。\(R\) 的任一真理想都包含于某个极大理想。因而非零交换含幺环至少有一个素理想；零环没有素理想。
\end{theorem}
\begin{proof}
固定真理想 \(I\)，在包含 \(I\) 的全部真理想中按包含排序。此集合非空。非空链的并仍为理想，且不含一：若一在并中，它已在某个链成员中，矛盾。空链由 \(I\) 提供上界。因此第一卷附录A定理A.1.4的 Zorn 引理给出极大元，它也是 \(R\) 的极大理想。取 \(I=0\) 得非零环的结论。零环没有真理想，故两个定义都不适用。
\end{proof}

一个常用后果是：\(a\in R\) 可逆当且仅当它不属于任何极大理想。若不可逆，\((a)\) 为真理想，用\cref{b3:thm:maximal-existence}即可找到包含它的极大理想；反向则因为包含单位的理想必为全环。注意这里并不要求 \(R\) 为 Noether 环。

\subsection{素理想的存在与回避}
\label{b3:subsec:0101:03}

\begin{lemma}[避开乘法集]\label{b3:lem:prime-disjoint-multiplicative}
设 \(R\) 为交换含幺环，\(S\subseteq R\) 含一且对乘法封闭，\(I\) 为 \(R\) 中与 \(S\) 不交的理想。则存在与 \(S\) 不交的素理想 \(\mathfrak p\supseteq I\)。
\end{lemma}
\begin{proof}
考虑包含 \(I\) 且与 \(S\) 不交的理想。链的并仍与 \(S\) 不交，空链有上界 \(I\)，所以 Zorn 引理给出极大者 \(\mathfrak p\)。它不含一，因而为真理想。若 \(a,b\notin\mathfrak p\)，极大性使 \(\mathfrak p+(a)\)、\(\mathfrak p+(b)\) 分别含有 \(s=u+ra\in S\)、\(t=v+qb\in S\)。若 \(ab\in\mathfrak p\)，展开乘积即得 \(st\in\mathfrak p\)，与 \(st\in S\) 矛盾。因此 \(\mathfrak p\) 素。
\end{proof}

这里使用选择原理的地方，是在避开乘法集的理想族中取极大元；它的极大性只相对于这个理想族而言，未必是原环的极大理想。对具体多项式环，常可直接写出所需的素理想；上述论证则保证任意交换环中相应对象存在。可参阅松村英之的《Commutative Ring Theory》\cite[Theorems 1.1--1.2, pp. 2--3]{Matsumura1989CommutativeRingTheory}。

\begin{lemma}[有限素理想回避]\label{b3:lem:finite-prime-avoidance}
设 \(R\) 为交换含幺环，\(I\) 是 \(R\) 的理想，\(\mathfrak p_1,\ldots,\mathfrak p_r\) 是 \(R\) 的素理想，其中 \(r\ge1\)。若 \(I\subseteq\bigcup_{i=1}^r\mathfrak p_i\)，则 \(I\subseteq\mathfrak p_i\) 对某个 \(i\) 成立。
\end{lemma}
\begin{proof}
删去被其他素理想包含的成员不改变并集，可设余下者两两不可比。若 \(I\) 不包含于任何一个，选 \(a_i\in I\setminus\mathfrak p_i\)。对 \(j\ne i\)，不可比性使我们可以选 \(b_{ij}\in\mathfrak p_j\setminus\mathfrak p_i\)。令
\[
c_i=a_i\prod_{j\ne i}b_{ij},\qquad c=\sum_i c_i.
\]
素性使 \(c_i\notin\mathfrak p_i\)，而 \(c_i\in I\) 且属于每个 \(\mathfrak p_j\)、\(j\ne i\)。对固定 \(i\)，和式中其他各项都在 \(\mathfrak p_i\) 内，所以 \(c\notin\mathfrak p_i\)。于是 \(c\in I\) 却不在所给并集中，矛盾。只有一个成员时把空乘积取为一，同一论证仍成立。
\end{proof}

这称为\term[sulixiang-huibi]{素理想回避引理}[prime avoidance lemma]的有限素理想版本。它常用来选择同时避开若干素理想的元素：若理想 \(I\) 不包含于任何一个 \(\mathfrak p_i\)，则 \(I\) 中至少有一个元素不属于它们的并，即
\[
I\nsubseteq\mathfrak p_i\quad(1\le i\le r)
\quad\Longrightarrow\quad
I\setminus\bigcup_{i=1}^r\mathfrak p_i\ne\varnothing.
\]

% !TeX root = ../../Book3.tex
% 纲要标题：### 1.2 根与约化
% 纲要位置：工作记录/计划纲要.md，第 1920 行。
% 纲要细目（写作范围）：
% * 幂零根和根理想；用有限混合乘积与无限变量例子区分逐元素幂零、统一指数和理想幂零
% * 根理想作为素理想的交；把不含指定元素转为避开其幂组成的乘法集
% * 约化环与零维例子；无限域乘积的论证对全部素理想成立

\section{根与约化}
\label{b3:sec:0102}

给定交换含幺环 \(R\) 的理想 \(I\)，一个元素 \(a\) 属于所有包含 \(I\) 的素理想，当且仅当它在商环 \(R/I\) 中的像是幂零元，也就是某次幂 \(a^n\) 落在 \(I\) 中。这一事实把素理想的包含关系与商环中的幂零现象联系起来。不过，这些素理想并不记录最小的幂次 \(n\)，也不能恢复原来的理想 \(I\)。将 \(I\) 换成它的根，保留的正是这组素理想包含信息。

\subsection{幂零根与根理想}
\label{b3:subsec:0102:01}

\begin{definition}\label{b3:def:radical-reduced}
设 \(R\) 为交换含幺环，\(I\) 为 \(R\) 的理想。定义 \(I\) 的\term[lixiang-de-gen]{根}[radical of an ideal]为
\[
\sqrt I=\{\,a\in R:\text{存在 }n\ge1\text{ 使 }a^n\in I\,\}.
\]
若 \(I=\sqrt I\)，则称 \(I\) 为\term[genlixiang]{根理想}[radical ideal]。\(\sqrt{(0)}\) 称为 \(R\) 的\term[miling-gen]{幂零根}[nilradical]，记为 \(\operatorname{Nil}(R)\)。若 \(\operatorname{Nil}(R)=0\)，则称 \(R\) 为\term[yuehua-huan]{约化环}[reduced ring]。
\end{definition}
\symidx{\(\sqrt I\)：理想的根}
\symidx{\(\operatorname{Nil}(R)\)：环的幂零根}

\(\sqrt I\) 确实是理想。若 \(a^m,b^n\in I\)，展开 \((a+b)^{m+n-1}\)，每项至少含 \(a^m\) 或 \(b^n\)，故该幂属于 \(I\)；负号和环元素倍乘也保持条件。还直接得到
\[
I\subseteq\sqrt I,\qquad \sqrt{\sqrt I}=\sqrt I,
\qquad \sqrt{I\cap J}=\sqrt{IJ}=\sqrt I\cap\sqrt J.
\]
最后一式中，若 \(a^m\in I,a^n\in J\)，则 \(a^{m+n}\in IJ\)；另一方向由 \(IJ\subseteq I\cap J\) 给出。对理想根迭代时，只需把两个存在的正指数相乘，便回到 \(I\)。

\ProseParagraphBreak
商环的幂零根为 \(\sqrt I/I\)，所以 \(R/I\) 约化当且仅当 \(I\) 是根理想。特别地，\(R_{\mathrm{red}}=R/\operatorname{Nil}(R)\) 约化，称为 \(R\) 的约化商。它把全部幂零元素送到零，而不要求剩下的非零元素都可逆。
\symidx{\(R_{\mathrm{red}}\)：环的约化商}

\begin{definition}\label{b3:def:nilpotent-ideal}
设 \(R\) 为交换含幺环，\(J\) 为 \(R\) 的理想。对正整数 \(N\)，\(J^N\) 是全部 \(N\) 个 \(J\) 中元素的乘积所生成的理想。若存在正整数 \(N\) 使 \(J^N=0\)，则称 \(J\) 为\term[mi-ling-lixiang]{幂零理想}[nilpotent ideal]；满足该式的最小正整数称为 \(J\) 的\term[miling-zhishu]{幂零指数}[nilpotency index]。
\end{definition}

理想幂零要求任意 \(N\) 个理想元素的混合乘积都为零，强于只对同一个元素取幂。对理想 \(J\)，下面三个条件逐项变强：
\[
\begin{aligned}
&\forall u\in J\ \exists n(u)\ge1, &&u^{n(u)}=0;\\
&\exists N\ge1\ \forall u\in J, &&u^N=0;\\
&\exists N\ge1, &&J^N=0.
\end{aligned}
\]
第一个条件在英文文献中称为 nil ideal，只要求每个元素各自幂零；第二个条件要求元素幂零有统一上界；第三个条件才是理想幂零。最后一个条件蕴含前两个，反向一般不成立。

\ProseParagraphBreak
例如，设 \(k\) 为域，在商环
\[
R=k[\varepsilon,\delta]/(\varepsilon^2,\delta^2)
\]
中仍用 \(\varepsilon,\delta\) 表示变量的像，令 \(J=(\varepsilon,\delta)\)。这个商环有 \(k\)-基 \(1,\varepsilon,\delta,\varepsilon\delta\)，所以混合乘积 \(\varepsilon\delta\ne0\)。因此
\[
J^2=(\varepsilon\delta)\ne0,\qquad J^3=0,
\]
后一式是因为任意三个生成元的乘积都会重复 \(\varepsilon\) 或 \(\delta\)。若 \(\operatorname{char}k=2\)，每个 \(u\in J\) 甚至都满足 \(u^2=0\)，而 \(J^2\) 仍不为零。这说明同一指数下，单个元素的幂与不同元素的混合乘积不能混为一谈。

\begin{proposition}\label{b3:prop:finite-nil-ideal}
设 \(R\) 为交换含幺环，\(J=(a_1,\ldots,a_m)\) 为 \(R\) 的理想，其中 \(m\ge0\)。若 \(a_i^{n_i}=0\)，且 \(n_i\ge1\)（\(1\le i\le m\)），则
\[
J^N=0,\qquad N=1+\sum_{i=1}^m(n_i-1).
\]
因此 Noether 环的幂零根是幂零理想。
\end{proposition}
\begin{proof}
\(m=0\) 时，\(J=0\)、\(N=1\)，结论成立。设 \(m\ge1\)。\(J^N\) 由全部 \(N\) 个生成元的乘积线性生成。若每个 \(a_i\) 出现次数都小于 \(n_i\)，总个数至多为 \(N-1\)，矛盾。所以每个乘积含某个 \(a_i^{n_i}\)，均为零。Noether 环的幂零根作为理想有限生成，每个生成元又幂零，故适用第一项。
\end{proof}

没有有限生成性时，一个理想即使每个元素都幂零，也未必是幂零理想。取域 \(k\)，考虑
\[
R=k[x_1,x_2,\ldots]/(x_1^2,x_2^2,\ldots).
\]
令 \(J=(\bar x_1,\bar x_2,\ldots)\)。每个 \(u\in J\) 只用有限多个变量，设有 \(r\) 个；任何 \(r+1\) 个正次单项式的乘积都会重复一个变量，所以 \(u^{r+1}=0\)。但 \(\bar x_1\cdots\bar x_N\ne0\) 对任意 \(N\) 成立，因为平方关系生成的理想由含某个平方因子的单项式线性生成，不能包含这个不含平方因子的单项式。又因商 \(R/J=k\) 约化，全部幂零元都在 \(J\) 中。因此 \(J=\operatorname{Nil}(R)\) 而 \(J^N\ne0\) 对全部 \(N\) 成立。

\ProseParagraphBreak
若 \(\operatorname{char}k=2\)，这个例子还有更强的性质：每个 \(u\in J\) 都满足 \(u^2=0\)。因为 \(u\) 是有限个正次单项式的线性组合，特征二时平方把有限和变成各项平方之和，而每个正次单项式的平方都为零。所以元素幂零即使有统一上界，也不能保证理想的某个幂为零。

\subsection{根理想的素理想交表示}
\label{b3:subsec:0102:02}

\begin{theorem}\label{b3:thm:radical-prime-intersection}
设 \(R\) 为交换含幺环。任意理想 \(I\subseteq R\) 满足
\begin{equation}\label{b3:eq:radical-prime-intersection}
\sqrt I=\bigcap_{\mathfrak p\supseteq I,\ \mathfrak p\text{ 素}}\mathfrak p.
\end{equation}
空交取为 \(R\)，因而公式也适用于 \(I=R\)。
\end{theorem}
\begin{proof}
若 \(a^n\in I\subseteq\mathfrak p\)，反复使用素性得到 \(a\in\mathfrak p\)，故左边包含于右边。

反过来，对 \(a\notin\sqrt I\)，我们要找一个包含 \(I\) 却不含 \(a\) 的素理想。对素理想 \(\mathfrak p\)，不含 \(a\) 等价于不含 \(a\) 的任何非负整数次幂：正次幂属于 \(\mathfrak p\) 会由素性推出 \(a\in\mathfrak p\)，而 \(1\) 本就不属于真理想。因此取乘法集 \(S=\{1,a,a^2,\ldots\}\)，便将不含 \(a\) 转化为与 \(S\) 不交。由 \(a\notin\sqrt I\) 可知 \(S\cap I=\varnothing\)：没有正次幂落在 \(I\) 中，且 \(I\ne R\)，所以 \(1\notin I\)。\cref{b3:lem:prime-disjoint-multiplicative}给出包含 \(I\) 并避开 \(S\) 的素理想 \(\mathfrak p\)，特别是 \(a\notin\mathfrak p\)。因此 \(a\) 不在右边的交中，证明反向包含。
\end{proof}

取 \(I=0\)，可知一个元素幂零当且仅当它属于全部素理想。于是映射
\[
R\longrightarrow\prod_{\mathfrak p\text{ 素}}R/\mathfrak p,
\qquad a\longmapsto(a+\mathfrak p)_{\mathfrak p}
\]
的核恰为 \(\operatorname{Nil}(R)\)。约化环因而可以嵌入整环的乘积，但本身未必是整环：\(k\times k\) 已给出最简单的反例。

\ProseParagraphBreak
在一般环上，不能把上述交中的素理想直接换成极大理想。以 \(R=k[t]_{(t)}\) 为例，其唯一极大理想为 \((t)R\)，但 \(R\) 是整环，幂零根为零。因此所有极大理想的交，即 Jacobson 根，可以严格大于幂零根。第二章将详细解释这个局部化环。

\subsection{约化环与零维例子}
\label{b3:subsec:0102:03}

\(k[\varepsilon]/(\varepsilon^2)\) 只有一个素理想 \((\varepsilon)\)：素理想必须包含幂零元 \(\varepsilon\)，商为域 \(k\)。它的约化商是 \(k\)。因此只记录有哪些素理想，尚不能区分这两个环；它们的区别在于前者存在非零平方零元。

\ProseParagraphBreak
\(\mathbb Z/12\mathbb Z\) 的素理想是 \((\bar2)\)、\((\bar3)\)，其交为 \((\bar6)\)，故幂零根平方为零。约化商为 \(\mathbb Z/6\mathbb Z\cong\mathbb F_2\times\mathbb F_3\)。后者没有非零幂零元，却仍有零因子，例如 \(\bar2\bar3=0\)。约化性只排除幂零元，不排除两个不同非零元素的乘积为零。

\ProseParagraphBreak
设 \(R\) 为非零交换含幺环。若 \(R\) 的每个素理想都是极大理想，则称 \(R\) 为零维环；这等价于素理想之间不存在严格包含。一般 Krull 维数将在第九章定义。\footnote{克鲁尔（Wolfgang Krull，1899-1971），德国数学家，研究交换环中的素理想、局部化与维数。}上述双数环及有限个域的乘积都零维。

\ProseParagraphBreak
零维性却不能控制理想的降链。令 \(k\) 为域，取无限乘积环 \(R=\prod_{i\ge1}k\)，其中的理想
\[
J_n=\{\,(a_i)\in R\mid a_1=\cdots=a_n=0\,\}
\qquad(n\ge1)
\]
满足 \(J_{n+1}\subsetneq J_n\)：第 \(n+1\) 个坐标为一、其余坐标为零的元素属于 \(J_n\)，却不属于 \(J_{n+1}\)。因此 \(R\) 不是 Artin 环。

\ProseParagraphBreak
另一方面，\(R\) 的每个素理想都极大。这里不把素理想限定为坐标投影的核；无限乘积中还可以有其他素理想，下面的论证对任意素理想都适用。对任意元素 \(a=(a_i)\)，在非零坐标取 \(b_i=a_i^{-1}\)，在零坐标取 \(b_i=0\)，便有 \(a=a^2b\)。模去任一素理想后，所得商是整环；其中任意非零的 \(\bar a\) 都可从 \(\bar a=\bar a^2\bar b\) 消去一份 \(\bar a\)，得到 \(1=\bar a\bar b\)，所以这个商是域。这样，\(R\) 是零维环，却有无限严格理想降链。第三章将讨论 Artin 环及其与零维性的关系。

% !TeX root = ../../Book3.tex
% 纲要标题：### 1.3 素谱的拓扑
% 纲要位置：工作记录/计划纲要.md，第 1828 行。
% 纲要细目（写作范围）：
% * Zariski 闭集与基本开集
% * 拟紧性（quasi-compactness，不附加 Hausdorff 条件）
% * 不可约闭集、一般点与闭点
% * Noether 素谱与有限不可约分量

\section{素谱的拓扑}
\label{b3:sec:0103}

把素理想作为点，还要指定哪些点属于同一个闭集。理想的包含关系给出一种自然规则：把包含同一理想的素理想收集起来。这个拓扑通常不是 Hausdorff 拓扑\footnote{豪斯多夫（Felix Hausdorff，1868-1942），德国数学家，建立一般拓扑中的分离公理等基本概念。}，甚至不同点之间可能有特殊化关系；这些现象正是需要保留的代数信息。

\subsection{Zariski 闭集与基本开集}
\label{b3:subsec:0103:01}

\begin{definition}\label{b3:def:spectrum-zariski}
设 \(R\) 为交换含幺环。\(R\) 的\term[sulixiang-pu]{素理想谱}[prime spectrum]是集合
\[
\operatorname{Spec}R=\{\,\mathfrak p\subset R:\mathfrak p\text{ 是素理想}\,\}.
\]
对 \(R\) 的理想 \(I\)，记 \(V(I)=\{\,\mathfrak p\in\operatorname{Spec}R\mid I\subseteq\mathfrak p\,\}\)。以这些 \(V(I)\) 为闭集得到的拓扑称为\term[Zariski-tuopu]{Zariski 拓扑}[Zariski topology]。对单个元素或有限个元素，简记 \(V(f)=V\bigl((f)\bigr)\)、\(V(f_1,\ldots,f_r)=V\bigl((f_1,\ldots,f_r)\bigr)\)。对 \(f\in R\)，记
\[
D(f)=\operatorname{Spec}R\setminus V(f)
=\{\,\mathfrak p:f\notin\mathfrak p\,\},
\]
称为\term[jiben-kaiji]{基本开集}[basic open subset]，也称主开集\termidx[zhukaiji]{主开集}。
\end{definition}
\symidx{\(\operatorname{Spec}R\)：交换环的素理想谱}
\symidx{\(V(I)\)：包含理想 \(I\) 的素理想组成的闭集}
\symidx{\(D(f)\)：素谱中的基本开集}

这里的规则确实定义拓扑，因为
\begin{equation}\label{b3:eq:zariski-closed-identities}
\begin{gathered}
V(0)=\operatorname{Spec}R,\qquad V(R)=\varnothing,
\\
\bigcap_\lambda V(I_\lambda)=V\left(\sum_\lambda I_\lambda\right),
\qquad V(I)\cup V(J)=V(IJ).
\end{gathered}
\end{equation}
前三式来自包含关系。最后一式是因为素理想包含 \(IJ\) 当且仅当它包含 \(I\) 或 \(J\)：若两者都不包含，分别从 \(I,J\) 中选取不属于该素理想的元素，它们的乘积属于 \(IJ\)，与素性矛盾。任意理想之和中的每个元素仅为有限和，所以这也说明无限交闭集的条件成立。

\ProseParagraphBreak
由根的交表示，还有
\begin{equation}\label{b3:eq:zariski-radical-order}
V(I)\subseteq V(J)\quad\Longleftrightarrow\quad
\sqrt J\subseteq\sqrt I.
\end{equation}
因此 \(V(I)=V(\sqrt I)\)，根理想与闭子集反向一一对应。另一方面，\(D(f)\cap D(g)=D(fg)\)，且
\(\operatorname{Spec}R\setminus V(I)=\bigcup_{f\in I}D(f)\)，故基本开集构成开集的一组基。特别地，\(D(f)=\varnothing\) 当且仅当 \(f\) 幂零，\(D(f)=\operatorname{Spec}R\) 当且仅当 \(f\) 是单位。

\subsection{拟紧性}
\label{b3:subsec:0103:02}

设 \(X\) 为拓扑空间。若 \(X\) 的每个开覆盖都有有限子覆盖，则称 \(X\) 为\term[nijin-kongjian]{拟紧空间}[quasi-compact space]。本书不在此词中附加 Hausdorff 条件。

\begin{proposition}\label{b3:prop:spec-quasicompact}
设 \(R\) 为交换含幺环，\(f\in R\)。则 \(\operatorname{Spec}R\) 及其基本开集 \(D(f)\) 都拟紧。
\end{proposition}
\begin{proof}
先将任一开覆盖细化为基本开覆盖。若 \(\operatorname{Spec}R=\bigcup_\lambda D(f_\lambda)\)，则不存在包含全部 \(f_\lambda\) 的素理想。由\cref{b3:thm:maximal-existence}，理想 \((f_\lambda)_\lambda\) 必为 \(R\)，所以一是其中有限个元素的组合。这有限个基本开集已经覆盖全谱，再取它们所属的原开集便得有限子覆盖。

对 \(D(f)\) 的任意子空间开覆盖，可先细化成由基本开集与 \(D(f)\) 的交组成的覆盖；交仍是基本开集，因为 \(D(f)\cap D(g)=D(fg)\)。因此只需考虑 \(D(f)\subseteq\bigcup_\lambda D(g_\lambda)\)。补集的包含关系是 \(V\bigl((g_\lambda)_\lambda\bigr)\subseteq V(f)\)。由\eqref{b3:eq:zariski-radical-order}，某个 \(f^N\) 属于 \((g_\lambda)_\lambda\)，于是属于其中有限个 \(g_\lambda\) 生成的理想。取基本开集又得到有限子覆盖。\(D(f)\) 为空时同样成立。
\end{proof}

\subsection{不可约闭集、一般点与闭点}
\label{b3:subsec:0103:03}

基本开覆盖可以由有限条代数关系控制。闭集还有另一种结构问题：它是否由几个更小的闭集拼成？若不能，下面的判据会把这个闭集辨认为某个素理想的闭包。

\begin{definition}\label{b3:def:irreducible-space}
设 \(X\) 为非空拓扑空间。若 \(X\) 不能写成两个真闭子集之并，则称 \(X\) 为\term[bukeyue-kongjian]{不可约空间}[irreducible space]。闭子集的不可约性用其子空间拓扑理解。若 \(X\) 不能分成两个非空、不交的开集，则称 \(X\) 为连通；不可约性比连通性强，因为在不可约空间中任意两个非空开集必须相交。
\end{definition}

\begin{proposition}\label{b3:prop:spec-irreducible}
设 \(R\) 为交换含幺环，\(I\) 为 \(R\) 的理想。若闭集 \(V(I)\subseteq\operatorname{Spec}R\) 非空，则 \(V(I)\) 不可约当且仅当 \(\sqrt I\) 为素理想。
\end{proposition}
\begin{proof}
若 \(\sqrt I=\mathfrak p\) 素，且 \(V(I)=\bigl(V(I)\cap V(J)\bigr)\cup\bigl(V(I)\cap V(K)\bigr)\)，则 \(V(\mathfrak p)\subseteq V(JK)\)。由\eqref{b3:eq:zariski-radical-order}，\(JK\subseteq\mathfrak p\)，故 \(J\subseteq\mathfrak p\) 或 \(K\subseteq\mathfrak p\)。相应的一个闭子集就是整个 \(V(I)\)。

反过来，若 \(\sqrt I\) 不素，选 \(a,b\notin\sqrt I\)、\(ab\in\sqrt I\)。有 \(V(I)\subseteq V(ab)=V(a)\cup V(b)\)，但 \(V(I)\) 不包含于任何一边，否则根的交表示会使 \(a\) 或 \(b\) 属于 \(\sqrt I\)。这给出两个真闭子集之并，违反不可约性。
\end{proof}

设 \(X\) 为拓扑空间，\(x\in X\)。若单点集 \(\{x\}\) 在 \(X\) 中为闭集，则称 \(x\) 为\term[bidian]{闭点}[closed point]。素谱中还常有闭包比单点大的点；这种闭包可以直接由理想的包含关系读出。

\begin{proposition}\label{b3:prop:generic-and-closed-points}
设 \(R\) 为交换含幺环，\(\mathfrak p\) 为 \(R\) 的素理想。它在 \(\operatorname{Spec}R\) 中所对应点的闭包为 \(V(\mathfrak p)\)。因此它是闭点当且仅当 \(\mathfrak p\) 极大；每个不可约闭子集恰有一个闭包为它的点。
\end{proposition}
\begin{proof}
闭集 \(V(I)\) 包含点 \(\mathfrak p\)，当且仅当 \(I\subseteq\mathfrak p\)，这又保证 \(V(\mathfrak p)\subseteq V(I)\)。所以 \(V(\mathfrak p)\) 是包含该点的最小闭集。它只有该点，当且仅当没有更大的真素理想；利用极大理想存在性，这等价于 \(\mathfrak p\) 极大。若闭集不可约，\cref{b3:prop:spec-irreducible}把它唯一写成 \(V(\mathfrak p)\)，根理想与闭集的一一对应保证此 \(\mathfrak p\) 唯一。
\end{proof}

设 \(X\) 为拓扑空间，\(Z\subseteq X\) 为不可约闭子集。若点 \(\eta\in Z\) 满足 \(\overline{\{\eta\}}=Z\)，则称 \(\eta\) 为 \(Z\) 的\term[yiban-dian]{一般点}[generic point]。一般点是相对于指定闭集而言的，并不与闭点互斥：闭点也是其单点闭子集的一般点。例如整环的素谱有一般点 \((0)\)；在 \(\operatorname{Spec}\mathbb Z\) 中，它的闭包是全谱，所有 \((p)\) 则是闭点。不同点具有不同闭包，所以素谱满足 \(T_0\) 分离性质，却一般不满足每点都闭，更不必是 Hausdorff 空间。

\ProseParagraphBreak
设 \(R\) 为交换含幺环，\(\mathfrak p,\mathfrak q\) 为 \(R\) 的素理想。若 \(\mathfrak p\subseteq\mathfrak q\)，则称 \(\mathfrak q\) 是 \(\mathfrak p\) 的\term[teshuhua]{特殊化}[specialization]，并称 \(\mathfrak p\) 是 \(\mathfrak q\) 的\term[yibanhua]{一般化}[generalization]。由闭包公式，这等价于 \(\mathfrak q\in\overline{\{\mathfrak p\}}\)。素理想 \(\mathfrak p\) 作为 \(V(\mathfrak p)\) 的一般点，记录这个不可约闭子集共同满足的代数关系；但这个闭集未必在全谱中极大。例如域 \(k\) 上的 \(\operatorname{Spec}k[x,y]\) 本身不可约，而 \((x)\) 是其真闭子集 \(V(x)\) 的一般点。

\ProseParagraphBreak
一个只有两个点的谱已经能显示这种拓扑与通常坐标点的差别。设 \(k\) 为域，令 \(R=k[t]_{(t)}\)。任一非零元素都可写为 \(t^n u\)，其中 \(n\ge0\)、\(u\in R\) 为单位；这是把分子中的 \(t\) 因子提取出来，而分子余下的因子及分母都不在 \((t)\) 中。若该元素属于素理想，则 \(n\ge1\)，并由素性得该理想含 \(t\)。另一方面，取值同态 \(f/g\mapsto\dfrac{f(0)}{g(0)}\) 将 \(R\) 满射到 \(k\)，核为 \((t)R\)，所以 \((t)R\) 极大。\(R\) 是整环，故 \((0)\) 也是素理想。这就得到\cref{b3:fig:two-point-specialization}中的全部两点。

\begin{figure}[htbp]
\centering
\begin{tikzcd}[column sep=huge]
\eta=(0) \arrow[r,"\text{特殊化}"] & s=(t)R
\end{tikzcd}
\caption{两点素谱的特殊化关系}
\label{b3:fig:two-point-specialization}
\end{figure}

图中的箭头表示 \(s\) 属于 \(\eta\) 的闭包，不表示距离或路径。具体地，
\[
\overline{\{\eta\}}=\{\eta,s\},\qquad
\overline{\{s\}}=\{s\}.
\]
所以三个开集恰为 \(\varnothing\)、\(\{\eta\}\) 和 \(\{\eta,s\}\)。其中 \(\{\eta\}=D(t)\) 非空，却不含全谱的唯一闭点 \(s\)。在一般素谱中，闭点的剩余域也可能比基域大；第六、七章将在代数闭基域的有限型情形把闭点解释为通常的坐标点。

\subsection{Noether 素谱与有限不可约分量}
\label{b3:subsec:0103:04}

不可约闭子集可能严格包含在另一个不可约闭子集中。例如上面的 \(V(x)\) 就包含在不可约的 \(\operatorname{Spec}k[x,y]\) 内。设 \(X\) 为拓扑空间，若不可约闭子集 \(Z\subseteq X\) 不严格包含于任何更大的不可约闭子集，则称 \(Z\) 为 \(X\) 的\term[bukeyue-fenliang]{不可约分量}[irreducible component]。

\ProseParagraphBreak
设 \(R\) 为交换含幺环，若素理想 \(\mathfrak p\) 不严格包含 \(R\) 的任何其他素理想，则称 \(\mathfrak p\) 为\term[jixiao-sulixiang]{极小素理想}[minimal prime ideal]。由\cref{b3:prop:spec-irreducible}与\eqref{b3:eq:zariski-radical-order}，\(V(\mathfrak p)\) 是全谱的不可约分量，当且仅当 \(\mathfrak p\) 是极小素理想：任一更大的不可约闭集都形如 \(V(\mathfrak q)\)，其包含关系恰对应 \(\mathfrak q\subseteq\mathfrak p\)。因此所有素理想都给出各自闭包的一般点，只有极小素理想给出全谱不可约分量的一般点。

\ProseParagraphBreak
不可约分量是否只有有限个，则是另一个问题。设 \(X\) 为拓扑空间，若 \(X\) 的每条闭集降链最终稳定，则称 \(X\) 为\term[Noether-tuopu-kongjian]{Noether 拓扑空间}[Noetherian topological space]。在素谱中，闭集降链对应根理想升链；环的 Noether 条件使这些升链稳定，因而保证素谱是 Noether 拓扑空间。

\begin{proposition}\label{b3:prop:noether-spectrum-components}
设 \(R\) 为交换含幺 Noether 环。则 \(\operatorname{Spec}R\) 是 Noether 拓扑空间，并且是有限多个不可约分量之并。它们恰为 \(V(\mathfrak p_1),\ldots,V(\mathfrak p_r)\)，其中 \(\mathfrak p_i\) 是 \(R\) 的全部极小素理想；每个素理想都包含其中至少一个。零环对应空分解。
\end{proposition}
\begin{proof}
若 \(R=0\)，素谱为空，结论成立。以下设 \(R\ne0\)。

闭集降链由\eqref{b3:eq:zariski-radical-order}对应根理想升链，Noether 升链条件使其稳定。若存在不能表示为有限个不可约闭集之并的闭集，在这些闭集中取一个关于包含关系的极小元 \(Y\)。降链条件保证极小元存在，否则从任一成员出发连续选择更小者，便得到无限严格降链。于是 \(Y\) 的每个真闭子集都有有限不可约分解。\(Y\) 非空且不可能不可约，否则它自己就是一项分解，故可写成两个真闭子集 \(Y_1,Y_2\) 的并。分别分解 \(Y_1,Y_2\)，再将所得有限个不可约闭集合并，就得到 \(Y\) 的分解，矛盾。

因此全谱有有限不可约闭分解。删去重复成员和被其他成员包含者，得到两两不包含的 \(V(\mathfrak p_1),\ldots,V(\mathfrak p_r)\)，其中各 \(\mathfrak p_i\) 由\cref{b3:prop:spec-irreducible}为素理想。若非空不可约闭集 \(Z\) 被闭集 \(F_1,\ldots,F_r\) 覆盖，则
\[
Z=\bigcup_{i=1}^r(Z\cap F_i).
\]
对 \(r\) 归纳可知某个 \(Z\cap F_i=Z\)：\(r=1\) 时显然；\(r>1\) 时，将第一项与余下各项的并看作 \(Z\) 的两个闭子集。不可约性使其中一边等于 \(Z\)；若是后者，再对 \(r-1\) 项使用归纳。因此任一不可约闭集都包含于某个 \(V(\mathfrak p_i)\)。若不可约闭集 \(Z\) 包含 \(V(\mathfrak p_i)\)，便有
\(V(\mathfrak p_i)\subseteq Z\subseteq V(\mathfrak p_j)\)
对某个 \(j\) 成立；两两不包含性迫使 \(j=i\)，故 \(Z=V(\mathfrak p_i)\)。这说明所列成员就是全部不可约分量，也就是全部极小素理想对应的闭集。

最后，对任一素理想 \(\mathfrak q\)，不可约闭集 \(V(\mathfrak q)\) 包含于某个 \(V(\mathfrak p_i)\)，即 \(\mathfrak p_i\subseteq\mathfrak q\)。
\end{proof}

不同不可约分量可以相交。例如域 \(k\) 上 \(k[x,y]/(xy)\) 的两个极小素理想为 \((\bar x)\) 和 \((\bar y)\)，因为包含 \(xy\) 的素理想必含 \(x\) 或 \(y\)，而 \((x)\)、\((y)\) 本身为互不包含的素理想。两条坐标轴分量 \(V(\bar x)\)、\(V(\bar y)\) 在 \((\bar x,\bar y)\) 处相交，所以存在多个不可约分量并不意味着空间不连通。

% !TeX root = ../../Book3.tex
% 纲要标题：### 1.4 环同态与素谱
% 纲要位置：工作记录/计划纲要.md，第 1834 行。
% 纲要细目（写作范围）：
% * 素理想收缩与连续映射
% * 商环对应闭子空间
% * 主开集 \(D(f)\) 对应 \(R_f\)；一般局部化的谱对应与乘法集不交的素理想子空间，未必是开子空间
% ---

\section{环同态与素谱}
\label{b3:sec:0104}

环同态将元素沿一个方向送出，素理想则沿收缩反向移动。商环与局部化是理解这个反向关系的两种基本操作：前者给出闭子空间，后者保留避开指定分母的点。关于元素 \(f\) 的各次幂的局部化称为单元素局部化，它对应基本开集 \(D(f)\)；一般乘法集给出的子空间却未必是开集。

\subsection{素理想收缩与谱上的连续映射}
\label{b3:subsec:0104:01}

\begin{samepage}
\begin{proposition}\label{b3:prop:spectrum-functor}
设 \(R,S\) 为交换含幺环，\(\varphi:R\rightarrow S\) 为保单位的环同态。它给出映射
\[
\varphi^*:\operatorname{Spec}S\longrightarrow\operatorname{Spec}R,
\qquad\mathfrak q\longmapsto\varphi^{-1}(\mathfrak q).
\]
它连续。对每个理想 \(I\subseteq R\) 和每个元素 \(f\in R\)，有
\begin{equation}\label{b3:eq:spectrum-pullback}
(\varphi^*)^{-1}\bigl(V(I)\bigr)=V(IS),\qquad
(\varphi^*)^{-1}\bigl(D(f)\bigr)=D\bigl(\varphi(f)\bigr),
\end{equation}
其中 \(IS\) 是由全部 \(\varphi(a)\)、\(a\in I\)，在 \(S\) 中生成的理想。恒等同态诱导恒等映射。若 \(A\) 也是交换含幺环，\(\psi:S\rightarrow A\) 为保单位的环同态，则 \((\psi\circ\varphi)^*=\varphi^*\circ\psi^*\)。
\end{proposition}
\end{samepage}
\begin{proof}
因 \(1\) 不在 \(\mathfrak q\) 中，收缩是真理想；若 \(ab\) 的像在 \(\mathfrak q\)，素性使 \(\varphi(a)\) 或 \(\varphi(b)\) 在其中，所以收缩为素理想。包含 \(I\subseteq\varphi^{-1}(\mathfrak q)\) 等价于 \(IS\subseteq\mathfrak q\)，给出第一式和连续性。对主理想取补集得到第二式。最后两项由原像的恒等及复合规则直接得到。
\end{proof}

极大理想的收缩未必极大。例如 \(\mathbb Z\hookrightarrow\mathbb Q\) 把域的唯一素理想 \((0)\) 收缩为 \(\mathbb Z\) 的零理想；后者不是极大理想。这也是素谱采用全部素理想而不只采用极大理想的一个理由。

\subsection{商环与闭子空间}
\label{b3:subsec:0104:02}

\begin{theorem}\label{b3:thm:spec-quotient}
设 \(R\) 为交换含幺环，\(I\) 为 \(R\) 的理想。商映射 \(R\rightarrow R/I\) 诱导同胚
\[
\operatorname{Spec}(R/I)\longrightarrow V(I)\subseteq\operatorname{Spec}R,
\qquad\overline{\mathfrak p}\longmapsto\mathfrak p.
\]
这里 \(\mathfrak p\) 是 \(\overline{\mathfrak p}\) 的原像，逆映射为 \(\mathfrak p\mapsto\mathfrak p/I\)。
\end{theorem}
\begin{proof}
理想的商对应和商整环判据给出所列互逆双射。对包含 \(I\) 的理想 \(J\)，\(R/I\) 的闭集 \(V(J/I)\) 对应为 \(V(J)\)，它已包含于 \(V(I)\)；另一方面，\(V(I)\) 中任一子空间闭集可写为 \(V(I)\cap V(K)=V(I+K)\)，也如此得到。故双射及其逆都连续。
\end{proof}

取 \(I=\operatorname{Nil}(R)\)，所有素理想都包含它，因而得到 \(\operatorname{Spec}R_{\mathrm{red}}\cong\operatorname{Spec}R\)。底拓扑空间看不见被取商消去的幂零元。这不表示两个环相同；后面讨论加厚点与纤维时，会反复用到这一区别。

\subsection{主开集与一般局部化的谱}
\label{b3:subsec:0104:03}

设 \(R\) 为交换含幺环，\(T\subseteq R\) 为含一的乘法集，允许 \(0\in T\)，此时局部化为零环。第一卷定理13.2.2构造了环 \(T^{-1}R\)，命题13.2.3给出零判据：\(a/t=0\) 当且仅当存在 \(u\in T\) 使 \(ua=0\)。下面把第一卷定理13.4.3的局部化素理想对应补上拓扑说明。

\begin{theorem}\label{b3:thm:spec-localization}
设 \(R\) 为交换含幺环，\(T\subseteq R\) 为含一的乘法集。局部化映射 \(\iota:R\rightarrow T^{-1}R\) 诱导同胚
\[
\operatorname{Spec}(T^{-1}R)\longrightarrow
X_T=\{\,\mathfrak p\in\operatorname{Spec}R:\mathfrak p\cap T=\varnothing\,\},
\]
右侧赋予子空间拓扑。对应的逆为 \(\mathfrak p\mapsto T^{-1}\mathfrak p\)。特别地，对每个 \(f\in R\)，令 \(R_f=\{1,f,f^2,\ldots\}^{-1}R\)，则 \(\operatorname{Spec}R_f\cong D(f)\)。
\end{theorem}
\begin{proof}
局部化中的素理想不能包含可逆分母的像，故其收缩与 \(T\) 不交。若 \(\mathfrak p\cap T=\varnothing\)，记 \(\overline T\) 为 \(T\) 在 \(R/\mathfrak p\) 中的像。由第一卷定理13.2.4的局部化泛性质，商映射诱导满同态
\[
\theta:T^{-1}R\longrightarrow\overline T^{-1}(R/\mathfrak p),
\qquad a/t\longmapsto\overline a/\overline t.
\]
满射性来自分子、分母都能取原像。局部化的零判据说明，\(\theta(a/t)=0\) 当且仅当某个 \(u\in T\) 使 \(ua\in\mathfrak p\)；此时 \(a/t=(ua)/(ut)\in T^{-1}\mathfrak p\)。反过来，\(\theta\) 把生成扩张理想的每个 \(b/1\)、\(b\in\mathfrak p\)，都送到零。因此 \(\ker\theta=T^{-1}\mathfrak p\)，由环的第一同构定理得到
\[
(T^{-1}R)/(T^{-1}\mathfrak p)\cong\overline T^{-1}(R/\mathfrak p).
\]
右端是非零整环，因为 \(R/\mathfrak p\) 是整环且 \(\overline T\) 不含零。所以 \(T^{-1}\mathfrak p\) 为素理想。前面的核计算还说明，若 \(a/1\in T^{-1}\mathfrak p\)，则某个 \(ua\in\mathfrak p\)，其中 \(u\in T\)；因 \(u\notin\mathfrak p\)，素性迫使 \(a\in\mathfrak p\)。因此扩张后的收缩恢复 \(\mathfrak p\)。

反过来，设 \(\mathfrak q\) 为 \(T^{-1}R\) 的理想，并令 \(J=\iota^{-1}(\mathfrak q)\)。若 \(a/t\in\mathfrak q\)，则 \(a/1=(t/1)(a/t)\in\mathfrak q\)，所以 \(a\in J\)，从而 \(a/t\in T^{-1}J\)。若 \(a\in J\)、\(t\in T\)，则 \(a/t=(1/t)(a/1)\in\mathfrak q\)。故 \(T^{-1}J=\mathfrak q\)，证明了收缩后扩张也恢复原理想。这验证了素理想之间的双射。

由\cref{b3:prop:spectrum-functor}，该映射连续。局部化中的基本开集 \(D(a/t)\) 在上述双射下对应 \(D(a)\cap X_T\)：素理想是否含 \(a/t\)，只取决于它是否含分子。这些正是子空间的一组开集基，故逆映射也连续。取 \(T=\{1,f,f^2,\ldots\}\) 时，素理想与 \(T\) 不交恰好等价于不含 \(f\)，得到主开集的断言。
\end{proof}

一般的 \(X_T=\bigcap_{t\in T}D(t)\) 未必开。取 \(R=\mathbb Z\)、\(T=\mathbb Z\setminus\{0\}\)，其像只有一般点 \((0)\)。任何包含它的基本开集 \(D(n)\)、\(n\ne0\)，都还包含一个 \((p)\)，其中素数 \(p\) 不整除 \(n\)；这样的素数存在，因为有限多个素数不能穷尽全部素数。因此单点 \(\bigl\{(0)\bigr\}\) 不是开集，尽管它与 \(\operatorname{Spec}\mathbb Q\) 同胚。

\ProseParagraphBreak
\cref{b3:tab:spectrum-operations}将上述操作并列。表中 \(R\) 为交换含幺环，\(I\) 为理想，\(f\in R\)，\(\mathfrak p\) 为素理想，\(T\) 为含一的乘法集；各行都经收缩把新环的谱视为原谱的子空间。在 \(\mathfrak p\) 处局部化保留包含于 \(\mathfrak p\) 的素理想，取商 \(R/\mathfrak p\) 则保留包含 \(\mathfrak p\) 的素理想，两个方向相反。

\begin{table}[htbp]
\centering
\caption{取商与局部化在素谱中保留的点}
\label{b3:tab:spectrum-operations}
\begin{tabularx}{\textwidth}{@{}l>{\raggedright\arraybackslash}p{.29\textwidth}>{\raggedright\arraybackslash}X@{}}
\toprule
环的操作 & 原环中保留的素理想 \(\mathfrak q\) & 在 \(\operatorname{Spec}R\) 中的像 \\
\midrule
\(R\rightarrow R/I\) & \(I\subseteq\mathfrak q\) & \(V(I)\)，闭子空间 \\
\(R\rightarrow R_f\) & \(f\notin\mathfrak q\) & \(D(f)\)，基本开集 \\
\(R\rightarrow R_{\mathfrak p}\) & \(\mathfrak q\subseteq\mathfrak p\) & \(\mathfrak p\) 及其一般化组成的子空间，未必开 \\
\(R\rightarrow T^{-1}R\) & \(\mathfrak q\cap T=\varnothing\) & \(\bigcap_{t\in T}D(t)\)，未必开 \\
\bottomrule
\end{tabularx}
\end{table}

第二章将把这些环上的分式构造推广到模。


\begin{chaptersummary}
素性与极大性分别使商环成为整环与域。避开乘法集的极大性论证不仅给出素理想存在性，还证明 \(\sqrt I\) 等于包含 \(I\) 的全部素理想之交。这个交精确刻画 \(R/I\) 中哪些元素幂零，却不能恢复原理想 \(I\)，也不记录元素或幂零理想的具体幂零指数。有限生成的逐元素幂零理想必为幂零理想；没有有限生成条件时，即使每个元素的平方都为零，理想的各次幂仍可能都非零。

\ProseParagraphBreak
Zariski 闭集对应根理想，基本开集对应元素不被素理想包含。全谱和每个基本开集都拟紧；每个素理想都是其闭包的一般点，极大理想恰好给出闭点，而极小素理想给出全谱的不可约分量。Noether 性使不可约分量有限，却不能从底拓扑空间反推出原环一定是 Noether 环。\(k[t]/(t^m)\) 对不同 \(m\ge1\) 总有相同的单点谱，商环中 \(\bar t\) 的幂零指数却不同。

\ProseParagraphBreak
环同态使素理想反向收缩。取商给闭子空间，单元素局部化给基本开集，一般局部化则给避开全部分母的子空间，它未必开。下一章把这种局部观察推广到模，并说明怎样从各点处的结果恢复全局结论。
\end{chaptersummary}

\BookExercises

\begin{exercise}\label{b3:ex:01-axes-primes}
设 \(k\) 为域，\(R=k[x,y]/(xy)\)。证明 \((\bar x)\)、\((\bar y)\) 素而不极大，\((\bar x,\bar y)\) 极大。证明每个素理想至少包含 \(\bar x,\bar y\) 中一个。
\end{exercise}
\begin{solution}
前两个商分别是 \(k[y],k[x]\)，为整环而非域；第三个商为 \(k\)，由\cref{b3:prop:prime-maximal-quotient}得到结论。又 \(\bar x\bar y=0\) 属于任何素理想，素性迫使至少一个因子在其中。因此这两个不可约分量已经覆盖全谱。
\end{solution}

\begin{exercise}\label{b3:ex:01-residue-field-four}
求 \(\mathbb F_2[t]/(t^2+t+1)\) 的元素与非零元素的逆，说明对应理想是极大理想，但剩余域不是原系数域。
\end{exercise}
\begin{solution}
多项式在零、一处都不为零，故这个二次多项式不可约。商以 \(1,\alpha\) 为基，有四个元素 \(0,1,\alpha,1+\alpha\)，其中 \(\alpha^2=\alpha+1\)。等式 \(\alpha(1+\alpha)=1\) 给出后二者互逆，一自逆，所以商为四元域。商域判据给出极大性，元素个数则说明它不是 \(\mathbb F_2\)。
\end{solution}

\begin{exercise}\label{b3:ex:01-radical-powers}
设 \(k\) 为域。在 \(k[x,y]\) 中证明 \(\sqrt{(x^2,y^3)}=(x,y)\)，并求商环的幂零根的最小正幂零指数。
\end{exercise}
\begin{solution}
两个变量都有幂在原理想中，故 \((x,y)\) 包含于根；反向因原理想包含于极大理想 \((x,y)\)，其根也包含于其中。商中的幂零根 \(N=(\bar x,\bar y)\)。任一四次单项式都含 \(x^2\) 或 \(y^3\)，故 \(N^4=0\)；但 \(\bar x\bar y^2\in N^3\) 非零，因为 \((x^2,y^3)\) 中每个单项式都含 \(x^2\) 或 \(y^3\)，故 \(N^3\ne0\)。最小正指数为四。
\end{solution}

\begin{exercise}\label{b3:ex:01-radical-sum}
设 \(R\) 为交换含幺环，\(I,J\) 为 \(R\) 的理想。证明 \(\sqrt{I+J}=\sqrt{\sqrt I+\sqrt J}\)。再取域 \(k\)，以 \(I=(y),J=(y-x^2)\subseteq k[x,y]\) 说明两个根理想之和未必为根理想。
\end{exercise}
\begin{solution}
取根保持理想的包含关系，所以 \(I+J\subseteq\sqrt I+\sqrt J\) 给出
\[
\sqrt{I+J}\subseteq\sqrt{\sqrt I+\sqrt J}.
\]
另一方面，由 \(I,J\subseteq I+J\) 得 \(\sqrt I,\sqrt J\subseteq\sqrt{I+J}\)，故 \(\sqrt I+\sqrt J\subseteq\sqrt{I+J}\)。再取根，并用 \(\sqrt{I+J}\) 已是根理想，得到
\[
\sqrt{\sqrt I+\sqrt J}\subseteq\sqrt{I+J}.
\]
两式给出所求等式；这里取根保持包含，而闭集运算 \(V\) 才反转包含。例中两个商都为 \(k[x]\)，所以 \(I,J\) 为素理想，特别是根理想；但其和为 \((y,x^2)\)，其中 \(x^2\) 属于它而 \(x\) 不属于它，故和不为根理想。
\end{solution}

\begin{exercise}\label{b3:ex:01-units-nil-quotient}
设 \(R\) 为交换含幺环，\(J\subseteq R\) 为理想，且 \(J\) 的每个元素都幂零。证明 \(a\in R\) 可逆当且仅当 \(a+J\) 在 \(R/J\) 中可逆；这里不要求 \(J\) 本身为幂零理想。
\end{exercise}
\begin{solution}
只有反向需要说明。若 \(ab-1=u\in J\)，选 \(N\) 使 \(u^N=0\)，则 \(1+u\) 的逆为 \(\sum_{j=0}^{N-1}(-u)^j\)。所以 \(b\sum_{j=0}^{N-1}(-u)^j\) 是 \(a\) 的逆。只用了这一个元素 \(u\) 的幂零性，不需要对全部 \(J\) 同时选一个指数。
\end{solution}

\begin{exercise}\label{b3:ex:01-integer-radical}
求 \(\mathbb Z/18\mathbb Z\) 的全部素理想、幂零根和约化商。它的谱有几个点，哪些点闭？
\end{exercise}
\begin{solution}
由商环对应，素理想来自 \(\mathbb Z\) 中包含十八的素理想，即 \((2),(3)\)。故幂零根为两像理想的交 \((\bar6)\)，包含 \(0,\bar6,\bar{12}\)，其平方为零。约化商为 \(\mathbb Z/6\mathbb Z\)，谱有两个点，都是极大理想对应的闭点。两个单点的补集也闭，因此这个有限谱是离散的。
\end{solution}

\begin{exercise}\label{b3:ex:01-product-spectrum}
设 \(A,B\) 为交换含幺环。证明 \(\operatorname{Spec}(A\times B)\) 是 \(\operatorname{Spec}A\) 与 \(\operatorname{Spec}B\) 的不交并，且这两个部分都是开闭子集。
\end{exercise}
\begin{solution}
令 \(e=(1,0)\)、\(f=(0,1)\)。因 \(ef=0,e+f=1\)，素理想恰好包含二者之一。若包含 \(f\)，它形如 \(\mathfrak p\times B\)，且商为 \(A/\mathfrak p\)，故 \(\mathfrak p\) 必素；另一类为 \(A\times\mathfrak q\)。两部分分别是 \(D(e)=V(f)\)、\(D(f)=V(e)\)，均开闭。各部分的拓扑由商环同胚\cref{b3:thm:spec-quotient}识别为相应因子的谱，所以这不仅是集合双射。
\end{solution}

\begin{exercise}\label{b3:ex:01-infinite-prime-avoidance}
设 \(k\) 为域，令 \(I=(x,y)\subseteq k[x,y]\)。对全部常数项为零的不可约多项式 \(h\)，考察素理想 \((h)\)。证明它们的并包含 \(I\)，但其中没有一个包含整个 \(I\)。
\end{exercise}
\begin{solution}
域上的多项式环为唯一分解整环，这是第一卷定理12.2.6逐次添加两个变量的结论。对 \(0\ne f\in I\)，不可约分解的常数项乘积为零，所以某个不可约因子 \(h\) 的常数项为零；于是 \(f\in(h)\)。零也在这些理想中。若 \(I\subseteq(h)\)，则 \(h\) 同时整除 \(x,y\)，只能为单位，矛盾。对 \(n\ge1\)，商环 \(k[x,y]/(x-y^n)\cong k[y]\) 是整环，故 \(x-y^n\) 为素元，特别是不可约元；它们已经给出无限多个不同的此类不可约元。这个例子说明有限素理想回避不能直接推广到无限并。
\end{solution}

\begin{exercise}\label{b3:ex:01-closed-integers}
在 \(\operatorname{Spec}\mathbb Z\) 中求 \(V(30)\)、\(D(6)\) 及点集 \(\bigl\{(2),(3)\bigr\}\) 的闭包。再证明任意无限多个不同的闭点组成的集合稠密。
\end{exercise}
\begin{solution}
\(V(30)=\bigl\{(2),(3),(5)\bigr\}\)，\(D(6)\) 包含一般点 \((0)\) 及除 \((2),(3)\) 外的全部闭点。所给两个点的集合本身等于 \(V(6)\)，所以已经闭。若一个闭集 \(V(n)\) 包含无限多个不同闭点，则整数 \(n\) 被无限多个素数整除，只能 \(n=0\)。因此包含这些点的唯一闭集是全谱，它们的闭包即为全谱。
\end{solution}

\begin{exercise}\label{b3:ex:01-affine-line-spectrum}
设 \(k\) 为代数闭域。描述 \(\operatorname{Spec}k[t]\) 的点及全部真闭集，并证明两个非空开集必相交。
\end{exercise}
\begin{solution}
一元多项式环的非零素理想由不可约多项式生成；代数闭性使这些多项式都是一次式。因此点为 \((0)\) 及 \((t-a)\)、\(a\in k\)。任何真闭集形如 \(V(f)\)、\(f\ne0\)，恰为这个多项式的有限多个根对应的闭点。反过来，任意有限闭点集 \(\{(t-a_1),\ldots,(t-a_r)\}\) 都等于
\[
V\!\left(\prod_{i=1}^{r}(t-a_i)\right).
\]
空集对应空乘积 \(1\)，即 \(V(1)=\varnothing\)。一般点 \((0)\) 不属于任何真闭集，所以非空开集必含 \((0)\)，任意两个这样的开集因而相交。其不可约性也可由整环的零理想为素理想得到。
\end{solution}

\begin{exercise}\label{b3:ex:01-local-open-no-closed-point}
设 \(k\) 为域，\(R=k[t]_{(t)}\)。分别计算 \(R/(t^2)\) 与 \(R[1/t]\) 的素谱，以及它们在 \(\operatorname{Spec}R\) 中的像。判断这两个像是否包含 \(\operatorname{Spec}R\) 的闭点，说明它们虽同为单点空间，为何处于原谱的不同位置。
\end{exercise}
\begin{solution}
\(\operatorname{Spec}R\) 的两个点为 \(\eta=(0)\)、\(s=(t)R\)，其中只有 \(s\) 为闭点。由\cref{b3:thm:spec-quotient}，\(\operatorname{Spec}(R/(t^2))\) 的像是 \(V(t^2)=V(t)=\{s\}\)，所以它含有原谱的闭点。

由\cref{b3:thm:spec-localization}，\(\operatorname{Spec}(R[1/t])\) 的像为 \(D(t)=\{\eta\}\)，不含原谱的闭点。事实上，\(R\) 中每个非零元素都可写成 \(t^n u\)，其中 \(n\ge0\)、\(u\) 为单位；再使 \(t\) 可逆便使全部非零元素可逆，得到 \(R[1/t]=k(t)\)。两个源谱都是单点空间，但商映射保留的是闭点，局部化保留的是一般点。\(\eta\) 在开子空间 \(\{\eta\}\) 中是闭点，在原谱中却不是；闭点的含义取决于所在空间。
\end{solution}

\begin{exercise}\label{b3:ex:01-nonquasicompact-open}
设 \(k\) 为域，在 \(R=k[x_1,x_2,\ldots]\) 中令 \(U=\bigcup_{i\ge1}D(x_i)\)。证明 \(U\) 不拟紧，尽管全谱及每个 \(D(x_i)\) 都拟紧。
\end{exercise}
\begin{solution}
若有限个 \(D(x_i)\) 覆盖 \(U\)，取未在这有限指标中出现的 \(j\)。把所选变量送到零、\(x_j\) 送到一、其余变量送到零，得到满同态 \(R\rightarrow k\)。它的核为极大理想，属于 \(D(x_j)\subseteq U\)，却不在有限个所选基本开集内，矛盾。所以这个开覆盖没有有限子覆盖。
\end{solution}

\begin{exercise}\label{b3:ex:01-noether-topology-not-ring}
设 \(k\) 为域，\(V\) 为无限维 \(k\) 线性空间，在 \(R=k\oplus V\) 上规定 \((a,v)(b,w)=(ab,aw+bv)\)。证明 \(R\) 的谱只有一个点，但 \(R\) 不是 Noether 环。
\end{exercise}
\begin{solution}
对第三个元素 \((c,u)\)，乘法的两种结合次序都得到 \((abc,ab u+ac w+bc v)\)，单位为 \((1,0)\)。理想 \(0\oplus V\) 平方为零，商为 \(k\)，所以它是唯一素理想。谱因而是单点 Noether 空间。选无限线性无关序列 \(v_i\)，其有限初段所张成的 \(0\oplus\operatorname{span}_k\{v_1,\ldots,v_n\}\) 都是理想，并组成严格升链。故 \(R\) 不是 Noether 环；底拓扑空间无法检测全部环的有限性。
\end{solution}

\begin{exercise}\label{b3:ex:01-components-plane-line}
设 \(k\) 为域。在 \(k[x,y,z]\) 中证明 \((xy,xz)=(x)\cap(y,z)\)，并求 \(V(xy,xz)\) 的全部不可约分量及其一般点。
\end{exercise}
\begin{solution}
左边显然包含于右边。若 \(xh\in(y,z)\)，模去 \((y,z)\) 得 \(x\cdot h(x,0,0)=0\) 于整环 \(k[x]\)，所以 \(h(x,0,0)=0\)，即 \(h\in(y,z)\)。故反向包含成立。两个理想的商分别为 \(k[y,z]\) 与 \(k[x]\)，所以都素，且互不包含。因此不可约分量为 \(V(x)\)、\(V(y,z)\)，一般点分别为 \((x)\)、\((y,z)\)。它们的交为 \(V(x,y,z)\)，非空；不可约分量不是必然互不相交的部分。
\end{solution}

\begin{exercise}\label{b3:ex:01-thickened-point}
设 \(k\) 为域，\(m\ge1\) 为整数。比较环 \(k[t]/(t^m)\) 的素谱、\(k\) 维数和幂零根的幂零指数，说明谱同胚不能决定环同构。
\end{exercise}
\begin{solution}
每个素理想都含 \(\bar t\)，商为域，所以谱总为一个点。环以 \(1,\bar t,\ldots,\bar t^{m-1}\) 为基，维数为 \(m\)；\((\bar t)^m=0\)，若 \(m>1\)，其 \(m-1\) 次幂非零，因此最小正指数为 \(m\)。不同 \(m\) 的这些 \(k\) 代数因维数或幂零指数不同而不同构，但底谱相同。
\end{solution}

\begin{exercise}\label{b3:ex:01-maps-direction}
在交换含幺环及保单位环同态中，分别给出反例，说明环同态单射未必使谱映射单射，以及谱映射同胚未必使环同态单射或满射。
\end{exercise}
\begin{solution}
取域 \(k\)。包含 \(k\hookrightarrow k[t]\) 单射，谱上的映射却把多个点全部送到 \(\operatorname{Spec}k\) 的唯一点。商映射 \(k[\varepsilon]/(\varepsilon^2)\rightarrow k\) 不是单射，却在谱上给同胚。一个真域扩张 \(k\hookrightarrow K\) 不是满射，但两端的谱都是单点，也给同胚。环的元素及剩余域信息并不能由这个底空间映射单独恢复。
\end{solution}

\begin{exercise}\label{b3:ex:01-two-localizations}
描述 \(\operatorname{Spec}\mathbb Z[1/6]\) 和 \(\operatorname{Spec}\mathbb Z_{(2)}\) 在 \(\operatorname{Spec}\mathbb Z\) 中的像，判断两者是否开。
\end{exercise}
\begin{solution}
前者的像是 \(D(6)\)，为开集。后者只保留包含于 \((2)\) 的素理想，即 \(\bigl\{(0),(2)\bigr\}\)。任何包含一般点的开集都包含某个 \(D(n)\)、\(n\ne0\)，而这个基本开集还含无限多个其他素数点。因此后者不是开集。两种像均按\cref{b3:thm:spec-localization}具有正确的子空间拓扑。
\end{solution}

\begin{exercise}\label{b3:ex:01-square-map}
设 \(k\) 为域，\(b\in k\)，令 \(\varphi:k[s]\rightarrow k[t]\) 为满足 \(\varphi(s)=t^2\) 的 \(k\)-代数同态。计算闭点 \((s-b)\) 在 \(\varphi^*\) 下的原像及其补集。分别在 \(b\) 不是 \(k\) 中的平方、\(b=0\)、\(b=a^2\ne0\) 的情形下，求原像的点数及这些点的剩余域，并判断 \(k[t]/(t^2-b)\) 是否约化；最后一种情形须区分 \(\operatorname{char}k=2\) 与 \(\operatorname{char}k\ne2\)。
\end{exercise}
\begin{solution}
\((s-b)\) 是极大理想，所以 \(V(s-b)=\{(s-b)\}\)。由\eqref{b3:eq:spectrum-pullback}，所求原像及其补集分别为 \(V(t^2-b)\) 与 \(D(t^2-b)\)。包含 \(t^2-b\) 的素理想收缩后包含 \((s-b)\)，而收缩仍是真理想，故恰好收缩为这个闭点。

若 \(b\) 不是平方，则二次多项式 \(t^2-b\) 没有 \(k\) 中的根，因而不可约。原像只有 \((t^2-b)\) 一个点，剩余域是二次扩张 \(k[t]/(t^2-b)\)；这个商是域，所以约化。

若 \(b=0\)，原像只有 \((t)\)，剩余域为 \(k\)，但商 \(k[t]/(t^2)\) 中 \(\bar t\ne0\) 且 \(\bar t^2=0\)，因此不约化。若 \(b=a^2\ne0\) 且特征不为二，则
\[
t^2-b=(t-a)(t+a),\qquad a\ne-a.
\]
原像有 \((t-a)\)、\((t+a)\) 两个点，剩余域都为 \(k\)。两个因子互素，中国剩余定理给出商环 \(k\times k\)，所以约化。特征为二时，\(t^2-b=(t-a)^2\)，原像只有 \((t-a)\)，剩余域为 \(k\)，商中 \(t-a\) 的像非零而平方为零，所以不约化。原像同为单点时，剩余域和商环的幂零元仍可能不同。
\end{solution}

\begin{exercise}\label{b3:ex:01-reduction-universal}
设 \(R,S\) 为交换含幺环，且 \(S\) 约化。证明每个保单位环同态 \(R\rightarrow S\)，都唯一经 \(R_{\mathrm{red}}\) 分解。说明为何仅有谱同胚不能反向推出这个同态的泛性质。
\end{exercise}
\begin{solution}
幂零元的像仍幂零，而 \(S\) 约化，所以同态零化 \(\operatorname{Nil}(R)\)。定义 \(r+\operatorname{Nil}(R)\mapsto\varphi(r)\) 即得到良定义的环同态；原商映射满射，故分解唯一。这个泛性质谈的是到每个约化环的全部同态。例如 \(\mathbb Q\hookrightarrow\mathbb Q(i)\) 的谱同样是单点间同胚，若恒等映射 \(\mathbb Q\rightarrow\mathbb Q\) 能经它分解，就须有保单位环同态 \(\psi\) 使
\[
\mathbb Q\hookrightarrow\mathbb Q(i)\xrightarrow{\psi}\mathbb Q
\]
的复合为恒等映射。第一支箭头已经存在，第二支却不存在：由 \(i^2=-1\) 将得到 \(\psi(i)^2=-1\)，而有理数中没有这样的元素。因此底拓扑同胚没有提供这一代数泛性质。
\end{solution}

\begin{exercise}\label{b3:ex:01-irreducible-not-domain}
设 \(R\) 为交换含幺环。证明 \(\operatorname{Spec}R\) 不可约当且仅当 \(\operatorname{Nil}(R)\) 为素理想。给出一个非整环但谱不可约的非零环，并说明加入约化条件后能得到什么。
\end{exercise}
\begin{solution}
零环的谱为空，不可约性与幂零根为素理想这两项均不成立。对非零环，结论是\cref{b3:prop:spec-irreducible}在 \(I=0\) 时的应用。取域 \(k\)，双数环 \(k[\varepsilon]/(\varepsilon^2)\) 的幂零根 \((\varepsilon)\) 为极大理想，故谱不可约，但 \(\varepsilon\ne0\) 且平方为零，环不是整环。若另有约化性，幂零根为零，所以零理想素，商整环判据说明 \(R\) 本身是整环。
\end{solution}

\begin{exercise}\label{b3:ex:01-basic-open-certificate}
设 \(R\) 为交换含幺环，\(f\in R\)，\((g_\lambda)_{\lambda\in\Lambda}\) 为 \(R\) 中的一族元素。证明
\[
D(f)\subseteq\bigcup_{\lambda\in\Lambda}D(g_\lambda)
\]
当且仅当存在有限子集 \(F\subseteq\Lambda\)、整数 \(N\ge1\) 和元素 \(a_\lambda\in R\)（\(\lambda\in F\)），使
\[
f^N=\sum_{\lambda\in F}a_\lambda g_\lambda.
\]
再在 \(R=\mathbb Z\) 中取 \(f=6\)、\(g_1=8\)、\(g_2=20\)，求这样的最小 \(N\)，并写出一组系数。
\end{exercise}
\begin{solution}
令 \(J=(g_\lambda)_{\lambda\in\Lambda}\)。对所给开集包含关系取补集，得到 \(V(J)\subseteq V(f)\)。由\eqref{b3:eq:zariski-radical-order}，这等价于 \(f\in\sqrt J\)，即某个正次幂 \(f^N\) 属于 \(J\)。理想生成中的每个表达式只含有限项，因而恰好得到题设的有限等式。反过来，若该等式成立，任何包含全部 \(g_\lambda\) 的素理想都包含 \(f^N\)，从而包含 \(f\)，所以 \(V(J)\subseteq V(f)\)，再取补集便得开集包含关系。\(F\) 为空时，等式为 \(f^N=0\)，此时 \(D(f)=\varnothing\)。

整数例中 \((8,20)=(4)\)。\(6\) 不是四的倍数，而 \(6^2=36\) 是，所以最小指数为 \(N=2\)。等式
\[
36=2\cdot8+20
\]
给出系数 \(a_1=2\)、\(a_2=1\)，也给出无需逐个检查素理想的包含证书。
\end{solution}

\begin{exercise}\label{b3:ex:01-spectrum-image-closure}
设 \(R,S\) 为交换含幺环，\(\varphi:R\rightarrow S\) 为保单位的环同态，\(\varphi^*:\operatorname{Spec}S\rightarrow\operatorname{Spec}R\) 为收缩映射。证明
\[
\overline{\varphi^*(\operatorname{Spec}S)}=V(\ker\varphi).
\]
\end{exercise}
\begin{solution}
对 \(X=\operatorname{Spec}R\) 的任意点集 \(E\)，包含 \(E\) 的闭集 \(V(I)\) 恰好满足 \(I\subseteq\bigcap_{\mathfrak p\in E}\mathfrak p\)。因此
\[
\overline E=V\!\left(\bigcap_{\mathfrak p\in E}\mathfrak p\right),
\]
其中空交取为 \(R\)，也给出 \(\overline\varnothing=V(R)=\varnothing\)。取 \(E=\varphi^*(\operatorname{Spec}S)\)，由\cref{b3:thm:radical-prime-intersection}，
\[
\bigcap_{\mathfrak p\in E}\mathfrak p
=\bigcap_{\mathfrak q\in\operatorname{Spec}S}\varphi^{-1}(\mathfrak q)
=\varphi^{-1}\bigl(\operatorname{Nil}(S)\bigr).
\]
一个元素 \(a\in R\) 属于最后的理想，当且仅当某个 \(\varphi(a)^n=0\)，也就是 \(a^n\in\ker\varphi\)。故该理想等于 \(\sqrt{\ker\varphi}\)，而取根不改变零集，得到所需等式。
\end{solution}

% !TeX root = ../../Book3.tex
% 纲要标题：## 第2章　模的局部化、局部环与局部判据
% 纲要位置：工作记录/计划纲要.md，第 1941 行。

\chapter{模的局部化、局部环与局部判据}
\label{b3:ch:02}
\BookChapterNumber{b3:ch:02}

\begin{chapterabstract}
把允许的分母作用在模上，既会消去部分元素，也能把全局问题分解到各个素理想处。本章构造模的局部化，证明它保持正合序列，并据此逐点检测元素的消失、子模成员关系以及给定同态的单射性与满射性。有限生成与有限展示分别控制生成元和关系的共同分母；相容的局部元素则能在有限基本开覆盖上黏合。
\end{chapterabstract}

\begin{learninggoals}
\begin{itemize}
\item 用分式零判据处理模局部化中的核、商与子模。
\item 区分乘法集饱和、理想幂饱和、局部模及剩余向量空间。
\item 正确使用支集与 Hom 的局部化公式，并指出所需的有限性条件。
\item 用局部环检测给定同态，说明从点处结论扩展到邻域需要哪些有限性条件。
\item 对有限基本开覆盖检验局部模元素的相容性，并由此构造唯一的全局元素。
\end{itemize}
\end{learninggoals}

取底环 \(R=\mathbb Z\)，把 \(M=\mathbb Z/12\mathbb Z\) 看成 \(R\) 模。元素 \(\bar3\) 被 \(4\) 消去，因为 \(4\bar3=0\)。若在 \(R\) 的素理想 \((5)\) 处局部化，\(4\) 可以作为分母，\(\bar3\) 于是成为零；在 \((2)\) 处，允许的分母都是奇数，而零化 \(\bar3\) 的整数恰是 \(4\) 的倍数，所以 \(\bar3/1\) 仍非零。这里的素理想取自底环 \(R\)，不是从商环 \(\mathbb Z/12\mathbb Z\) 中选取。元素是否消失，由是否存在一个允许的分母将它消去来判定。

\ProseParagraphBreak

从一个元素转向整个模，量词便发生变化：每个生成元分别有消去它的分母，未必有一个分母同时消去所有元素。若模有限生成，把有限多个分母相乘就能得到这样的共同分母；控制全部关系则常需有限展示。支集记录哪些点仍看见这个模。另一方面，若一个给定同态在各点处都单射或都满射，局部化的正合性就能把这些结论还原为全局结论。

\ProseParagraphBreak

本章需要第一章的环局部化和素谱，以及第二卷第2章的模、核与商、第6章的张量积。除明确说明外，不假设底环是 Noether 环，也不假设所讨论的模有限生成。给单个元素清分母，与给全部生成元或全部关系选取共同分母，是两种不同的任务；这项区别决定了多数定理的有限性假设。

\ProseParagraphBreak
进一步阅读可参阅 Matsumura 的《Commutative Ring Theory》\cite[Theorems 4.4--4.6, pp. 26--27]{Matsumura1989CommutativeRingTheory}，对照局部化、张量积与局部消失的证明。该书\cite[Theorem 4.10, pp. 28--29]{Matsumura1989CommutativeRingTheory}还证明：对有限生成模及固定整数 \(r\ge0\)，局部模能由至多 \(r\) 个元素生成的点组成开集；对有限展示模，局部模为自由模的点也组成开集。

% !TeX root = ../../Book3.tex
% 纲要标题：### 2.1 模局部化
% 纲要位置：工作记录/计划纲要.md，第 1844 行。
% 纲要细目（写作范围）：
% * 分式模的构造
% * 与环局部化的张量描述
% * 局部化的正合性

\section{模局部化}
\label{b3:sec:0201}

环的局部化允许把指定元素当作分母；模也可以作同样的变换。不过，模中通常没有乘法，分式相等必须用标量作用来定义。以下始终设 \(R\) 为交换含幺环，\(S\subseteq R\) 含 \(1\) 且对乘法封闭，\(M\) 为幺性 \(R\) 模。允许 \(0\in S\)，此时局部化环及局部化模都为零。第二卷命题6.4.5已给出模局部化的张量描述；本节重新展开分母的处理，使后续局部判据能够直接使用。

\subsection{分式模的构造}
\label{b3:subsec:0201:01}

\begin{definition}\label{b3:def:module-localization}
设 \(R\) 为交换含幺环，\(S\subseteq R\) 为含一的乘法集，\(M\) 为 \(R\) 模。在 \(M\times S\) 上规定
\[
(m,s)\sim(n,t)
\quad\Longleftrightarrow\quad
\text{存在 }u\in S\text{ 使 }u(tm-sn)=0.
\]
所得等价类记为 \(m/s\)，全部等价类组成的集合记为 \(S^{-1}M\)，称为 \(M\) 关于 \(S\) 的\term[mo-jubuhua]{局部化}[localization of a module]。若 \(\mathfrak p\) 是 \(R\) 的素理想，则 \(S=R\setminus\mathfrak p\) 时记为 \(M_{\mathfrak p}\)；若 \(f\in R\)，则 \(S=\{1,f,f^2,\ldots\}\) 时记为 \(M_f\)。
\end{definition}
\symidx{\(S^{-1}M\)、\(M_{\mathfrak p}\)、\(M_f\)：模的局部化}

先核对这确实是等价关系。自反性取 \(u=1\)，对称性把等式取负。若 \(u(tm-sn)=0\)、\(v(rn-tp)=0\)，则
\[
uvt(rm-sp)=uv\bigl(r(tm-sn)+s(rn-tp)\bigr)=0,
\]
所以传递性也成立。关系中的额外乘子 \(u\) 不能省去：若某个分母已经零化了 \(m\)，局部化后 \(m\) 就应消失。

\begin{proposition}\label{b3:prop:module-fraction-operations}
设 \(R\) 为交换含幺环，\(S\subseteq R\) 为含一的乘法集，\(M\) 为 \(R\) 模。下列公式使 \(S^{-1}M\) 成为 \(S^{-1}R\) 模：
\[
\frac ms+\frac nt=\frac{tm+sn}{st},\qquad
\frac as\frac mt=\frac{am}{st}.
\]
自然映射 \(\iota_M:M\rightarrow S^{-1}M\)、\(m\mapsto m/1\) 为 \(R\) 线性映射，且
\begin{equation}\label{b3:eq:localized-element-zero}
\frac ms=0\quad\Longleftrightarrow\quad
um=0\text{ 对某个 }u\in S\text{ 成立}.
\end{equation}
若 \(N\) 是 \(S^{-1}R\) 模，任意 \(R\) 线性映射 \(h:M\rightarrow N\) 唯一延拓为 \(S^{-1}R\) 线性映射
\[
\widetilde h:S^{-1}M\longrightarrow N,\qquad
\widetilde h(m/s)=(s/1)^{-1}h(m).
\]
这唯一的 \(S^{-1}R\) 线性延拓使下图交换：
\[
\begin{tikzcd}[column sep=large,row sep=large]
M\arrow[r,"\iota_M"]\arrow[dr,"h"']
&S^{-1}M\arrow[d,dashed,"\widetilde h"]\\
&N
\end{tikzcd}
\]
\end{proposition}
\begin{proof}
若 \(m/s=m'/s'\)，取 \(u\in S\) 使 \(u(s'm-sm')=0\)。加上同一个 \(n/t\) 后，两分子的交叉差为 \(t^2(s'm-sm')\)，仍被 \(u\) 零化；乘上同一个 \(a/t\) 后，交叉差为 \(at(s'm-sm')\)，也被 \(u\) 零化。改变第二个加数的代表元同理。若 \(a/t=a'/t'\)，相应交叉差为 \(s(t'a-ta')m\)，由环分式的等价关系零化。因此运算良定义。把有限个分式通分后，结合律、分配律、零元 \(0/1\)、负元 \(-m/s\) 以及单位作用，都化为 \(M\) 中的相应公理。零判据直接把 \(m/s\) 与 \(0/1\) 比较。

若 \(u(tm-sn)=0\)，在 \(N\) 中 \(u,s,t\) 均可逆，施加 \(h\) 再消去这些单位，即得 \(s^{-1}\cdot h(m)=t^{-1}\cdot h(n)\)。这证明延拓良定义；通分验证线性。因 \(m/s=(1/s)(m/1)\)，任何这样的延拓都只能采用所给公式，故唯一。
\end{proof}

例如 \(R=\mathbb Z\)、\(M=\mathbb Z/12\mathbb Z\)，在 \((2)\) 处局部化时 \(3\) 可逆，原来阶为 \(3\) 的部分消失，得到 \(M_{(2)}\cong\mathbb Z/4\mathbb Z\)。确切的同构由 \(\bar a/s\mapsto \bar a\bar s^{-1}\) 给出，其中 \(s\) 为奇数；其核为零可由 \(4\mid a\Rightarrow12\mid3a\) 核验。

\subsection{模局部化的张量描述}
\label{b3:subsec:0201:02}

\begin{proposition}\label{b3:prop:localization-tensor}
设 \(R\) 为交换含幺环，\(S\subseteq R\) 为含一的乘法集，\(M\) 为 \(R\) 模。存在关于 \(M\) 自然的 \(S^{-1}R\) 模同构
\[
(S^{-1}R)\otimes_R M\longrightarrow S^{-1}M,
\qquad \frac as\otimes m\longmapsto\frac{am}s.
\]
其逆映射为 \(m/s\mapsto(1/s)\otimes m\)。
\end{proposition}
\begin{proof}
正向公式对两个变量可加，且 \((ar/s)\otimes m\) 与 \((a/s)\otimes rm\) 具有相同的像，故由张量积的泛性质给出线性映射。若 \(u(tm-sn)=0\)，则
\[
\frac1s\otimes m-\frac1t\otimes n
=\frac1{ust}\otimes u(tm-sn)=0.
\]
因此逆向公式良定义，通分后可直接验证它线性。两种公式的复合在 \(m/s\) 及纯张量 \((a/s)\otimes m\) 上均为恒等，故互为逆映射。模同态 \(f:M\rightarrow N\) 在两边都把分子 \(m\) 变为 \(f(m)\)，这就是自然性。
\end{proof}

局部化因而是沿 \(R\rightarrow S^{-1}R\) 的标量扩张。它允许把原来的生成元乘以新的分式标量，并消去被某个分母零化的元素。例如对任意理想 \(I\)，该描述给出
\[
S^{-1}(M/IM)\cong(S^{-1}M)/(S^{-1}I)(S^{-1}M).
\]
这里也可以把分式 \((m+IM)/s\) 直接送到 \(m/s\) 的商类；其良定义和核将在下一小节的正合性中一并得到。

\subsection{局部化的正合性}
\label{b3:subsec:0201:03}

\begin{theorem}\label{b3:thm:localization-exact}
设 \(R\) 为交换含幺环，\(S\subseteq R\) 为含一的乘法集，\(L,M,N\) 为 \(R\) 模。若 \(f:M\rightarrow N\) 是 \(R\) 模同态，则它诱导映射 \(S^{-1}f(m/s)=f(m)/s\)。这一定义给出正合函子：若
\[
0\longrightarrow L\xrightarrow{i}M\xrightarrow{q}N\longrightarrow0
\]
正合，则局部化后的序列仍正合。因此局部化与核、像、余核及有限交相容。
\end{theorem}
\begin{proof}
若两分式相等，对它们的交叉差施加 \(f\) 即可证明诱导映射良定义；恒等映射与复合的相容性在分子上直接成立。

若 \(i(l)/s=0\)，由\eqref{b3:eq:localized-element-zero}有 \(u\cdot i(l)=i(ul)=0\)。原映射 \(i\) 单射，故 \(ul=0\)，于是 \(l/s=0\)。若 \(q(m)/s=0\)，取 \(u\in S\) 使 \(q(um)=0\)，再取 \(l\in L\) 使 \(i(l)=um\)。则
\[
\frac ms=\frac{um}{us}=S^{-1}i\left(\frac l{us}\right).
\]
反过来，\(qi=0\) 保证像包含于核。最后，对 \(n/s\) 选 \(m\) 满足 \(q(m)=n\)，便有原像 \(m/s\)，证明满射性。

把任意同态分别分解为到其像的满射及像的包含，便得到关于核、像和余核的断言。对子模 \(L_1,L_2\subseteq M\)，把交写为映射 \(L_1\oplus L_2\rightarrow M\)、\((x,y)\mapsto x-y\) 的核，再用 \(S^{-1}(L_1\oplus L_2)\cong S^{-1}L_1\oplus S^{-1}L_2\)，得到有限交的相容性。
\end{proof}

第二卷命题7.3.10用有限关系判据证明 \(S^{-1}R\) 是平坦 \(R\) 模；结合\cref{b3:prop:localization-tensor}，那也能推出正合性。这里的分式证明则明确指出：落在局部核中不意味着原分子已经在核中，而是要先乘上某个分母。这个区别将在理想的收缩和模的支集中反复使用。

% !TeX root = ../../Book3.tex
% 纲要标题：### 2.2 理想与局部化
% 纲要位置：工作记录/计划纲要.md，第 1949 行。
% 纲要细目（写作范围）：
% * 扩张、收缩和饱和；有向并与单元素饱和的几何实例
% * 局部化后的素理想对应
% * 剩余域与纤维；局部模和剩余向量空间、空纤维及非约化纤维

\section{理想与局部化}
\label{b3:sec:0202}

局部化可能使一个非零理想变成零，也可能使一个真理想变成单位理想。要准确追踪理想，必须同时使用扩张与收缩；回到原环后增加的元素，正是能被某个允许的分母送入原理想的元素。本节还把第一章的素谱对应用于剩余域及纤维。

\subsection{扩张、收缩与饱和}
\label{b3:subsec:0202:01}

设 \(R\) 为交换含幺环，\(S\subseteq R\) 为含一的乘法集，\(I\) 为 \(R\) 的理想。由\cref{b3:thm:localization-exact}，包含映射 \(I\hookrightarrow R\) 局部化后仍为单射，因而可把模的局部化 \(S^{-1}I\) 识别为 \(S^{-1}R\) 中由分式 \(a/s\)、\(a\in I\)、\(s\in S\) 组成的子模。这些分式也对乘以任意环分式封闭，而且每个 \(a/s=(a/1)(1/s)\)，所以其像正是扩张理想 \(I(S^{-1}R)\)。以下用 \(S^{-1}I\) 同时表示这个扩张理想。对 \(S^{-1}R\) 的理想 \(J\)，其收缩是 \(J^c=\{r\in R:r/1\in J\}\)。

\begin{definition}\label{b3:def:multiplicative-saturation}
设 \(R\) 为交换含幺环，\(S\subseteq R\) 为含一的乘法集，\(I\) 为 \(R\) 的理想。对 \(s\in R\)，记 \((I:s)=\{\,r\in R\mid sr\in I\,\}\)。\(I\) 关于 \(S\) 的\term[chengfaji-baohe]{饱和}[saturation with respect to a multiplicative set]为
\[
\operatorname{Sat}_S(I)\coloneqq
\{\,r\in R\mid sr\in I\text{ 对某个 }s\in S\text{ 成立}\,\}
=\bigcup_{s\in S}(I:s).
\]
若 \(\operatorname{Sat}_S(I)=I\)，则称 \(I\) 关于 \(S\) 饱和。
\end{definition}
\symidx{\(\operatorname{Sat}_S(I)\)：理想关于乘法集的饱和}

理想的任意并一般不是理想，但这里的并具有额外的相容性。对 \(s,t\in S\)，有
\[
(I:s)\subseteq(I:st),\qquad (I:t)\subseteq(I:st).
\]
因为 \(st\in S\)，任意两个成员都包含于这族理想中的同一个成员，即这族理想在包含关系下有向。因此并中的任意两个元素可以在同一理想内相加，每个元素乘以 \(R\) 的任意标量也仍在并中，故 \(\operatorname{Sat}_S(I)\) 确实是理想。

\begin{proposition}\label{b3:prop:ideal-localization-saturation}
设 \(R\) 为交换含幺环，\(S\subseteq R\) 为含一的乘法集，\(I\) 为 \(R\) 的理想，\(J\) 为 \(S^{-1}R\) 的理想。记 \(S^{-1}I\) 为扩张理想，\(J^c=\{r\in R:r/1\in J\}\) 为收缩理想。则
\[
(S^{-1}I)^c=\operatorname{Sat}_S(I),\qquad S^{-1}(J^c)=J.
\]
因此 \(S^{-1}R\) 的理想与 \(R\) 中关于 \(S\) 饱和的理想一一对应。此外，\(S^{-1}I=S^{-1}R\) 当且仅当 \(I\cap S\ne\varnothing\)。
\end{proposition}
\begin{proof}
若 \(r/1=a/s\)、\(a\in I\)，则存在 \(u\in S\) 使 \(u(sr-a)=0\)，所以 \((us)r=ua\in I\)。反过来，若 \(sr\in I\)，则 \(r/1=(sr)/s\in S^{-1}I\)。这证明第一个等式。

对 \(a/s\in J\)，有 \(a/1=(s/1)(a/s)\in J\)，因而 \(a\in J^c\)；反向包含由 \(J\) 是理想得到。两个等式还表明扩张收缩在饱和理想上互逆。最后，\(1\in(S^{-1}I)^c\) 等价于存在 \(s\in S\) 属于 \(I\)，这就是单位理想的判据。
\end{proof}

设 \(k\) 为域。在 \(R=k[x,y]\) 中，取 \(S=\{1,x,x^2,\ldots\}\)、\(I=(xy)\)，则
\[
\operatorname{Sat}_S(I)=(y),\qquad I R_x=(y)R_x.
\]
确实，\(x(y)\subseteq I\) 给出一个包含；若 \(x^n f\in(xy)\)，在整环 \(k[x,y]/(y)=k[x]\) 中约去非零的 \(x^n\)，便得 \(f\in(y)\)。同一个 \(I\) 在 \(x,y\) 都可逆的环中扩张成单位理想。

\ProseParagraphBreak
在这个例子中，\(V(xy)=V(x)\cup V(y)\)。把 \(x\) 变成单位相当于限制到 \(D(x)\)，其中不再有 \(V(x)\) 上的点，剩下的是 \(V(y)\cap D(x)\)。这个集合包含 \(V(y)\) 的一般点 \((y)\)，因为 \(x\notin(y)\)；所以它在原谱中的闭包正是 \(V(y)\)。所算出的饱和理想 \((y)\) 是根理想，恰好对应这个闭包。

\ProseParagraphBreak
若 \(S\) 由单个 \(f\in R\) 的幂组成，则 \(\operatorname{Sat}_S(I)=I:f^\infty=\bigcup_{n\ge0}(I:f^n)\)。它与关于多生成元理想 \(J\) 的 \(I:J^\infty\) 不同：后者要求 \(J^n r\subseteq I\)，即该幂理想的每个元素都满足条件；乘法集的饱和只要求存在一个分母。理想商与饱和的系统计算将在\chapref{b3:ch:08}继续讨论。

\subsection{局部化中的素理想对应}
\label{b3:subsec:0202:02}

\begin{proposition}\label{b3:prop:localized-prime-ideals}
设 \(R\) 为交换含幺环，\(S\subseteq R\) 为含一的乘法集。扩张与收缩给出素理想之间的对应
\[
\{\,\mathfrak p\in\operatorname{Spec}R\mid\mathfrak p\cap S=\varnothing\,\}
\longleftrightarrow\operatorname{Spec}(S^{-1}R),
\qquad \mathfrak p\longmapsto S^{-1}\mathfrak p.
\]
对 \(R\) 的任一素理想 \(\mathfrak p\)，在 \(R_{\mathfrak p}\) 中唯一的极大理想为 \(\mathfrak pR_{\mathfrak p}\)，其素理想对应于 \(R\) 中包含在 \(\mathfrak p\) 内的素理想。
\end{proposition}
\begin{proof}
第一章的\cref{b3:thm:spec-localization}已证明这一对应及其拓扑相容性；这里说明代数上怎样使用。若 \(sr\in\mathfrak p\)、\(s\notin\mathfrak p\)，素性给出 \(r\in\mathfrak p\)，故 \(\mathfrak p\) 饱和。商环
\[
(S^{-1}R)/(S^{-1}\mathfrak p)\cong\bar S^{-1}(R/\mathfrak p)
\]
是非零整环，因此扩张仍为素理想。相反，局部化环中所有 \(s/1\) 都是单位，不可能落入真素理想；其收缩是与 \(S\) 不交的素理想。由\cref{b3:prop:ideal-localization-saturation}两种操作互逆。

取 \(S=R\setminus\mathfrak p\)。一个素理想与 \(S\) 不交正是它包含于 \(\mathfrak p\)，故最大的一个为 \(\mathfrak p\)。也可直接看出，若 \(a/s\notin\mathfrak pR_{\mathfrak p}\)，则 \(a\notin\mathfrak p\)，其逆为 \(s/a\)；所以剩下的全部元素都是单位，极大理想唯一。
\end{proof}

\subsection{剩余域与纤维}
\label{b3:subsec:0202:03}

设 \(R,A\) 为交换含幺环。保单位环同态 \(\varphi:R\rightarrow A\) 把 \(A\) 的素理想 \(\mathfrak q\) 送到收缩 \(\varphi^{-1}(\mathfrak q)\)。固定 \(R\) 的一个素理想 \(\mathfrak p\)，哪些点恰好落在它上面？要使收缩等于 \(\mathfrak p\)，可以先把 \(\varphi(\mathfrak p)\) 中的元素模掉，再把 \(\varphi(R\setminus\mathfrak p)\) 中的元素变成单位。这两种操作合起来，就是从 \(R\) 到 \(\mathfrak p\) 处的剩余域作标量变换。

\begin{definition}\label{b3:def:residue-field-fibre}
设 \(R\) 为交换含幺环，\(\mathfrak p\) 为 \(R\) 的素理想。\(\mathfrak p\) 处的剩余域为
\[
\kappa(\mathfrak p)\coloneqq R_{\mathfrak p}/\mathfrak pR_{\mathfrak p}
\cong\operatorname{Frac}(R/\mathfrak p).
\]
给定交换含幺环 \(A\) 及保单位环同态 \(\varphi:R\rightarrow A\)，称 \(A\otimes_R\kappa(\mathfrak p)\) 为 \(\varphi\) 在 \(\mathfrak p\) 处的\term[xianwei-daishu]{纤维代数}[fibre algebra]。对任一 \(R\) 模 \(M\)，它在 \(\mathfrak p\) 处的剩余向量空间为 \(M\otimes_R\kappa(\mathfrak p)\cong M_{\mathfrak p}/\mathfrak pM_{\mathfrak p}\)。
\end{definition}
\symidx{\(\kappa(\mathfrak p)\)：素理想处的剩余域}

剩余域的同构可从分式直接看出：模掉 \(\mathfrak p\) 后，\(R\setminus\mathfrak p\) 的像恰是整环 \(R/\mathfrak p\) 的全部非零元素。两个张量描述则依次使用\cref{b3:prop:localization-tensor}及商模的张量公式。

\ProseParagraphBreak
局部模 \(M_{\mathfrak p}\) 只把 \(R\setminus\mathfrak p\) 中的标量变成单位；剩余向量空间还进一步模掉 \(\mathfrak pM_{\mathfrak p}\)，因此两者保留的信息不同。例如设 \(k\) 为域，取 \(R=k[t]_{(t)}\)、\(\mathfrak p=(t)R\)、\(M=R/(t^2)\)。\(R\setminus\mathfrak p\) 中的元素已经是单位，故
\[
M_{\mathfrak p}\cong R/(t^2),\qquad
M\otimes_R\kappa(\mathfrak p)\cong M\otimes_R k\cong k.
\]
在局部模中，\(t\) 的商类 \(\bar t\) 非零，却满足 \(t\bar t=0\)；在剩余向量空间中，\(\bar t\otimes1=\bar1\otimes(t\cdot1)=0\)。所以取纤维会丢掉这个局部模中仍然可见的非零元素。

\begin{proposition}\label{b3:prop:fibre-spectrum}
设 \(R,A\) 为交换含幺环，\(\varphi:R\to A\) 为保单位环同态，\(\mathfrak p\) 为 \(R\) 的素理想。纤维代数 \(A\otimes_R\kappa(\mathfrak p)\) 的素理想自然对应于 \(A\) 中满足 \(\varphi^{-1}(\mathfrak q)=\mathfrak p\) 的素理想 \(\mathfrak q\)。
\end{proposition}
\begin{proof}
令 \(S=R\setminus\mathfrak p\)。通过 \(\varphi\) 赋予 \(A\) 一个 \(R\) 代数结构，并记
\[
S^{-1}A\coloneqq\varphi(S)^{-1}A.
\]
这里 \(S^{-1}A\) 是把 \(\varphi(S)\) 中的元素变成单位所得的环；\(\mathfrak pS^{-1}A\) 表示 \(\varphi(\mathfrak p)\) 在该环中生成的理想。局部化及取商的张量公式给出
\[
A\otimes_R\kappa(\mathfrak p)
\cong S^{-1}A/\mathfrak pS^{-1}A.
\]
商环的素理想对应于 \(S^{-1}A\) 中包含 \(\mathfrak pS^{-1}A\) 的素理想；局部化对应又把它们识别为 \(A\) 中包含 \(\varphi(\mathfrak p)\) 且与 \(\varphi(S)\) 不交的素理想。这两个条件给出两边的包含：
\[
\begin{aligned}
\varphi(\mathfrak p)\subseteq\mathfrak q
&\ \Longrightarrow\ \mathfrak p\subseteq\varphi^{-1}(\mathfrak q),\\
\varphi(R\setminus\mathfrak p)\cap\mathfrak q=\varnothing
&\ \Longrightarrow\ \varphi^{-1}(\mathfrak q)\subseteq\mathfrak p.
\end{aligned}
\]
第二个方向是因为：若 \(r\in\varphi^{-1}(\mathfrak q)\) 却不属于 \(\mathfrak p\)，则 \(r\in S\)，于是 \(\varphi(r)\) 同时属于 \(\varphi(S)\) 和 \(\mathfrak q\)，矛盾。反过来，若 \(\varphi^{-1}(\mathfrak q)=\mathfrak p\)，则上述两个条件都成立。因此它们合起来恰好是所述等式。
\end{proof}

\begin{figure}[htbp]
\centering
\begin{tikzcd}[column sep=large,row sep=large]
\operatorname{Spec}\bigl(A\otimes_R\kappa(\mathfrak p)\bigr)
\arrow[r,"\iota"] \arrow[d]
& \operatorname{Spec}A \arrow[d,"\varphi^*"] \\
\{\mathfrak p\} \arrow[r,hook]
& \operatorname{Spec}R
\end{tikzcd}
\caption{谱映射在一个点上的纤维}
\label{b3:fig:spectrum-fibre}
\end{figure}

\cref{b3:fig:spectrum-fibre}中的 \(\iota\) 由 \(a\mapsto a\otimes1\) 诱导。上面的命题说明，它的像恰为 \(\{\,\mathfrak q\in\operatorname{Spec}A\mid\varphi^{-1}(\mathfrak q)=\mathfrak p\,\}\)，也就是 \(\varphi^*\) 在 \(\mathfrak p\) 上的集合论纤维；这解释了纤维代数的名称。

\ProseParagraphBreak
纤维也可能为空。例如 \(\mathbb Z\hookrightarrow\mathbb Q\) 在 \((2)\) 处的纤维代数为 \(\mathbb Q\otimes_{\mathbb Z}\mathbb F_2\)。它的单位元满足 \(1\otimes\bar1=(1/2)\otimes2\bar1=0\)，所以纤维代数是零环，素谱为空。

\ProseParagraphBreak
例如，设 \(k\) 为域，\(\varphi:k[t]\rightarrow k[x]\) 为满足 \(\varphi(t)=x^2\) 的 \(k\) 代数同态。对 \(a\in k\)，\(\varphi\) 在 \((t-a)\) 处的纤维代数为 \(k[x]/(x^2-a)\)。当 \(a=0\) 时，其中 \(x\) 的像非零却平方为零；素谱只有一个点，其剩余域为 \(k\)，但纤维代数仍保存这项幂零信息。当 \(a=b^2\ne0\)、\(b\in k\) 且 \(\operatorname{char}k\ne2\) 时，\(x-b\) 与 \(x+b\) 互素，中国剩余定理给出 \(k[x]/(x^2-a)\cong k\times k\)，有两个点，剩余域都是 \(k\)。若 \(\operatorname{char}k=2\) 且 \(a=b^2\ne0\)，则 \(x^2-a=(x-b)^2\)；令 \(\varepsilon\) 为 \(x-b\) 的商类，纤维代数同构于 \(k[\varepsilon]/(\varepsilon^2)\)。这仍只有一个剩余域为 \(k\) 的点，且 \(\varepsilon\ne0\) 是幂零元，即一个非约化双重点。若 \(a\) 在 \(k\) 中没有平方根，则 \(x^2-a\) 不可约，纤维代数是二次域扩张，只有一个点，其剩余域就是该二次扩张，而且没有非零幂零元。这些情形说明，点数、点的剩余域和纤维代数的约化性是三种不同的信息；同样只有一个点，既可能有幂零元，也可能是域。

% !TeX root = ../../Book3.tex
% 纲要标题：### 2.3 模的支集
% 纲要位置：工作记录/计划纲要.md，第 1955 行。
% 纲要细目（写作范围）：
% * 支集、零化子与有限生成性；循环模支集等于理想闭集
% * 局部化与 Hom 的交换条件；有限生成不能代替有限展示的循环模反例
% * 支集与短正合列

\section{模的支集}
\label{b3:sec:0203}

局部化把一个模分散到素谱的各点；支集记录它在哪些点仍然非零。有限生成性使这些局部信息受到一个理想的统一控制。处理同态时，还要进一步区分“有限个生成元”与“有限个生成关系”，因为清分母必须一次覆盖全部有关数据。

\subsection{支集、零化子与有限生成性}
\label{b3:subsec:0203:01}

\begin{definition}\label{b3:def:support-annihilator}
设 \(R\) 为交换含幺环，\(M\) 为 \(R\) 模。\(M\) 的\term[zhiji]{支集}[support]与\term[linghuazi]{零化子}[annihilator]分别为
\[
\begin{aligned}
\operatorname{Supp}_R M
&=\{\,\mathfrak p\in\operatorname{Spec}R\mid M_{\mathfrak p}\ne0\,\},\\
\operatorname{Ann}_R M
&=\{\,a\in R\mid am=0\text{ 对所有 }m\in M\,\}.
\end{aligned}
\]
零化子也称湮灭子。\termidx[yinmiezi]{湮灭子}对单个元素，简记
\(\operatorname{Ann}_R(m)=\{a\in R:am=0\}\)；它等于循环子模 \(Rm\) 的零化子。
\end{definition}
\symidx{\(\operatorname{Supp}_R M\)：模的支集}
\symidx{\(\operatorname{Ann}_R M\)：模的零化子}

\begin{theorem}\label{b3:thm:support-finite}
设 \(R\) 为交换含幺环，\(M\) 为 \(R\) 模。则
\[
\operatorname{Supp}_R M\subseteq V(\operatorname{Ann}_R M).
\]
若 \(M\) 有限生成，则等号成立。更具体地，对每个 \(R\) 的素理想 \(\mathfrak p\)，有限生成模 \(M\) 满足
\[
M_{\mathfrak p}=0
\quad\Longleftrightarrow\quad
sM=0\text{ 对某个 }s\notin\mathfrak p\text{ 成立}.
\]
\end{theorem}
\begin{proof}
若 \(a\in\operatorname{Ann}_R M\setminus\mathfrak p\)，在 \(R_{\mathfrak p}\) 中 \(a\) 可逆，又零化 \(M_{\mathfrak p}\)，故该模为零。这证明一般包含关系。

设 \(M=Rm_1+\cdots+Rm_r\) 且 \(M_{\mathfrak p}=0\)。由\eqref{b3:eq:localized-element-zero}，对每个 \(i\) 有 \(s_i\notin\mathfrak p\) 使 \(s_im_i=0\)。取 \(s=s_1\cdots s_r\)，则 \(s\notin\mathfrak p\) 且 \(sM=0\)。反向由 \(s\) 局部可逆得到。于是 \(M_{\mathfrak p}=0\) 等价于 \(\operatorname{Ann}_R M\not\subseteq\mathfrak p\)，即所述等式。对零模取空生成集及 \(s=1\)，论证仍成立。
\end{proof}

因而有限生成模的支集是闭集。特别地，对 \(R\) 的任意理想 \(I\)，循环模 \(R/I\) 有限生成，且 \(\operatorname{Ann}_R(R/I)=I\)：零化它的元素必须零化 \(1+I\)，反向则由商模的运算得到。因此
\[
\operatorname{Supp}_R(R/I)=V(I).
\]
例如，对域 \(k\)，有
\[
\operatorname{Supp}_{k[x,y]} k[x,y]/(xy)=V(xy)=V(x)\cup V(y).
\]
模掉 \(xy\) 得到的模同时分布在两条坐标轴上；若进一步在 \(x\) 处局部化，则只留下 \(y=0\) 且 \(x\ne0\) 的部分。

\ProseParagraphBreak
有限生成条件不能省去。整数模 \(\mathbb Q/\mathbb Z\) 的零化子为零：对非零整数 \(n\)，选不能整除 \(n\) 的正整数 \(d\)，则 \(n(1/d+\mathbb Z)\ne0\)。但在 \((0)\) 处局部化后，每个元素都被某个非零整数零化，故整个模变成零。另一方面，对任一素数 \(p\)，元素 \(1/p+\mathbb Z\) 在 \((p)\) 处局部化后仍不为零：若它变成零，就有某个 \(s\in\mathbb Z\setminus(p)\) 使 \(s/p\in\mathbb Z\)，与 \(p\nmid s\) 矛盾。由于 \(\mathbb Z\) 的素理想只有 \((0)\) 及各个 \((p)\)，这就完整算出
\[
\operatorname{Supp}_{\mathbb Z}(\mathbb Q/\mathbb Z)
=\{\,(p)\mid p\text{ 为素数}\,\},\qquad
V\bigl(\operatorname{Ann}_{\mathbb Z}(\mathbb Q/\mathbb Z)\bigr)
=\operatorname{Spec}\mathbb Z.
\]
非零整数不可能被所有素数整除，故 \(\bigcap_{p\text{ 为素数}}(p)=(0)\)。若闭集 \(V(I)\) 包含上述支集，则 \(I\) 包含于这个交，只能为零。因此支集在 \(\operatorname{Spec}\mathbb Z\) 中稠密；它又缺少 \((0)\)，因而不是闭集。此例中每个元素各有一个零化它的非零乘子，却没有一个非零整数同时零化全部元素。

\subsection{局部化与 Hom 的交换条件}
\label{b3:subsec:0203:02}

设 \(R\) 为交换含幺环，\(S\subseteq R\) 为含一的乘法集，\(M,N\) 为 \(R\) 模。局部化后的同态可以对分式取值，因而有自然映射
\begin{equation}\label{b3:eq:hom-localization-map}
S^{-1}\operatorname{Hom}_R(M,N)
\longrightarrow
\operatorname{Hom}_{S^{-1}R}(S^{-1}M,S^{-1}N),
\qquad
\frac fs\longmapsto
\left(\frac mt\longmapsto\frac{f(m)}{st}\right).
\end{equation}
若两侧的分式代表元改变，交叉差仍被某个分母零化；因此该公式良定义。也可先在 \(M\) 上取 \(m\mapsto f(m)/s\)，再用\cref{b3:prop:module-fraction-operations}的延拓泛性质。

\begin{theorem}\label{b3:thm:hom-localization}
设 \(R\) 为交换含幺环，\(S\subseteq R\) 为含一的乘法集，\(M,N\) 为 \(R\) 模。考虑自然映射
\[
S^{-1}\operatorname{Hom}_R(M,N)\longrightarrow
\operatorname{Hom}_{S^{-1}R}(S^{-1}M,S^{-1}N),\qquad
f/s\longmapsto\bigl(m/t\longmapsto f(m)/(st)\bigr).
\]
若 \(M\) 有限生成，则此映射为单射。若 \(M\) 有限展示，即存在非负整数 \(a,b\) 及 \(R\) 模正合列
\[
R^a\longrightarrow R^b\longrightarrow M\longrightarrow0,
\]
则该映射为同构。两种结论都不要求 \(N\) 有限生成。
\end{theorem}
\begin{proof}
先只假设 \(M\) 有限生成，取生成元 \(m_1,\ldots,m_r\)。单射性要求：一个同态若在局部化后消失，就能找到一个分母零化这个同态的全部值。若 \(f/s\) 在右侧为零，则每个 \(f(m_i)/s=0\)，故可取 \(u_i\in S\) 使 \(u_i\cdot f(m_i)=0\)。因为生成元只有有限个，乘积 \(u=\prod_{i=1}^r u_i\) 零化 \(f\) 在全部生成元上的值，因而 \(uf=0\)，即 \(f/s=0\)。这证明单射性；当 \(M=0\) 时可取空生成集及 \(u=1\)。

再设 \(M\) 有上述有限展示。要把局部化后的同态写成 \(f/s\)，仅给生成元的像通分还不够；通分后的像必须在 \(N\) 中满足全部关系，才能定义 \(f:M\to N\)。记 \(e_1,\ldots,e_b\) 在 \(M\) 中的像为 \(m_i\)。展示的第一个映射的像由有限个关系
\(\sum_i a_{ij}m_i=0\)、\(1\le j\le a\) 生成。给定右侧的同态 \(g\)，把有限个 \(g(m_i/1)\) 通分为 \(n_i/s\)。每个关系给出
\[
\frac{\sum_i a_{ij}n_i}{s}=0.
\]
故可选一个共同的 \(u\in S\)，使全部 \(\sum_i a_{ij}u n_i=0\) 在 \(N\) 中成立。于是 \(m_i\mapsto u n_i\) 满足全部关系，诱导 \(f:M\rightarrow N\)。此时 \(f/(us)\) 在\eqref{b3:eq:hom-localization-map}下的像在 \(m_i/1\) 上等于 \(g\)，因而就是 \(g\)。这证明满射性。
\end{proof}

有限生成不能代替有限展示来保证满射性。下面的循环模只有一个生成元，却无法同时清除所有关系中的分母。

\begin{proposition}\label{b3:prop:hom-localization-fg-counterexample}
设 \(k\) 为域，\(R=\prod_{n\ge1}k\)，\(I=\bigoplus_{n\ge1}k\subseteq R\) 为只有有限多个非零坐标的元素组成的理想。对有限集 \(F\subseteq\mathbb Z_{\ge1}\)，令 \(e_F\) 在 \(F\) 的坐标上等于 \(1\)，其余坐标上等于 \(0\)，并取 \(S=\{1-e_F:F\text{ 有限}\}\)。则 \(S\) 是含一而不含零的乘法集，\(M=R/I\) 是有限生成但非有限展示的 \(R\) 模，且自然映射
\[
S^{-1}\operatorname{Hom}_R(M,R)\longrightarrow
\operatorname{Hom}_{S^{-1}R}(S^{-1}M,S^{-1}R)
\]
的源为零，目标非零。还有自然同构 \(S^{-1}R\cong R/I\)，把 \(a/s\) 送到 \(a+I\)。
\end{proposition}
\begin{proof}
因为 \((1-e_F)(1-e_G)=1-e_{F\cup G}\)，\(S\) 对乘法封闭；取空集得到 \(1\in S\)，而有限集不可能包含所有正整数，所以 \(0\notin S\)。\(M\) 由 \(1+I\) 生成。一个同态 \(M\to R\) 在该生成元上的值 \(a\) 必须满足 \(Ia=0\)。用每个单坐标幂等元 \(e_{\{n\}}\in I\) 相乘得到 \(a_n=0\)，故 \(a=0\)，即
\[
\operatorname{Hom}_R(M,R)=\operatorname{Ann}_R I=0.
\]
另一方面，每个 \(x\in I\) 只有有限多个非零坐标；若这些坐标包含于 \(F\)，则 \((1-e_F)x=0\)。由分式零判据，\(S^{-1}I=0\)，局部化正合性于是给出 \(S^{-1}M\cong S^{-1}R\)。环 \(S^{-1}R\) 非零，因为 \(1/1=0\) 将要求 \(S\) 含有零元。因此目标包含非零恒等同态，而源为零，映射不满射。由\cref{b3:thm:hom-localization}，\(M\) 不可能有限展示。

对每个 \(s=1-e_F\in S\)，其在 \(R/I\) 中的像等于 \(1\)，故商映射延拓为 \(S^{-1}R\to R/I\)、\(a/s\mapsto a+I\)。其逆为 \(a+I\mapsto a/1\)：\(S^{-1}I=0\) 保证该映射良定义，且 \(e_F/1=0\) 使 \(s/1=1\)，从而 \(a/s=a/1\)。
\end{proof}

若连有限生成也不假设，取 \(R=\mathbb Z\)、\(S=\mathbb Z\setminus\{0\}\)、\(M=\mathbb Q\)、\(N=\mathbb Z\)。对任意 \(\mathbb Z\) 模同态 \(f:\mathbb Q\rightarrow\mathbb Z\)，有
\[
f(q)=n\cdot f(q/n)\qquad(q\in\mathbb Q,\ n\ge1).
\]
所以整数 \(f(q)\) 被每个正整数整除，只能为零。因而 \(\operatorname{Hom}_{\mathbb Z}(\mathbb Q,\mathbb Z)=0\)，而局部化后两模都成为 \(\mathbb Q\)，右侧为 \(\operatorname{Hom}_{\mathbb Q}(\mathbb Q,\mathbb Q)=\mathbb Q\)。因此一般的 Hom 不能无条件移过局部化。

\subsection{短正合列中的支集}
\label{b3:subsec:0203:03}

\begin{proposition}\label{b3:prop:support-exact}
设 \(R\) 为交换含幺环，\(L,M,N\) 为 \(R\) 模。若 \(0\rightarrow L\rightarrow M\rightarrow N\rightarrow0\) 正合，则
\[
\operatorname{Supp}_R M
=\operatorname{Supp}_R L\cup\operatorname{Supp}_R N.
\]
该结论对任意 \(R\) 模成立。此外，任意一族 \(R\) 模 \((M_\lambda)_\lambda\) 满足
\[
\operatorname{Supp}_R\left(\bigoplus_{\lambda}M_\lambda\right)
=\bigcup_{\lambda}\operatorname{Supp}_R M_\lambda.
\]
\end{proposition}
\begin{proof}
由\cref{b3:thm:localization-exact}，在每个 \(\mathfrak p\) 处仍有短正合列。中间项为零当且仅当两端同时为零，取补集即得第一式。

局部化与直和相容：把分式 \((m_\lambda)_\lambda/s\) 送到 \((m_\lambda/s)_\lambda\)。其像仍只有有限多个非零分量；反过来，对这些非零分量的分母取乘积，就能通分并得到逆映射。这里的有限性来自直和的定义，与素谱上模的支集是不同的概念。直和为零当且仅当各分量为零，因而得到第二式。
\end{proof}

第一式说明支集忽略了短正合列是否分裂。例如，设 \(k\) 为域，在 \(R=k[x]\) 上有
\[
0\longrightarrow(x)/(x^2)\longrightarrow R/(x^2)
\longrightarrow R/(x)\longrightarrow0
\]
的三个非零模支集都为 \(\bigl\{(x)\bigr\}\)，但中间项中的 \(x\) 作用非零而平方为零，两端的 \(x\) 作用均为零。这种厚度信息并未由支集单独记录，后面还要用长度、关联素理想及准素分解作更细的分析。

% !TeX root = ../../Book3.tex
% 纲要标题：### 2.4 局部—整体原则
% 纲要位置：工作记录/计划纲要.md，第 1862 行。
% 纲要细目（写作范围）：
% * 元素和模的消失判据
% * 单射、满射及同构的局部检测
% * 有限基本开覆盖与邻域有限性
% * 有限基本开覆盖上的相容局部元素黏合

\section{局部—整体原则}
\label{b3:sec:0204}

在一个素理想处看不见某个元素，只说明有一个不属于该素理想的乘子零化它；若在每个极大理想处都看不见它，这些乘子生成的理想便不能包含于任何极大理想，只能是单位理想。由此可以检测元素是否为零、是否属于给定子模，以及给定同态是否单射、满射或同构。这些判据不要求模有限生成；把点处的有限生成元延伸到开邻域，则需要另行考察有限性。

\subsection{元素与模的消失判据}
\label{b3:subsec:0204:01}

\begin{theorem}\label{b3:thm:local-zero-detection}
设 \(R\) 为交换含幺环，\(M\) 为 \(R\) 模，\(m\in M\)。下列条件等价：
\begin{enumerate}
\item \(m=0\)；
\item 对每个素理想 \(\mathfrak p\)，\(m/1=0\) 在 \(M_{\mathfrak p}\) 中成立；
\item 对每个极大理想 \(\mathfrak m\)，\(m/1=0\) 在 \(M_{\mathfrak m}\) 中成立。
\end{enumerate}
因此 \(M=0\) 当且仅当所有 \(M_{\mathfrak m}\) 为零，也当且仅当所有 \(M_{\mathfrak p}\) 为零。
\end{theorem}
\begin{proof}
\((1)\Rightarrow(2)\Rightarrow(3)\) 直接成立，其中第二次蕴含用到极大理想是素理想。为证明 \((3)\Rightarrow(1)\)，设 \(m\ne0\)，则其零化子 \(\operatorname{Ann}_R(m)\) 是真理想，故包含于某个极大理想 \(\mathfrak m\)。若 \(m/1=0\) 在 \(M_{\mathfrak m}\) 中成立，\eqref{b3:eq:localized-element-zero}给出 \(s\notin\mathfrak m\) 且 \(sm=0\)，与 \(\operatorname{Ann}_R(m)\subseteq\mathfrak m\) 矛盾。因此第三个条件推出第一个。对 \(M\) 的每个元素应用该结论，即得模的断言。若 \(R\) 为零环，幺性模只能为零，空集上的检验也与结论相符。
\end{proof}

\begin{corollary}\label{b3:cor:submodule-local-detection}
设 \(R\) 为交换含幺环，\(M\) 为 \(R\) 模，\(N\subseteq M\) 为子模，\(m\in M\)。则
\[
m\in N\quad\Longleftrightarrow\quad
m/1\in N_{\mathfrak m}\text{ 对每个极大理想 }\mathfrak m.
\]
特别地，两个子模 \(N,L\subseteq M\) 相等，当且仅当它们在每个极大理想处局部化后相等。
\end{corollary}
\begin{proof}
由\cref{b3:thm:localization-exact}，\((M/N)_{\mathfrak m}\cong M_{\mathfrak m}/N_{\mathfrak m}\)。所以条件恰好表示 \(m+N\) 在每个局部化中为零；应用\cref{b3:thm:local-zero-detection}即得 \(m+N=0\)。对子模的每个元素分别应用包含判据，得到相等判据。
\end{proof}

例如整环 \(R\) 在分式域 \(K\) 内满足
\[
R=\bigcap_{\mathfrak m\text{ 极大}}R_{\mathfrak m}.
\]
确实，对 \(z\in K\)，分母理想 \(I_z=\{a\in R:az\in R\}\) 若为真理想，选包含它的极大理想 \(\mathfrak m\)。若 \(z\in R_{\mathfrak m}\)，则 \(z=b/s\)、\(s\notin\mathfrak m\)，于是 \(s\in I_z\)，矛盾。故处处属于局部环便迫使 \(I_z=R\)、\(z\in R\)。

\subsection{单射、满射与同构的局部检测}
\label{b3:subsec:0204:02}

\begin{theorem}\label{b3:thm:local-morphism-detection}
设 \(R\) 为交换含幺环，\(M,N\) 为 \(R\) 模。对任意 \(R\) 模同态 \(f:M\rightarrow N\)，单射、满射及同构这三种性质，均可在所有极大理想处的 \(f_{\mathfrak m}\) 上检测，也可在所有素理想处检测。更一般地，若 \(L,M,N\) 为 \(R\) 模，则复形
\[
L\xrightarrow{u}M\xrightarrow{v}N,\qquad vu=0,
\]
在 \(M\) 处正合，当且仅当每个极大理想处的局部化复形在中间项正合。
\end{theorem}
\begin{proof}
正合性定理给出
\[
(\ker f)_{\mathfrak m}\cong\ker f_{\mathfrak m},\qquad
(\operatorname{coker}f)_{\mathfrak m}\cong
\operatorname{coker}f_{\mathfrak m}.
\]
用\cref{b3:thm:local-zero-detection}分别检测这两个模是否为零，便得到单射与满射的判据；两者同时成立等价于同构。对复形同样考虑
\(H=(\ker v)/\operatorname{im}u\)。由局部化与核、像及商相容，其局部化恰为局部复形的中间同调。因此所有局部复形正合等价于所有 \(H_{\mathfrak m}=0\)，也就等价于 \(H=0\)。改用所有素理想的证明完全相同。
\end{proof}

上述结论检测的是一个已经给定的映射。若只知道每个 \(M_{\mathfrak m}\) 与 \(N_{\mathfrak m}\) 各自同构，却没有一个共同的全局映射，就不能直接应用该定理。局部同构之间是否相容，是另一项需要处理的数据。

\ProseParagraphBreak
给定一组元素 \(m_1,\ldots,m_r\in M\)，它们生成 \(M\)，当且仅当其像生成每个 \(M_{\mathfrak m}\)：把判据应用于 \(R^r\rightarrow M\)、\(e_i\mapsto m_i\) 的满射性即可。这里生成元是全局预先选定的；不能偷换成“每个点都能另选有限个局部生成元”。

\subsection{有限基本开覆盖与邻域有限性}
\label{b3:subsec:0204:03}

\begin{proposition}\label{b3:prop:principal-cover-finiteness}
设 \(R\) 为交换含幺环，\(M\) 为 \(R\) 模，\(f_1,\ldots,f_r\in R\)（\(r\ge1\)）生成单位理想。
\begin{enumerate}
\item 若 \(M_{f_i}=0\) 对每个 \(i\) 成立，则 \(M=0\)。
\item 若每个 \(M_{f_i}\) 都是有限生成 \(R_{f_i}\) 模，则 \(M\) 有限生成。
\end{enumerate}
\end{proposition}
\begin{proof}
对 \(m\in M\)，第一项的假设给出整数 \(n_i\ge1\) 使 \(f_i^{n_i}m=0\)。理想 \((f_1^{n_1},\ldots,f_r^{n_r})\) 仍为单位理想：否则它包含于某个极大理想，该极大理想由素性包含每个 \(f_i\)，与原假设矛盾。将 \(1\) 写成这些幂的线性组合，便得 \(m=0\)。

对第二项，把每个局部模的有限生成元写成 \(m_{ij}/f_i^{a_{ij}}\)。取由全部有限个分子 \(m_{ij}\) 生成的子模 \(N\)。由于 \(f_i\) 已可逆，\(N_{f_i}=M_{f_i}\)；正合性给出 \((M/N)_{f_i}=0\)。第一项应用于 \(M/N\)，即得 \(M=N\)。
\end{proof}

第一章中 \(D(f_i)\) 覆盖素谱恰好等价于 \((f_1,\ldots,f_r)=R\)。所以该命题说明，在一组有限基本开覆盖上给出有限生成性，足以得到全局有限生成性。若每个点有这样的基本开邻域，还可由素谱拟紧性抽取有限子覆盖。

\ProseParagraphBreak
仅要求每个茎 \(M_{\mathfrak p}\) 有限生成则不够。例如
\[
M=\bigoplus_{p\text{ 为素数}}\mathbb Z/p\mathbb Z
\]
不是有限生成整数模，因为有限个元素只涉及有限个直和分量；但 \(M_{(p)}\cong\mathbb Z/p\mathbb Z\)，而 \(M_{(0)}=0\)。另一方面，对每个非零整数 \(f\)，
\[
M_f\cong\bigoplus_{p\nmid f}\mathbb Z/p\mathbb Z.
\]
右边仍有无限多个非零分量，因而不是有限生成 \(\mathbb Z_f\) 模。因此所有茎都有限生成，并不能保证存在局部模有限生成的基本开邻域。这与\cref{b3:thm:support-finite}对有限生成性的要求相呼应。

\ProseParagraphBreak
此外，“在局部环上检验”不能随意改成“张量到剩余域上检验”。取域 \(k\)。在 \(R=k[t]_{(t)}\) 上，商映射 \(R/(t^2)\rightarrow R/(t)\) 不是同构，但张量到唯一的剩余域 \(k\) 后成为 \(k\) 上的恒等映射。下一章的 Nakayama 引理能在有限生成条件下用剩余向量空间检测生成性；它并不使剩余域张量变成一般的正合操作。

\subsection{相容局部元素的黏合}
\label{b3:subsec:0204:04}

前面的检测定理从已经给定的全局元素或同态出发。若先在有限个基本开集上给出模元素，还须检查它们在交集上相等，才能反过来构造全局元素。这个构造仍只需局部化的分式零判据与单位理想条件。

\begin{lemma}\label{b3:lem:localization-sheaf-gluing}
设 \(A\) 为交换含幺环，\(M\) 为 \(A\) 模，\(r\ge1\)，\(f_1,\ldots,f_r\in A\)，且 \(D(f_1),\ldots,D(f_r)\) 覆盖 \(\operatorname{Spec}A\)。在下列映射中，\(M\rightarrow\prod_{i=1}^rM_{f_i}\) 将 \(m\) 送到 \((m/1)_i\)；定义 \(\delta((s_i))_{ij}\) 时，先通过自然局部化映射将 \(s_i\in M_{f_i}\) 与 \(s_j\in M_{f_j}\) 送入 \(M_{f_if_j}\)，再取两者之差。则
\[
0\longrightarrow M\longrightarrow\prod_{i=1}^rM_{f_i}
\xrightarrow{\delta}\prod_{1\le i,j\le r}M_{f_if_j},
\qquad \delta((s_i))_{ij}=s_i-s_j
\]
在 \(M\) 及 \(\prod_{i=1}^rM_{f_i}\) 处正合。换言之，在各交集上相等的一组局部元素来自唯一的 \(M\) 中元素。
\end{lemma}
\begin{proof}
若 \(m\) 在各 \(M_{f_i}\) 中为零，\cref{b3:prop:principal-cover-finiteness}的第一项便给出 \(m=0\)，所以第一个映射单射。它的像在各交集上的限制相等，因而包含于 \(\ker\delta\)。

现设 \((s_i)\in\ker\delta\)。把每个 \(s_i\) 写成 \(u_i/f_i^{n_i}\)，其中 \(u_i\in M\)、\(n_i\ge0\)。因为指标 \(i\) 只有有限多个，可以选整数 \(N\ge1\)，使 \(N\ge n_i\) 对每个 \(i\) 成立。令 \(t_i=f_i^{N-n_i}u_i\)，则
\[
t_i/1=f_i^Ns_i\quad\text{在 }M_{f_i}\text{ 中成立}.
\]
在 \(M_{f_j}\) 中考虑 \(t_i/1-f_i^Ns_j\)。它进一步在 \(f_i\) 处局部化后为零，因为 \((M_{f_j})_{f_i}\cong M_{f_if_j}\)，且 \(s_i,s_j\) 在这一交集上的像相等。由分式零判据，对每对 \((i,j)\) 存在整数 \(L_{ij}\ge0\)，使 \(f_i^{L_{ij}}(t_i/1-f_i^Ns_j)=0\)。指标对也只有有限多个，故可选 \(L\ge L_{ij}\) 对所有 \((i,j)\) 成立。于是
\[
f_i^Lt_i/1=f_i^{N+L}s_j\quad\text{在 }M_{f_j}\text{ 中成立}.
\]
令 \(N'=N+L\)、\(t_i'=f_i^Lt_i\)，便得到 \(t_i'/1=f_i^{N'}s_j\) 在每个 \(M_{f_j}\) 中同时成立。此时 \((f_1^{N'},\ldots,f_r^{N'})=A\)：若这个理想包含于某个极大理想 \(\mathfrak m\)，由 \(N'\ge1\) 和素性可知所有 \(f_i\in\mathfrak m\)，这与 \(D(f_i)\) 覆盖全谱矛盾。因此可以选 \(a_1,\ldots,a_r\in A\)，使 \(\sum_i a_if_i^{N'}=1\)。取
\[
m=\sum_{i=1}^r a_it_i'\in M.
\]
对每个 \(j\)，在 \(M_{f_j}\) 中有
\[
m/1=\sum_{i=1}^r a_i(t_i'/1)
=\left(\sum_{i=1}^r a_if_i^{N'}\right)s_j=s_j.
\]
所以 \((s_i)\) 属于第一个映射的像，唯一性则由该映射的单射性得到。
\end{proof}

第二十四章构造仿射概形的结构层时，将用这条引理把基本开集上的相容截面拼成一个截面。

% !TeX root = ../../Book3.tex
% 纲要标题：### 2.5 局部环与局部同态
% 纲要位置：工作记录/计划纲要.md，第 1968 行。
% 纲要细目（写作范围）：
% * 局部环 \((R,\mathfrak m,\kappa)\) 的单位判据、极大理想与剩余域；局部性不蕴含 Noether 性
% * 局部同态及诱导的剩余域映射；有限型仿射代数的点处局部环
% * \(k[x_1,\ldots,x_n]_{(x_1,\ldots,x_n)}\)、离散赋值环及 Artin 局部环的对照
% * 先区分局部化、Hensel 化与完备化的任务；后续第13、23章分别建立相应理论
% ---

\section{局部环与局部同态}
\label{b3:sec:0205}

在 \(\mathfrak p\) 处局部化后，所有不属于 \(\mathfrak p\) 的元素都成为单位，留下的非单位恰好组成一个极大理想。这种情形值得独立命名，因为接下来的生成元判据和维数理论都将在这样的环上表述。局部性只说明单位与非单位的分界，不包含任何隐含的有限性假设。

\subsection{局部环的单位、极大理想与剩余域}
\label{b3:subsec:0205:01}

\begin{definition}\label{b3:def:local-ring}
设 \(R\) 为非零交换含幺环。若 \(R\) 恰有一个极大理想 \(\mathfrak m\)，则称 \(R\) 为\term[jubu-huan]{局部环}[local ring]。写作 \((R,\mathfrak m)\)，或连同剩余域 \(\kappa=R/\mathfrak m\) 写作 \((R,\mathfrak m,\kappa)\)。
\end{definition}
\symidx{\((R,\mathfrak m,\kappa)\)：局部环、极大理想及剩余域}

\begin{proposition}\label{b3:prop:local-ring-units}
设 \(R\) 为非零交换含幺环。下列条件等价：
\begin{enumerate}
\item \(R\) 为局部环；
\item 全部非单位组成一个理想；
\item 存在真理想 \(\mathfrak m\)，使 \(R\setminus\mathfrak m=R^\times\)。
\end{enumerate}
在这些条件下，该非单位理想就是唯一的极大理想；特别地，\(a\in\mathfrak m\) 时 \(1-a\) 是单位。
\end{proposition}
\begin{proof}
一个元素 \(a\) 非单位，当且仅当主理想 \((a)\) 为真理想，也就当且仅当 \(a\) 属于某个极大理想。若极大理想唯一，全部非单位恰为它，得到第二项与第三项。反过来，若非单位集是理想，则它包含每个真理想中的元素，又不含 \(1\)，所以是唯一的极大理想。第三项也给出同样结论。最后，若 \(a,1-a\) 都属于 \(\mathfrak m\)，则 \(1\in\mathfrak m\)，矛盾。
\end{proof}

局部环不必是 Noether 环。取域 \(k\)，令
\[
R=k[x_1,x_2,\ldots]_{(x_1,x_2,\ldots)}.
\]
这是由\cref{b3:prop:localized-prime-ideals}得到的局部环，但理想链
\((x_1)\subsetneq(x_1,x_2)\subsetneq\cdots\) 仍严格上升。若 \(x_{n+1}\) 属于前一项，清分母后有
\(s x_{n+1}\in(x_1,\ldots,x_n)\)，其中 \(s\) 的常数项非零；令前 \(n\) 个变量为零，在剩余变量的多项式整环中得到一个非零多项式乘 \(x_{n+1}\) 为零，矛盾。

\subsection{局部同态与点处局部环}
\label{b3:subsec:0205:02}

\begin{definition}\label{b3:def:local-homomorphism}
设 \((R,\mathfrak m)\) 与 \((A,\mathfrak n)\) 为局部环，\(\varphi:R\to A\) 为保单位环同态。若 \(\varphi(\mathfrak m)\subseteq\mathfrak n\)，则称 \(\varphi\) 为\term[jubu-tongtai]{局部同态}[local homomorphism]。
\end{definition}

任意含幺同态都把单位送到单位，所以总有
\(\varphi^{-1}(\mathfrak n)\subseteq\mathfrak m\)。因此局部同态的条件等价于
\(\varphi^{-1}(\mathfrak n)=\mathfrak m\)，也等价于非单位不会变成单位。它诱导一个含幺的剩余域同态
\[
R/\mathfrak m\longrightarrow A/\mathfrak n,
\]
该映射为单射，因为域到非零环的含幺同态核为零。

\begin{proposition}\label{b3:prop:induced-local-map}
设 \(R,A\) 为交换含幺环，\(\varphi:R\rightarrow A\) 为保单位环同态，\(\mathfrak q\in\operatorname{Spec}A\)，并令
\(\mathfrak p=\varphi^{-1}(\mathfrak q)\)，则有自然的局部同态
\[
R_{\mathfrak p}\longrightarrow A_{\mathfrak q},
\qquad r/s\longmapsto\varphi(r)/\varphi(s).
\]
它诱导剩余域的嵌入 \(\kappa(\mathfrak p)\rightarrow\kappa(\mathfrak q)\)。
\end{proposition}
\begin{proof}
\(s\notin\mathfrak p\) 意味着 \(\varphi(s)\notin\mathfrak q\)，因而目标分母合法；局部化的泛性质给出同态。若 \(r\in\mathfrak p\)，其像落在 \(\mathfrak qA_{\mathfrak q}\)，故该同态为局部同态。再取剩余域即得所述嵌入。
\end{proof}

例如，设 \(k\) 为域，\(I\) 为 \(k[x_1,\ldots,x_n]\) 的理想，则有限型仿射代数 \(A=k[x_1,\ldots,x_n]/I\) 在素理想 \(\mathfrak p\) 处的局部环为 \(A_{\mathfrak p}\)。对一个 \(k\) 有理点 \(a=(a_1,\ldots,a_n)\)，即所有 \(f\in I\) 在 \(a\) 处取零，代入映射的核 \(\mathfrak m_a\) 为极大理想，剩余域为 \(k\)。局部环中的元素可写成 \(f/g\)，要求 \(g(a)\ne0\)。这个描述表示在该点附近允许分母不消失，而不要求 \(g\) 在所有点上都非零。

\ProseParagraphBreak
局部环之间的任意含幺同态未必局部。例如 \(\mathbb Z_{(p)}\hookrightarrow\mathbb Q\) 把非单位 \(p\) 变为单位，其极大理想不能映入 \(\mathbb Q\) 的极大理想 \(0\)。所以使用剩余域映射前必须检查局部条件。

\subsection{多项式局部环、离散赋值环与 Artin 局部环}
\label{b3:subsec:0205:03}

设 \(k\) 为域，\(n\ge1\)。在
\[
R=k[x_1,\ldots,x_n]_{(x_1,\ldots,x_n)}
\]
中，唯一极大理想由 \(x_1,\ldots,x_n\) 生成，剩余域为 \(k\)。当 \(n=1\) 时，每个非零元素都唯一写成 \(x^r u\)，其中 \(r\ge0\)、\(u\) 为单位。因为多项式分子可提取唯一的 \(x\) 次幂，而局部化的分母常数项非零。非零理想中取这样的最小 \(r\)，便知该理想等于 \((x^r)\)。

这里的指数 \(r\) 测量函数在原点的零点阶数。将它推广到分式域 \(k(x)\) 时，若非零多项式 \(f=x^a f_0\)、\(g=x^b g_0\)，且 \(f_0(0),g_0(0)\ne0\)，就令
\[
v(f/g)=a-b.
\]
这是分子与分母零点阶数之差；若另一个分式表示同一元素，交叉相乘使这两个差相等，所以它不依赖表示。乘积使阶数相加，非零和的阶数则至少为两项阶数的较小者，因为通分后最低次项可能相消。于是 \(v(x^r u)=r\)，局部环正好由阶数非负的有理函数及零组成，而其中的单位恰好阶数为零。离散赋值环将这种整数阶数及其运算规则抽象出来。

\begin{definition}\label{b3:def:discrete-valuation-example}
设 \(K\) 为域，满射 \(v:K^\times\rightarrow\mathbb Z\) 满足
\[
\begin{aligned}
v(ab)&=v(a)+v(b) &&(a,b\in K^\times),\\
v(a+b)&\ge\min\bigl\{v(a),v(b)\bigr\} &&(a,b\in K^\times,\ a+b\ne0).
\end{aligned}
\]
这样的 \(v\) 称为一个\term[lisan-fuzhi]{离散赋值}[discrete valuation]；相应子环
\(V=\{0\}\cup\bigl\{\,a\in K^\times\mid v(a)\ge0\,\bigr\}\) 称为\term[lisan-fuzhi-huan]{离散赋值环}[discrete valuation ring]，简称 DVR。
\end{definition}

由赋值条件，\(V\) 确实对加法及乘法封闭；其单位恰为赋值为零的元素。若选 \(\pi\) 满足 \(v(\pi)=1\)，每个非零元素为 \(\pi^r u\)，并且非单位组成 \((\pi)\)。因此它是非域的局部整环，所有非零理想均为 \((\pi^r)\)，证明与上一段的最小赋值论证相同。取素数 \(p\)，则 \(\mathbb Z_{(p)}\) 对应有理数的 \(p\) 指数；\(k[x]_{(x)}\) 对应有理函数在 \(x=0\) 处的零点阶数。DVR 的其他内在刻画见\chapref{b3:ch:10}。

\ProseParagraphBreak
另一类局部环来自有限维代数。仍设 \(k\) 为域，对整数 \(r\ge2\)，取
\[
B=k[\epsilon]/(\epsilon^r),\qquad
C=k[x,y]/(x,y)^2.
\]
在 \(B\) 中，一个元素为单位当且仅当常数项非零：把它写成 \(a(1+n)\)，其中 \(a\in k^\times\)、\(n^r=0\)，逆为
\(a^{-1}\bigl(1-n+\cdots+(-n)^{r-1}\bigr)\)。唯一极大理想 \((\bar\epsilon)\) 的 \(r\) 次幂为零，而 \(\bar\epsilon\ne0\)，所以 \(B\) 不是整环；若取 \(r=1\)，则商环就是域 \(k\)。在 \(C\) 中同理，唯一极大理想 \(\mathfrak m=(\bar x,\bar y)\) 的平方为零。极大理想 \(\mathfrak m\) 作为 \(k\) 向量空间以 \(\bar x,\bar y\) 为基，而整个环 \(C\) 以 \(1,\bar x,\bar y\) 为基。\(C\) 对 \(\mathfrak m\) 的作用只通过常数项，故 \(\mathfrak m\) 中任何一个元素所生成的理想至多一维；这个极大理想必须用两个元素生成。这些环作为 \(k\) 向量空间有限维，故理想降链必稳定，都是 Artin 局部环。

\ProseParagraphBreak
\cref{b3:tab:local-ring-examples}列出三种情形。极大理想只有一个，并不决定该理想是否主生成、是否幂零，也不决定环是否为整环。DVR 的极大理想由一个非零元生成，但它的任何次幂都非零；后两类环则具有非零幂零元。

\begin{table}[htbp]
\centering
\caption{局部环的三个例子}
\label{b3:tab:local-ring-examples}
\begin{tabularx}{\textwidth}{@{}>{\raggedright\arraybackslash}p{.30\textwidth}>{\raggedright\arraybackslash}p{.23\textwidth}*{3}{>{\raggedright\arraybackslash}X}@{}}
\toprule
环 & 极大理想 \(\mathfrak m\) & 整环 & \(\mathfrak m\) 主生成 & \(\mathfrak m\) 幂零 \\
\midrule
\(k[x]_{(x)}\) & \((x)\) & 是 & 是 & 否 \\
\(k[\epsilon]/(\epsilon^r)\)，\(r\ge2\) & \((\bar\epsilon)\) & 否 & 是 & 是 \\
\(k[x,y]/(x,y)^2\) & \((\bar x,\bar y)\) & 否 & 否 & 是 \\
\bottomrule
\end{tabularx}
\end{table}

\subsection{局部化、Hensel 化与完备化的比较}
\label{b3:subsec:0205:04}

设 \(k\) 为域，\(R=k[x]_{(x)}\)。在 \(x=0\) 附近作分式计算，只需允许常数项非零的多项式作分母。例如
\(\dfrac{1}{1-x}\in R\)，但一个任意给定的形式幂级数未必能写成这样的有理函数。局部化保留的是分式计算，而不是所有无穷次近似。

\ProseParagraphBreak
若希望同时保存模掉 \(x,x^2,\ldots\) 的相容近似，就要考虑序列 \((a_n)_{n\ge1}\)，其中 \(a_n\in R/(x^n)\)，且 \(a_{n+1}\) 模 \(x^n\) 的像为 \(a_n\)。逐项相加、相乘使这些序列成为一个环。对一般局部环 \((A,\mathfrak m)\)，以同样方式由 \(A/\mathfrak m^n\) 构成的环记为
\[
\widehat A=\varprojlim_{n\ge1} A/\mathfrak m^n,
\]
称为 \(A\) 的\term[wanbeihua]{完备化}[completion]。对上面的 \(R\)，每个商同构于 \(k[x]/(x^n)\)：常数项非零的多项式在这个商中已经可逆。相容的截断多项式逐项确定一个形式幂级数，反过来每个级数也给出这样一组截断，因此 \(\widehat R\cong k[\![x]\!]\)。例如，\(\dfrac{1}{1-x}\) 对应各阶截断 \(1+x+\cdots+x^{n-1}\)，也就是级数 \(\sum_{j\ge0}x^j\)。完备化与正合性、原环到完备化的映射等问题见\chapref{b3:ch:13}。

\ProseParagraphBreak
还有一种要求来自解方程。再设 \(\operatorname{char}k\ne2\)，则方程 \(T^2-(1+x)=0\) 模掉 \(x\) 后在 \(T=1\) 处有单根；但它在 \(R\) 中没有根，因为有理函数的平方在不可约多项式 \(1+x\) 处具有偶数阶，而 \(1+x\) 的阶为一。在 \(k[\![x]\!]\) 中，若要求根的常数项为 \(1\)，比较各次系数时新系数总乘以 \(2\)，便可逐次唯一确定这个根。

\ProseParagraphBreak
第23章的\cref{b3:thm:henselization-existence}将给出 \(R\) 的 Hensel 化 \(R^{\mathrm h}\)。它保证剩余域中的单根或互素因子能够提升，因而其中也有上述方程的根，且其剩余类为 \(1\)。两种扩张都能解决这个方程，但所加入的数据不同：完备化接纳所有相容的有限阶近似，Hensel 化则以\cref{b3:def:henselization}的泛性质规定解决这类局部代数提升问题所需的扩张，不要求接纳任意给定的形式幂级数。这也是在完备化之外另行构造 Hensel 化的原因；其与完备化的关系将在第23章具体说明。


\begin{chaptersummary}
局部化中的等式允许再乘一个分母后成立；这一零判据同时解释了分式模的构造、局部化的正合性和理想收缩时的饱和。从局部模 \(M_{\mathfrak p}\) 取剩余向量空间，还要模去 \(\mathfrak pM_{\mathfrak p}\)，得到纤维 \(M\otimes_R\kappa(\mathfrak p)\)。这次取商会丢失局部模中的信息，也未必保持单射。

\ProseParagraphBreak
对任意模，元素是否为零、是否属于给定子模、两个子模是否相等，以及给定同态是否单射、满射或同构，都能在所有极大理想处检测；给定复形的正合性也如此。这些判据不要求有限生成。有限生成使支集等于零化子的闭集，并使点处的一组全局生成元能在某个基本开邻域上继续生成；有限展示进一步使有限个关系可以统一清分母，从而得到 Hom 的局部化同构。仅有每个局部模有限生成，仍不足以保证邻域上的有限生成性。

\ProseParagraphBreak
有限基本开覆盖上的局部模元素若在两两交集上一致，清除有限多个分母后就能拼合成唯一的全局元素。这里先构造的是元素；从各处分别给出的抽象同构构造全局对象，还须处理同构之间的相容性。

\ProseParagraphBreak
局部环把非单位集中在唯一极大理想中，为下一章的 Nakayama 引理准备了单位判据。多项式局部环、DVR 及 Artin 局部环已经展示不同的极大理想结构；完备化与 Hensel 化各自增加的内容将在后续章节另行建立。
\end{chaptersummary}

\BookExercises

\begin{exercise}\label{b3:ex:02-cyclic-integers}
设 \(n\ge1\)，\(p\) 为素数，写 \(n=p^a d\)、\(p\nmid d\)。求
\((\mathbb Z/n\mathbb Z)_{(p)}\)，并说明 \(a=0\) 时的结果。
\end{exercise}
\begin{solution}
在 \(\mathbb Z_{(p)}\) 中 \(d\) 是单位，所以
\(n\mathbb Z_{(p)}=p^a\mathbb Z_{(p)}\)。由商与局部化相容，
\[
(\mathbb Z/n\mathbb Z)_{(p)}
\cong\mathbb Z_{(p)}/p^a\mathbb Z_{(p)}
\cong\mathbb Z/p^a\mathbb Z.
\]
末个同构把 \(b/s\) 送到 \(b s^{-1}\) 模 \(p^a\) 的类，\(p\nmid s\) 保证逆存在；核恰为 \(p^a\mathbb Z_{(p)}\)。若 \(a=0\)，\(n\) 已是单位，商为零。
\end{solution}

\begin{exercise}\label{b3:ex:02-fraction-kernel}
设 \(k\) 为域，令 \(R=k[x,y]/(xy)\)，仍用 \(x,y\) 记剩余类。求 \(R\rightarrow R_x\) 的核，并显式确定 \(R_x\)。
\end{exercise}
\begin{solution}
\(x\) 可逆后，由 \(xy=0\) 得 \(y=0\)，所以 \(R_x\cong k[x,x^{-1}]\)。每个 \(R\) 中的元素唯一写成
\(a(x)+y\cdot b(y)\)。其在局部化中的像为 \(a(x)\)，由 \(k[x]\) 到 Laurent 多项式环的嵌入可知像为零当且仅当 \(a=0\)。故核为 \((y)\)，恰好是被某个 \(x\) 幂零化的元素组成的理想。
\end{solution}

\begin{exercise}\label{b3:ex:02-infinite-product}
固定素数 \(p\)，取 \(M=\prod_{n\ge1}\mathbb Z/p^n\mathbb Z\)，把 \(p\) 变为单位。证明每个分量的局部化均为零，而 \(M[1/p]\ne0\)。这与局部化正合有无矛盾？
\end{exercise}
\begin{solution}
第 \(n\) 个分量被 \(p^n\) 零化，故局部化为零。取 \(m=(1\bmod p^n)_n\)。对任何 \(r\)，其第 \(r+1\) 个分量满足
\(p^r\not\equiv0\pmod{p^{r+1}}\)，所以 \(p^rm\ne0\)。由零判据，\(m/1\ne0\)。
因此局部化不必与无限直积交换。正合性讨论的是核、像和有限的正合条件，不包含任意无限直积相容性，故无矛盾。
\end{solution}

\begin{exercise}\label{b3:ex:02-saturation-computation}
设 \(k\) 为域。在 \(R=k[x,y]\) 中，令 \(I=(x^2,xy)\)。求 \(I:y^\infty\)、\(I:x^\infty\)，并写出在 \(R_y,R_x\) 中的扩张理想。
\end{exercise}
\begin{solution}
因 \(yx\in I\)，有 \((x)\subseteq I:y^\infty\)。若 \(y^n f\in I\subseteq(x)\)，模掉 \(x\) 后在整环 \(k[y]\) 中得到 \(y^n\bar f=0\)，故 \(\bar f=0\)，即 \(f\in(x)\)。所以 \(I:y^\infty=(x)\)，相应扩张为 \((x)R_y\)。
又 \(x^2\in I\)，故 \(1\in I:x^\infty\)，得到 \(I:x^\infty=R\)、\(IR_x=R_x\)。
\end{solution}

\begin{exercise}\label{b3:ex:02-two-saturations}
设 \(k\) 为域，在 \(k[x,y]\) 中取 \(I=(xy)\)、\(J=(x,y)\)，并令 \(S\) 为 \(x,y\) 生成的乘法集。比较 \(I:J^\infty\) 与 \(\operatorname{Sat}_S(I)\)。
\end{exercise}
\begin{solution}
若 \(fJ^n\subseteq(xy)\)，则 \(fx^n\) 和 \(fy^n\) 都被 \(xy\) 整除。第一式模掉 \(y\) 后迫使 \(f\in(y)\)，第二式模掉 \(x\) 后迫使 \(f\in(x)\)。多项式的单项式展开给出
\((x)\cap(y)=(xy)\)，故 \(f\in I\)。反向显然，因而 \(I:J^\infty=I\)。
另一方面 \(xy\in I\cap S\)，所以 \(\operatorname{Sat}_S(I)=R\)。前者要求一整个幂理想都送入 \(I\)，后者只要求找到一个乘子。
\end{solution}

\begin{exercise}\label{b3:ex:02-primes-invert-six}
列出 \(\mathbb Z[1/6]\) 的素理想，并求
\((\mathbb Z/42\mathbb Z)\otimes_{\mathbb Z}\mathbb Z[1/6]\) 的支集。
\end{exercise}
\begin{solution}
素理想对应于 \(\mathbb Z\) 中与 \(\{6^n:n\ge0\}\) 不交的素理想，故为
\(0\) 及 \((p)\mathbb Z[1/6]\)，其中 \(p\ne2,3\) 为素数。局部化后 \(42=6\cdot7\) 与 \(7\) 生成同一理想，所以该模为 \(\mathbb Z[1/6]/(7)\cong\mathbb F_7\)。其支集只有极大理想 \((7)\mathbb Z[1/6]\)。
\end{solution}

\begin{exercise}\label{b3:ex:02-residue-generic}
设 \(k\) 为域。在 \(R=k[x,y]\) 中取 \(\mathfrak p=(x)\)、\(\mathfrak m=(x,y)\)。分别求剩余域，说明为什么 \(R_{\mathfrak p}\) 中 \(y\) 可逆，而 \(R_{\mathfrak m}\) 中不可逆。
\end{exercise}
\begin{solution}
\(R/\mathfrak p=k[y]\)，故 \(\kappa(\mathfrak p)=k(y)\)；\(R/\mathfrak m=k\)，故
\(\kappa(\mathfrak m)=k\)。因 \(y\notin\mathfrak p\)，它是 \(R_{\mathfrak p}\) 的允许分母；而 \(y\in\mathfrak m\)，在 \(R_{\mathfrak m}\) 中属于唯一极大理想，不能为单位。素理想处的剩余域不必是基域，也不必由一个闭点给出。
\end{solution}

\begin{exercise}\label{b3:ex:02-hyperbola-fibres}
设 \(k\) 为域，\(\varphi:k[t]\rightarrow k[x,y]\) 为满足 \(\varphi(t)=xy\) 的 \(k\)-代数同态。求 \(\varphi\) 在素理想 \((t-a)\)、\(a\in k\)，处的纤维代数，并分别讨论 \(a=0\) 与 \(a\ne0\)。
\end{exercise}
\begin{solution}
纤维代数为 \(k[x,y]/(xy-a)\)。若 \(a=0\)，素谱是 \(V(x)\cup V(y)\)，且
\(x,y\) 的非零剩余类乘积为零。若 \(a\ne0\)，关系说明 \(x\) 可逆，逆为 \(a^{-1}y\)；映射
\[
k[x,y]/(xy-a)\longrightarrow k[x,x^{-1}],\qquad
x\longmapsto x,\quad y\longmapsto ax^{-1}
\]
的逆把 \(x^{-1}\) 送到 \(a^{-1}y\)，故为同构。此时纤维环是整环，且不包含 \(x=0\) 的点。
\end{solution}

\begin{exercise}\label{b3:ex:02-prufer-support}
设 \(p\) 为素数。求整数模 \(\mathbb Z[1/p]/\mathbb Z\) 的零化子与支集，并据此说明无限生成模的支集与零化子闭集可以不同。
\end{exercise}
\begin{solution}
没有非零整数能同时零化全部 \(p^{-n}+\mathbb Z\)，所以零化子为零。在 \((q)\)、\(q\ne p\)，处，\(p\) 已可逆，每个元素被某个 \(p\) 幂零化，故局部模为零；在 \((0)\) 处同样为零。在 \((p)\) 处，任何不被 \(p\) 整除的整数作用在每个有限循环子群上都是可逆的，因此自然局部化不消去非零元素。支集恰为 \(\bigl\{(p)\bigr\}\)，而 \(V(0)=\operatorname{Spec}\mathbb Z\)。
\end{solution}

\begin{exercise}\label{b3:ex:02-zero-neighbourhood}
设 \(R\) 为交换含幺环，\(M\) 为有限生成 \(R\) 模，\(\mathfrak p\) 为 \(R\) 的素理想。若 \(M_{\mathfrak p}=0\)，证明存在 \(f\notin\mathfrak p\) 使 \(M_f=0\)。更一般地，若 \(m_1,\ldots,m_r\in M\) 的像生成 \(M_{\mathfrak p}\)，证明存在 \(f\notin\mathfrak p\)，使 \(m_1/1,\ldots,m_r/1\) 生成 \(M_f\)。
\end{exercise}
\begin{solution}
由\cref{b3:thm:support-finite}，存在 \(f\notin\mathfrak p\) 使 \(fM=0\)，因而 \(M_f=0\)。对第二问，令 \(N\) 为这些全局元素生成的子模。有限生成商模 \(Q=M/N\) 满足 \(Q_{\mathfrak p}=0\)，故存在 \(f\notin\mathfrak p\) 使 \(Q_f=0\)。由局部化正合性，\(N_f=M_f\)，正是所需结论。
\end{solution}

\begin{exercise}\label{b3:ex:02-hom-finite-generated-insufficient}
设 \(k\) 为域，令 \(R=\prod_{n\ge1}k\)，\(I=\bigoplus_{n\ge1}k\subseteq R\)。对有限集 \(F\subseteq\mathbb Z_{\ge1}\)，令 \(e_F\) 在 \(F\) 的坐标上等于 \(1\)，其余坐标上等于 \(0\)，并取 \(S=\{1-e_F:F\;\text{有限}\}\)。再令 \(e\in R\) 在偶数坐标上等于 \(1\)，奇数坐标上等于 \(0\)，\(J=eI\)、\(M=R/I\)、\(N=R/J\)。计算自然映射
\[
S^{-1}\operatorname{Hom}_R(M,N)\longrightarrow
\operatorname{Hom}_{S^{-1}R}(S^{-1}M,S^{-1}N)
\]
的像与余核，说明本题中源与像都非零，但映射仍不满射。
\end{exercise}
\begin{solution}
记 \(A=S^{-1}R\)。由\cref{b3:prop:hom-localization-fg-counterexample}中的同构，\(A\cong R/I\)，自然映射把 \(a/s\) 送到 \(a+I\)。由于 \(J\subseteq I\)，有 \(S^{-1}J=0\)；局部化商模的正合列给出 \(S^{-1}M\cong A\)、\(S^{-1}N\cong A\)。因此目标同态模由同态在 \(1\) 上的值识别为 \(A\)。

一个同态 \(M\to N\) 由 \(1+I\) 的像 \(a+J\) 决定，存在的条件是 \(Ia\subseteq J\)。用各个单坐标幂等元相乘可知，这等价于 \(a\) 的所有奇数坐标均为零：偶数坐标产生的乘积已属于 \(J\)，奇数坐标产生的乘积只能为零。所以
\[
\operatorname{Hom}_R(M,N)\cong eR/J.
\]
局部化后，\(S^{-1}(eR/J)\cong eA\)。自然映射把在 \(1+I\) 上取值为 \(a+J\) 的同态送到在 \(1\) 上取值为 \(a/1\) 的同态；每个 \(s\in S\) 在 \(A\) 中等于 \(1\)，所以在上述识别下它就是包含 \(eA\hookrightarrow A\)。其像为 \(eA\)，余核为
\[
A/eA\cong(1-e)A.
\]
最后，\(e\) 与 \(1-e\) 都有无限多个非零坐标，均不属于 \(I\)，所以它们在 \(A\cong R/I\) 中的像均非零。因而源与像非零，余核也非零。例如，\(A\to A\)、\(a\mapsto(1-e)a\) 就不能由局部化前的同态统一清分母得到。
\end{solution}

\begin{exercise}\label{b3:ex:02-unbounded-hom-denominators}
取 \(M=\bigoplus_{n\ge1}\mathbb Z e_n\)、\(N=\mathbb Z\)，把一个素数 \(p\) 变为单位。证明局部化后由 \(e_n\mapsto p^{-n}\) 定义的同态，不来自
\(\operatorname{Hom}_{\mathbb Z}(M,\mathbb Z)[1/p]\)。
\end{exercise}
\begin{solution}
局部化后 \(M[1/p]\) 是具有基 \(e_n\) 的自由 \(\mathbb Z[1/p]\) 模，因此指定这些像确实定义同态。若它来自 \(f/p^r\)，则对每个 \(n\) 有
\(f(e_n)=p^{r-n}\) 作为有理数成立。但 \(n>r\) 时右侧不是整数，矛盾。这个例子的困难已经发生在无限个生成元上，无需再考虑生成关系。
\end{solution}

\begin{exercise}\label{b3:ex:02-intersection-calculation}
设 \(k\) 为域。在 \(k[x,y]\) 中，先求 \((x)\cap(y)\)，再在 \(y\) 处局部化，核对结果等于两个扩张理想的交。另取素数 \(p\)，比较整数环中 \(\bigcap_{n\ge1}(p^n)\) 的局部化与各个 \((p^n)\) 在 \(\mathbb Z[1/p]\) 中的扩张理想之交。
\end{exercise}
\begin{solution}
属于 \((x)\) 的多项式的每个非零单项式都含 \(x\)，属于 \((y)\) 的则都含 \(y\)。因此交为 \((xy)\)。在 \(R_y\) 中，\((y)R_y=R_y\)，而
\((xy)R_y=(x)R_y\)，故
\[
\bigl((x)\cap(y)\bigr)R_y=(x)R_y=(x)R_y\cap(y)R_y.
\]
这与有限交的局部化相容性一致。对于第二问，\(\bigcap_{n\ge1}p^n\mathbb Z=0\)，因为一个非零整数不能被任意高次的 \(p\) 幂整除；但 \(p\) 可逆后，每个 \(p^n\mathbb Z[1/p]=\mathbb Z[1/p]\)。于是
\[
\left(\bigcap_{n\ge1}p^n\mathbb Z\right)[1/p]=0,
\qquad
\bigcap_{n\ge1}p^n\mathbb Z[1/p]=\mathbb Z[1/p]\ne0.
\]
因此局部化与有限交相容，却不必与无限交相容。
\end{solution}

\begin{exercise}\label{b3:ex:02-omit-prime}
固定素数 \(q\)。在 \(\mathbb Q\) 内计算 \(\bigcap_{p\ne q}\mathbb Z_{(p)}\)，其中 \(p\) 遍历其余素数。
\end{exercise}
\begin{solution}
把一个有理数写成互素分子分母 \(a/b\)、\(b>0\)。它属于 \(\mathbb Z_{(p)}\) 当且仅当 \(p\nmid b\)。要求对所有 \(p\ne q\) 成立，就等价于 \(b\) 只有素因子 \(q\)。因此交恰为 \(\mathbb Z[1/q]\)。删去一个极大理想处的检测，确实可能失去信息。
\end{solution}

\begin{exercise}\label{b3:ex:02-two-open-membership}
设 \(R\) 为交换含幺环，\(M\) 为 \(R\) 模，\(N\subseteq M\) 为子模，\(m\in M\)，且 \(f,g\in R\) 满足 \(fR+gR=R\)。若 \(m/1\in N_f\) 且 \(m/1\in N_g\)，直接用分母证明 \(m\in N\)，不要只引用全部极大理想处的判据。
\end{exercise}
\begin{solution}
在 \(M/N\) 中，\(m+N\) 分别在 \(f,g\) 处消失，所以存在 \(r,s\ge1\) 使
\(f^rm,g^sm\in N\)。理想 \((f^r,g^s)\) 为单位理想：否则一个包含它的极大理想将同时包含 \(f,g\)。取 \(af^r+bg^s=1\)，得到
\(m=a(f^rm)+b(g^sm)\in N\)。
\end{solution}

\begin{exercise}\label{b3:ex:02-stalk-versus-neighbourhood}
令 \(M=\bigoplus_p\mathbb Z/p\mathbb Z\)。证明对每个非零整数 \(n\)，\(M[1/n]\) 仍不是有限生成 \(\mathbb Z[1/n]\) 模。由此解释为什么 \(M_{(0)}=0\) 不能在 \((0)\) 的一个基本开邻域上给出有限生成性。
\end{exercise}
\begin{solution}
若 \(p\mid n\)，第 \(p\) 个分量被局部化消去；若 \(p\nmid n\)，\(n\) 在
\(\mathbb Z/p\mathbb Z\) 上已可逆，该分量保留。故
\(M[1/n]\cong\bigoplus_{p\nmid n}\mathbb Z/p\mathbb Z\)。
仍有无限多个非零分量，而有限个元素的 \(\mathbb Z[1/n]\) 线性组合只能涉及有限个分量，所以不有限生成。包含 \((0)\) 的基本开集恰好是非零 \(n\) 给出的 \(D(n)\)，因此没有一个这样的邻域满足所需有限性。
\end{solution}

\begin{exercise}\label{b3:ex:02-fibre-misses-kernel}
设 \(k\) 为域，\(R=k[t]_{(t)}\)。计算 \(R/(t^3)\rightarrow R/(t)\) 的核，并证明张量到剩余域后该核的自然像为零，而张量后的商映射是同构。
\end{exercise}
\begin{solution}
核为 \((t)/(t^3)\)，其中 \(t+(t^3)\ne0\)。张量到 \(k=R/(t)\) 时，两个商模都变成 \(k\)，诱导映射为恒等。核的张量为
\[
\bigl((t)/(t^3)\bigr)\otimes_R k
\cong(t)/(t^2)\cong k,
\]
但它到 \(\bigl(R/(t^3)\bigr)\otimes_R k=k\) 的自然映射把 \(t\) 的类送到零。于是原序列的单射在张量后不再单射；这恰是剩余域张量不能替代局部化正合性的原因。
\end{solution}

\begin{exercise}\label{b3:ex:02-local-map-examples}
设 \(k\) 为域。证明 \(k[x]_{(x)}\rightarrow k[x,y]_{(x,y)}\) 的自然同态为局部同态，并求其剩余域映射。再说明 \(k[x]_{(x)}\hookrightarrow k(x)\) 为什么不是局部同态。
\end{exercise}
\begin{solution}
来源中允许的分母在 \(x=0\) 处非零，进入目标后在 \((0,0)\) 处仍非零，所以同态存在。极大理想 \((x)\) 映入 \((x,y)\)，故同态局部；两边剩余域都是 \(k\)，诱导映射为恒等。到分式域的嵌入把非单位 \(x\) 变为单位，或者说 \((x)\) 的像不包含于目标极大理想 \(0\)，所以不局部。
\end{solution}

\begin{exercise}\label{b3:ex:02-dvr-ideals}
设 \(k\) 为域。在 \(R=k[x]_{(x)}\) 中求理想 \(\bigl(x^3(1+x),x^5\bigr)\)，并求商环的所有理想及其作为 \(k\) 向量空间的维数。
\end{exercise}
\begin{solution}
\(1+x\) 为单位，所以原理想等于 \((x^3)\)。\(R\) 的每个非零理想为一个
\((x^r)\)，因此商环 \(R/(x^3)\) 的理想为
\(0,(\bar x^2),(\bar x),R/(x^3)\)。该商同构于 \(k[x]/(x^3)\)，基为
\(1,\bar x,\bar x^2\)，维数为 \(3\)；其中三个非零理想的维数分别为 \(1,2,3\)。
\end{solution}

\begin{exercise}\label{b3:ex:02-square-zero-local}
设 \(k\) 为域。对 \(A=k[x,y]/(x,y)^2\)，写出一般单位的逆，并证明其极大理想不能由一个元素生成。
\end{exercise}
\begin{solution}
一般元素为 \(a+b\bar x+c\bar y\)。若 \(a\ne0\)，令 \(n=b\bar x+c\bar y\)，有
\(n^2=0\)，故逆为 \(a^{-1}-a^{-2}n\)。若 \(a=0\)，该元素平方为零，不能为单位。所以极大理想为 \(\mathfrak m=k\bar x\oplus k\bar y\)。因 \(\mathfrak m^2=0\)，一个元素 \(u\in\mathfrak m\) 生成的理想只有 \(ku\)，至多一维，不能等于二维的 \(\mathfrak m\)。
\end{solution}

\begin{exercise}\label{b3:ex:02-completion-compatible}
设 \(k\) 为域。直接从相容序列定义构造
\(\varprojlim_{n\ge1} k[x]/(x^n)\cong k[\![x]\!]\)，并说明 \(1/(1-x)\) 在该同构下的像。再证明自然映射 \(k[x]_{(x)}\rightarrow k[\![x]\!]\) 单射但不满射，构造一个不在其像中的级数。
\end{exercise}
\begin{solution}
一个相容序列可唯一写成
\(a_0+a_1x+\cdots+a_{n-1}x^{n-1}\pmod{x^n}\)，因为相邻项的低次系数必须相同。把它送到 \(\sum_{j\ge0}a_jx^j\) 得到双射；每个固定次数的加法与乘法只涉及有限个系数，故双射保留环运算。
在每个商环中，\(1-x\) 的逆为 \(1+x+\cdots+x^{n-1}\)，因此相容序列对应
\(\sum_{j\ge0}x^j\)。这是形式系数的等式，不涉及实数或复数意义的收敛。

常数项非零的多项式在 \(k[\![x]\!]\) 中可逆，因此把 \(f/g\in k[x]_{(x)}\) 送到 \(f\) 与 \(g\) 的形式逆的乘积，就得到所述自然映射。若其像为零，乘以 \(g\) 得 \(f=0\)，故映射单射。取
\[
h(x)=\sum_{j\ge1}x^{j!}.
\]
若 \(h=f/g\)，其中 \(f,g\in k[x]\)、\(g(0)\ne0\)，则 \(gh=f\)。选充分大的 \(j\)，使 \(j!>\deg f\) 且 \(j!-(j-1)!>\deg g\)。在 \(h\) 中，\(x^{j!}\) 前面的 \(\deg g\) 个系数全为零，故 \(gh\) 的 \(x^{j!}\) 系数等于 \(g(0)\ne0\)，而 \(f\) 的该系数为零，矛盾。所以 \(h\) 不在像中，映射不满射。
\end{solution}

\begin{exercise}\label{b3:ex:02-root-lifting-warning}
设 \(k\) 为特征不为 \(2\) 的域。证明 \(T^2-(1+x)\) 在 \(k[\![x]\!]\) 中有唯一常数项为 \(1\) 的根，却在 \(k(x)\) 中没有根。
\end{exercise}
\begin{solution}
令根为 \(1+\sum_{n\ge1}a_nx^n\)。比较一次项得 \(2a_1=1\)；对 \(n\ge2\) 比较系数得
\[
2a_n+\sum_{i=1}^{n-1}a_i a_{n-i}=0.
\]
因 \(2\) 可逆，这些等式递归且唯一地确定全部系数，并使平方恰为 \(1+x\)。
若有有理函数根 \(f/g\)，约成互素多项式后由 \(f^2=(1+x)g^2\)，不可约多项式 \(1+x\) 在左侧的指数为偶数，在右侧为奇数，矛盾。因此完备化中可以出现原局部环及其分式域都没有的代数根。
\end{solution}

\begin{exercise}\label{b3:ex:02-two-open-gluing-calculation}
令 \(R=\mathbb Z/12\mathbb Z\)，\(f=\bar2\)、\(g=\bar3\)。在 \(R_f\) 和 \(R_g\) 中分别给出 \(s_f=1/f\)、\(s_g=1/g\)。证明这两个元素在 \(R_{fg}\) 中相容，并直接清分母构造其唯一的全局元素。具体要求找出共同指数 \(N\)、元素 \(u,v\in R\) 和系数 \(a,b\in R\)，使
\[
u/1=f^Ns_f,\qquad v/1=g^Ns_g,
\qquad af^N+bg^N=1,
\]
且 \(m=au+bv\) 的两个局部像分别为 \(s_f,s_g\)；不能仅引用黏合引理或中国剩余定理。
\end{exercise}
\begin{solution}
由 \(-f+g=1\)，\(D(f),D(g)\) 覆盖全谱。又 \((fg)^2=\bar{36}=0\)，故 \(R_{fg}=0\)，两个局部元素在交叠上相等。

取 \(N=3\)、\(u=f^2=\bar4\)、\(v=g^2=\bar9\)，便有所要求的两个分母等式。整数恒等式 \(-10\cdot8+3\cdot27=1\) 给出 \(a=\overline{-10}\)、\(b=\bar3\)。在 \(R_f\) 中，由 \(f^2g=0\) 且 \(f\) 可逆，得 \(g=0\)，所以 \(v=0\)；把系数等式除以 \(f\)，得到 \(m/1=af^2=1/f\)。在 \(R_g\) 中，由 \(f^2g^2=0\) 且 \(g\) 可逆，得 \(f^2=0\)，所以 \(u=0\)；把系数等式除以 \(g\)，同理得到 \(m/1=bg^2=1/g\)。这样构造的元素是
\[
m=\overline{-10\cdot4+3\cdot9}=\overline{11}.
\]
若另一个元素具有相同局部像，二者之差 \(h\) 满足 \(f^rh=g^sh=0\)，其中可取 \(r,s\ge1\)。理想 \((f^r,g^s)\) 为单位理想，否则包含它的极大理想将同时包含 \(f,g\)。因此 \(h=0\)，证明唯一性。
\end{solution}

% !TeX root = ../../Book3.tex
% 纲要标题：## 第3章　Noether 性、Artin 性与 Nakayama 引理
% 纲要位置：工作记录/计划纲要.md，第 1977 行。

\chapter{Noether 性、Artin 性与 Nakayama 引理}
\label{b3:ch:03}
\BookChapterNumber{b3:ch:03}

\begin{chapterabstract}
有限生成与有限展示使模能够用有限数据描述，Noether 条件进一步保证子模仍有这样的描述。本章以行列式技巧证明交换环上的 Nakayama 引理，进而用剩余向量空间计算局部模的最少生成元数。对交换 Artin 环，则先证明根的幂零性与有限长度，再得到局部直积分解和零维 Noether 环的刻画。最后比较商、局部化、有限代数与子环中的有限性行为。
\end{chapterabstract}

\begin{learninggoals}
\begin{itemize}
\item 区分有限生成、有限展示、Noether 性、Artin 性与有限长度。
\item 能在具体证明中指出行列式技巧和 Nakayama 引理的假设。
\item 用 \(M/\mathfrak mM\) 选择极小生成组，并避免把生成组误认为自由基。
\item 分解交换 Artin 环，并用反例说明零维刻画为何需要 Noether 条件。
\end{itemize}
\end{learninggoals}

固定素数 \(p\)，设 \(R=\mathbb Z_{(p)}\)，把有理数域 \(\mathbb Q\) 当作 \(R\) 模。它非零，却满足 \(p\mathbb Q=\mathbb Q\)，所以 \(\mathbb Q/p\mathbb Q=0\)。因此“模掉极大理想后为零”不能单独推出原模为零；有限生成是 Nakayama 引理不可少的条件。反过来，对一个有限生成局部模，剩余域上的向量空间又准确提示最少需要多少生成元。

\ProseParagraphBreak

有限性还会改变理想与子模的行为。若一切理想都有限生成，任意有限生成模的子模也能继续由有限数据控制；若升链必定停止，反复取理想或子模就不会无限产生新的条件。对 Artin 环，还要把这种停止性和局部商的结构合起来，才能说明它如何分解成有限个局部部分。

\ProseParagraphBreak

预备知识为前两章的理想、素谱、局部化与局部环，以及第二卷第12.2节的有限长度与 Jordan–Hölder 定理\footnote{若尔当（Marie Ennemond Camille Jordan，1838-1922），法国数学家，研究群论与 Galois 理论。赫尔德（Ludwig Otto Hölder，1859-1937），德国数学家，研究群的合成列与实数的公理化。}。Hilbert 基定理的主要证明已在第一卷定理12.4.10给出，本章据此发展模论后果。

\ProseParagraphBreak
Matsumura 的《Commutative Ring Theory》\cite[Theorems 2.1--2.3, pp. 7--9; Theorem 3.2, p. 16]{Matsumura1989CommutativeRingTheory}可用于对照行列式技巧、局部生成元和交换 Artin 环的证明。

% !TeX root = ../../Book3.tex
% 纲要标题：### 3.1 有限性条件
% 纲要位置：工作记录/计划纲要.md，第 1979 行。
% 纲要细目（写作范围）：
% * 有限生成与有限展示；基本关系生成全部关系及有限展示对生成组的独立性
% * Noether 环上的有限生成模；升链条件与子模有限生成的等价性
% * Hilbert 基定理的模论后果

\section{有限性条件}
\label{b3:sec:0301}

有限个生成元给出模中元素的表达方式，却不保证生成元之间的关系能由有限组关系生成。再要求每个子模都有限生成，就得到 Noether 条件。这些条件控制不同的有限性，不能混用；Hilbert 基定理则把底环的 Noether 性传给有限变量多项式环，使多项式环上的有限生成模及其子模也具有有限描述。

\subsection{有限生成与有限展示}
\label{b3:subsec:0301:01}

\begin{definition}\label{b3:def:finite-presentation}
设 \(R\) 为交换含幺环，\(M\) 为 \(R\) 模。若存在非负整数 \(b\) 及元素 \(m_1,\ldots,m_b\in M\)，使每个 \(m\in M\) 都能写成
\(\sum_{i=1}^b a_i m_i\)，其中 \(a_i\in R\)，则称 \(M\) 为\term[youxian-shengcheng-mo]{有限生成模}[finitely generated module]。若存在非负整数 \(a,b\) 及 \(R\) 模正合列
\[
R^a\xrightarrow{A}R^b\longrightarrow M\longrightarrow0,
\]
则称 \(M\) 为\term[youxian-zhanshi-mo]{有限展示模}[finitely presented module]。该正合列给出一个有限展示：满射把 \(R^b\) 的标准基送到 \(M\) 的生成元，其核由矩阵 \(A\) 的列生成，称为这组生成元的关系模。
\end{definition}
有限展示模也称有限呈示模\termidx[youxian-chengshi-mo]{有限呈示模}；本书以“展示”为主称。“有限生成”限制的是生成元数目，而非底集合的大小：例如 \(\mathbb Z\) 作为整数模由 \(1\) 生成，任意非零有限维实向量空间也有限生成，二者的底集合却都无限。

\ProseParagraphBreak
若 \(m_1,\ldots,m_b\) 生成 \(M\)，自由模 \(R^b\) 的标准基 \(e_i\) 可以分别送到 \(m_i\)，得到满射 \(p:R^b\to M\)。向量 \(\sum_i r_i e_i\) 属于 \(\ker p\)，恰好表示关系 \(\sum_i r_i m_i=0\)，所以这个核记录全部关系。例如，对 \(f,g\in R\)，有有限展示
\[
R^2\xrightarrow{(f\ g)}R\longrightarrow R/(f,g)\longrightarrow0.
\]
其第一个映射把 \((a,b)\) 送到 \(af+bg\)，第二个映射为自然商映射。模 \(R/(f,g)\) 由 \(e=1+(f,g)\) 生成，\(fe=0\)、\(ge=0\) 是两条基本关系；任一关系 \(re=0\) 都由 \(r\in(f,g)\) 得到，因而是这两条基本关系的 \(R\) 线性组合。这里的有限性限制基本关系的数目，不要求全部关系只有有限条。

\ProseParagraphBreak
换一组有限生成元后，需要考虑相应的自由模满射及其核。下面的命题说明，一旦某组有限生成元的关系模有限生成，其他有限生成组也有这一性质；因此有限展示性不依赖生成组的选择。

\begin{proposition}\label{b3:prop:finite-presentation-kernel}
设 \(R\) 为交换含幺环，\(M\) 为 \(R\) 模。\(M\) 有限展示，当且仅当它有限生成且存在满射 \(R^b\rightarrow M\)，其中 \(b\ge0\)，其核有限生成。此时，对每个非负整数 \(n\)，满射 \(R^n\rightarrow M\) 的核也都有限生成。
\end{proposition}
\begin{proof}
第一句只是把关系矩阵的像写成核；反过来，对核选有限个生成元即可构成有限展示。证明最后一句时，设 \(p:R^b\rightarrow M\) 的核 \(L\) 有限生成，另有满射 \(q:R^n\rightarrow M\)。对 \(R^b\) 的每个标准基向量 \(e_i\)，由 \(q\) 满射选取 \(y_i\in R^n\)，使 \(q(y_i)=p(e_i)\)；规定 \(u(e_i)=y_i\)，便得到 \(u:R^b\rightarrow R^n\) 且 \(qu=p\)。反过来，对 \(R^n\) 的标准基向量 \(f_j\)，选取 \(z_j\in R^b\)，使 \(p(z_j)=q(f_j)\)；规定 \(v(f_j)=z_j\)，得到 \(v:R^n\rightarrow R^b\) 且 \(pv=q\)。这两个映射分别把一组生成元用另一组来表达。若 \(K=\ker q\)，则 \(u(L)\subseteq K\)，且
\(q(1-uv)=q-pv=0\)。对 \(x\in K\)，有 \(vx\in L\)，所以
\[
x=u(vx)+(1-uv)x.
\]
因此 \(K=u(L)+(1-uv)(R^n)\)，由 \(L\) 的有限生成元的像和有限个 \((1-uv)f_j\) 生成。
\end{proof}

循环模 \(R/I\) 总是有限生成；由这个命题，它有限展示当且仅当理想 \(I\) 有限生成。取
域 \(k\)，令 \(R=k[x_1,x_2,\ldots]\)、\(I=(x_1,x_2,\ldots)\)，便得到有限生成而非有限展示的模。若 \(I\) 由有限个多项式生成，把其中出现的变量集中在前 \(n\) 个内，再将这 \(n\) 个变量取零；所有生成元的像为零，\(x_{n+1}\) 的像却非零，矛盾。

\subsection{Noether 环上的有限生成模}
\label{b3:subsec:0301:02}

\begin{definition}\label{b3:def:noether-artin-module}
设 \(R\) 为交换含幺环，\(M\) 为 \(R\) 模。若 \(M\) 的每条子模升链 \(N_1\subseteq N_2\subseteq\cdots\) 最终稳定，则称 \(M\) 为 \term[Noether-mo]{Noether 模}[Noetherian module]。若每条子模降链最终稳定，则称 \(M\) 为 \term[Artin-mo]{Artin 模}[Artinian module]。若 \(R\) 作为自身的模满足相应条件，则分别称 \(R\) 为 \term[Noether-huan]{Noether 环}[Noetherian ring]与 \term[Artin-huan]{Artin 环}[Artinian ring]。
\end{definition}

一个模是 Noether 模，当且仅当它的每个子模有限生成。若某个子模不有限生成，逐次添入不能由已有元素生成的元素，就得到严格升链。反过来，若每个子模有限生成，对一条升链 \(N_1\subseteq N_2\subseteq\cdots\)，其并 \(N\) 是子模。若 \(N=0\)，链各项已全为零；否则选有限组生成元 \(m_1,\ldots,m_r\)，分别取 \(i_j\) 使 \(m_j\in N_{i_j}\)。令 \(n=\max\{i_1,\ldots,i_r\}\)，链的嵌套关系保证全部生成元都在 \(N_n\) 中，所以 \(N_n=N\)，以后各项也只能等于 \(N\)。

\begin{proposition}\label{b3:prop:noether-finite-modules}
设 \(R\) 为交换含幺环，\(0\rightarrow L\rightarrow M\rightarrow N\rightarrow0\) 为 \(R\) 模的短正合列。中间项 \(M\) 为 Noether 模，当且仅当两端 \(L,N\) 均为 Noether 模；Artin 模也满足同样结论。特别地，若 \(R\) 为 Noether 环，则每个有限生成 \(R\) 模及其子模都为 Noether 模，并且都有限展示。
\end{proposition}
\begin{proof}
子模的链也是原模的链，商模的链可以取逆像，所以中间项的链条件传给两端。反向设两端具有升链条件。对 \(M\) 的升链 \((M_i)\)，交链 \(M_i\cap L\) 及其在 \(N\) 中的像链都稳定。若 \(M_i\subseteq M_j\)，且交与像相同，对 \(x\in M_j\) 可取 \(y\in M_i\) 有相同的像，则 \(x-y\in M_j\cap L=M_i\cap L\)，所以 \(x\in M_i\)。因此原链稳定。降链在共同的稳定起点后，对反向包含使用同一个论证，即得 Artin 情形。

若 \(R\) 为 Noether 环，对 \(n\ge1\)，在向量末尾补一个零给出 \(R^{n-1}\to R^n\) 的单射，取末坐标则给出 \(R^n\to R\) 的满射，其核恰为前一个映射的像。因此有短正合列 \(0\rightarrow R^{n-1}\rightarrow R^n\rightarrow R\rightarrow0\)，从零模 \(R^0\) 开始归纳得到有限自由模为 Noether 模；其商即任意有限生成模仍为 Noether 模，所有子模也如此。对这些模选有限自由满射，核作为有限自由 Noether 模的子模有限生成，再用\cref{b3:prop:finite-presentation-kernel}得到有限展示。
\end{proof}

这个命题把第二卷第12.1节的链条件论证落实到交换环上。它还解释了\cref{b3:thm:hom-localization}在 Noether 环上为何只需假设源模有限生成：有限展示此时自动成立。若底环不满足 Noether 条件，就不能省略那项假设。

\subsection{Hilbert 基定理的模论后果}
\label{b3:subsec:0301:03}

第一卷定理12.4.10，即 Hilbert 基定理，已经证明：若 \(R\) 为 Noether 环，则有限变量多项式环 \(R[x_1,\ldots,x_n]\) 仍为 Noether 环。这里不重复其主证明，而将它与\cref{b3:prop:noether-finite-modules}合用。

\begin{corollary}\label{b3:cor:hilbert-module-consequences}
设 \(R\) 为交换含幺 Noether 环，\(A=R[a_1,\ldots,a_n]\) 为有限型交换 \(R\) 代数，其中 \(n\ge0\)。则 \(A\) 为 Noether 环。每个有限生成 \(A\) 模有有限展示，每个子模有限生成；有限矩阵给出的核、像及余核也都是有限生成 \(A\) 模。
\end{corollary}
\begin{proof}
代入同态 \(R[x_1,\ldots,x_n]\rightarrow A\)、\(x_i\mapsto a_i\) 为满射。由 Hilbert 基定理，来源是 Noether 环；商环的理想链取逆像，故商环 \(A\) 仍为 Noether 环。现在对 \(A\) 使用\cref{b3:prop:noether-finite-modules}，得到关于有限生成模和子模的断言。矩阵的核是有限自由模的子模，像与余核分别是有限自由模的同态像及商，也满足这些结论。
\end{proof}

例如，取域 \(k\)，在 \(A=k[x,y]/(y^2-x^3)\) 上，理想、理想之间的关系以及有限矩阵的解模，都能用有限数据表示。这个结论保证有限描述存在，却没有给出具体的生成元算法；当系数域上可以作精确计算时，第一卷附录F的 Gröbner 基方法和本卷\chapref{b3:ch:08}将进一步处理算法问题。

\ProseParagraphBreak
必须分清有限型代数与有限生成模。对域 \(k\)，多项式环 \(k[x]\) 由一个元素作为 \(k\) 代数生成，却有无限多个线性无关的单项式，因而不是有限生成的 \(k\) 模。上面的推论说有限生成的 \(A\) 模是 Noether 模，没有说它作为 \(R\) 模也有限生成。

\ProseParagraphBreak
模的有限长度控制另一种有限性：它要求存在有限合成列
\(0=M_0\subsetneq M_1\subsetneq\cdots\subsetneq M_\ell=M\)，其中每个相邻商 \(M_i/M_{i-1}\) 都是简单模；零模对应 \(\ell=0\) 的空合成列。第二卷定理12.2.4证明，有限长度等价于同时满足 Noether 与 Artin 条件。\cref{b3:tab:module-finiteness}将这些条件放在一起比较。

\begin{table}[htbp]
\centering
\caption[模的有限性条件]{模的有限性条件。设 \(R\) 为交换含幺环，\(M\) 为 \(R\) 模；除最后一行外，不附加底环的链条件。}
\label{b3:tab:module-finiteness}
\begin{tabularx}{\textwidth}{@{}l>{\raggedright\arraybackslash}X>{\raggedright\arraybackslash}p{0.29\textwidth}@{}}
\toprule
条件 & 所要求的有限性 & 与其他条件的关系 \\
\midrule
有限生成 & 有有限生成集 & 不一定有限展示或为 Noether 模 \\
有限展示 & 有有限生成集，全部关系均为有限组基本关系的 \(R\) 线性组合 & 推出有限生成；不一定为 Noether 模 \\
Noether & 每条子模升链最终稳定；等价地，每个子模都有限生成 & 推出有限生成；不一定为 Artin 模 \\
Artin & 每条子模降链最终稳定 & 不一定有限生成 \\
有限长度 & 有有限合成列 & 等价于 Noether 且 Artin；因而有限生成 \\
\midrule
\(R\) 为 Noether 环 & 对任意 \(R\) 模 \(M\)，有限生成、有限展示与 Noether 条件等价 & 不要求 \(M\) 为 Artin 模 \\
\bottomrule
\end{tabularx}
\end{table}

有限长度在一般底环上也不保证有限展示。上面的无限变量例子中，\(R/I\cong k\) 是简单 \(R\) 模，长度为 \(1\)，但理想 \(I\) 不有限生成，故 \(R/I\) 没有有限展示。

% !TeX root = ../../Book3.tex
% 纲要标题：### 3.2 Nakayama 引理
% 纲要位置：工作记录/计划纲要.md，第 1985 行。
% 纲要细目（写作范围）：
% * 行列式技巧；生成元组上的矩阵关系及标量作用的像
% * Jacobson 根版本与局部环版本
% * 极小生成元及剩余向量空间；区别极小、最少与自由基

\section{Nakayama 引理}
\label{b3:sec:0302}

在局部环上，将有限生成模按极大理想取商，就得到剩余域上的有限维向量空间。Nakayama 引理说明，向量空间中的生成元可以提升为原模的生成元。这里的有限性十分关键；证明它的行列式技巧，还将在整依赖中给出首一方程。

\subsection{行列式技巧}
\label{b3:subsec:0302:01}

考虑有限生成 \(R\) 模 \(M\) 满足 \(M=IM\) 的情形，其中 \(R\) 为交换含幺环、\(I\) 为理想。选取生成元 \(m_1,\ldots,m_n\)，每个生成元都能写成 \(m_i=\sum_j a_{ij}m_j\)、\(a_{ij}\in I\)。把这些关系合在一起，就有
\[
(I_n-A)\begin{pmatrix}m_1\\ \vdots\\m_n\end{pmatrix}=0,
\qquad A=(a_{ij}).
\]
于是可以利用伴随矩阵恒等式，寻找同时零化全部生成元的一个环元素，而不再分别处理这些关系。对于满足 \(u(M)\subseteq IM\) 的一般自同态，同一个方法也能产生零化多项式。

\begin{theorem}[行列式技巧]\label{b3:thm:determinant-trick}
设 \(R\) 为交换含幺环，\(M\) 为由 \(n\ge1\) 个元素生成的 \(R\) 模，\(I\subseteq R\) 为理想，
\(u:M\rightarrow M\) 为 \(R\) 线性自同态，并且 \(u(M)\subseteq IM\)。则存在
\[
P(T)=T^n+c_1T^{n-1}+\cdots+c_n,\qquad c_i\in I^i,
\]
使 \(P(u)=0\) 作为 \(M\) 的自同态成立。
\end{theorem}
\begin{proof}
选生成元 \(m_1,\ldots,m_n\)，写
\(u(m_i)=\sum_j a_{ij}m_j\)、\(a_{ij}\in I\)。令 \(A=(a_{ij})\in\operatorname{Mat}_n(R)\)。在自同态环 \(\operatorname{End}_R(M)\) 中，\(R\) 的标量作用像与 \(u\) 生成一个交换子环，简记为 \(R[u]\)。这里使用的是映射 \(R\to\operatorname{End}_R(M)\) 的像，不要求它单射，也不要求整个自同态环交换。将每个 \(a_{ij}\) 理解为对应的标量自同态，上述等式写成
\[
(uI_n-A)\begin{pmatrix}m_1\\ \vdots\\m_n\end{pmatrix}=0.
\]
这是 \(M^n\) 中的等式，列中的 \(m_i\) 是 \(M\) 的元素，并不是相对于一组自由基的坐标。矩阵的每个自同态作用在相应分量上，因此可以继续左乘矩阵。对 \(B=uI_n-A\)，在交换环 \(R[u]\) 上使用 \(\operatorname{adj}(B)B=\det(B)I_n\)，得到
\[
\det(uI_n-A)m_i=0\qquad(1\le i\le n).
\]
令 \(P(T)=\det(TI_n-A)\in R[T]\)。行列式的展开中，\(T^{n-i}\) 的系数每项都含 \(i\) 个 \(A\) 的元素，故 \(c_i\in I^i\)。由于 \(P(u)\) 零化全部生成元，它零化整个 \(M\)。
\end{proof}

这称为\term[hanglieshi-jiqiao]{行列式技巧}[determinant trick]。矩阵 \(A\) 未必唯一，所产生的首一多项式也不必唯一，但它确实零化给定自同态。若 \(M\) 自由且所取生成元是一组基，便得到熟悉的 Cayley–Hamilton 恒等式\footnote{凯莱（Arthur Cayley，1821-1895），英国数学家，研究矩阵、不变量和代数几何。哈密顿（William Rowan Hamilton，1805-1865），爱尔兰数学家及物理学家，创立四元数理论。}。上述一般版本可对照 Matsumura 的《Commutative Ring Theory》\cite[Theorem 2.1, pp. 7--8]{Matsumura1989CommutativeRingTheory}。

\ProseParagraphBreak
若 \(M=IM\)，对 \(u=\operatorname{id}_M\) 使用该定理，得到
\[
(1+c_1+\cdots+c_n)M=0.
\]
由于 \(c_i\in I^i\subseteq I\) 对 \(i\ge1\) 成立，令 \(a\coloneqq-(c_1+\cdots+c_n)\)，便有 \(a\in I\) 且 \((1-a)M=0\)。有限生成将逐个生成元的关系合成了一个统一的零化元。

\subsection{Jacobson 根版本与局部环版本}
\label{b3:subsec:0302:02}

\begin{definition}\label{b3:def:jacobson-radical}
设 \(R\) 为交换含幺环。\(R\) 的\term[Jacobson-gen]{Jacobson 根}[Jacobson radical]为
\[
J(R)=\bigcap_{\mathfrak m\;\text{极大}}\mathfrak m.
\]
零环采用空交为整个环的约定。
\end{definition}
\symidx{\(J(R)\)：交换环的 Jacobson 根}

若 \(a\in J(R)\)，则对每个 \(r\in R\) 都有 \(ra\in J(R)\)，所以 \(1-ra\) 不属于任何极大理想，因而是单位。反过来，若 \(a\notin J(R)\)，由交的定义可取极大理想 \(\mathfrak m\) 使 \(a\notin\mathfrak m\)。于是 \(\bar a\) 在域 \(R/\mathfrak m\) 中非零，取其逆并提升为某个 \(r\in R\)，便有 \(1-ra\in\mathfrak m\)，所以 \(1-ra\) 不是单位。因此 Jacobson 根的单位判据为
\[
a\in J(R)\quad\Longleftrightarrow\quad
1-ra\in R^\times\;\text{对每个}\;r\in R\;\text{成立}.
\]
在局部环上，\(J(R)\) 就是唯一的极大理想。

\begin{theorem}[Nakayama 引理]\label{b3:thm:nakayama}
设 \(R\) 为交换含幺环，\(M\) 为有限生成 \(R\) 模，\(I\) 为 \(R\) 中满足 \(I\subseteq J(R)\) 的理想，\(N\subseteq M\) 为子模。
\begin{enumerate}
\item 若 \(IM=M\)，则 \(M=0\)。
\item 若 \(N\subseteq M\) 且 \(M=N+IM\)，则 \(M=N\)。
\end{enumerate}
第二项实际上只需 \(M/N\) 有限生成，而不必要求 \(M\) 本身有限生成。特别地，在局部环 \((R,\mathfrak m,\kappa)\) 上，有限生成模满足
\(M/\mathfrak mM=0\) 当且仅当 \(M=0\)。
\end{theorem}
\begin{proof}
第一项中，若 \(M\ne0\)，选有限个生成元。行列式技巧给出 \(a\in I\) 使 \((1-a)M=0\)。由 \(I\subseteq J(R)\)，\(1-a\) 为单位，乘其逆得到 \(M=0\)，矛盾；零模情形本来成立。第二项对有限生成商模 \(Q=M/N\) 使用第一项，因为原等式恰好表示 \(IQ=Q\)。局部版本取 \(I=\mathfrak m\)。
\end{proof}

\term[Nakayama-yinli]{Nakayama 引理}[Nakayama's lemma]的这一交换环证明，也可见\cite[Theorem 2.2 and Corollary, p. 8]{Matsumura1989CommutativeRingTheory}。第二卷定理14.2.1则用 Jacobson 根的单位判据直接处理非交换模，没有把交换行列式用于非交换系数。

\ProseParagraphBreak
第二项中的 \(M=N+IM\) 恰好表示自然映射 \(N\to M/IM\) 满射。因此，当 \(M\) 有限生成且 \(I\subseteq J(R)\) 时，在 \(M/IM\) 中选出一组生成元，再任取它们到 \(M\) 的提升，这些提升便生成 \(M\)。在局部环上取 \(I=\mathfrak m\)，只需在剩余向量空间中选择生成元，就能得到原模的生成元。

\ProseParagraphBreak
两项假设都有明确作用。取素数 \(p\)，对局部环 \(R=\mathbb Z_{(p)}\) 的模 \(M=\mathbb Q\)，有 \(pM=M\ne0\)，因为 \(M\) 不有限生成：有限个有理数只能给出有统一上界的负 \(p\) 指数。若省去 \(I\subseteq J(R)\)，则 \(R=\mathbb Z\)、\(I=(2)\)、\(M=\mathbb Z/3\mathbb Z\) 满足 \(IM=M\ne0\)。有限生成与 Jacobson 根条件不能相互替代。

\begin{corollary}\label{b3:cor:finite-surjective-endomorphism}
设 \(R\) 为交换含幺环，\(M\) 为有限生成 \(R\) 模。则 \(M\) 的每个满 \(R\) 线性自同态都是同构。
\end{corollary}
\begin{proof}
设 \(u:M\rightarrow M\) 满。满射性是 \(u(M)=M\)；把 \(u\) 编码成一个标量的作用，就能将它改写成前面使用的 \(IM=M\)。为此赋予 \(M\) 一个 \(R[T]\) 模结构，令 \(T\) 作用为 \(u\)；原来的有限个 \(R\) 生成元也生成此 \(R[T]\) 模，并且 \((T)M=u(M)=M\)。对交换环 \(R[T]\) 及其理想 \((T)\) 使用上面的统一零化元结论，得到
\(\bigl(1-T\cdot g(T)\bigr)M=0\)，其中 \(g(T)\in R[T]\)。所以
\(u\,g(u)=g(u)\,u=\operatorname{id}_M\)，给出了双侧逆。
\end{proof}

\subsection{极小生成元与剩余向量空间}
\label{b3:subsec:0302:03}

一个生成组不可再删去元素，只说明其中每个元素对这组生成元都不可缺少；元素个数最少则需要与所有生成组比较。生成组是否为自由基，还要检查是否存在系数不全为零的线性关系。在局部环上的有限生成模中，前两个条件由剩余向量空间联系起来。

\begin{corollary}\label{b3:cor:local-generators}
设 \((R,\mathfrak m,\kappa)\) 为交换局部环，其中 \(\kappa=R/\mathfrak m\)，\(M\) 为有限生成 \(R\) 模，\(m_1,\ldots,m_r\in M\)。这些元素生成 \(M\)，当且仅当其剩余类生成
\(\kappa\) 向量空间 \(M/\mathfrak mM\)。它们构成不可再删去元素的极小生成组，当且仅当剩余类构成该向量空间的一组基。因此所有极小生成组的元素个数都等于
\[
\mu_R(M)\coloneqq\dim_\kappa(M/\mathfrak mM).
\]
\end{corollary}
\begin{proof}
正向由取商得到。反向令 \(N=\sum_i Rm_i\)，剩余类生成意味着 \(M=N+\mathfrak mM\)，由\cref{b3:thm:nakayama}有 \(M=N\)。若剩余类为基，删去任一项便不再生成剩余向量空间，故不再生成 \(M\)。若某个极小生成组的剩余类线性相关，可从中删去至少一项而保持对剩余向量空间的生成性；刚证明的反向蕴含说明原模仍被其余元素生成，与极小性矛盾。
\end{proof}
\symidx{\(\mu_R(M)\)：局部环上有限生成模的最少生成元个数}

任意生成组的剩余类都必须生成 \(M/\mathfrak mM\)，所以其元素个数至少是 \(\mu_R(M)\)。而提升一组剩余向量空间的基就能达到这个数目。因此在局部环上，不可删去元素的极小生成组也具有最小可能的元素个数；非局部环则可能有不可删去、但个数并非最少的生成组。

\ProseParagraphBreak
例如，设 \(k\) 为域。在 \(R=k[x,y]_{(x,y)}\) 中，\(\mathfrak m=(x,y)R\) 满足
\(\mathfrak m/\mathfrak m^2\cong k^2\)，基为 \(x,y\) 的类。确实一次项给出这两个独立坐标，极大理想平方恰好没有一次项；局部化的分母常数项非零，只会把一次项乘以一个非零标量。因此 \(\mathfrak m\) 需要两个生成元，不是主理想。同样，
\((x,y)^r/(x,y)^{r+1}\) 的基由全次数恰为 \(r\) 的 \(r+1\) 个单项式给出，故 \(\mu_R(\mathfrak m^r)=r+1\)。

\ProseParagraphBreak
极小生成组不一定是自由基。例如 \(R=k[\epsilon]/(\epsilon^2)\) 上的模
\(M=R/(\epsilon)\) 只需一个生成元，但该生成元被非零的 \(\epsilon\) 零化。剩余向量空间的维数测量生成元数，并没有排除原模的关系。上述极小生成组的结论也可对照\cite[Theorem 2.3, pp. 8--9]{Matsumura1989CommutativeRingTheory}。

% !TeX root = ../../Book3.tex
% 纲要标题：### 3.3 Artin 环
% 纲要位置：工作记录/计划纲要.md，第 1991 行。
% 纲要细目（写作范围）：
% * 降链条件与有限长度；通过简单商模证明根幂零及商层的域坐标分解
% * 交换 Artin 环的局部直积分解；满足根幂为零的指数及由极大理想恢复因子
% * 幂零根及零维 Noether 环

\section{Artin 环}
\label{b3:sec:0303}

降链条件对一般模弱于有限长度，但对交换环本身有更强的后果：它迫使极大理想只有有限多个，迫使 Jacobson 根幂零，最终使 \(R\) 作为自身的模具有有限长度。这些性质又把环分成有限个 Artin 局部环。

\subsection{降链条件与有限长度}
\label{b3:subsec:0303:01}

设 \(R\) 为交换含幺环。若非零 \(R\) 模 \(M\) 只有 \(0\) 和 \(M\) 两个子模，则称 \(M\) 为\term[jiandan-mo]{简单模}[simple module]。若 \(R\) 模 \(M\) 存在有限子模链
\[
0=M_0\subsetneq M_1\subsetneq\cdots\subsetneq M_r=M,
\]
且每个 \(M_i/M_{i-1}\) 都简单，则称 \(M\) 为\term[youxian-changdu-mo]{有限长度模}[module of finite length]。上述链称为\term[hecheng-lie]{合成列}[composition series]，各商称为合成因子。第二卷定理12.2.3，即模的 Jordan–Hölder 定理，已经证明其层数与所取链无关，称为模的\term[mo-changdu]{长度}[length of a module]，记作 \(\ell_R(M)\)。第二卷定理12.2.4进一步证明：有限长度等价于同时满足升链与降链条件。对环 \(R\) 而言，这一判据作用于 \(R\) 作为自身的模；任意 Artin 模则未必有限生成。
零模的合成列没有合成因子，长度为零。
\symidx{\(\ell_R(M)\)：模的长度}

\begin{lemma}\label{b3:lem:artin-primes-radical}
设 \(R\ne0\) 为交换含幺 Artin 环。则每个素理想都极大，极大理想只有有限多个，而且 \(J(R)\) 幂零。
\end{lemma}
\begin{proof}
若 \(\mathfrak p\) 素，则整环 \(R/\mathfrak p\) 仍满足降链条件。要证明这个商是域，只需证明其中每个非零元可逆。对这样的元素 \(a\)，主理想链
\((a)\supseteq(a^2)\supseteq\cdots\) 稳定，所以 \(a^n=a^{n+1}b\) 对某个 \(n,b\) 成立。约去非零的 \(a^n\)，得到 \(ab=1\)。该整环是域，故 \(\mathfrak p\) 极大。

若极大理想有无穷多个，从中选取一列两两不同的极大理想 \(\mathfrak m_1,\mathfrak m_2,\ldots\)，令
\(I_n=\mathfrak m_1\cdots\mathfrak m_n\)。不同极大理想之间不能有包含关系，否则较小者的极大性会迫使二者相等。因此对每个 \(i\le n\)，可以选
\(a_i\in\mathfrak m_i\setminus\mathfrak m_{n+1}\)，则
\(\prod_i a_i\in I_n\setminus\mathfrak m_{n+1}\)，因为 \(\mathfrak m_{n+1}\) 素。另一方面
\(I_{n+1}\subseteq I_n\cap\mathfrak m_{n+1}\)，故
\(I_{n+1}\subsetneq I_n\)，与降链条件矛盾。因此极大理想只有
\(\mathfrak m_1,\ldots,\mathfrak m_r\)。

记 \(J=J(R)\)。降链条件给出 \(J^n=J^{n+1}\)，其中 \(n\ge1\)，但稳定并不直接意味着 \(J^n=0\)。此时还不知道 \(J^n\) 有限生成，不能直接对它使用 Nakayama 引理。改看其零化子
\(K=\operatorname{Ann}_R(J^n)\)：证明 \(K=R\) 就等价于证明 \(J^n=0\)，而幂的稳定性给出
\[
(K:J)=\operatorname{Ann}_R(J^{n+1})=K.
\]
若 \(K\ne R\)，严格包含 \(K\) 的理想至少有 \(R\)，降链条件允许从中选一个极小者 \(K'\)。由于 \(K\) 与 \(K'\) 之间没有中间理想，商模 \(Q=K'/K\) 是非零简单模。取 \(x\in K'\setminus K\)，它的像生成 \(Q\)，也就是 \(K'=K+Rx\)，所以 \(Q\) 循环且有限生成。\cref{b3:thm:nakayama}于是给出 \(JQ\ne Q\)；否则这个非零模就会被迫为零。简单性使 \(JQ\) 只能是零或整个 \(Q\)，因此 \(JQ=0\)。这表示 \(Jx\subseteq K\)，于是 \(x\in(K:J)=K\)，与所取 \(x\notin K\) 矛盾。所以 \(K=R\)，即 \(J^n=0\)。
\end{proof}

这个证明用降链条件找到简单商模 \(K'/K\)，把 Nakayama 引理所需的有限生成性落实在这个循环模上。上面的处理可参阅\cite[Theorem 3.2, p. 16]{Matsumura1989CommutativeRingTheory}。

\begin{theorem}\label{b3:thm:artin-finite-length}
设 \(R\) 为交换含幺 Artin 环。则 \(R\) 作为自身的模具有有限长度，因而是 Noether 环；每个有限生成 \(R\) 模也具有有限长度。
\end{theorem}
\begin{proof}
零环的结论直接成立，以下设 \(R\ne0\)。沿用\cref{b3:lem:artin-primes-radical}中极大理想及其幂零交的记号，不同极大理想两两互素，中国剩余定理给出
\[
R/J\cong\prod_{i=1}^r R/\mathfrak m_i.
\]
考虑有限滤过 \(R\supseteq J\supseteq\cdots\supseteq J^n=0\)。对 \(0\le j<n\)，商层 \(L_j=J^j/J^{j+1}\) 被 \(J\) 零化，因而是 \(R/J\) 模；子模和商模又继承降链条件，所以每个 \(L_j\) 是 Artin 模。这是降链条件在商层上提供的有限性，尚未用到环的 Noether 性。

记 \(K_i=R/\mathfrak m_i\)，在 \(R/J\cong\prod_{i=1}^rK_i\) 中，令 \(e_i=(0,\ldots,1,\ldots,0)\)，其中只有第 \(i\) 个坐标为一。由 \(\sum_i e_i=1\) 及 \(e_ie_h=0\)、\(i\ne h\)，有
\[
L_j=e_1L_j\oplus\cdots\oplus e_rL_j.
\]
每个 \(e_iL_j\) 的标量作用只经过第 \(i\) 个域 \(K_i\)，其 \(R\) 子模正是 \(K_i\) 线性子空间。它作为 \(L_j\) 的子模满足降链条件，故对应向量空间也满足降链条件。这样的向量空间必有限维：若含无限线性无关序列，将其扩充为基，逐次删去其中前几个向量所得的线性生成空间就构成严格降链。有限维向量空间中，将基向量逐个加入，就得到相邻商一维的子空间链；各一维商的 \(R\) 子模只有零和自身，所以是简单模。因此各 \(e_iL_j\)，进而每个 \(L_j\)，都具有有限长度。把各商层的合成列取逆像并连接，得到 \(R\) 的有限合成列。

由第二卷定理12.2.4，有限长度推出 Noether 性。若 \(M\) 由 \(s\) 个元素生成，它是 \(R^s\) 的商；有限直和与商都保留有限长度，因此 \(M\) 也具有有限长度。
\end{proof}

\subsection{交换 Artin 环的局部直积分解}
\label{b3:subsec:0303:02}

中国剩余定理已给出一个具体模型：
\[
\mathbb Z/72\mathbb Z\cong\mathbb Z/8\mathbb Z\times\mathbb Z/9\mathbb Z.
\]
右侧两个因子都是局部环，分别对应原环的极大理想 \((\bar2)\)、\((\bar3)\)。剩余类 \(e_1=9\)、\(e_2=64\) 在这个同构下分别对应 \((1,0)\)、\((0,1)\)，所以二者和为一，乘积为零，且各自平方等于自身。任意 \(\mathbb Z/72\mathbb Z\) 模 \(M\) 自动分成 \(9M\oplus64M\)，分别保留 \(2\) 主部分与 \(3\) 主部分，不需要为模另选生成元。一般交换 Artin 环也能由各极大理想恢复这样的局部因子。

\begin{theorem}\label{b3:thm:artin-decomposition}
设 \(R\ne0\) 为交换含幺 Artin 环，极大理想为
\(\mathfrak m_1,\ldots,\mathfrak m_r\)。对任意满足 \(J(R)^n=0\) 的整数 \(n\ge1\)，自然映射给出
\[
R\cong\prod_{i=1}^r R/\mathfrak m_i^n
\cong\prod_{i=1}^r R_{\mathfrak m_i}.
\]
每个因子都是非零 Artin 局部环。相应地，\(1=e_1+\cdots+e_r\) 分解为两两正交的非零幂等元，每个 \(R\) 模都有
\(M=e_1M\oplus\cdots\oplus e_rM\)。这些局部因子由各极大理想唯一确定，分解仅差排列。
\end{theorem}
\begin{proof}
固定满足 \(J(R)^n=0\) 的整数 \(n\ge1\)。要使商映射给出直积分解，需要各 \(\mathfrak m_i^n\) 两两互素，且它们的交为零。对不同的 \(i,j\)，有
\(\mathfrak m_i+\mathfrak m_j=R\)。写 \(a+b=1\)、\(a\in\mathfrak m_i\)、\(b\in\mathfrak m_j\)，展开 \((a+b)^{2n-1}\) 后每一项至少含 \(a^n\) 或 \(b^n\)，故
\(\mathfrak m_i^n+\mathfrak m_j^n=R\)。于是
\[
\bigcap_i\mathfrak m_i^n
=\prod_i\mathfrak m_i^n
=\left(\prod_i\mathfrak m_i\right)^n=J(R)^n=0.
\]
中国剩余定理给出第一个同构。

这些商因子都是局部环，因为包含 \(\mathfrak m_i^n\) 的极大理想必包含 \(\mathfrak m_i\)：对每个 \(a\in\mathfrak m_i\)，有 \(a^n\in\mathfrak m_i^n\)，再用素性即可。因 \(\mathfrak m_i\) 本身极大，这样的极大理想只能是 \(\mathfrak m_i\)。又因 \(\mathfrak m_i^n\subseteq\mathfrak m_i\subsetneq R\)，商环非零；它继承 Artin 性，所以是非零 Artin 局部环。

在上述直积中，令 \(e_i\) 对应第 \(i\) 个坐标为 \(1\)、其余为 \(0\) 的元素。它们非零、两两正交且和为 \(1\)。对每个 \(R\) 模 \(M\)，等式 \(m=\sum_i e_i m\) 给出所述分解；若 \(\sum_i m_i=0\)、\(m_i\in e_iM\)，逐次乘 \(e_j\) 得 \(m_j=0\)，故和是直和。

局部因子还可直接从原环的极大理想恢复，这将使它们不依赖所取的 \(n\)。考虑任意有限直积分解
\[
R\cong\prod_{j=1}^s B_j,
\]
其中 \(B_j\) 是非零局部环，唯一极大理想记作 \(\mathfrak n_j\)。沿此同构识别 \(R\) 与直积。若 \(\mathfrak m\) 是 \(R\) 的极大理想，各坐标幂等元在域 \(R/\mathfrak m\) 中只能取 \(0\) 或 \(1\)，且由正交性及其和为 \(1\)，恰有一个坐标幂等元取值 \(1\)。设它对应第 \(j\) 个坐标，则其他坐标因子都在 \(\mathfrak m\) 中，而第 \(j\) 个坐标上的商必须是域。因此
\[
\mathfrak m=\mathfrak m_j
=B_1\times\cdots\times\mathfrak n_j\times\cdots\times B_s.
\]
反过来，每个 \(\mathfrak m_j\) 的商是域 \(B_j/\mathfrak n_j\)，故它们恰好穷尽 \(R\) 的极大理想。

第 \(j\) 个坐标投影将 \(R\setminus\mathfrak m_j\) 映到 \(B_j\) 的单位，因而诱导映射
\[
R_{\mathfrak m_j}\longrightarrow B_j,
\qquad a/u\longmapsto a_j u_j^{-1}.
\]
它是满射，因为坐标投影已经满射。若其值为零，则 \(a_j=0\)；第 \(j\) 个坐标幂等元 \(e_j\notin\mathfrak m_j\) 满足 \(e_j a=0\)，所以 \(a/u=0\)。映射的核为零，从而 \(R_{\mathfrak m_j}\cong B_j\)。应用于中国剩余定理所得的分解，就得到 \(R_{\mathfrak m_i}\cong R/\mathfrak m_i^n\)。对任意另一组非零局部因子，极大理想集合仍恢复同一组局部化，故因子只差排列和同构。
\end{proof}

零环没有极大理想，采用空直积为零环的约定后，它对应空的局部分解。零环上的幺性模只有零模，相应的模分解也是空直和。

\subsection{幂零根与零维 Noether 环}
\label{b3:subsec:0303:03}

本小节把“零维”用于每个素理想都极大的非零环；一般 Krull 维数的定义和基本定理留在\chapref{b3:ch:09}。零环作为单独的平凡情形处理。

\begin{theorem}\label{b3:thm:noether-dimension-zero}
设 \(R\) 为非零交换含幺环。下列条件等价：
\begin{equivalences}
\item \(R\) 为 Artin 环；
\item \(R\) 为 Noether 环，且每个素理想都极大。
\end{equivalences}
此时幂零根与 Jacobson 根相等，均为幂零理想。\(R\) 约化当且仅当它是有限个域的直积。
\end{theorem}
\begin{proof}
由\cref{b3:lem:artin-primes-radical}及\cref{b3:thm:artin-finite-length}，第一项推出第二项。

反过来，设 \(R\) 为 Noether 环且每个素理想极大。第一章的\cref{b3:prop:noether-spectrum-components}表明其极小素理想只有有限多个，且每个素理想包含一个极小素理想。现在包含在两个素理想之间只能是相等，所以整个素谱就是有限个极大理想
\(\mathfrak m_1,\ldots,\mathfrak m_r\)。由\cref{b3:thm:radical-prime-intersection}，
\[
\sqrt{(0)}=\bigcap_i\mathfrak m_i=J(R).
\]
Noether 条件使幂零根有限生成；它的每个元素幂零，故由\cref{b3:prop:finite-nil-ideal}，该理想的某个幂为零。

记 \(J=J(R)\)，中国剩余定理仍给出 \(R/J\cong\prod_i R/\mathfrak m_i\)。这里各理想 \(J^j\) 因 Noether 性而有限生成，所以商层 \(J^j/J^{j+1}\) 是有限生成 \(R\) 模。它被 \(J\) 零化，因而也是有限生成的 \(R/J\) 模。与前面从降链条件得到 Artin 商层不同，这个方向依靠有限生成性：按坐标幂等元分解以后，每个域上的向量空间分量由原来有限个生成元的像生成，故有限维。因此各商层具有有限长度；沿有限的 \(J\) 幂滤过连接合成列，得到 \(R\) 具有有限长度，故满足降链条件。

最后，若 \(R\) 约化，上述公共根为零，直接得到 \(R\cong\prod_i R/\mathfrak m_i\)。有限个域的直积没有非零幂零元，反向也成立。
\end{proof}

上述两种有限性给出不同的蕴含关系。在 \(R\) 为 Artin 环的假设下，有
\[
\left.
\begin{gathered}
\#\operatorname{Max}R<\infty,\quad J^n=0,\\
J^i/J^{i+1}\;\text{满足降链条件}
\end{gathered}
\right\}
\Longrightarrow\ell_R(R)<\infty
\Longrightarrow R\;\text{为 Noether 环}.
\]
而在 \(R\) 为 Noether 环且每个素理想都极大的假设下，有
\[
\left.
\begin{gathered}
\operatorname{Spec}R\;\text{有限},\quad J=\sqrt{(0)},\quad J^n=0,\\
J^i/J^{i+1}\;\text{有限生成}
\end{gathered}
\right\}
\Longrightarrow\ell_R(R)<\infty
\Longrightarrow R\;\text{为 Artin 环}.
\]
\ProseParagraphBreak
去掉 Noether 条件，零维便不足以推出 Artin 性。\secref{b3:sec:0102}中的无限域直积 \(R=\prod_{n\ge1}k\) 已经给出反例：它的每个素理想都极大，而使前 \(n\) 个坐标为零的理想构成严格降链。

% !TeX root = ../../Book3.tex
% 纲要标题：### 3.4 有限性在变换下的行为
% 纲要位置：工作记录/计划纲要.md，第 1897 行。
% 纲要细目（写作范围）：
% * 商、局部化和有限代数
% * Noether 性并非任意子环继承
% * 有限性假设的常见误用
% ---

\section{有限性在变换下的行为}
\label{b3:sec:0304}

商、局部化和有限代数都保留相当多的有限性，而取子环一般不保留。把这些结论放在一起，可以避免在具体计算中因“环变小了”或“只剩一个点了”就误判其理想结构。本节还区分茎上的条件与基本开邻域上的条件。

\subsection{商、局部化与有限代数}
\label{b3:subsec:0304:01}

\begin{proposition}\label{b3:prop:chain-conditions-localization}
设 \(R\) 为交换含幺环，\(I\subseteq R\) 为理想，\(S\subseteq R\) 为含一的乘法集。若 \(R\) 为 Noether 环，则 \(R/I\) 与 \(S^{-1}R\) 仍为 Noether 环；若 \(R\) 为 Artin 环，则二者仍为 Artin 环。对任意 \(R\) 模 \(M\)，若 \(M\) 有限生成或有限展示，则 \(S^{-1}M\) 作为 \(S^{-1}R\) 模仍分别满足相应条件。
\end{proposition}
\begin{proof}
商环的理想链取逆像，严格包含仍严格，故两种链条件都传给商。局部化环中每个理想都是其收缩的扩张，见\cref{b3:prop:ideal-localization-saturation}。因此一条严格理想链收缩后仍严格；原环的相应链条件排除了无限严格升链或降链。

模的生成元取分式 \(m_i/1\) 就生成局部化模。对有限展示
\(R^a\rightarrow R^b\rightarrow M\rightarrow0\)，使用\cref{b3:thm:localization-exact}，得到
\((S^{-1}R)^a\rightarrow(S^{-1}R)^b\rightarrow S^{-1}M\rightarrow0\)，故有限展示也被保持。
\end{proof}

\begin{proposition}\label{b3:prop:finite-algebra-noether}
设 \(R\) 为交换含幺环，\(A\) 为交换含幺 \(R\) 代数。若 \(A\) 作为 \(R\) 模有限生成，则称 \(A\) 为有限 \(R\) 代数。设 \(A\) 满足这一条件。
\begin{enumerate}
\item 若 \(R\) 为 Noether 环，则 \(A\) 为 Noether 环。
\item 若 \(R\) 为 Artin 环，则 \(A\) 为 Artin 环。
\item 若 \(M\) 为有限生成 \(A\) 模，则它也是有限生成 \(R\) 模。
\end{enumerate}
\end{proposition}
\begin{proof}
第一项中，\(A\) 作为有限生成 \(R\) 模满足升链条件，而每个 \(A\) 理想都是 \(R\) 子模，所以理想升链稳定。第二项同理：由\cref{b3:thm:artin-finite-length}，\(A\) 作为 \(R\) 模具有有限长度，其 \(R\) 子模降链稳定，特别地 \(A\) 理想降链稳定。第三项中，若 \(a_1,\ldots,a_s\) 生成 \(A\) 作为 \(R\) 模，\(m_1,\ldots,m_t\) 生成 \(M\) 作为 \(A\) 模，则全部 \(a_i m_j\) 生成 \(M\) 作为 \(R\) 模。
\end{proof}

这里的“有限”比有限型严格得多。\(k[x]\) 是域 \(k\) 上的有限型代数，却不是 Artin 环，因为 \((x)\supsetneq(x^2)\supsetneq\cdots\)。Noether 底环上的有限型代数保持 Noether 性，依据是 Hilbert 基定理；Artin 底环上的有限代数保持 Artin 性，依据则是模的有限长度。

\begin{proposition}\label{b3:prop:finite-presentation-open-cover}
设 \(R\) 为交换含幺环，\(M\) 为 \(R\) 模，\(f_1,\ldots,f_r\in R\)（\(r\ge1\)）生成单位理想。若每个 \(M_{f_i}\) 是有限展示
\(R_{f_i}\) 模，则 \(M\) 有限展示。若每个 \(R_{f_i}\) 是 Noether 环，则 \(R\) 为 Noether 环。
\end{proposition}
\begin{proof}
由\cref{b3:prop:principal-cover-finiteness}，\(M\) 有限生成，取满射
\(R^n\rightarrow M\)，核为 \(K\)。局部化后仍正合，且目标有限展示，所以由
\cref{b3:prop:finite-presentation-kernel}，每个 \(K_{f_i}\) 有限生成。再次使用有限基本开覆盖上的有限生成性判据，得到 \(K\) 有限生成，因而 \(M\) 有限展示。

对第二句，任取理想 \(I\subseteq R\)。其局部化是 Noether 环 \(R_{f_i}\) 的理想，故有限生成。对 \(I\) 使用同一有限生成性判据，得到每个理想都有限生成，所以 \(R\) 为 Noether 环。
\end{proof}

\subsection{Noether 环的非 Noether 子环}
\label{b3:subsec:0304:02}

设 \(k\) 为域，取多项式环 \(B=k[x,y]\) 的子环
\[
A=k+xk[x,y]
=k[x,xy,xy^2,\ldots].
\]
等号右侧表示常数项之外的每个单项式都至少含一个 \(x\)；它对乘法封闭，并且与 \(B\) 具有相同的单位元。环 \(B\) 由 Hilbert 基定理是 Noether 环，\(A\) 却不是。

\ProseParagraphBreak
为此令 \(I=xk[x,y]\)，它是 \(A\) 的理想，\(A/I\cong k\)，而
\(I^2=x^2k[x,y]\)。于是
\[
I/I^2
=\bigoplus_{j\ge0}k\cdot(xy^j+I^2)
\]
为无限维 \(k\) 向量空间。如果 \(I\) 由有限个元素作为 \(A\) 理想生成，模掉 \(I^2\) 后，这些元素的剩余类就会在 \(A/I=k\) 上生成 \(I/I^2\)，与无限维性矛盾。故 \(I\) 不有限生成，\(A\) 不是 Noether 环。

\ProseParagraphBreak
这个例子也显示问题所在：在较大的环 \(B\) 中，\(I\) 由一个元素 \(x\) 生成，因为可以任意乘以 \(y\)；在 \(A\) 中没有这个乘子。取子环虽然减少了元素，却也减少了生成理想时允许使用的系数，可能使有限生成性失败。

\subsection{有限性假设的适用范围}
\label{b3:subsec:0304:03}

有限个生成元只描述模本身，并不保证它的所有子模也能如此描述。设 \(k\) 为域，环 \(k[x_1,x_2,\ldots]\) 作为自身的模由 \(1\) 生成，其理想 \((x_1,x_2,\ldots)\) 却不有限生成：任意有限组候选生成元只涉及有限多个变量，令这些变量全为零后，它们生成的理想也映为零，而其余变量的像非零。在 Noether 底环上，\cref{b3:prop:noether-finite-modules}保证有限生成模的每个子模都有限生成。

\ProseParagraphBreak
即使每个点处都能找到有限组生成元，也未必已经得到有限的全局描述。\subsecref{b3:subsec:0204:03}中的
\(\bigoplus_p\mathbb Z/p\mathbb Z\) 在每个素理想处都有限生成，但全局不是有限生成模。\cref{b3:prop:principal-cover-finiteness}及\cref{b3:prop:finite-presentation-open-cover}要求的是一组基本开邻域上的有限数据；仅检查各个茎不足以自动得到这些邻域。

\ProseParagraphBreak
环的 Noether 性也有相同的区别。\secref{b3:sec:0102}中的无限域直积 \(R=\prod_{n\ge1}k\) 满足：每个元素 \(a\) 都可写成 \(a=a^2b\)。若 \(a\in\mathfrak p\)，则 \(1-ab\notin\mathfrak p\)，而 \(a(1-ab)=0\)，所以 \(a/1=0\) 于 \(R_{\mathfrak p}\)。于是 \(\mathfrak pR_{\mathfrak p}=0\)，每个局部环 \(R_{\mathfrak p}\) 都是域。然而，令 \(e_n\) 为第 \(n\) 个坐标为一、其余坐标为零的元素，理想 \(Re_1+\cdots+Re_n\) 组成严格升链，故 \(R\) 不是 Noether 环。各点处的 Noether 性也不能代替基本开邻域上的条件。

\ProseParagraphBreak
降链条件也不能一般地保证模有限生成。设 \(p\) 为素数，考虑整数模
\[
C_{p^\infty}=\mathbb Z[1/p]/\mathbb Z
\]
记 \(C_n=\langle p^{-n}+\mathbb Z\rangle\)（\(n\ge0\)），则 \(C_0=0\)，\(C_n\) 是阶为 \(p^n\) 的循环群，且
\[
C_0\subsetneq C_1\subsetneq C_2\subsetneq\cdots,
\qquad C_{p^\infty}=\bigcup_{n\ge0}C_n.
\]
任意阶为 \(p^n\) 的元素都能写成 \(a/p^n+\mathbb Z\)，其中 \(p\nmid a\)。取整数 \(u,v\) 使 \(ua+vp^n=1\)，就有
\(u(a/p^n+\mathbb Z)=p^{-n}+\mathbb Z\)，所以该元素恰好生成 \(C_n\)。此外，被 \(p^n\) 零化的元素恰为 \(C_n\) 的元素。故任意阶为 \(p^n\) 的子群都包含于 \(C_n\)，比较元素个数可知它等于 \(C_n\)：每一阶的子群只有一个。

\ProseParagraphBreak
现在设 \(H\) 为子模。若 \(H\) 中元素的阶无界，则对每个 \(n\) 都能取到阶为 \(p^r\)、\(r\ge n\) 的元素；它生成 \(C_r\)，因而 \(C_n\subseteq H\)，于是 \(H=C_{p^\infty}\)。若阶有界，则 \(H\subseteq C_n\) 对某个 \(n\) 成立。在有限群 \(H\) 中选阶最大的元素，设其阶为 \(p^r\)，则它生成 \(C_r\)，而 \(H\) 中其余元素都被 \(p^r\) 零化，所以 \(H=C_r\)。这就证明真子模恰为各个有限循环群 \(C_r\)。任何子模降链一旦离开全模便进入有限群，此后稳定，故 \(C_{p^\infty}\) 是 Artin 模；任意有限组元素都落在某个 \(C_n\) 中，不能生成全模，故它不是有限生成模。

\ProseParagraphBreak
\cref{b3:tab:finiteness-counterexamples}汇集了这些反例，以及保证相应结论成立的充分条件。

\begin{table}[htbp]
\centering
\caption{有限性推论的反例与适用条件}
\label{b3:tab:finiteness-counterexamples}
\begin{tabularx}{\textwidth}{@{}>{\raggedright\arraybackslash}p{.25\textwidth}>{\raggedright\arraybackslash}p{.30\textwidth}>{\raggedright\arraybackslash}X@{}}
\toprule
一般不成立的推论 & 反例 & 保证相应结论的条件\\
\midrule
有限生成模的子模都有限生成 & 无限变量多项式环的变量理想 & 底环是 Noether 环\\
各个素理想处有限生成，便全局有限生成 & \(\bigoplus_{p\text{ 素数}}\mathbb Z/p\mathbb Z\) & 一组覆盖素谱的基本开邻域上有限生成\\
Artin 模必有限生成 & 整数模 \(C_{p^\infty}\) & 模同时满足 Noether 条件，此时具有有限长度\\
\bottomrule
\end{tabularx}
\end{table}

有限生成使有限组生成元可以共用分母，也是 Nakayama 引理对模的要求；有限展示进一步要求关系模有限生成。Noether 条件控制所有子模的有限生成性，Artin 条件控制降链；对交换环本身，后者还给出\cref{b3:thm:artin-decomposition}中的局部直积分解。


\begin{chaptersummary}
有限展示要求模有有限个生成元，且全部关系由有限多个基本关系线性生成，并不把关系总数限制为有限。Noether 环使每个有限生成模的子模仍有限生成。Hilbert 基定理保证 Noether 环上的有限型代数仍是 Noether 环，因而其中的理想与有限矩阵的关系模具有有限描述。

\ProseParagraphBreak
行列式技巧不要求生成元线性无关，却能产生一个统一的首一零化方程。对恒等映射使用它，再利用 Jacobson 根的单位判据，便得到 Nakayama 引理。在局部环上，极小生成元数因此等于剩余向量空间的维数；这项维数不负责判断关系是否消失。

\ProseParagraphBreak
交换 Artin 环具有有限长度，并分解为各极大理想处的局部环。反向刻画需要 Noether 性与零维同时成立。有限代数、商及局部化保留适当的链条件，任意子环却不保留；后续整扩张与维数论中的有限性假设将继续依赖这些区别。
\end{chaptersummary}

\BookExercises

\begin{exercise}\label{b3:ex:03-cyclic-presentation}
设 \(k\) 为域，\(R=k[x,y]\)、\(\mathfrak m=(x,y)\)、\(I=(x^2,xy,y^2)\)。给出循环模 \(R/I\) 的有限展示，并说明三个基本关系 \(x^2\cdot(1+I)=xy\cdot(1+I)=y^2\cdot(1+I)=0\) 都不能直接删去。再通过 \(I/\mathfrak mI\) 证明：针对生成元 \(1+I\)，其关系模不能由两个元素生成。
\end{exercise}
\begin{solution}
关系模是商映射 \(R\to R/I\) 的核 \(I\)，故有有限展示
\[
R^3\xrightarrow{\begin{pmatrix}x^2&xy&y^2\end{pmatrix}}R
\longrightarrow R/I\longrightarrow0.
\]
一个模生成元为 \(1+I\)，全部关系是这三个基本关系的 \(R\) 线性组合。若删去某个二次单项式，它不能由另两个生成：乘以多项式后的二次项只能是这两个单项式的 \(k\) 线性组合。

更一般地，\(\mathfrak mI=(x,y)^3\) 只含全次数至少为三的项，所以 \(I/\mathfrak mI\) 以 \(x^2,xy,y^2\) 的类为 \(k\) 基，维数为三。由于 \(\mathfrak m\) 在这个商上作用为零，\(R\) 的作用经过 \(R/\mathfrak m=k\)。若任意两个元素生成 \(I\)，它们的剩余类就生成该三维向量空间，矛盾。因此不是换两条巧选的基本关系就能代替这三条。这里计算的是关系模的生成元数，模 \(R/I\) 本身仍只有一个生成元。
\end{solution}

\begin{exercise}\label{b3:ex:03-two-variable-relation}
设 \(k\) 为域，令 \(R=k[x,y]\)。求映射 \(R^2\rightarrow(x,y)\)、\((a,b)\mapsto ax+by\) 的核，并给出 \((x,y)\) 的有限展示。
\end{exercise}
\begin{solution}
若 \(ax+by=0\)，模掉 \(y\) 后在 \(k[x]\) 中有 \(\bar a x=0\)，故 \(a=yc\)。代回得到 \(y(cx+b)=0\)，整环性给出 \(b=-xc\)。因此核由 \((y,-x)\) 生成，得到正合列
\[
0\longrightarrow R\xrightarrow{c\mapsto(yc,-xc)}R^2
\longrightarrow(x,y)\longrightarrow0.
\]
第一个映射单射，因为 \(y\ne0\) 且 \(R\) 为整环。
\end{solution}

\begin{exercise}\label{b3:ex:03-finite-extension-modules}
设 \(R\) 为交换含幺环，\(0\rightarrow L\rightarrow M\rightarrow N\rightarrow0\) 为 \(R\) 模的正合列。证明若 \(L,N\) 有限生成，则 \(M\) 有限生成；同时给出 \(M\) 有限生成而 \(L\) 不有限生成的例子。
\end{exercise}
\begin{solution}
取 \(L\) 的有限生成元，并提升 \(N\) 的有限生成元到 \(M\)。对任意 \(m\in M\)，先减去提升元的适当线性组合，使差落入 \(L\)，再用 \(L\) 的生成元表示，故这些有限个元素生成 \(M\)。
反向取域 \(k\)，令 \(R=k[x_1,x_2,\ldots]\)、\(M=R\)、\(L=(x_1,x_2,\ldots)\)。\(M\) 循环，但 \(L\) 不有限生成：任意有限组候选生成元只涉及有限个变量，将这些变量置零后，它们全变为零，而一个未出现的变量仍非零。此处底环不是 Noether 环。
\end{solution}

\begin{exercise}\label{b3:ex:03-cusp-finite}
设 \(k\) 为域，令 \(A=k[t^2,t^3]\subseteq k[t]\)、\(R=k[t^2]\)。证明 \(A\) 是有限 \(R\) 代数并为 Noether 环；比较 \(k[t,y]\) 作为 \(k[t]\) 代数的有限型性与有限性。
\end{exercise}
\begin{solution}
每个 \(A\) 中的单项式次数都为偶数，或为至少 \(3\) 的奇数，故
\(A=R+Rt^3\)。所以 \(A\) 作为 \(R\) 模由 \(1,t^3\) 生成。\(R\) 同构于一元多项式环，是 Noether 环，故由\cref{b3:prop:finite-algebra-noether}，\(A\) 为 Noether 环。
另一方面 \(k[t,y]\) 作为 \(k[t]\) 代数由 \(y\) 生成，是有限型代数；若它作为模由有限个多项式生成，生成元的 \(y\) 次数有共同上界，而 \(k[t]\) 系数的组合不能超过这个上界，不能得到任意高次的 \(y^n\)。所以它不是有限代数。
\end{solution}

\begin{exercise}\label{b3:ex:03-determinant-integrality}
设 \(R\) 为整环，\(K\) 为其分式域，\(0\ne M\subseteq K\) 为有限生成 \(R\) 子模。若 \(z\in K\) 满足 \(zM\subseteq M\)，证明 \(z\) 满足一个 \(R\) 系数的首一多项式。
\end{exercise}
\begin{solution}
乘 \(z\) 给出 \(M\) 的 \(R\) 线性自同态 \(m\mapsto zm\)，其中“自同态”指从 \(M\) 到自身的同态，不要求可逆。对 \(I=R\) 应用\cref{b3:thm:determinant-trick}，得到首一 \(P(T)\in R[T]\)，使
\(P(z)\cdot m=0\) 对所有 \(m\in M\) 成立。取一个非零 \(m\in M\subseteq K\)，在域中约去它，得到 \(P(z)=0\)。非零子模的条件负责最后的约去；若 \(M=0\)，这个论证不能给 \(z\) 任何限制。
\end{solution}

\begin{exercise}\label{b3:ex:03-general-nakayama-factor}
在 \(R=\mathbb Z\)、\(I=(2)\)、\(M=\mathbb Z/3\mathbb Z\) 中，找出 \(a\in I\) 使 \((1-a)M=0\)。解释为什么这不意味着 \(M=0\)。
\end{exercise}
\begin{solution}
取 \(a=4\)，则 \(1-a=-3\) 零化 \(M\)，而 \(4\in(2)\)。这里 \(-3\) 不是整数环的单位，不能从 \((-3)M=0\) 约去 \(-3\)。\((2)\) 也不包含于
\(J(\mathbb Z)=0\)，因为所有极大理想 \((p)\) 的交为零。所以满足的是行列式技巧给出的统一零化元结论，没有满足 Nakayama 的 Jacobson 根假设。
\end{solution}

\begin{exercise}\label{b3:ex:03-local-matrix}
设 \(R\) 为交换含幺局部环，极大理想为 \(\mathfrak m\)，剩余域为 \(\kappa=R/\mathfrak m\)，\(n\ge1\)，\(A\in M_n(R)\)。证明 \(A\) 可逆当且仅当其模 \(\mathfrak m\) 的矩阵 \(\bar A\in M_n(\kappa)\) 可逆。再设 \(M\) 为有限生成 \(R\) 模，\(m_1,\ldots,m_n\) 与 \(m'_1,\ldots,m'_n\) 是它的两组极小生成元。证明任意满足
\[
m'_i=\sum_{j=1}^n b_{ij}m_j\qquad(1\le i\le n)
\]
的转换矩阵 \(B=(b_{ij})\in M_n(R)\) 都可逆。
\end{exercise}
\begin{solution}
若 \(A\) 可逆，取商后仍可逆。反向若 \(\bar A\) 可逆，则
\(\det A\bmod\mathfrak m=\det\bar A\ne0\)，故 \(\det A\) 为 \(R\) 的单位，伴随矩阵公式给出 \(A^{-1}=(\det A)^{-1}\operatorname{adj}(A)\)。
由\cref{b3:cor:local-generators}，两组极小生成元模掉 \(\mathfrak mM\) 后都是 \(M/\mathfrak mM\) 的 \(\kappa\) 基。因此上述等式约化后给出两组基之间的转换矩阵 \(\bar B\)，它必可逆；由第一部分，\(B\) 在 \(R\) 上可逆。生成元未必是自由基，所以表示它们的矩阵未必唯一；这里证明的是每个这样的矩阵都可逆。
\end{solution}

\begin{exercise}\label{b3:ex:03-minimal-power-generators}
设 \(k\) 为域，\(d\ge1\)，令 \(R=k[x_1,\ldots,x_d]_{(x_1,\ldots,x_d)}\)、\(\mathfrak m=(x_1,\ldots,x_d)R\)。求 \(\mu_R(\mathfrak m^r)\)、\(r\ge0\)。
\end{exercise}
\begin{solution}
\(\mathfrak m^r/\mathfrak m^{r+1}\) 以全次数 \(r\) 的单项式为 \(k\) 基。局部化分母的常数项非零，在这个商中只影响其非零常数倍，不引入新的线性关系。因此
\[
\mu_R(\mathfrak m^r)=\dim_k(\mathfrak m^r/\mathfrak m^{r+1})
=\binom{r+d-1}{d-1}.
\]
末个数是在 \(d\) 个非负整数中分配总和 \(r\) 的方案数。\(r=0\) 时即 \(R/\mathfrak m\) 的维数 \(1\)，与 \(R\) 作为自身的模循环一致。
\end{solution}

\begin{exercise}\label{b3:ex:03-nonlocal-minimal}
证明整数模 \(\mathbb Z\) 同时有极小生成组 \(\{1\}\) 和 \(\{2,3\}\)。说明这不违反局部环上的极小生成元结论。
\end{exercise}
\begin{solution}
\(\{1\}\) 显然生成；\(\{2,3\}\) 也生成，因为 \(3-2=1\)，但删去任意一个只得到 \(2\mathbb Z\) 或 \(3\mathbb Z\)，所以它也极小。“极小”是不能删去元素，并不是在所有生成组中个数最少。
整数环不是局部环，因此不能把不同极大理想的剩余信息统一成同一个剩余向量空间。局部结论所用的 \(\mathfrak m\) 唯一性在此缺失。
\end{solution}

\begin{exercise}\label{b3:ex:03-surjective-endomorphism}
设 \(R\) 为交换含幺环，\(M,N\) 为 \(R\) 模，其中 \(M\) 有限生成。若 \(M\cong M\oplus N\) 为 \(R\) 模同构，证明 \(N=0\)。再用整数模 \(\mathbb Z\) 的一个单自同态说明“满自同态可逆”不能改成“单自同态可逆”。
\end{exercise}
\begin{solution}
将给定同构 \(M\rightarrow M\oplus N\) 与第一投影复合，得到 \(M\) 的满自同态。由\cref{b3:cor:finite-surjective-endomorphism}它单射，而其核在原同构下恰为 \(0\oplus N\)，故 \(N=0\)。
整数模的乘二映射为单射但不是满射，因 \(1\) 没有原像，所以不具有逆映射。
\end{solution}

\begin{exercise}\label{b3:ex:03-idempotent-ideal}
设 \(R\) 为交换含幺环，理想 \(I\subseteq R\) 有限生成且 \(I^2=I\)。证明存在幂等元 \(e\in R\) 使 \(I=eR\)。若 \(R\) 局部，说明只能有 \(I=0\) 或 \(I=R\)。
\end{exercise}
\begin{solution}
把 \(I\) 看作有限生成模，行列式技巧应用于 \(I=II\)，得到 \(e\in I\) 使
\((1-e)I=0\)。于是 \(x=ex\) 对所有 \(x\in I\) 成立，故 \(I\subseteq eR\)；另一包含由 \(e\in I\) 得到。再取 \(x=e\)，便有 \(e=e^2\)。
在局部环中，\(e\) 与 \(1-e\) 至少一个是单位，因为二者之和为 \(1\)。由
\(e(1-e)=0\)，若 \(e\) 可逆则 \(e=1\)，若 \(1-e\) 可逆则 \(e=0\)。因此 \(I\) 只能是所述两个理想。
\end{solution}

\begin{exercise}\label{b3:ex:03-integer-artin-product}
显式分解 \(\mathbb Z/360\mathbb Z\) 为 Artin 局部环的直积，并找出三个对应的正交幂等元。
\end{exercise}
\begin{solution}
\(360=8\cdot9\cdot5\)，三因子两两互素，所以
\[
\mathbb Z/360\mathbb Z\cong
\mathbb Z/8\mathbb Z\times\mathbb Z/9\mathbb Z\times\mathbb Z/5\mathbb Z.
\]
可取 \(e_1=225\)、\(e_2=280\)、\(e_3=216\) 的剩余类。它们在对应因子中取 \(1\)，在另外两个因子中取 \(0\)；由中国剩余同构可知它们幂等、两两正交，且和为 \(1\)。也可直接核对 \(225+280+216=721\equiv1\pmod{360}\)。
\end{solution}

\begin{exercise}\label{b3:ex:03-artin-square-example}
设 \(k\) 为域，对 \(A=k[x,y]/(x^2,y^2)\)，求唯一极大理想、其幂零指数、环的长度，以及零化该极大理想的元素组成的理想。
\end{exercise}
\begin{solution}
\(A\) 的 \(k\) 基为 \(1,\bar x,\bar y,\bar x\bar y\)。非零常数项的元素可用有限几何级数求逆，零常数项的元素幂零，所以唯一极大理想为
\(\mathfrak m=(\bar x,\bar y)\)。有
\(\mathfrak m^2=k\bar x\bar y\ne0\)、\(\mathfrak m^3=0\)，幂零指数为 \(3\)。
商层 \(A/\mathfrak m\)、\(\mathfrak m/\mathfrak m^2\)、\(\mathfrak m^2\) 的维数分别为 \(1,2,1\)，将中间二维层细分即得长度 \(4\)。
若 \(a+b\bar x+c\bar y+d\bar x\bar y\) 被 \(\bar x,\bar y\) 同时零化，分别相乘得到 \(a=c=0\) 与 \(a=b=0\)，故零化子为 \(k\bar x\bar y\)。
\end{solution}

\begin{exercise}\label{b3:ex:03-polynomial-artin-product}
设 \(k\) 为域。分解 \(A=k[t]/\bigl(t^2(t-1)^3\bigr)\) 为局部因子，求其长度及约化商。
\end{exercise}
\begin{solution}
\((t^2)\) 与 \(\bigl((t-1)^3\bigr)\) 互素，所以
\[
A\cong k[t]/(t^2)\times k[t]/\bigl((t-1)^3\bigr).
\]
两个因子分别是以 \((\bar t)\) 与 \((\overline{t-1})\) 为唯一极大理想的 Artin 局部环，具有长度 \(2,3\)，因此 \(A\) 的长度为 \(5\)。幂零根在两个因子中分别为这两个极大理想，约化商为 \(k\times k\)。
\end{solution}

\begin{exercise}\label{b3:ex:03-nilpotent-maximal-not-artin}
设 \(k\) 为域，\(V\) 为无限维 \(k\) 向量空间，在 \(R=k\oplus V\) 上规定
\((a,u)(b,v)=(ab,av+bu)\)。证明它是局部环，极大理想平方为零，但它不是 Artin 环。
\end{exercise}
\begin{solution}
该乘法直接满足结合律及分配律，单位为 \((1,0)\)。若 \(a\ne0\)，
\((a,u)^{-1}=(a^{-1},-a^{-2}u)\)；若 \(a=0\)，元素平方为零，所以非单位集恰为 \(\mathfrak m=0\oplus V\)，且 \(\mathfrak m^2=0\)。
\(V\) 的每个 \(k\) 子空间都是 \(R\) 理想，因为 \(V\) 对它作用为零，标量 \(k\) 按通常方式作用。选一列线性无关向量并扩充为基，逐次删去前面的向量所得子空间构成严格降链。因此 \(R\) 不是 Artin 环。极大理想幂零本身不能替代有限性。
\end{solution}

\begin{exercise}\label{b3:ex:03-zero-dimensional-reduced}
设 \(k\) 为域。证明有限维约化交换含幺 \(k\) 代数是有限个有限域扩张的直积（零代数对应空乘积）。指出当 \(k\) 代数闭时这些因子是什么。
\end{exercise}
\begin{solution}
设该代数为 \(A\)。若 \(A=0\)，按约定它就是空乘积。以下设 \(A\ne0\)。有限维性使理想升链、降链都稳定，故 \(A\) 为 Artin 环；约化性结合
\cref{b3:thm:noether-dimension-zero}给出有限个域的直积。各域为原代数的商，因此作为 \(k\) 向量空间仍有限维，是有限域扩张。若 \(k\) 代数闭，每个有限扩张中的元素都有 \(k\) 上的最小多项式，该不可约多项式只能为一次，故各因子均等于 \(k\)。
\end{solution}

\begin{exercise}\label{b3:ex:03-artin-local-length}
设 \(R\) 为交换含幺 Artin 局部环，极大理想为 \(\mathfrak m\)，剩余域为 \(\kappa=R/\mathfrak m\)，\(M\) 为有限生成 \(R\) 模。证明
\[
\ell_R(M)=\sum_{i\ge0}
\dim_\kappa(\mathfrak m^iM/\mathfrak m^{i+1}M).
\]
推导 \(\ell_R(M)\ge\mu_R(M)\)，并刻画等号成立的情形。
\end{exercise}
\begin{solution}
由\cref{b3:thm:artin-finite-length}，\(R\) 是 Noether 环；再由\cref{b3:prop:noether-finite-modules}，有限生成模 \(M\) 的每个子模 \(\mathfrak m^iM\) 都有限生成。因此各商层也是有限生成 \(R\) 模；它们被 \(\mathfrak m\) 零化，遂成为有限维 \(\kappa\) 向量空间。又因 \(\mathfrak m\) 幂零，和中只有有限个非零项。按各层一组基的前若干项构造向量空间旗标。这些商层上的 \(R\) 作用都经过 \(R/\mathfrak m=\kappa\)，所以其 \(R\) 子模恰是 \(\kappa\) 线性子空间；每个一维商只有零子模与自身，因而是简单 \(R\) 模 \(\kappa\)。将各层的旗标取逆像并连接，就得到总层数为右侧的合成列。
第一项为 \(\dim_\kappa M/\mathfrak mM=\mu_R(M)\)，其他项非负，因此有不等式。若等号成立，则 \(\mathfrak mM/\mathfrak m^2M=0\)。上面已经证明 \(\mathfrak mM\) 有限生成，对它应用\cref{b3:thm:nakayama}，得到 \(\mathfrak mM=0\)。反向若 \(\mathfrak mM=0\)，只有第一项非零，当然取等号。
\end{solution}

\begin{exercise}\label{b3:ex:03-artin-localized-length}
设 \(R\) 为交换含幺 Artin 环，极大理想为 \(\mathfrak m_1,\ldots,\mathfrak m_r\)，\(M\) 为有限生成 \(R\) 模。证明
\[
\ell_R(M)=\sum_{i=1}^r\ell_{R_{\mathfrak m_i}}(M_{\mathfrak m_i}).
\]
\end{exercise}
\begin{solution}
由\cref{b3:thm:artin-decomposition}，\(R=\prod_i e_iR\)，且
\(M=\bigoplus_i e_iM\)、\(e_iM\cong M_{\mathfrak m_i}\)。
一个 \(e_iM\) 的 \(R\) 子模恰为其 \(e_iR\) 子模，因为其他坐标作用为零；因此两种底环下的合成列及长度相同。分别取各分量的合成列，再在直和中依次添入各分量，便得到长度为各长度之和的 \(M\) 合成列，证明所需等式。
\end{solution}

\begin{exercise}\label{b3:ex:03-nonnoether-subring-local}
设 \(k\) 为域，令 \(A=k+xk[x,y]\)、\(I=xk[x,y]\)。证明局部环 \(A_I\) 仍不是 Noether 环。
\end{exercise}
\begin{solution}
\(A/I=k\)，故 \(I\) 极大，\(A_I\) 的极大理想为 \(IA_I\)。由局部化正合性及乘积的分式公式，
\[
IA_I/(IA_I)^2\cong (I/I^2)_I.
\]
\(I/I^2\) 被 \(I\) 零化，且 \(A\setminus I\) 的作用只是非零常数的作用，已经可逆，所以右侧仍为以 \(xy^j+I^2\)、\(j\ge0\)，为基的无限维 \(k\) 向量空间。若 \(IA_I\) 有限生成，其模极大理想后的向量空间就有限维，矛盾。因此该局部环不是 Noether 环。
\end{solution}

\begin{exercise}\label{b3:ex:03-finite-gaussian-algebra}
证明 \(\mathbb Z[i]\) 是 Noether 环，并分解 \(\mathbb Z[i]/(5)\) 为 Artin 局部环的直积。
\end{exercise}
\begin{solution}
\(\mathbb Z[i]\) 由 \(1,i\) 作为 \(\mathbb Z\) 模生成，而整数环为 Noether 环，所以它是有限代数，从而是 Noether 环。模掉 \(5\) 后，
\[
\mathbb Z[i]/(5)\cong\mathbb F_5[T]/(T^2+1)
=\mathbb F_5[T]/\bigl((T-2)(T+2)\bigr).
\]
两个一次因子互素，由中国剩余定理得到 \(\mathbb F_5\times\mathbb F_5\)。这两个域就是局部因子，商环约化且长度为 \(2\)。
\end{solution}

\begin{exercise}\label{b3:ex:03-pointwise-noether}
设 \(k\) 为域，令 \(R=\prod_{n\ge1}k\)。对 \(f=(f_n)\in R\)，记 \(E=\{n\ge1:f_n\ne0\}\)。计算 \(R_f\)，并证明它是 Noether 环当且仅当 \(E\) 有限（允许 \(E\) 为空）。由此证明，不存在由满足 \(R_{f_i}\) 为 Noether 环的基本开集 \(D(f_i)\) 组成的有限覆盖。解释这与正文中对任意素理想 \(\mathfrak p\) 所证的 \(R_{\mathfrak p}\) 都是域为何相容。
\end{exercise}
\begin{solution}
向 \(E\) 坐标的投影把 \(f\) 送到可逆元素，逆逐坐标为 \(f_n^{-1}\)，故诱导同态
\[
R_f\longrightarrow\prod_{n\in E}k,\qquad
a/f^r\longmapsto(a_n/f_n^r)_{n\in E}.
\]
它满射，因为任意 \(E\) 坐标序列都能在其余坐标补零，成为 \(R\) 中的元素。若其值为零，\(a\) 的所有 \(E\) 坐标均为零，从而 \(fa=0\)，分式零判据给出 \(a/f^r=0\)。因此该映射是同构。\(E=\varnothing\) 时 \(f=0\)，两侧都是零环。

\(E\) 有限时，右侧是有限个域的直积，为 Artin 环，因而 Noether；\(E\) 无限时，选两两不同的坐标 \(n_1,n_2,\ldots\in E\)，由前 \(j\) 个单坐标幂等元生成的理想组成严格升链，故不是 Noether 环。

若 \(D(f_1),\ldots,D(f_s)\) 覆盖全谱，则 \((f_1,\ldots,f_s)=R\)。在每个坐标 \(n\) 上，至少一个 \(f_i\) 的值必须非零，否则这几个元素的任何线性组合在该坐标都为零，不可能等于 \(1\)。因此各 \(E_i\) 的并必须是全部正整数，有限个有限集不能满足此条件。于是 Noether 的基本开局部化不能给出这样的有限覆盖。若每个点都有这样的 Noether 基本开邻域，全谱的拟紧性就会给出一个有限子覆盖，与上面的结论矛盾。正文的域结论对每个素理想都成立，却不能保证各点有 Noether 的基本开邻域；这正是\cref{b3:prop:finite-presentation-open-cover}所需的额外条件。
\end{solution}

\begin{exercise}\label{b3:ex:03-noether-artin-distinction}
分别给出下列例子并验证：Noether 而非 Artin 的环；Artin 而非 Noether 的模；有限生成但非有限展示的模。最后说明为什么不存在交换 Artin 而非 Noether 的环。
\end{exercise}
\begin{solution}
\(\mathbb Z\) 的每个理想为主理想，故是 Noether 环；理想链
\(2^n\mathbb Z\) 严格下降，故不是 Artin 环。
取素数 \(p\)，整数模 \(\mathbb Z[1/p]/\mathbb Z\) 的子模分类已在\subsecref{b3:subsec:0304:03}证明：每个真子模有限循环，所以降链稳定；但它含有阶依次为 \(p^n\) 的严格上升子群链，所以不是 Noether 模。
取域 \(k\)，令 \(R=k[x_1,x_2,\ldots]\)、\(I=(x_1,x_2,\ldots)\)，则 \(R/I\) 循环而有限生成；若它有限展示，\cref{b3:prop:finite-presentation-kernel}会迫使 \(I\) 有限生成，与变量数的论证矛盾。
最后一种环由\cref{b3:thm:artin-finite-length}排除：交换 Artin 环 \(R\) 作为自身的模具有有限长度，必满足升链条件。
\end{solution}

% !TeX root = ../../Book3.tex
% 纲要标题：## 第4章　关联素理想与准素分解
% 纲要位置：工作记录/计划纲要.md，第 2005 行。

\chapter{关联素理想与准素分解}
\label{b3:ch:04}
\BookChapterNumber{b3:ch:04}

\begin{chapterabstract}
本章用模中元素的零化子定义关联素理想，证明其存在性、有限性与局部化公式，并由此描述零因子。随后证明 Noether 环中的理想，以及这类环上有限生成模中的子模，均存在准素分解，并区分孤立准素分量与嵌入准素分量的唯一性。
\end{chapterabstract}

\begin{learninggoals}
\begin{itemize}
\item 由实际元素的零化子寻找关联素理想，说明存在性与有限性各需什么假设。
\item 用素理想滤过和局部化研究有限生成模，以关联素理想的并描述零因子，并辨认支集中的极小点。
\item 证明准素分解存在性，准确表述两类唯一性，给出素理想幂不准素及嵌入准素分量不唯一的例子。
\item 对理想、子模与矩阵余核作准素分析，区分非约化结构与嵌入关联素理想。
\end{itemize}
\end{learninggoals}

设 \(k\) 为域。在 \(R=k[x,y]/(xy)\) 中，\(\bar x\bar y=0\)，因而两个非零元素都是零因子。这里可以直接从两条分支理解：\(\bar x\) 在 \(y\) 轴上消失，\(\bar y\) 在 \(x\) 轴上消失。但取 \(I=(x^2,xy)\)、\(A=k[x,y]/I\)，有 \(\sqrt I=(x)\)，所以 \(A\) 作为 \(k[x,y]\) 模的支集只是 \(V(x)\)。同时，\(A\) 中的非零元 \(\bar x\) 满足 \(x\bar x=y\bar x=0\)，其零化子是 \((x,y)\)。约化商 \(A/(\bar x)\cong k[y]\) 有同一个底闭集，却没有这份原点处的额外零化信息；支集本身不能辨认这种区别。

\ProseParagraphBreak

要找出支集没有区分的信息，应当反过来问：哪些标量杀死了给定的非零模元素？得到的零化子中若有素理想，便形成关联素理想。它们不仅标出不可约分支的一般点，也可能检测到分支内部特殊点处的额外模结构。准素分解把这些素理想及其相应的商模结构落实为理想的交表示；随后还须分清哪些素理想由原理想确定，哪些分量会随分解的选择改变。

\ProseParagraphBreak

本章使用\chapref{b3:ch:01}的素谱、根理想及有限素理想回避，\chapref{b3:ch:02}的模局部化和支集，以及\chapref{b3:ch:03}的 Noether 模与有限性。理想商的基础运算可结合第一卷第10.3节复习。

\ProseParagraphBreak
除这条带有嵌入信息的直线外，环 \(k[x,y,z]/(xy-z^2)\) 还提供了另一项检验：即使环境是整环，素理想的普通幂也未必准素。它要求我们直接检查准素定义中的乘法条件，不能只从理想的根推测准素性。

% !TeX root = ../../Book3.tex
% 纲要标题：### 4.1 关联素理想
% 纲要位置：工作记录/计划纲要.md，第 2007 行。
% 纲要细目（写作范围）：
% * 元素零化子的素性；循环子模与素商模的嵌入、短正合列的非对称性
% * 关联素理想的存在和有限性；非典范素理想滤过、局部化中的统一清分母
% * 零因子作为关联素理想的并；支集极小点与孤立、嵌入关联素理想

\section{关联素理想}
\label{b3:sec:0401}

同一个元素可以被许多环元素零化，但它的全部零化子组成一个确定的理想。对交换含幺环 \(R\) 的模 \(M\) 及元素 \(m\in M\)，映射 \(a\mapsto am\) 给出 \(R\) 模同构 \(Rm\cong R/\operatorname{Ann}_R(m)\)。当这个零化子是素理想 \(\mathfrak p\) 时，模中便含有一个与 \(R/\mathfrak p\) 同构的循环子模。这给出了从模中提取素理想的方法。在 Noether 环上，每个非零模都有这样的素理想，有限生成模的关联素理想还只有有限个。前一个结论不要求模有限生成，而这些假设也不是每个具体模存在关联素理想的必要条件。

\subsection{元素零化子的素性}
\label{b3:subsec:0401:01}

\begin{definition}\label{b3:def:associated-prime}
设 \(R\) 为交换含幺环，\(M\) 为 \(R\) 模，\(m\in M\)。元素 \(m\) 的零化子为
\[
\operatorname{Ann}_R(m)=\{\,a\in R\mid am=0\,\}.
\]
若 \(R\) 的素理想 \(\mathfrak p\) 等于 \(M\) 中某个非零元素的零化子，则称 \(\mathfrak p\) 为 \(M\) 的\term[guanlian-sulixiang]{关联素理想}[associated prime]。所有这样的素理想组成集合 \(\operatorname{Ass}_R(M)\)，并约定 \(\operatorname{Ass}_R(0)=\varnothing\)。
\symidx{\(\operatorname{Ass}_R(M)\)：模的关联素理想集合}
\end{definition}

因此 \(\mathfrak p\in\operatorname{Ass}_R(M)\) 等价于存在 \(R\) 模单射 \(R/\mathfrak p\hookrightarrow M\)。定义说的是某个元素的零化子，不是整个模的零化子必须为素理想。

\begin{lemma}\label{b3:lem:maximal-element-annihilator-prime}
设 \(R\) 为交换含幺环，\(M\) 为 \(R\) 模，\(0\ne m\in M\)。若 \(\operatorname{Ann}_R(m)\) 在集合 \(\{\operatorname{Ann}_R(x):0\ne x\in M\}\) 中按包含关系极大，则它是 \(R\) 的素理想。
\end{lemma}
\begin{proof}
记 \(I=\operatorname{Ann}_R(m)\)。因 \(m\ne0\)，\(I\ne R\)。若 \(ab\in I\) 且 \(b\notin I\)，则 \(bm\ne0\)。乘以 \(b\) 后，原来零化 \(m\) 的元素仍零化 \(bm\)，所以 \(I\subseteq\operatorname{Ann}_R(bm)\)。由于 \(bm\) 仍是非零元素，极大性不允许这个零化子严格增大，只能有 \(\operatorname{Ann}_R(bm)=I\)。此时 \(a(bm)=0\) 说明 \(a\) 原来就属于 \(I\)，这正是素性。
\end{proof}

这里的“极大”指元素零化子这一族中的极大元，并不表示 \(I\) 是环的极大理想。例如整环 \(R\) 作为自身上的模时，每个非零元素的零化子都是 \((0)\)，于是 \(\operatorname{Ass}_R(R)=\bigl\{(0)\bigr\}\)，即使 \(R\) 不是域也成立。

\begin{proposition}\label{b3:prop:associated-exact-sequence}
设 \(R\) 为交换含幺环，\(0\longrightarrow N\longrightarrow M\longrightarrow P\longrightarrow0\) 为 \(R\) 模的短正合列。则
\[
\operatorname{Ass}_R(N)\subseteq\operatorname{Ass}_R(M)
 \subseteq\operatorname{Ass}_R(N)\cup\operatorname{Ass}_R(P).
\]
特别地，有限直和的关联素理想集合是各因子对应集合的并。对一般短正合列，却不能保证 \(\operatorname{Ass}_R(P)\subseteq\operatorname{Ass}_R(M)\)。
\end{proposition}
\begin{proof}
子模中的元素在 \(N\) 与 \(M\) 内有同一个零化子，得第一个包含。给定 \(\mathfrak p\in\operatorname{Ass}_R(M)\)，在 \(M\) 内取 \(L\cong R/\mathfrak p\)。\(L\) 中每个非零元素的零化子都是 \(\mathfrak p\)，因为 \(R/\mathfrak p\) 是整环。如果 \(L\cap N\ne0\)，取交内非零元素即得 \(\mathfrak p\in\operatorname{Ass}_R(N)\)；如果交为零，\(L\) 单射到 \(P\)，故 \(\mathfrak p\in\operatorname{Ass}_R(P)\)。直和结论由分裂正合列和两因子各自的嵌入得到，再对因子数归纳。

商模可以产生原模没有的关联素理想。例如，设 \(k\) 为域，\(R=k[t]\)，考虑短正合列
\[
0\longrightarrow tR\longrightarrow R\longrightarrow R/(t)\longrightarrow0.
\]
环 \(R\) 是整环，故 \(\operatorname{Ass}_R(R)=\{(0)\}\)；而 \(R/(t)\cong k\) 的每个非零元素的零化子均为 \((t)\)，所以 \(\operatorname{Ass}_R(R/(t))=\{(t)\}\)。因此商模的关联素理想不必属于中间模的关联素理想集合。
\end{proof}

\subsection{关联素理想的存在性与有限性}
\label{b3:subsec:0401:02}

\begin{theorem}\label{b3:thm:associated-existence}
若 \(R\) 是交换含幺 Noether 环，\(M\ne0\) 是 \(R\) 模，则 \(\operatorname{Ass}_R(M)\ne\varnothing\)。这里不要求 \(M\) 有限生成。
\end{theorem}
\begin{proof}
所有非零元素零化子构成非空理想族。Noether 升链条件保证其中有极大元，再用\cref{b3:lem:maximal-element-annihilator-prime}。
\end{proof}

Noether 环上的有限生成模可以逐层拆成素商模，再由短正合列的包含关系控制它的关联素理想。每一步的素商子模需要作出选择，所得滤过通常不唯一，其中出现的素理想只给关联素理想集合一个上界。

\begin{theorem}\label{b3:thm:associated-finite}
设 \(R\) 是交换含幺 Noether 环，\(M\) 是有限生成 \(R\) 模。存在有限子模链
\[
0=M_0\subsetneq M_1\subsetneq\cdots\subsetneq M_r=M,
\qquad M_i/M_{i-1}\cong R/\mathfrak p_i,
\]
其中各 \(\mathfrak p_i\) 为素理想；这种链称为\term[sulixiang-lvguo]{素理想滤过}[prime filtration]。并且
\[
\operatorname{Ass}_R(M)\subseteq\{\mathfrak p_1,\ldots,\mathfrak p_r\},
\]
所以关联素理想只有有限个。\(M=0\) 时采用空滤过。
\end{theorem}
\begin{proof}
若已构造 \(M_i\ne M\)，对非零商模 \(M/M_i\) 应用\cref{b3:thm:associated-existence}，取其中同构于 \(R/\mathfrak p_{i+1}\) 的子模，并令 \(M_{i+1}\) 为其原像。由\cref{b3:prop:noether-finite-modules}，有限生成模 \(M\) 是 Noether 模，故这个严格上升过程必在有限步停止。对短正合列
\(0\rightarrow M_{i-1}\rightarrow M_i\rightarrow R/\mathfrak p_i\rightarrow0\)
逐次应用\cref{b3:prop:associated-exact-sequence}，再用 \(\operatorname{Ass}_R(R/\mathfrak p_i)=\{\mathfrak p_i\}\)，得到所需包含。
\end{proof}

例如 \(\mathbb Z\supset2\mathbb Z\supset0\) 的两个因子为 \(\mathbb Z/2\mathbb Z\) 与 \(\mathbb Z\)，但 \(\operatorname{Ass}_{\mathbb Z}(\mathbb Z)=\bigl\{(0)\bigr\}\)，滤过中的 \((2)\) 是额外出现的素理想。此外，有限生成性不能删除：\(\bigoplus_{p\;\text{为素数}}\mathbb Z/p\mathbb Z\) 的关联素理想包括所有 \((p)\)，因而是无限集。

\begin{proposition}\label{b3:prop:associated-localization}
设 \(R\) 为交换含幺 Noether 环，\(S\subseteq R\) 为含一的乘法集，\(M\) 为任意 \(R\) 模。则
\[
\operatorname{Ass}_{S^{-1}R}(S^{-1}M)
=\bigl\{\,S^{-1}\mathfrak p\mid \mathfrak p\in\operatorname{Ass}_R(M),\quad
                       \mathfrak p\cap S=\varnothing\,\bigr\}.
\]
\end{proposition}
\begin{proof}
若 \(\operatorname{Ann}_R(m)=\mathfrak p\) 且 \(\mathfrak p\cap S=\varnothing\)，则 \(m/1\ne0\)。当 \((a/s)(m/1)=0\) 时，有 \(t\in S\) 使 \(tam=0\)；由 \(ta\in\mathfrak p\)、\(t\notin\mathfrak p\) 得 \(a\in\mathfrak p\)。故其零化子恰为 \(S^{-1}\mathfrak p\)。

反之，设 \(\mathfrak q\) 是非零元素 \(m/s\) 的素零化子，令 \(\mathfrak p\) 为 \(\mathfrak q\) 在 \(R\) 中的收缩。由\cref{b3:thm:spec-localization}，\(\mathfrak p\) 是避开 \(S\) 的素理想，且 \(\mathfrak q=S^{-1}\mathfrak p\)。以单位 \(s/1\) 乘此元素后可改用 \(m/1\)。要在原模中找到零化子恰为 \(\mathfrak p\) 的元素，必须把局部成立的零化关系同时变成原模中的等式；这一步需要共同的分母。由于 \(R\) 是 Noether 环，理想 \(\mathfrak p\) 有有限生成集，写成 \(\mathfrak p=(a_1,\ldots,a_l)\)。因 \((a_i m)/1=0\)，可选 \(s_i\in S\) 使 \(s_i a_i m=0\)。令 \(t=s_1\cdots s_l\)，此时 \(\mathfrak p\subseteq\operatorname{Ann}_R(tm)\)；当 \(\mathfrak p=(0)\) 时采用空生成集并令 \(t=1\)。若 \(a(tm)=0\)，由于 \(t/1\) 可逆，局部化后 \(a/1\) 便零化 \(m/1\)，所以 \(a\in\mathfrak p\)。于是 \(\operatorname{Ann}_R(tm)=\mathfrak p\)，且 \(tm\ne0\)，否则 \(m/1=0\)。这就找到了原模中的关联元素，而所用的有限生成性来自理想 \(\mathfrak p\)，并未要求 \(M\) 有限生成。
\end{proof}

例如，设 \(k\) 为域，\(R=k[x,y]\)、\(M=R/(x)\)、\(\mathfrak m=(x,y)\)。因 \(R/(x)\cong k[y]\) 是整环，每个非零模元素的零化子都是 \((x)\)，故 \(\operatorname{Ass}_R(M)=\{(x)\}\)。局部化后有
\[
\begin{aligned}
M_{\mathfrak m}&\cong R_{\mathfrak m}/(x)R_{\mathfrak m}
\cong k[y]_{(y)},\\
\operatorname{Ass}_{R_{\mathfrak m}}(M_{\mathfrak m})
&=\{(x)R_{\mathfrak m}\}.
\end{aligned}
\]
这个唯一的关联素理想严格小于局部环的极大理想 \(\mathfrak mR_{\mathfrak m}\)：元素 \(y/1\) 属于后者，在商环 \(k[y]_{(y)}\) 中的像却非零，所以不属于前者。

\begin{proposition}\label{b3:prop:minimal-primes-associated}
设 \(R\) 为交换含幺 Noether 环，\(M\) 为非零有限生成 \(R\) 模。则 \(\operatorname{Supp}_R(M)\) 的每个素理想都包含某个关联素理想，且
\[
\min\operatorname{Supp}_R(M)=\min\operatorname{Ass}_R(M).
\]
尤其，对任意真理想 \(I\)，包含 \(I\) 的极小素理想都是 \(R/I\) 的关联素理想。
\end{proposition}
\begin{proof}
若 \(\mathfrak p\in\operatorname{Ass}_R(M)\)，上述局部化命题保证 \(M_{\mathfrak p}\ne0\)，故 \(\operatorname{Ass}_R(M)\subseteq\operatorname{Supp}_R(M)\)。给定 \(\mathfrak q\in\operatorname{Supp}_R(M)\)，\(R_{\mathfrak q}\) 是 Noether 环，非零模 \(M_{\mathfrak q}\) 有关联素理想；由\cref{b3:prop:associated-localization}，它来自某个 \(\mathfrak p\in\operatorname{Ass}_R(M)\)，且 \(\mathfrak p\subseteq\mathfrak q\)。因此两集合的极小元一致。对有限生成模，\cref{b3:thm:support-finite}给出 \(\operatorname{Supp}_R(M)=V\bigl(\operatorname{Ann}_R(M)\bigr)\)。取 \(M=R/I\) 就得到最后一句。
\end{proof}

设 \(k\) 为域，取 \(R=k[x,y]\)、\(M=R/(x^2,xy)\)。由 \(\sqrt{(x^2,xy)}=(x)\)，其支集为 \(V(x)\)，极小元只有 \((x)\)，所以 \((x)\) 也是极小关联素理想。另一方面，非零元素 \(\bar x\) 的零化子为 \((x,y)\)，具体计算见\subsecref{b3:subsec:0404:02}。因此 \((x,y)\) 也是关联素理想，却严格包含 \((x)\)。

\subsection{零因子的关联素理想描述}
\label{b3:subsec:0401:03}

设 \(R\) 为交换含幺环，\(M\) 为 \(R\) 模。若存在 \(0\ne m\in M\) 使 \(am=0\)，则称 \(a\in R\) 是 \(M\) 上的\term[mo-shang-lingyinzi]{零因子}[zero divisor on a module]；否则称它在 \(M\) 上是正则的。记这些零因子的集合为 \(Z_R(M)\)。这里允许 \(a=0\)：对非零模，零当然属于这个集合；对零模，集合为空。

\ProseParagraphBreak
标量 \(a\) 的乘法映射 \(M\to M\)、\(m\mapsto am\) 单射，当且仅当 \(a\notin Z_R(M)\)。因此零因子集合决定哪些标量在模上的作用失去单射性；在 Noether 环上，关联素理想恰好给出这个集合的描述。

\begin{theorem}\label{b3:thm:associated-zero-divisors}
设 \(R\) 为交换含幺 Noether 环，\(M\) 为 \(R\) 模。则
\[
Z_R(M)=\bigcup_{\mathfrak p\in\operatorname{Ass}_R(M)}\mathfrak p.
\]
若 \(M\) 有限生成，这就是有限个素理想的并。
\end{theorem}
\begin{proof}
关联素理想中的元素会零化定义它的非零元素，因此右端包含于左端。反向设 \(am=0\)、\(m\ne0\)。包含 \(\operatorname{Ann}_R(m)\) 的非零元素零化子构成非空理想族；由于 \(R\) 是 Noether 环，其中存在极大元 \(\operatorname{Ann}_R(n)\)。它在全部非零元素零化子中也是极大的，因任何更大零化子仍包含 \(\operatorname{Ann}_R(m)\)。\cref{b3:lem:maximal-element-annihilator-prime}说明它是关联素理想，而 \(a\) 在其中。
\end{proof}

例如，设 \(k\) 为域，在 \(R=k[x,y]/(xy)\) 中，\(\bar x\) 与 \(\bar y\) 的零化子分别为 \((\bar y)\) 和 \((\bar x)\)。每个元素都可唯一写为常数加一个无常数项的 \(x\) 多项式，再加一个无常数项的 \(y\) 多项式。限制到两轴给出单射 \(R\to k[x]\times k[y]\)；不在 \((\bar x)\cup(\bar y)\) 中的元素在两个坐标上都非零，故乘积为零时，两个整环中的消去律迫使另一乘子为零。而 \((\bar x)\)、\((\bar y)\) 中的元素分别被 \(\bar y\)、\(\bar x\) 零化，所以零因子恰为这两个素理想的并。\(\bar x+\bar y\) 不在其中，因而不是零因子，尽管它的两个加项都是零因子。这说明零因子集合不一定对加法封闭，因而不一定是理想。

\ProseParagraphBreak
更一般地，设 \(R\) 为交换含幺 Noether 环，若 \(R\) 的理想 \(J\) 的每个元素都是有限生成 \(R\) 模 \(M\) 上的零因子，由\cref{b3:thm:associated-finite}和\cref{b3:thm:associated-zero-divisors}及\cref{b3:lem:finite-prime-avoidance}的有限素理想回避性质，\(J\) 必包含于某一个关联素理想。这把“逐个检查乘法是否单射”的问题转成了有限个素理想的包含问题。以上存在性、局部化和滤过方法可对照 \cite[\S6, Theorems 6.1--6.5]{Matsumura1989CommutativeRingTheory}。

\ProseParagraphBreak
设 \(R\) 为交换含幺 Noether 环，\(M\) 为有限生成 \(R\) 模。按包含关系极小的关联素理想称为\term[guli-guanlian-sulixiang]{孤立关联素理想}[isolated associated prime]，其余称为\term[qianru-guanlian-sulixiang]{嵌入关联素理想}[embedded associated prime]。这里的“孤立”指素理想包含关系中的极小性，不要求它是 Zariski 拓扑中的孤立点，即不要求单点集合为开集。由\cref{b3:prop:minimal-primes-associated}，前者就是支集的极小元，后者则严格包含某个极小关联素理想；两类都由 \(M\) 本身确定。\cref{b3:tab:support-associated-primes}汇集了这些关系。

\begin{table}[htbp]
\centering
\caption[支集与关联素理想]{支集与关联素理想。设 \(R\) 为交换含幺 Noether 环，\(M\) 为有限生成 \(R\) 模；极小元均按素理想的包含关系理解。}
\label{b3:tab:support-associated-primes}
\begin{tabularx}{\textwidth}{@{}>{\raggedright\arraybackslash}p{0.25\textwidth}>{\raggedright\arraybackslash}X@{}}
\toprule
集合 & 与支集及零化子的关系 \\
\midrule
支集 & \(\operatorname{Supp}_R(M)=V\bigl(\operatorname{Ann}_R(M)\bigr)\)；每个支集中的素理想都包含某个关联素理想 \\
关联素理想 & \(\operatorname{Ass}_R(M)\subseteq\operatorname{Supp}_R(M)\)，且该集合有限 \\
孤立关联素理想 & \(\min\operatorname{Ass}_R(M)=\min\operatorname{Supp}_R(M)\)，即支集的极小点 \\
嵌入关联素理想 & \(\operatorname{Ass}_R(M)\setminus\min\operatorname{Supp}_R(M)\)，即支集中非极小的关联点 \\
\bottomrule
\end{tabularx}
\end{table}

% !TeX root = ../../Book3.tex
% 纲要标题：### 4.2 准素理想
% 纲要位置：工作记录/计划纲要.md，第 1913 行。
% 纲要细目（写作范围）：
% * 准素性与根
% * 素幂不必准素的注意事项
% * 局部化中的准素理想

\section{准素理想}
\label{b3:sec:0402}

素理想在商环中排除一切非零零因子，准素理想则允许零因子存在，但要求它们最终被某个幂消去。这个差别使准素理想既保留一个确定的素支集，又能描述支集上的幂零结构。先把准素性直接写成乘法条件，再研究它与局部化的相容性。

\subsection{准素性与根理想}
\label{b3:subsec:0402:01}

\begin{definition}\label{b3:def:primary-ideal}
设 \(R\) 为交换含幺环，\(Q\subsetneq R\) 为真理想。若对任意 \(a,b\in R\) 都有
\[
ab\in Q,\quad b\notin Q
\quad\Longrightarrow\quad a^n\in Q\ \text{对某个整数 }n\ge1.
\]
则称 \(Q\) 为\term[zhunsu-lixiang]{准素理想}[primary ideal]。若 \(\sqrt Q=\mathfrak p\)，也称 \(Q\) 是 \(\mathfrak p\) 准素理想。
\end{definition}

等价地，\(R/Q\) 中每个零因子都是幂零元。素理想满足更强的条件 \(n=1\)，因而总是准素理想；反之不成立，例如 \(\mathbb Z/(p^e)\) 中每个非单位都是 \(p\) 的倍数，故 \((p^e)\) 为 \((p)\) 准素理想。

\begin{proposition}\label{b3:prop:primary-radical-colon}
设 \(R\) 为交换含幺环，\(Q\) 为 \(R\) 的准素理想。则其根 \(\mathfrak p=\sqrt Q\) 是素理想。对 \(c\notin Q\)，理想商 \((Q:c)\) 仍为 \(\mathfrak p\) 准素理想；若还满足 \(c\notin\mathfrak p\)，则 \((Q:c)=Q\)。同一素理想所属的有限个准素理想的交仍为该素理想所属的准素理想。
\end{proposition}
\begin{proof}
若 \(ab\in\sqrt Q\)、\(b\notin\sqrt Q\)，则某个 \(a^m b^m\in Q\)，且 \(b^m\notin Q\)。准素性给出 \(a^{mn}\in Q\)，从而 \(a\in\sqrt Q\)。根是真理想，故为素理想。

若 \(ab\in(Q:c)\)、\(b\notin(Q:c)\)，即 \(a(bc)\in Q\)、\(bc\notin Q\)，则某个 \(a^n\in Q\subseteq(Q:c)\)。又 \(Q\subseteq(Q:c)\)，而 \(dc\in Q\)、\(c\notin Q\) 蕴涵 \(d\in\mathfrak p\)，所以两者根相同。若 \(c\notin\mathfrak p\)，从 \(dc\in Q\) 与 \(d\notin Q\) 将推出 \(c\in\mathfrak p\)，故 \(d\in Q\)。

最后令 \(Q=\bigcap_{i=1}^rQ_i\)，各 \(Q_i\) 的根均为 \(\mathfrak p\)。若 \(ab\in Q\)、\(b\notin Q\)，则至少一个因子给出 \(a\in\mathfrak p\)。于是每个 \(Q_i\) 都包含 \(a\) 的某个幂，取这些指数的最大值即得 \(a^n\in Q\)；同时 \(\sqrt Q=\bigcap_i\sqrt{Q_i}=\mathfrak p\)。
\end{proof}

\begin{proposition}\label{b3:prop:maximal-radical-primary}
设 \(R\) 为交换含幺环，\(Q\) 为 \(R\) 的真理想。若其根是极大理想 \(\mathfrak m\)，则 \(Q\) 是 \(\mathfrak m\) 准素理想。因此极大理想的任意正整数次幂均为准素理想。
\end{proposition}
\begin{proof}
\(R/Q\) 只有一个极大理想 \(\mathfrak m/Q\)。若 \(a\notin\mathfrak m\)，其像为单位，故 \(ab\in Q\) 必导致 \(b\in Q\)。取逆否命题即得到准素条件。又 \(\sqrt{\mathfrak m^n}=\mathfrak m\)，得最后一句。
\end{proof}

\begin{proposition}\label{b3:prop:primary-associated-singleton}
设 \(R\) 为交换含幺 Noether 环，\(Q\subsetneq R\) 为真理想。则 \(Q\) 是准素理想，当且仅当 \(\operatorname{Ass}_R(R/Q)\) 仅有一个元素；这唯一的元素就是 \(\sqrt Q\)。
\end{proposition}
\begin{proof}
若 \(Q\) 为准素理想，对每个非零 \(\bar b\in R/Q\)，\(\operatorname{Ann}_R(\bar b)=(Q:b)\) 的根由\cref{b3:prop:primary-radical-colon}等于 \(\sqrt Q\)。因此当零化子本身为素理想时，它只能是 \(\sqrt Q\)；存在性保证这样的零化子确实出现。

反之，设关联素理想只有 \(\mathfrak p\)。\cref{b3:prop:minimal-primes-associated}说明 \(Q\) 的极小素理想只有 \(\mathfrak p\)，故 \(\sqrt Q=\mathfrak p\)。由\cref{b3:thm:associated-zero-divisors}，商环中的每个零因子来自 \(\mathfrak p\)，所以它都是幂零元，\(Q\) 准素。
\end{proof}

\subsection{素理想幂不为准素理想的例子}
\label{b3:subsec:0402:02}

极大理想幂的结论不能直接推广到任意素理想。以下例子甚至发生在一个有限型整环中。

\begin{example}\label{b3:exm:prime-square-not-primary}
设 \(k\) 为任意域，
\[
R=k[x,y,z]/(xy-z^2),\qquad \mathfrak p=(\bar x,\bar z).
\]
则 \(R\) 是整环，\(\mathfrak p\) 是素理想，而 \(\mathfrak p^2\) 不是准素理想。
\end{example}
\begin{proof}
映射 \(\bar x\mapsto s^2\)、\(\bar y\mapsto t^2\)、\(\bar z\mapsto st\) 将 \(R\) 嵌入 \(k[s,t]\)。具体地，利用 \(z^2=xy\)，每个类可写为 \(a(x,y)+z\cdot b(x,y)\)；代入后的第一项只含偶数指数对，第二项只含奇数指数对，两类单项式不能抵消，因此映射单射，\(R\) 是整环。又 \(R/\mathfrak p\cong k[y]\)，故 \(\mathfrak p\) 素。

关系 \(xy-z^2\) 是全次数二的齐次式，因而商环仍按全次数分解，互不相同的次数不能相互抵消。理想 \(\mathfrak p^2=(\bar x^2,\bar x\bar z,\bar z^2)\) 由全次数二的齐次元生成，所以非零的一次齐次元 \(\bar x\) 不属于它。然而
\[
\bar y\bar x=\bar z^2\in\mathfrak p^2,
\qquad \bar y^n\notin\mathfrak p\supseteq\mathfrak p^2\quad(n\ge1),
\]
后一句由商环 \(k[y]\) 核验。这直接违反准素条件。
\end{proof}

失效的原因是 \(\bar y\) 虽不在素理想 \(\mathfrak p\) 中，却在 \(R/\mathfrak p^2\) 中成为零因子。只知道 \(\sqrt{\mathfrak p^2}=\mathfrak p\) 并不足以推出准素性；上一小节中的“根是极大理想”确有作用。

\subsection{局部化中的准素理想}
\label{b3:subsec:0402:03}

\begin{proposition}\label{b3:prop:primary-localization}
设 \(R\) 为交换含幺环，\(Q\) 为 \(R\) 的 \(\mathfrak p\) 准素理想，\(S\subseteq R\) 为含一的乘法集。若 \(S\cap\mathfrak p\ne\varnothing\)，则 \(S^{-1}Q=S^{-1}R\)；若 \(S\cap\mathfrak p=\varnothing\)，则 \(S^{-1}Q\) 为 \(S^{-1}\mathfrak p\) 准素理想，且其在 \(R\) 内的收缩是 \(Q\)。反之，\(S^{-1}R\) 中任一准素理想的收缩也是准素理想。
\end{proposition}
\begin{proof}
第一种情形取 \(s\in S\cap\mathfrak p\)，则某个 \(s^n\in Q\)，局部化后 \(Q\) 包含单位。第二种情形中，若 \(a/1\in S^{-1}Q\)，则某个 \(s\in S\) 满足 \(sa\in Q\)；因为 \(s\notin\mathfrak p\)，\cref{b3:prop:primary-radical-colon}给出 \(a\in Q\)，这证明收缩断言，也说明扩张仍是真理想。

若 \((a/s)(b/t)\in S^{-1}Q\)、\(b/t\notin S^{-1}Q\)，收缩断言给出 \(ab\in Q\)、\(b\notin Q\)。于是 \(a^n\in Q\)，进而 \((a/s)^n\in S^{-1}Q\)。又由已证的收缩性质，\((a/s)^n\in S^{-1}Q\) 当且仅当 \(a^n\in Q\)。某个这样的正整数 \(n\) 存在，恰等价于 \(a\in\sqrt Q=\mathfrak p\)，所以局部化后的根为 \(S^{-1}\mathfrak p\)。最后，对于收缩的理想，把 \(ab\)、\(b\) 送入局部化环并应用准素条件，所获 \((a/1)^n\) 属于该理想就意味着 \(a^n\) 属于收缩。
\end{proof}

\begin{definition}\label{b3:def:symbolic-power}
设 \(R\) 为交换含幺环，\(\mathfrak p\) 为 \(R\) 的素理想，\(n\ge1\) 为整数。定义
\[
\mathfrak p^{(n)}\coloneqq\mathfrak p^nR_{\mathfrak p}\cap R
\]
并称它为 \(\mathfrak p\) 的第 \(n\) 个\term[fuhao-mi]{符号幂}[symbolic power]；交表示沿 \(R\rightarrow R_{\mathfrak p}\) 的收缩，不要求该映射单射。
\symidx{\(\mathfrak p^{(n)}\)：素理想的符号幂}
\end{definition}

\(\mathfrak pR_{\mathfrak p}\) 是极大理想，所以由\cref{b3:prop:maximal-radical-primary}和\cref{b3:prop:primary-localization}，\(\mathfrak p^{(n)}\) 总是 \(\mathfrak p\) 准素理想，且包含普通幂 \(\mathfrak p^n\)。在\cref{b3:exm:prime-square-not-primary}中，\(\bar y\) 在 \(R_{\mathfrak p}\) 中可逆，故 \(\bar x=\bar z^2/\bar y\in\mathfrak p^2R_{\mathfrak p}\)，于是 \(\bar x\in\mathfrak p^{(2)}\setminus\mathfrak p^2\)。符号幂正是先在目标素理想处集中研究，再把结果收回原环所得的准素对象。

% !TeX root = ../../Book3.tex
% 纲要标题：### 4.3 准素分解定理
% 纲要位置：工作记录/计划纲要.md，第 1919 行。
% 纲要细目（写作范围）：
% * Noether 环中的存在性
% * 极小分解、孤立分量和嵌入分量
% * 哪些部分唯一、哪些部分不唯一

\section{准素分解定理}
\label{b3:sec:0403}

整数的素因数分解给出 \((12)=(4)\cap(3)\)。一般 Noether 环的理想也能分解为有限个准素理想的交，但“各因子都唯一”不再成立。极小分解中的根素理想集合由原理想决定；其中对应极小素理想的准素分量也唯一，嵌入准素分量却可能改变。

\subsection{Noether 环中的准素分解存在性}
\label{b3:subsec:0403:01}

设 \(R\) 为交换含幺环，\(I\subsetneq R\) 为真理想。若对 \(R\) 的任意理想 \(J,K\)，等式 \(I=J\cap K\) 总使 \(I=J\) 或 \(I=K\)，则称 \(I\) 为\term[buke-yue-lixiang]{不可约理想}[irreducible ideal]。这里讨论的是理想作为交的不可约性，不是元素的不可约性，也不是模的单性。

\begin{lemma}\label{b3:lem:irreducible-ideal-primary}
设 \(R\) 为交换含幺 Noether 环。则 \(R\) 的每个真理想是有限个不可约理想的交，而每个不可约理想都是准素理想。
\end{lemma}
\begin{proof}
若有真理想不能写成有限个不可约理想的交，在这些理想中取极大者 \(I\)。它不可能不可约，故 \(I=J\cap K\)，其中 \(J,K\) 都严格包含 \(I\)。极大性保证两者都已有这样的有限分解；将两个分解合并就得到 \(I\) 的分解，矛盾。

设 \(Q\) 不可约且 \(ab\in Q\)、\(b\notin Q\)。升链
\[
(Q:a)\subseteq(Q:a^2)\subseteq\cdots
\]
稳定；选 \(n\ge1\) 使 \((Q:a^n)=(Q:a^{n+1})\)。我们证明
\[
Q=\bigl(Q+(a^n)\bigr)\cap\bigl(Q+(b)\bigr).
\]
只需检查反向包含。若 \(w=q+ra^n=q'+sb\)，其中 \(q,q'\in Q\)，则乘以 \(a\) 后第二式给出 \(aw\in Q\)，于是第一式给出 \(ra^{n+1}\in Q\)。由理想商稳定性，\(ra^n\in Q\)，故 \(w\in Q\)。因为 \(b\notin Q\)，不可约性迫使 \(Q+(a^n)=Q\)，即 \(a^n\in Q\)。这正是准素条件。
\end{proof}

\begin{theorem}[准素分解存在性]\label{b3:thm:primary-decomposition}
设 \(R\) 为交换含幺 Noether 环，\(I\subsetneq R\) 为真理想。则 \(I\) 能写成
\begin{equation}\label{b3:eq:primary-decomposition}
I=Q_1\cap\cdots\cap Q_r,
\end{equation}
其中各 \(Q_i\) 为准素理想。可以进一步要求没有任何 \(Q_i\) 能删去，且各根 \(\mathfrak p_i=\sqrt{Q_i}\) 两两不同。
\end{theorem}
\begin{proof}
\cref{b3:lem:irreducible-ideal-primary}给出一个有限准素分解。把根相同的因子合并为一个交，由\cref{b3:prop:primary-radical-colon}它仍为同根准素理想。再删去冗余因子，即得附加条件。
\end{proof}

单项式理想的成员关系可以逐项判断。设 \(k\) 为域，在 \(k[x_1,\ldots,x_n]\) 中，由单项式 \(m_1,\ldots,m_s\) 生成的理想，作为 \(k\) 向量空间，恰由那些被某个 \(m_i\) 整除的单项式生成。事实上，把各多项式系数展开，就得到这一方向的包含；反向，每个这样的单项式本来就是某个 \(m_i\) 的倍数。由于不同单项式线性无关，一个多项式属于此理想，当且仅当它的每个非零单项都属于此理想。因此两个单项式理想取交时可以逐项检查，不需要预先运行 Gröbner 基算法。

\ProseParagraphBreak
准素理想不一定不可约。例如，设 \(k\) 为域，在 \(k[x,y]\) 内有
\[
(x,y)^2=(x^2,y)\cap(x,y^2).
\]
等式可逐个比较单项式：同时属于右端的单项式恰被 \(x^2,xy,y^2\) 之一整除。两边的理想都是 \((x,y)\) 准素理想，右边两个因子却都严格大于左边。因此不可约分解只是证明存在性的一种工具，并不是最后要保存的标准形式。

\subsection{极小分解、孤立分量与嵌入分量}
\label{b3:subsec:0403:02}

设 \(R\) 为交换含幺环，\(I\subsetneq R\) 为真理想。若 \(I=Q_1\cap\cdots\cap Q_r\) 的各因子均为准素理想、各根 \(\sqrt{Q_i}\) 两两不同，且删去任一因子都会改变交，则称它为\term[jixiao-zhunsu-fenjie]{极小准素分解}[minimal primary decomposition]，其中各个 \(Q_i\) 称为\term[zhunsu-fenliang]{准素分量}[primary component]。“极小”不声称每个准素因子在包含关系下极小。

\begin{theorem}\label{b3:thm:primary-associated-primes}
设 \(R\) 为交换含幺 Noether 环，\(I\subsetneq R\) 为真理想。若 \(I=\bigcap_{i=1}^rQ_i\) 是 \(R\) 中的极小准素分解，其中 \(r\ge1\)，各 \(Q_i\) 为准素理想，\(\sqrt{Q_i}=\mathfrak p_i\)，则
\[
\operatorname{Ass}_R(R/I)=\{\mathfrak p_1,\ldots,\mathfrak p_r\}.
\]
特别地，分解中出现哪些素理想与分解的选择无关。
\end{theorem}
\begin{proof}
对角映射给出单射
\[
R/I\longrightarrow\bigoplus_{i=1}^rR/Q_i,\qquad
\bar a\longmapsto(a+Q_1,\ldots,a+Q_r).
\]
由\cref{b3:prop:associated-exact-sequence}和\cref{b3:prop:primary-associated-singleton}，左端的关联素理想都在 \(\{\mathfrak p_i\}\) 中。

固定 \(i\)，令 \(J_i=\bigcap_{j\ne i}Q_j\)，在 \(r=1\) 时约定 \(J_i=R\)。无冗余保证 \(J_i/I\ne0\)，而映射 \(J_i/I\rightarrow R/Q_i\) 单射，因其核为 \((J_i\cap Q_i)/I=0\)。由\cref{b3:thm:associated-existence}，这个非零子模有至少一个关联素理想；由于它嵌入 \(R/Q_i\)，该素理想只能是 \(\mathfrak p_i\)。它同时是 \(R/I\) 的子模，故 \(\mathfrak p_i\in\operatorname{Ass}_R(R/I)\)。逐个取 \(i\) 得反向包含。
\end{proof}

按照前面对关联素理想的命名，若 \(\mathfrak p_i\) 为孤立关联素理想，则称相应 \(Q_i\) 为孤立准素分量；若 \(\mathfrak p_i\) 为嵌入关联素理想，则称 \(Q_i\) 为嵌入准素分量。由\cref{b3:prop:minimal-primes-associated}，孤立关联素理想恰是包含 \(I\) 的极小素理想。它们对应 \(V(I)\) 的不可约分量，而嵌入素理想对应的闭集位于某个不可约分量之内，并不另给出一个不可约分量。

\ProseParagraphBreak
例如，设 \(k\) 为域，理想 \(I=(x^2,xy)\subset k[x,y]\) 满足
\begin{equation}\label{b3:eq:embedded-line-decomposition}
I=(x)\cap(x^2,y).
\end{equation}
第二因子因根为极大理想 \((x,y)\) 而准素。要验证交，设 \(xf=x^2a+yb\)，模掉 \(x\) 后有 \(y\bar b=0\) 于 \(k[y]\)，故 \(b=xc\)，于是 \(xf=x^2a+xyc\in I\)。两因子均不可删，根之间的包含关系为
\[
\underbrace{(x)}_{\text{极小关联素理想}}
\subsetneq
\underbrace{(x,y)}_{\text{嵌入关联素理想}}.
\]
因此 \((x)\) 是孤立准素分量，\((x^2,y)\) 是嵌入准素分量；\(V(I)=V(x)\) 只有一个不可约分量，\(V(x,y)\) 是其中的一个闭点。

\subsection{准素分解的唯一性与非唯一性}
\label{b3:subsec:0403:03}

\begin{theorem}\label{b3:thm:primary-isolated-uniqueness}
设 \(R\) 为交换含幺 Noether 环，\(I\subsetneq R\) 为真理想，\(I=\bigcap_{i=1}^rQ_i\) 为极小准素分解，其中 \(r\ge1\)、\(\mathfrak p_i=\sqrt{Q_i}\)。若 \(\mathfrak p_i\) 为 \(\operatorname{Ass}_R(R/I)\) 中的极小元，则
\[
Q_i=IR_{\mathfrak p_i}\cap R.
\]
因此每个孤立准素分量由 \(I\) 与其根唯一决定。
\end{theorem}
\begin{proof}
有限交与局部化交换。若 \(j\ne i\)，由于 \(\mathfrak p_i\) 在所有根中极小且根互异，\(\mathfrak p_j\not\subseteq\mathfrak p_i\)。取 \(a\in\mathfrak p_j\setminus\mathfrak p_i\)，则某个 \(a^n\in Q_j\)，而 \(a\) 在 \(R_{\mathfrak p_i}\) 中成为单位，所以 \(Q_jR_{\mathfrak p_i}=R_{\mathfrak p_i}\)。于是 \(IR_{\mathfrak p_i}=Q_iR_{\mathfrak p_i}\)。\cref{b3:prop:primary-localization}又说明 \(Q_i\) 是其扩张的收缩，得到公式。
\end{proof}

嵌入准素分量一般不唯一。在\cref{b3:eq:embedded-line-decomposition}的同一个理想上，对每个 \(n\ge1\) 都有
\begin{equation}\label{b3:eq:nonunique-embedded-components}
(x^2,xy)=(x)\cap(x^2,xy,y^n).
\end{equation}
右边第二因子的根为 \((x,y)\)，故它准素。由单项式整除判据，第二因子中同时被 \(x\) 整除的单项式恰属于 \((x^2,xy)\)，所以交确实不变。\(x\) 不属于第二因子，\(y^n\) 不属于第一因子，因而各分解都极小；而 \(y^n\) 属于第 \(n\) 个第二因子，却不属于第 \(n+1\) 个，故这些准素分量互不相同。

\ProseParagraphBreak
极小准素分解中出现的根素理想集合唯一，并等于 \(\operatorname{Ass}_R(R/I)\)；孤立准素分量唯一，嵌入准素分量一般不唯一。\cref{b3:tab:primary-uniqueness-data}区分了这些数据。模的对应结论及另一种存在性证明见下一节；文献中的统一表述可对照 \cite[\S6, Theorem 6.8]{Matsumura1989CommutativeRingTheory}。

\begin{table}[htbp]
\centering
\caption[极小准素分解的唯一性数据]{极小准素分解的唯一性数据。设 \(R\) 为交换含幺 Noether 环，\(I\subsetneq R\) 为真理想。}
\label{b3:tab:primary-uniqueness-data}
\begin{tabularx}{\textwidth}{@{}>{\raggedright\arraybackslash}X>{\raggedright\arraybackslash}p{.27\textwidth}@{}}
\toprule
数据 & 是否由 \(I\) 唯一确定\\
\midrule
极小准素分解中的根素理想集合 \(\operatorname{Ass}_R(R/I)\) & 是\\
孤立关联素理想集合 & 是\\
对应于给定孤立关联素理想的准素分量 & 是\\
嵌入关联素理想集合 & 是\\
对应于给定嵌入关联素理想的准素分量 & 一般不唯一\\
\bottomrule
\end{tabularx}
\end{table}

% !TeX root = ../../Book3.tex
% 纲要标题：### 4.4 子模的准素分解与嵌入分量
% 纲要位置：工作记录/计划纲要.md，第 2026 行。
% 纲要细目（写作范围）：
% * 子模的准素分解；商模上的统一幂零作用及有限生成假设的无限直和反例
% * 嵌入关联素理想的显式例子；实际元素零化子及矩阵余核的生成关系
% * 与几何分支及非约化结构的联系；区分多分支、沿分量非约化与特殊点处的嵌入结构
% ---

\section{子模的准素分解与嵌入分量}
\label{b3:sec:0404}

理想分解其实是自由模 \(R\) 中的子模分解。若要研究一个矩阵给出的余核，或研究有限生成模的支集与局部结构，就需要允许一般有限生成模，并把准素条件写成商模上标量作用的性质。

\subsection{子模的准素分解}
\label{b3:subsec:0404:01}

\begin{definition}\label{b3:def:primary-submodule}
设 \(R\) 为交换含幺环，\(M\) 为 \(R\) 模，\(N\subsetneq M\) 为真子模。若对任意 \(a\in R\)、\(m\in M\) 都有
\[
am\in N,\quad m\notin N
\quad\Longrightarrow\quad a^nM\subseteq N
\ \text{对某个}\;n\ge1.
\]
则称 \(N\) 为\term[zhunsu-zimo]{准素子模}[primary submodule]。若 \(N\) 准素，且 \(\sqrt{\operatorname{Ann}_R(M/N)}\) 等于素理想 \(\mathfrak p\)，则称 \(N\) 为 \(\mathfrak p\) 准素子模。整个条件仅依赖商模 \(M/N\)，但“\(N\) 准素”总须指明它位于哪个模中。
\end{definition}

令 \(T=M/N\)。条件 \(am\in N\)、\(m\notin N\) 表示 \(a\) 零化了 \(T\) 中的某个非零元素，因而是 \(T\) 上的零因子；结论 \(a^nM\subseteq N\) 则表示 \(a^nT=0\)。所以准素性要求每个零因子的某次幂零化整个商模。指数可以随标量 \(a\) 改变，但选定 \(a\) 后，必须有同一个指数适用于 \(T\) 的所有元素。仅要求眼前的 \(m\) 被 \(a^n\) 零化并不能表达这一点，因为 \(am\in N\) 已经使它成立。若只知道对每个 \(t\in T\) 分别存在指数使 \(a^{n(t)}t=0\)，在 \(T\) 有限生成时，可对有限组生成元取指数的最大值，得到 \(a^nT=0\)。当 \(T=R/Q\) 时，\(a^nT=0\) 等价于 \(a^n\in Q\)，故此定义与准素理想的定义一致。

\begin{proposition}\label{b3:prop:primary-submodule-associated}
设 \(R\) 为交换含幺 Noether 环，\(M\) 为有限生成 \(R\) 模，\(N\subsetneq M\) 为真子模。\(N\) 准素，当且仅当 \(\operatorname{Ass}_R(M/N)\) 只有一个元素；此元素就是 \(\sqrt{\operatorname{Ann}_R(M/N)}\)，因而该根是素理想。\(M\) 的有限生成假设一般不能删去。
\end{proposition}
\begin{proof}
记 \(T=M/N\) 与 \(J=\operatorname{Ann}_R(T)\)。若 \(N\) 准素，取 \(\mathfrak p=\operatorname{Ann}_R(t)\in\operatorname{Ass}_R(T)\)。因 \(J\subseteq\mathfrak p\)，有 \(\sqrt J\subseteq\mathfrak p\)。另一方面，每个 \(a\in\mathfrak p\) 都零化非零元 \(t\)，准素条件给出 \(a^nT=0\)，即 \(a\in\sqrt J\)。故所有关联素理想都等于 \(\sqrt J\)，且存在性保证集合非空。

反向若 \(\operatorname{Ass}_R(T)=\{\mathfrak p\}\)，这里 \(T\) 的有限生成性保证 \(\operatorname{Supp}_R(T)=V(J)\)。由\cref{b3:prop:minimal-primes-associated}，\(V(J)\) 的极小素理想只有 \(\mathfrak p\)，且其中每个素理想都包含 \(\mathfrak p\)。于是包含 \(J\) 的全部素理想之交就是 \(\mathfrak p\)，即 \(\sqrt J=\mathfrak p\)。若 \(at=0\)、\(t\ne0\)，由\cref{b3:thm:associated-zero-divisors}得 \(a\in\mathfrak p=\sqrt J\)，故某个 \(a^n\in J\)，这就是 \(a^nM\subseteq N\)。

为说明有限生成假设的作用，取域 \(k\)，令 \(R=k[x]\) 与
\[
T=\bigoplus_{n\ge1}R/(x^n).
\]
每个元素只有有限个非零坐标，因此被 \(x\) 的某次幂零化。有限个元素总共只涉及有限个直和分量，不能生成全部 \(T\)，故 \(T\) 不是有限生成模。若非零元素 \(t\) 的零化子是素理想 \(\mathfrak q\)，由 \(x^nt=0\) 得 \(x\in\mathfrak q\)；因 \((x)\) 是极大理想，只能有 \(\mathfrak q=(x)\)。而第一坐标中 \(1+(x)\) 的零化子恰为 \((x)\)，故 \(\operatorname{Ass}_R(T)=\{(x)\}\)。\(x\) 零化这个非零元素，所以是 \(T\) 上的零因子；但对任意 \(n\ge1\)，第 \(n+1\) 个坐标中的 \(1+(x^{n+1})\) 不被 \(x^n\) 零化。因而不存在统一指数使 \(x^nT=0\)，\(0\subset T\) 不是准素子模。逐元素的指数在这里无法统一，且 \(\operatorname{Ann}_R(T)=\bigcap_{n\ge1}(x^n)=0\)，其根也不同于唯一的关联素理想 \((x)\)。
\end{proof}

\begin{theorem}\label{b3:thm:module-primary-decomposition}
设 \(R\) 为交换含幺 Noether 环，\(M\) 为有限生成 \(R\) 模。则 \(M\) 的每个真子模 \(N\) 都有极小准素分解
\[
N=N_1\cap\cdots\cap N_r,
\qquad \operatorname{Ass}_R(M/N_i)=\{\mathfrak p_i\},
\]
其中各 \(\mathfrak p_i\) 互异，且无因子可删。并且
\[
\operatorname{Ass}_R(M/N)=\{\mathfrak p_1,\ldots,\mathfrak p_r\}.
\]
若 \(\mathfrak p_i\) 为该集合的极小元，记 \(\lambda_i:M\rightarrow M_{\mathfrak p_i}\) 为局部化映射，则
\[
N_i=\lambda_i^{-1}(N_{\mathfrak p_i}).
\]
\end{theorem}
\begin{proof}
若真子模 \(N\subsetneq M\) 不能写成两个严格更大的子模的交，则称 \(N\) 在 \(M\) 中不可约。因为 \(M\) 是 Noether 模，\cref{b3:lem:irreducible-ideal-primary}第一段的极大反例证明逐字适用于子模，故 \(N\) 是有限个不可约子模的交。这个有限分解提供了候选因子，还需证明这些因子准素，并将其整理为定理要求的极小准素分解。

不可约子模必准素。否则由\cref{b3:prop:primary-submodule-associated}，非零模 \(M/N\) 至少有两个不同关联素理想 \(\mathfrak p,\mathfrak q\)。分别取子模 \(L\cong R/\mathfrak p\)、\(K\cong R/\mathfrak q\)。若 \(0\ne z\in L\cap K\)，在 \(L\cong R/\mathfrak p\) 中将 \(z\) 对应于非零剩余类 \(\bar b\)。因 \(R/\mathfrak p\) 为整环，\(az=0\) 等价于 \(\bar a\bar b=0\)，也就等价于 \(a\in\mathfrak p\)。所以 \(\operatorname{Ann}_R(z)=\mathfrak p\)；同理，由 \(z\in K\) 又得 \(\operatorname{Ann}_R(z)=\mathfrak q\)，矛盾。所以 \(L\cap K=0\)，两个非零子模在 \(M\) 中的原像便把 \(N\) 写成两个严格更大子模的交，与不可约性矛盾。

为合并根相同的因子，考虑任意非空有限族同根准素子模的交。其商模单射到各商模的直和，由\cref{b3:prop:associated-exact-sequence}，关联素理想只能属于这个单元素集合；商模非零，由\cref{b3:thm:associated-existence}恰有该关联素理想，再由\cref{b3:prop:primary-submodule-associated}得到准素性。于是合并同根因子并删除冗余后，就得到根两两不同、无因子可删的分解。

这些根确实都是 \(M/N\) 的关联素理想。对角单射 \(M/N\hookrightarrow\bigoplus_iM/N_i\) 先给出 \(\operatorname{Ass}_R(M/N)\subseteq\{\mathfrak p_1,\ldots,\mathfrak p_r\}\)。反向，去掉第 \(i\) 个因子会使交严格增大，故 \((\bigcap_{j\ne i}N_j)/N\) 非零；当 \(r=1\) 时，此处空族的交取为 \(M\)。它是 \(M/N\) 的子模，而映入 \(M/N_i\) 的核是 \((\bigcap_jN_j)/N=0\)，所以也嵌入 \(M/N_i\)。它的关联素理想非空且只能为 \(\mathfrak p_i\)，故每个 \(\mathfrak p_i\) 都在 \(\operatorname{Ass}_R(M/N)\) 中出现。

若 \(\mathfrak p_i\) 极小，在 \(\mathfrak p_i\) 处局部化可以使其余分量的商模消失。对 \(j\ne i\)，极小性与根互异保证 \(\mathfrak p_j\not\subseteq\mathfrak p_i\)，故可取 \(a\in\mathfrak p_j\setminus\mathfrak p_i\)。因为 \(\mathfrak p_j=\sqrt{\operatorname{Ann}_R(M/N_j)}\)，某个 \(a^n\) 零化 \(M/N_j\)。局部化后 \(a\) 是单位，故 \((N_j)_{\mathfrak p_i}=M_{\mathfrak p_i}\)。局部化与有限交相容，于是 \(N_{\mathfrak p_i}=(N_i)_{\mathfrak p_i}\)。再将这个局部子模收缩：若 \(m/1\in(N_i)_{\mathfrak p_i}\)，则某个 \(s\notin\mathfrak p_i\) 满足 \(sm\in N_i\)。若 \(m\notin N_i\)，准素性将给出 \(s\in\sqrt{\operatorname{Ann}_R(M/N_i)}=\mathfrak p_i\)，矛盾。所以收缩恰为 \(N_i\)，公式得证。
\end{proof}

\subsection{嵌入关联素理想的显式例子}
\label{b3:subsec:0404:02}

取域 \(k\)，回到 \(R=k[x,y]\)、\(I=(x^2,xy)\)。商模 \(R/I\) 中有
\[
\operatorname{Ann}_R(\bar x)=(x,y),\qquad
\operatorname{Ann}_R(\bar y)=(x).
\]
第一式因为 \(f\bar x=f(0,0)\bar x\)；第二式可将 \(fy\in(x^2,xy)\) 模掉 \(x\)，由 \(k[y]\) 是整环得到 \(f\in(x)\)。这不仅列出了素理想，还给出了定义它们的实际元素。它们正是\cref{b3:thm:primary-associated-primes}对\cref{b3:eq:embedded-line-decomposition}所给出的两个关联素理想。

\ProseParagraphBreak
一个矩阵也可以同时包含孤立与嵌入成分。令
\[
\Phi:R^3\longrightarrow R^2,\qquad
[\Phi]=\begin{pmatrix}x&y&0\\0&0&x\end{pmatrix},
\qquad T=\operatorname{coker}\Phi.
\]
设 \(e_1,e_2\) 为 \(R^2\) 的标准基，\(\bar e_1,\bar e_2\) 为它们在余核中的像。矩阵的三列给出关系
\[
x\bar e_1=0,\qquad y\bar e_1=0,\qquad x\bar e_2=0.
\]
像子模为 \(\operatorname{im}\Phi=(x,y)e_1\oplus(x)e_2\)，两组关系没有混合两个坐标，故 \(R\bar e_1\cong R/(x,y)\)、\(R\bar e_2\cong R/(x)\)，并且
\[
T\cong R/(x,y)\oplus R/(x).
\]
因此 \(\operatorname{Ass}_R(T)=\bigl\{(x,y),(x)\bigr\}\)。\(N=\operatorname{im}\Phi\) 的极小准素分解为
\[
N=\bigl(Re_1\oplus(x)e_2\bigr)
   \cap\bigl((x,y)e_1\oplus Re_2\bigr).
\]
第一个商模为 \(R/(x)\)，第二个为 \(R/(x,y)\)，故两因子分别是 \((x)\) 与 \((x,y)\) 准素分量。在 \((x)\) 处局部化时 \(y\) 成为单位，第一坐标的商消失，第二坐标留下 \(k(y)\)。这就是孤立准素分量的收缩公式在一个展示矩阵上的具体表现。

\subsection{几何分支与非约化结构}
\label{b3:subsec:0404:03}

设 \(R\) 为交换含幺 Noether 环，\(I\subsetneq R\) 为真理想。闭集 \(V(I)=V(\sqrt I)\) 的不可约分量由包含 \(I\) 的极小素理想决定，其中仍包含这些极小点的所有特殊化点，包括嵌入关联素理想对应的点。底拓扑却不能辨认一个非极小点是否为嵌入关联点：将 \(I\) 换为 \(\sqrt I\) 不改变这个闭集，却会除去商环中的幂零元，并可能改变关联素理想集合。以下直接在素谱中解释这种区别，无须预设基域代数闭。

\begin{proposition}\label{b3:prop:reduced-ring-no-embedded-primes}
设 \(R\) 为交换含幺约化环。则
\[
\operatorname{Ass}_R(R)\subseteq\min\operatorname{Spec}R.
\]
因此任意约化环都没有嵌入关联素理想。若 \(R\) 还为 Noether 环，则上述包含为等号。
\end{proposition}
\begin{proof}
若 \(R=0\)，两侧均为空集。以下设 \(R\ne0\)。取 \(\mathfrak p=\operatorname{Ann}_R(x)\in\operatorname{Ass}_R(R)\)，其中 \(x\ne0\)。若 \(\mathfrak p\) 不是极小素理想，则存在素理想 \(\mathfrak q\subsetneq\mathfrak p\)。取 \(a\in\mathfrak p\setminus\mathfrak q\)，由 \(ax=0\in\mathfrak q\) 和 \(a\notin\mathfrak q\)，素性给出 \(x\in\mathfrak q\subseteq\mathfrak p=\operatorname{Ann}_R(x)\)。于是 \(x^2=0\)，约化性迫使 \(x=0\)，矛盾。这证明了包含式。

若 \(R\) 为 Noether 环，对有限生成模 \(R\) 使用\cref{b3:prop:minimal-primes-associated}；因 \(\operatorname{Supp}_R(R)=\operatorname{Spec}R\)，每个极小素理想都属于 \(\operatorname{Ass}_R(R)\)，从而得到等号。
\end{proof}

取域 \(k\)，先看 \(C=k[x,y]/(xy)\)。因为 \((xy)=(x)\cap(y)\)，它是约化环，有两个极小素理想 \((\bar x)\)、\((\bar y)\)，对应两条坐标轴这两个不可约分量。由刚才的命题，它的关联素理想也恰为这两个极小素理想，没有嵌入关联素理想。这里的两个分支来自两个极小素理想，并不来自幂零结构。

\ProseParagraphBreak
即使在 Noether 环中，没有嵌入关联素理想也不保证约化。环 \(B=k[x,y]/(x^2)\) 非约化，却只有孤立关联素理想 \((\bar x)\)。每个元素唯一写成 \(a(y)+\bar x\cdot b(y)\)，其中 \(a,b\in k[y]\)。对 \(c,d\in k[y]\)，有
\[
(a+\bar x\cdot b)(c+\bar x\cdot d)=ac+\bar x\cdot(ad+bc).
\]
若 \(a\ne0\) 且乘积为零，比较两个 \(k[y]\) 分量得 \(ac=0\)、\(ad+bc=0\)。因为 \(k[y]\) 是整环，先得 \(c=0\)，再得 \(d=0\)；故 \(a+\bar x\cdot b\) 的乘法是单射。零因子因此恰为 \((\bar x)\) 中的元素，且都平方为零，说明 \((x^2)\subset k[x,y]\) 为 \((x)\) 准素理想。

\ProseParagraphBreak
在 \(B\) 中，幂零理想 \((\bar x)\) 作为 \(B\) 模同构于 \(B/(\bar x)\cong k[y]\)，其支集是整条直线 \(V(\bar x)\)，因而在这条线的每个点处局部化都不会消失。再看 \(A=k[x,y]/(x^2,xy)\)，其幂零理想 \((\bar x)\) 被 \((\bar x,\bar y)\) 零化，且作为 \(k\) 向量空间只有一维，所以支集只有原点；在开集 \(D(\bar y)\) 上它消失。上一小节计算的 \(\operatorname{Ann}_{k[x,y]}(\bar x)=(x,y)\) 正是这个原点处嵌入关联素理想的来源，它没有增加新的不可约分量。

\ProseParagraphBreak
\(A\) 与 \(B\) 的约化商都同构于 \(k[y]\)，底空间只有一条不可约直线，幂零理想的支集却不同。\(B\) 没有嵌入关联素理想，仍然非约化；\(A\) 在原点有嵌入关联素理想，仍然只有一个不可约分量。相较之下，约化环 \(C\) 有两个几何分支，却没有幂零结构。


\begin{chaptersummary}
关联素理想由循环子模 \(R/\mathfrak p\) 的实际嵌入定义。Noether 环保证非零模有这样的子模，有限生成模的升链条件进一步给出有限素理想滤过，因而关联素理想只有有限个。它们的并是全部零因子，它们的极小元与支集的极小元一致；局部化保留恰好那些不遇到分母集的关联素理想。

\ProseParagraphBreak
在 Noether 环中，理想及有限生成模中的子模可以写成有限个准素分量的交。极小分解中出现的根素理想集合由原商模唯一确定；孤立准素分量还可以用局部化及收缩唯一恢复，嵌入准素分量则一般不唯一。任意约化交换环的关联素理想都是极小素理想，Noether 情形下两者的集合相等；反过来，没有嵌入关联素理想的环仍可能带有幂零元。
\end{chaptersummary}

\BookExercises
\begin{exercise}\label{b3:ex:04-integer-module}
求 \(M=\mathbb Z\oplus\mathbb Z/12\mathbb Z\) 的关联素理想与零因子，并为每个关联素理想给出一个零化子恰为它的元素。
\end{exercise}
\begin{solution}
自由因子的非零元素零化子为 \((0)\)。第二因子内，\(\bar a\) 的零化子由 \(12/\gcd(12,a)\) 生成，成为素理想时只能是 \((2)\) 或 \((3)\)。因此关联素理想为 \((0),(2),(3)\)，可分别取 \((1,\bar0),(0,\bar6),(0,\bar4)\) 为见证。由\cref{b3:thm:associated-zero-divisors}，零因子为 \((0)\cup(2)\cup(3)=(2)\cup(3)\)，即 \(2\mathbb Z\cup3\mathbb Z\)。关联素理想 \((0)\) 包含于另两个理想，没有向并集增加元素；若只保留模 \(\mathbb Z/12\mathbb Z\)，关联集合便少了 \((0)\)，零因子集合却相同。因此仅由零因子的并集不能恢复全部关联素理想。
\end{solution}

\begin{exercise}\label{b3:ex:04-primary-polynomial-module}
设 \(k\) 为域。在 \(R=k[x,y]\) 中，直接计算 \(R/(x^3)\) 的零因子及关联素理想，并求元素 \(\bar x^2\) 的零化子。
\end{exercise}
\begin{solution}
每个类唯一写成 \(a_0(y)+a_1(y)\bar x+a_2(y)\bar x^2\)。若 \(a_0\ne0\)，将其与另一个这样的式子相乘，逐次比较 \(1,\bar x,\bar x^2\) 的系数，利用 \(k[y]\) 为整环可知乘法单射。若 \(a_0=0\)，该元素是幂零元，且会零化非零元 \(\bar x^2\)。所以零因子来自 \((x)\)，\((x^3)\) 为 \((x)\) 准素理想，关联素理想集合为 \(\bigl\{(x)\bigr\}\)。计算 \(f\bar x^2=a_0(y)\bar x^2\) 给出零化子 \((x)\)。
\end{solution}

\begin{exercise}\label{b3:ex:04-exact-quotient-warning}
设 \(k\) 为域，\(R=k[t]\)，令 \(M=R/(t^2)\oplus R\)、\(N=R/(t^2)\oplus t(t-1)R\subset M\)、\(P=M/N\)。计算这三个模的关联素理想，并验证\cref{b3:prop:associated-exact-sequence}中的两个包含。说明第二个包含可以严格，且商模既可以产生新的关联素理想，也可以失去原有的关联素理想。
\end{exercise}
\begin{solution}
因 \(t(t-1)R\cong R\)，有限直和公式给出
\[
\operatorname{Ass}_R(N)=\operatorname{Ass}_R(M)=\{(0),(t)\}.
\]
商模为 \(P\cong R/\bigl(t(t-1)\bigr)\)。理想 \((t)\)、\((t-1)\) 互素，中国剩余定理给出 \(P\cong R/(t)\oplus R/(t-1)\)，所以
\[
\operatorname{Ass}_R(P)=\{(t),(t-1)\}.
\]
第一个包含在此取等号，而第二个为
\(\{(0),(t)\}\subsetneq\{(0),(t),(t-1)\}\)。商模增加了 \((t-1)\)，同时失去 \((0)\)；两个方向都不能从原模的关联集合直接推定。
\end{solution}

\begin{exercise}\label{b3:ex:04-infinite-associated}
求 \(\mathbb Q/\mathbb Z\) 作为 \(\mathbb Z\) 模的关联素理想。说明它的例子不违反关联素理想有限性定理。
\end{exercise}
\begin{solution}
非零元 \(a/n+\mathbb Z\) 在 \(\gcd(a,n)=1\) 时零化子为 \((n)\)，故素零化子恰为所有素数 \(p\) 对应的 \((p)\)，由 \(1/p+\mathbb Z\) 实现。不存在零化子为零的元素，因为每个元素都有有限加法阶。所得集合无限，而 \(\mathbb Q/\mathbb Z\) 不有限生成：有限多个有理数的分母有公倍数 \(d\)，它们生成的模被 \(d\) 零化，却不能包含所有 \(1/p+\mathbb Z\)。有限性定理的模假设不满足。
\end{solution}

\begin{exercise}\label{b3:ex:04-localization-primes}
设 \(k\) 为域，令 \(R=k[x,y]\)、\(M=R/(x)\oplus R/(x,y)\)，取 \(\mathfrak p=(x)\)、\(\mathfrak m_0=(x,y)\)、\(\mathfrak m_1=(x,y-1)\)。计算这三个素理想处的局部模、关联素理想，以及 \(M_{\mathfrak q}/\mathfrak qM_{\mathfrak q}\) 在剩余域 \(\kappa(\mathfrak q)=R_{\mathfrak q}/\mathfrak qR_{\mathfrak q}\) 上的维数，其中 \(\mathfrak q\) 分别取上述三个理想。说明局部环的极大理想何时是该模的关联素理想。
\end{exercise}
\begin{solution}
原关联素理想为 \((x),(x,y)\)。在 \(\mathfrak p\) 处，\(y\) 成为单位，第二因子消失，第一因子将 \(k[y]\) 中所有非零元素倒置，故 \(M_{\mathfrak p}\cong k(y)\)。关联集合为 \(\{(x)R_{\mathfrak p}\}\)，该理想就是局部环的极大理想；模已被它零化，剩余向量空间在 \(k(y)\) 上维数为一。

在 \(\mathfrak m_0\) 处，两因子都保留，得到 \(M_{\mathfrak m_0}\cong k[y]_{(y)}\oplus k\)。关联集合为 \(\{(x)R_{\mathfrak m_0},\mathfrak m_0R_{\mathfrak m_0}\}\)，其中极大理想来自第二因子。对两因子分别模掉 \(\mathfrak m_0\)，各得到一个 \(k\)，故剩余向量空间维数为二。

在 \(\mathfrak m_1\) 处，\(y\notin\mathfrak m_1\)，第二因子消失，第一因子为 \(k[y]_{(y-1)}\)。关联集合为 \(\{(x)R_{\mathfrak m_1}\}\)，而 \(y-1\) 不在 \((x)R_{\mathfrak m_1}\) 中，所以这个关联素理想严格小于局部环的极大理想。剩余向量空间仍是一个 \(k\)，维数为一；维数非零并不说明极大理想一定关联。
\end{solution}

\begin{exercise}\label{b3:ex:04-prime-filtration-explicit}
设 \(k\) 为域。对 \(R=k[x,y]\)、\(M=R/(x^2,xy)\)，构造长度为二的素理想滤过，写出两个因子的具体同构。
\end{exercise}
\begin{solution}
取 \(M_1=R\bar x\)。因 \(\operatorname{Ann}_R(\bar x)=(x,y)\)，映射 \(R/(x,y)\rightarrow M_1\)、\(\bar a\mapsto a\bar x\) 同构。商掉 \(M_1\) 相当于再加关系 \(x=0\)，故 \(M/M_1\cong R/(x)\)。链 \(0\subsetneq M_1\subsetneq M\) 即所求；两个商都非零，且都是相应素商环。
\end{solution}

\begin{exercise}\label{b3:ex:04-primary-colon-basis}
设 \(k\) 为域，在 \(R=k[x,y]\) 中令 \(Q=(x^2,y^3)\)。证明 \(Q\) 准素，写出 \(R/Q\) 的一组 \(k\) 基，并计算 \((Q:x)\)、\((Q:y^2)\)。
\end{exercise}
\begin{solution}
\(\sqrt Q=(x,y)\) 极大，所以由\cref{b3:prop:maximal-radical-primary}，\(Q\) 为 \((x,y)\) 准素理想。标准单项式为 \(1,y,y^2,x,xy,xy^2\)，它们生成商空间且因原理想为单项式理想而线性无关。逐个判断乘积所属关系得 \((Q:x)=(x,y^3)\)、\((Q:y^2)=(x^2,y)\)。两理想的根仍为 \((x,y)\)，符合准素理想商命题。
\end{solution}

\begin{exercise}\label{b3:ex:04-prime-radical-insufficient}
设 \(k\) 为域，令 \(R=k[x,y,z]\)、\(J=(x^2,xz)\)。求其根，并用两个明确元素证明 \(J\) 不准素。指出这说明关于根的哪一个逆命题错误。
\end{exercise}
\begin{solution}
\(x^2\in J\subseteq(x)\)，故 \(\sqrt J=(x)\)，为素理想。但 \(zx\in J\)、\(x\notin J\)，而 \(z^n\notin(x)\supseteq J\) 对所有 \(n\ge1\) 成立，所以 \(J\) 不准素。因而“准素理想的根是素理想”的逆命题错误；极大根准素命题的“极大”不能换成“素”。
\end{solution}

\begin{exercise}\label{b3:ex:04-primary-localization-calculation}
设 \(k\) 为域，在 \(R=k[x,y]\) 中令 \(Q=(x^2,y^3)\)。分别倒置 \(1+x\) 与 \(x\)，计算 \(Q\) 的扩张及收缩；说明两种答案的差别。
\end{exercise}
\begin{solution}
\(1+x\notin(x,y)=\sqrt Q\)，因此在 \(R_{1+x}\) 中，扩张是由 \(x^2,y^3\) 生成的真准素理想，收缩仍为 \(Q\)。商环中 \((1+x)(1-x)=1\)，所以 \(1+x\) 原已可逆，得到 \(R_{1+x}/QR_{1+x}\cong R/Q\)。不变的是这个商环；环境环 \(R_{1+x}\) 则严格包含 \(R\)，因为 \(1/(1+x)\) 不是多项式。倒置 \(x\) 时 \(x^2\in Q\) 成为单位，扩张是全环，收缩是 \(R\)。区别恰在分母集是否与根相交。
\end{solution}

\begin{exercise}\label{b3:ex:04-symbolic-square-cone}
设 \(k\) 为域，\(R=k[x,y,z]/(xy-z^2)\)，\(\mathfrak p=(\bar x,\bar z)\)。证明符号平方 \(\mathfrak p^{(2)}=(\bar x)\)。
\end{exercise}
\begin{solution}
商环 \(R/(\bar x)\cong k[y,z]/(z^2)\) 中，每个元素为 \(a(y)+z\cdot b(y)\)。若 \(a\ne0\)，比较乘积的两个系数可见其为非零因子；若 \(a=0\)，它平方为零。故 \((\bar x)\) 为 \(\mathfrak p\) 准素理想。局部化于 \(\mathfrak p\) 后 \(\bar y\) 可逆，\(\bar x=\bar z^2/\bar y\)，且 \(\mathfrak pR_{\mathfrak p}=(\bar z)\)，于是 \(\mathfrak p^2R_{\mathfrak p}=(\bar x)R_{\mathfrak p}\)。利用准素理想的收缩性质得到 \(\mathfrak p^{(2)}=(\bar x)\)。这也再次显示 \(\mathfrak p^2\subsetneq\mathfrak p^{(2)}\)。
\end{solution}

\begin{exercise}\label{b3:ex:04-two-isolated-components}
设 \(k\) 为域。在 \(R=k[x,y,z]\) 中分解 \(I=(xy,xz)\)，求 \(\operatorname{Ass}_R(R/I)\)，并解释其准素分量为何全部唯一。
\end{exercise}
\begin{solution}
\(I=(x)\cap(y,z)\)：若 \(xf\in(y,z)\)，模掉 \((y,z)\) 后在 \(k[x]\) 中有 \(x\bar f=0\)，故 \(f\in(y,z)\)。两个因子都是素理想，互不包含，且 \(x\)、\(y\) 分别见证无冗余。因此关联素理想恰为 \((x),(y,z)\)，两者都极小，准素分量由\cref{b3:thm:primary-isolated-uniqueness}唯一决定。
\end{solution}

\begin{exercise}\label{b3:ex:04-three-components}
设 \(k\) 为域。证明 \(I=(x^2y,xy^2)\subset k[x,y]\) 有极小准素分解
\[
I=(x)\cap(y)\cap(x,y)^3.
\]
求关联素理想，并给出一个零化子为嵌入素理想的元素。
\end{exercise}
\begin{solution}
前两个理想的交为 \((xy)\)，其中全次数至少三的单项式恰被 \(x^2y\) 或 \(xy^2\) 整除，故再与 \((x,y)^3\) 取交得到 \(I\)。第三因子因根极大而准素。删去第三因子会增加 \(xy\)，删去第一因子会增加 \(y^3\)，删去第二因子会增加 \(x^3\)，因此分解极小。关联素理想为 \((x),(y),(x,y)\)，最后一个嵌入。非零元 \(\overline{xy}\) 的零化子为 \((x,y)\)：若 \(fxy\in(x^2y,xy^2)\)，可写成 \(fxy=xy(xa+yb)\)，在整环 \(k[x,y]\) 中约去非零元 \(xy\) 得 \(f=xa+yb\in(x,y)\)；反向直接相乘即可。
\end{solution}

\begin{exercise}\label{b3:ex:04-saturation-components}
设 \(k\) 为域。对 \(I=(x^2,xy)\subset R=k[x,y]\)，计算 \(I:y^\infty\) 与 \(I:x^\infty\)，并用准素分解解释被保留与被消去的准素分量。
\end{exercise}
\begin{solution}
对每个 \(n\ge1\)，\(fy^n\in I\) 当且仅当 \(f\in(x)\)，故 \(I:y^\infty=(x)\)。又 \(x^2\in I\)，所以 \(I:x^\infty=R\)。分解 \(I=(x)\cap(x^2,y)\) 中，倒置 \(y\) 保留根为 \((x)\) 的孤立准素分量，消去根含 \(y\) 的嵌入准素分量；倒置 \(x\) 则同时消去两个准素分量。饱和是相应局部化扩张的收缩，因此与直接计算一致。
\end{solution}

\begin{exercise}\label{b3:ex:04-maximal-components-crt}
设 \(R\) 为交换含幺 Noether 环，\(I\subsetneq R\) 为真理想，且 \(R/I\) 的全部关联素理想都是极大理想。证明 \(I\) 的极小准素分解除分量排列外唯一；若记作 \(I=\bigcap_{i=1}^r Q_i\)，则 \(R/I\cong\prod_{i=1}^rR/Q_i\)。
\end{exercise}
\begin{solution}
由\cref{b3:thm:primary-associated-primes}，极小准素分解的根素理想恰为 \(\operatorname{Ass}_R(R/I)\)。不同极大理想互不包含，所以这些关联素理想都孤立，\cref{b3:thm:primary-isolated-uniqueness}使各对应准素分量唯一。因此分解仅差分量的排列。若 \(Q_i+Q_j\ne R\)，取包含此和的极大理想 \(\mathfrak m\)，它包含 \(\sqrt{Q_i}\) 与 \(\sqrt{Q_j}\)，故两者都等于 \(\mathfrak m\)，与根互异矛盾。因此各 \(Q_i\) 两两互素，中国剩余定理给出所述乘积分解，其核正是 \(\bigcap_iQ_i=I\)。
\end{solution}

\begin{exercise}\label{b3:ex:04-parameter-embedded-components}
设 \(k\) 为域。对任意 \(c\in k\)，令 \(Q_c=(x^2,y-cx)\subset R=k[x,y]\)。证明 \((x)\cap Q_c=(x^2,xy)\)，并证明 \(c\ne d\) 时 \(Q_c\ne Q_d\)。
\end{exercise}
\begin{solution}
\(R/Q_c\cong k[x]/(x^2)\)，其中 \(y=cx\)，所以 \(Q_c\) 的根为 \((x,y)\)，它准素。若 \(xf\in Q_c\)，在此商环中 \(xf=x\cdot f(0,0)\)，而 \(x\ne0\)，故 \(f(0,0)=0\)，即 \(f\in(x,y)\)。因此交包含于 \((x^2,xy)\)，反向由 \(xy=x(y-cx)+cx^2\) 立即验证。若 \(c\ne d\)，\(y-cx\) 在 \(R/Q_d\) 中等于 \((d-c)x\ne0\)，所以不属于 \(Q_d\)。这给出另一类明确不同的嵌入准素分量。
\end{solution}

\begin{exercise}\label{b3:ex:04-diagonal-presentation}
设 \(k\) 为域，\(R=k[t]\)、\(M=R^2\)，\(N\) 是矩阵 \(\operatorname{diag}\bigl(t^2,(t-1)^3\bigr)\) 的像。给出 \(N\) 的极小准素分解、\(M/N\) 的关联素理想，以及商模的 \(k\) 维数。
\end{exercise}
\begin{solution}
令 \(N_0=t^2Re_1\oplus Re_2\)、\(N_1=Re_1\oplus(t-1)^3Re_2\)，则逐坐标取交得 \(N=N_0\cap N_1\)。两个商模分别为 \(R/(t^2)\) 与 \(R/\bigl((t-1)^3\bigr)\)，因而两因子分别属于 \((t)\)、\((t-1)\)，且无冗余。关联素理想就是这两个极大理想。商模同构于上述两个循环模的直和，维数为 \(2+3=5\)。
\end{solution}

\begin{exercise}\label{b3:ex:04-finiteness-primary-module}
设 \(k\) 为域，令 \(R=k[x]\)、\(P=\prod_{n\ge1}R/(x^n)\)。计算 \(\operatorname{Ann}_R(P)\)、\(\operatorname{Ass}_R(P)\) 与 \(Z_R(P)\)，并与\cref{b3:prop:primary-submodule-associated}证明中的无限直和反例比较。说明为什么每个坐标因子都被 \(x\) 的某个幂零化，直积里仍能出现零化子为 \((0)\) 的元素。
\end{exercise}
\begin{solution}
\(\operatorname{Ann}_R(P)=\bigcap_{n\ge1}(x^n)=(0)\)。对任意非零坐标元，它的零化子为某个 \((x^r)\)、\(r\ge1\)。一个非零序列的零化子是其非零坐标零化子的交：若这些指数有界，就取其中最大值，得到 \((x^r)\)；若指数无界，交为 \((0)\)。所以素零化子只能是 \((x)\) 或 \((0)\)。仅第一个坐标为 \(\bar1\) 的序列实现 \((x)\)，而所有坐标都为 \(\bar1\) 的序列实现 \((0)\)，故
\[
\operatorname{Ass}_R(P)=\{(0),(x)\}.
\]
若 \(f\notin(x)\)，它在每个 \(R/(x^n)\) 中都是单位，因此乘 \(f\) 在直积上也是单射；若 \(f\in(x)\)，它会零化仅第一个坐标非零的序列。所以 \(Z_R(P)=(x)\)，关联素理想 \((0)\) 没有向此并集增加元素。

无限直和中的元素只用有限个坐标，每个元素仍被某个 \(x^r\) 零化；无限直积没有这种限制，序列 \((\bar1,\bar1,\ldots)\) 的坐标需要无界的幂，故没有非零多项式能同时零化它。两个模本身的零化子都为 \((0)\)，零因子集合都为 \((x)\)，关联素理想集合却不同。
\end{solution}

\begin{exercise}\label{b3:ex:04-primary-not-irreducible-module}
设 \(R\) 为交换含幺环，\(\mathfrak p\) 为 \(R\) 的素理想，\(M=R/\mathfrak p\oplus R/\mathfrak p\)。直接证明零子模准素但可约。
\end{exercise}
\begin{solution}
若 \(a(m_1,m_2)=0\) 且元组非零，至少一个坐标在整环 \(R/\mathfrak p\) 中非零，故 \(a\in\mathfrak p\)。于是 \(aM=0\)，准素条件甚至以指数一成立。另一方面，两个坐标子模 \(R/\mathfrak p\oplus0\) 与 \(0\oplus R/\mathfrak p\) 都非零，交为零，所以零子模可约。子模准素与不能分解为交是不同性质。
\end{solution}

\begin{exercise}\label{b3:ex:04-triangular-presentation}
设 \(k\) 为域，在 \(R=k[x,y]\) 中令
\[
\Phi=\begin{pmatrix}x&y\\0&x\end{pmatrix},\qquad
N=\operatorname{im}\Phi\subset R^2.
\]
证明 \(R^2/N\) 的关联素理想只有 \((x)\)，因而 \(N\) 是准素子模。
\end{exercise}
\begin{solution}
记 \(\Psi=\begin{pmatrix}x&-y\\0&x\end{pmatrix}\)，则 \(\Phi\Psi=\Psi\Phi=x^2I_2\)。映射 \(v\mapsto\Psi v\bmod x^2R^2\) 诱导
\[
R^2/N\longrightarrow\bigl(R/(x^2)\bigr)^2.
\]
它是单射：若 \(\Psi v=x^2w\)，乘以 \(\Phi\) 得 \(x^2v=x^2\Phi w\)，而 \(R\) 是整环，所以 \(v=\Phi w\)。目标的关联素理想只有 \((x)\)，故源的关联集合包含于它。源非零，例如模掉 \((x,y)\) 后矩阵为零而余核为 \(k^2\)，所以关联素理想存在，必须恰为 \((x)\)。再用\cref{b3:prop:primary-submodule-associated}得 \(N\) 准素。
\end{solution}

\begin{exercise}\label{b3:ex:04-three-axes}
设 \(k\) 为域，在 \(R=k[x,y,z]\) 中令 \(I=(xy,xz,yz)\)。证明 \(I\) 为根理想，求 \(R/I\) 的关联素理想，并判断 \(\overline{x+y+z}\) 是否为零因子。
\end{exercise}
\begin{solution}
由单项式整除判据，\(I=(x,y)\cap(x,z)\cap(y,z)\)。它是素理想的交，故为根理想；三个因子互不包含且无冗余，因此关联素理想为这三个理想，或在商环中写作它们的像。\(x+y+z\) 模这三个理想分别为 \(z,y,x\)，均非零，故不在任何关联素理想中。由\cref{b3:thm:associated-zero-divisors}，其像为非零因子。
\end{solution}

\begin{exercise}\label{b3:ex:04-embedded-thickness}
设 \(k\) 为域，\(r\ge1\)，令 \(A=k[x,y]/(x^2,xy^r)\)。求幂零理想的 \(k\) 维数、它的支集，以及 \(A\) 的关联素理想。
\end{exercise}
\begin{solution}
根为 \((x)\)，所以幂零理想为 \((\bar x)\)，其 \(k\) 基为 \(\bar x,\bar x\bar y,\ldots,\bar x\bar y^{r-1}\)，维数 \(r\)。零化子为 \((\bar x,\bar y^r)\)，故支集只有 \((\bar x,\bar y)\) 这一点。直接验证
\[
(x^2,xy^r)=(x)\cap(x^2,y^r)
\]
给出极小准素分解，两个关联素理想为 \((\bar x)\) 与 \((\bar x,\bar y)\)。后者也由元素 \(\bar x\bar y^{r-1}\) 的零化子实现。倒置 \(y\) 后幂零理想消失，而不同的 \(r\) 给出不同维数的原点幂零部分。
\end{solution}

\begin{exercise}\label{b3:ex:04-annihilator-radical}
设 \(R\) 为交换含幺 Noether 环，\(M\ne0\) 为有限生成 \(R\) 模，\(\operatorname{Ass}_R(M)=\{\mathfrak p_1,\ldots,\mathfrak p_r\}\)。证明
\[
\sqrt{\operatorname{Ann}_R(M)}=\bigcap_{i=1}^r\mathfrak p_i,
\]
并证明某个 \(n\ge1\) 满足 \((\mathfrak p_1\cdots\mathfrak p_r)^nM=0\)。
\end{exercise}
\begin{solution}
关联素理想集合非空且有限，所以每个 \(\mathfrak p\in\operatorname{Ass}_R(M)\) 都包含某个极小关联素理想 \(\mathfrak q\)。因此，属于全部极小关联素理想的元素也属于每个关联素理想；反向包含则显然。由\cref{b3:prop:minimal-primes-associated}及\cref{b3:thm:radical-prime-intersection}，得到
\[
\begin{aligned}
\bigcap_{\mathfrak p\in\operatorname{Ass}_R(M)}\mathfrak p
&=\bigcap_{\mathfrak q\in\min\operatorname{Ass}_R(M)}\mathfrak q\\
&=\bigcap_{\mathfrak q\in\min V(\operatorname{Ann}_R(M))}\mathfrak q
=\sqrt{\operatorname{Ann}_R(M)}.
\end{aligned}
\]
记这个交为 \(J=(a_1,\ldots,a_s)\)，有限生成来自 Noether 性。对每个 \(a_i\) 选 \(n_i\ge1\) 使 \(a_i^{n_i}M=0\)。令 \(n=1+\sum_i(n_i-1)\)，则任意 \(n\) 次生成元乘积中某个 \(a_i\) 至少出现 \(n_i\) 次，所以 \(J^nM=0\)。因为 \(\mathfrak p_1\cdots\mathfrak p_r\subseteq J\)，所需结论成立。若 \(J=0\)，直接取 \(n=1\)。
\end{solution}

% !TeX root = ../../Book3.tex
% 纲要标题：## 第5章　整扩张、整闭性与 Noether 正规化
% 纲要位置：工作记录/计划纲要.md，第 2034 行。

\chapter{整扩张、整闭性与 Noether 正规化}\label{b3:ch:05}
\BookChapterNumber{b3:ch:05}

\begin{chapterabstract}
整性把首一方程与有限生成模联系起来，使环扩张中的元素和素理想都能受到控制。本章从有限生成模判据出发，建立整闭包、素理想的提升与不可比性，完整证明整闭基整环上的 Going-down 定理。最后用适用于任意域的加权代换证明 Noether 正规化，把有限型代数组织成多项式子环上的有限生成模。
\end{chapterabstract}

\begin{learninggoals}
\begin{itemize}
\item 用有限生成模判据证明整性，区分有限型、有限与整扩张。
\item 明确整闭包的环境，检验整闭性及局部化的作用。
\item 在指定端点上提升素理想链，并核对 Going-up 与 Going-down 的不同假设。
\item 构造使原代数成为有限生成模的多项式子环，说明有限域上线性代换可能失败及加权代换的作用。
\end{itemize}
\end{learninggoals}

设 \(k\) 为域，考虑子环 \(A=k[t^2,t^3]\subset k[t]\)。元素 \(t\) 不属于 \(A\)，却满足首一方程 \(T^2-t^2=0\)，其系数都在 \(A\) 中。每个多项式都可写成 \(a(t^2)+t\,b(t^2)\)，故 \(k[t]=A\cdot1+A\cdot t\) 是有限生成的 \(A\) 模。这是一组模生成元，却不是模基：系数 \(t^2,t^3\) 都在 \(A\) 中，并给出非平凡关系 \(t^2\cdot t-t^3\cdot1=0\)。整性的有限生成模判据只需要生成性，不要求生成元线性无关。首一方程把添入的新元素的高次幂限制在有限个低次幂的线性生成空间中，这是整性的最初模型。

\ProseParagraphBreak

对于一般的有限型 \(k\) 代数 \(k[x_1,\ldots,x_n]/I\)，原来的生成元通常满足代数关系。能否在这个代数中选出若干代数独立元素 \(y_1,\ldots,y_d\)，使整个代数成为 \(k[y_1,\ldots,y_d]\) 上的有限生成模？这就是 Noether 正规化要回答的问题。首一方程不仅限制元素的幂次，也约束扩张前后素理想的行为；理解这些约束，才能把所选坐标与素谱中的点和包含关系联系起来。

\ProseParagraphBreak

本章用到环局部化、素理想对应和有限生成模的行列式技巧。最小多项式与代数独立沿用第一卷的域论定义；Noether 正规化的对象包括含零因子的有限型代数，证明也在这一范围内进行。素理想链的比较先按包含关系处理，维数术语留待下一编。

\ProseParagraphBreak
Matsumura 的《Commutative Ring Theory》\cite[§9, Theorems 9.1--9.4, pp. 64--68]{Matsumura1989CommutativeRingTheory}适合对照整性与素理想提升；其中 Going-down 的主证明使用正规扩域中的素理想共轭，本章改用最小多项式系数的局部化论证。\cite[§33, Lemma 2, p. 262]{Matsumura1989CommutativeRingTheory}给出无限域上的线性代换证明，本章同时处理有限域，避免把非零多项式必有非零取值的结论用于有限域。

% !TeX root = ../../Book3.tex
% 纲要标题：### 5.1 整元素
% 纲要位置：工作记录/计划纲要.md，第 1935 行。
% 纲要细目（写作范围）：
% * 首一方程及有限生成模判据
% * 整性的传递性
% * 整闭包及正规整环

\section{整元素}\label{b3:sec:0501}

有理数 \(1/2\) 满足整数系数方程 \(2T-1=0\)，却不满足整数系数的首一方程。首项系数能否取为一，决定了高次幂能否用低次幂表示，也决定了加入一个元素以后是否仍能用有限个模生成元控制运算。本节始终约定环交换且有单位，环包含保持单位。

\subsection{首一方程与有限生成模判据}\label{b3:subsec:0501:01}

\begin{definition}\label{b3:def:integral-element}
设 \(A\subseteq B\) 为交换含幺环的包含，\(b\in B\)。若存在整数 \(n\ge1\) 和 \(a_1,\ldots,a_n\in A\)，使 \(b\) 满足
\[
b^n+a_1b^{n-1}+\cdots+a_n=0,
\]
则称 \(b\) 在 \(A\) 上\term[zheng-yuansu]{整}[integral]。若 \(B\) 的每个元素都在 \(A\) 上整，则称 \(B/A\) 为\term[zheng-kuozhang]{整扩张}[integral extension]。若 \(B\) 作为 \(A\) 模有限生成，则称 \(B/A\) 为\term[youxian-huan-kuozhang]{有限扩张}[finite extension]。
\end{definition}

这里的“有限”是模的有限生成性，不是底集合有限，也不是仅指代数生成元有限。例如 \(A[T]\) 是有限型 \(A\) 代数，但当 \(A\ne0\) 时不可能是有限生成的 \(A\) 模。

\begin{proposition}[整性的有限生成模判据]\label{b3:prop:integral-finite-criterion}
设 \(A\subseteq B\) 为交换含幺环的包含，\(b\in B\)。下列条件等价：
\begin{enumerate}[label=\Roman*.]
\item \(b\) 在 \(A\) 上整；
\item \(A[b]\) 是有限生成的 \(A\) 模；
\item 存在有限生成 \(A\) 模 \(M\)，它同时是忠实 \(A[b]\) 模，且两种 \(A\) 作用一致。
\end{enumerate}
因此，有限环扩张必为整扩张；若 \(B=A[b_1,\ldots,b_r]\)，则 \(B/A\) 有限当且仅当每个 \(b_i\) 在 \(A\) 上整。
\end{proposition}
\begin{proof}
首一方程可把 \(b^n\) 及更高幂递次化为 \(1,b,\ldots,b^{n-1}\) 的 \(A\) 线性组合，故第一项推出第二项。取 \(M=A[b]\)，其中一使乘法作用忠实，得到第三项。

设第三项成立，选 \(A\) 模生成元 \(m_1,\ldots,m_s\)，写 \(bm_j=\sum_i c_{ij}m_i\)，其中 \(c_{ij}\in A\)。\cref{b3:thm:determinant-trick}说明首一多项式 \(\det(TI_s-C)\) 在乘法算子 \(b\) 上取值为零。忠实性再说明这个取值作为 \(A[b]\) 元素也为零，所以 \(b\) 整。

对有限扩张取 \(M=B\)，乘法作用同样忠实，便得到每个元素整。反过来，若 \(b_i\) 满足次数 \(n_i\) 的首一方程，则所有单项式
\[
b_1^{e_1}\cdots b_r^{e_r},\qquad 0\le e_i<n_i,
\]
生成整个 \(B\)，因为可依次降低每个变量的指数。
\end{proof}

例如 \(\mathbb Z[\sqrt2,\sqrt3]\) 由 \(1,\sqrt2,\sqrt3,\sqrt6\) 作为整数模生成。每个有限代数扩域都是整扩张；一般代数扩域也整，却未必有限。例如 \(\overline{\mathbb F}_p/\mathbb F_p\) 包含任意大的有限子域，不是有限维扩张。

\subsection{整性的传递性}\label{b3:subsec:0501:02}

\begin{proposition}[整性的传递性]\label{b3:prop:integral-transitivity}
设 \(A\subseteq B\subseteq C\) 为交换含幺环的包含。若 \(B/A\) 和 \(C/B\) 都整，则 \(C/A\) 整。
\end{proposition}
\begin{proof}
对 \(c\in C\)，取其在 \(B\) 上的首一方程，设其中的有限个系数为 \(b_1,\ldots,b_r\)。由\cref{b3:prop:integral-finite-criterion}，\(B_0=A[b_1,\ldots,b_r]\) 是有限生成的 \(A\) 模，\(B_0[c]\) 是有限生成的 \(B_0\) 模。若两层分别由 \(u_i\) 与 \(v_j\) 生成，则乘积 \(u_iv_j\) 作为 \(A\) 模生成 \(B_0[c]\)。再次应用有限生成模判据，得到 \(c\) 在 \(A\) 上整。
\end{proof}

\begin{proposition}\label{b3:prop:integral-quotient-localization}
设 \(A\subseteq B\) 为交换含幺环的整扩张。对理想 \(J\subseteq B\)，包含 \(A/(J\cap A)\subseteq B/J\) 整；对 \(A\) 中含一的乘法集 \(S\)，包含 \(S^{-1}A\subseteq S^{-1}B\) 整。
\end{proposition}
\begin{proof}
商环中的首一方程由原方程逐项取像得到。局部化中，对 \(b/s\) 把 \(b\) 的次数为 \(n\) 的首一方程除以 \(s^n\)，得到系数 \(a_i/s^i\in S^{-1}A\) 的首一方程。局部化的包含仍是单射：若 \(a/s\) 在 \(S^{-1}B\) 中为零，则某个 \(t\in S\) 满足 \(ta=0\)，该等式也在 \(A\) 中成立。
\end{proof}

这一命题允许把整性问题移到某个素理想附近，或移到某个闭集上。它不声称任意 \(B\) 中分母的局部化仍在 \(A\) 上整；例如把整数中的所有非零元素倒置便得到非整扩张 \(\mathbb Z\subseteq\mathbb Q\)。

\subsection{整闭包与正规整环}\label{b3:subsec:0501:03}

\begin{proposition}\label{b3:prop:integral-closure-ring}
设 \(A\subseteq B\) 为交换含幺环的包含。\(B\) 中所有在 \(A\) 上整的元素组成一个包含 \(A\) 的子环 \(\overline A^{\,B}\)。若 \(c\in B\) 在此子环上整，则 \(c\in\overline A^{\,B}\)。
\end{proposition}
\begin{proof}
若 \(x,y\) 整，则 \(A[x,y]\) 是有限生成的 \(A\) 模，其中的 \(x+y\)、\(xy\)、\(-x\) 均由\cref{b3:prop:integral-finite-criterion}在 \(A\) 上整。这给出子环性。最后一项由\cref{b3:prop:integral-transitivity}得到。
\end{proof}

\begin{definition}\label{b3:def:integral-closure}
设 \(A\subseteq B\) 为交换含幺环的包含，\(\overline A^{\,B}\) 为 \(B\) 中所有在 \(A\) 上整的元素组成的子环。称 \(\overline A^{\,B}\) 为 \(A\) 在 \(B\) 中的\term[zhengbi-bao]{整闭包}[integral closure]。若 \(\overline A^{\,B}=A\)，则称 \(A\) 在 \(B\) 中整闭。若整环 \(A\) 在其分式域 \(\operatorname{Frac}(A)\) 中整闭，则称 \(A\) 为\term[zhengbi-zhenghuan]{整闭整环}[integrally closed domain]，也称\term[zhenggui-zhenghuan]{正规整环}[normal domain]。
\end{definition}

本章的正规整环不把 Noether 性包含在定义中。整闭性必须带有环境：\(\mathbb Z\) 在 \(\mathbb Q\) 中整闭，却不在 \(\mathbb Q(\sqrt2)\) 中整闭。

\begin{proposition}\label{b3:prop:ufd-integrally-closed}
设 \(A\) 为交换含幺整环。若 \(A\) 是唯一分解整环，则 \(A\) 整闭。若 \(A\) 整闭，且 \(S\subseteq A\) 是含一、不含零的乘法集，则 \(S^{-1}A\) 仍是整闭整环。
\end{proposition}
\begin{proof}
设 \(A\) 是唯一分解整环，\(a/b\) 为互素分子分母表示。若它满足次数 \(n\) 的首一方程，乘以 \(b^n\) 得 \(b\mid a^n\)。唯一分解及互素性迫使 \(b\) 为单位，故 \(a/b\in A\)。

设 \(A\) 整闭、\(0\notin S\)，\(z\in\operatorname{Frac}(A)\) 在 \(S^{-1}A\) 上整。把其方程的有限个系数通分为 \(a_i/s\)，其中 \(s\in S\)。令 \(w=sz\)，乘原方程以 \(s^n\)，得
\[
w^n+a_1w^{n-1}+a_2s w^{n-2}+\cdots+a_ns^{n-1}=0.
\]
于是 \(w\in A\)，从而 \(z=w/s\in S^{-1}A\)。
\end{proof}

特别地，域上的多项式环是整闭整环。相反，\(A=k[t^2,t^3]\subseteq k[t]\) 的分式域含有 \(t=t^3/t^2\)，而 \(t\) 满足 \(T^2-t^2=0\) 却不在 \(A\) 中，因为 \(A\) 的每个多项式的一次项系数为零。这是后面正规化计算的基本例子。

% !TeX root = ../../Book3.tex
% 纲要标题：### 5.2 素理想的提升
% 纲要位置：工作记录/计划纲要.md，第 2042 行。
% 纲要细目（写作范围）：
% * Lying-over 定理；局部化与商域判据所确定的收缩
% * Incomparability 定理；转入分式域及同一纤维的不可比性
% * Going-up 定理；取商保证指定下端

\section{素理想的提升}\label{b3:sec:0502}

环包含 \(A\subseteq B\) 给出素谱上的收缩映射 \(\pi:\operatorname{Spec}B\to\operatorname{Spec}A\)、\(\mathfrak q\mapsto\mathfrak q\cap A\)，且收缩总保持素理想间的包含关系。整性还给出三项约束：每个底层素理想都有原像；收缩到同一个底层素理想的两个素理想若可比较，便必须相等；给定 \(\mathfrak p\subseteq\mathfrak p'\) 以及 \(\mathfrak p\) 上方的 \(\mathfrak q\)，还可找到包含 \(\mathfrak q\) 且收缩为 \(\mathfrak p'\) 的素理想。

\subsection{Lying-over 定理}\label{b3:subsec:0502:01}

在整扩张 \(A\subseteq B\) 中，若 \(\mathfrak q\) 极大，则 \(B/\mathfrak q\) 是域。要证明其收缩也极大，就要由整性推出下方的商环 \(A/(\mathfrak q\cap A)\) 也是域。下面的域判据提供了这一推理，随后将用于局部化后的扩张。

\begin{lemma}\label{b3:lem:integral-field-criterion}
设 \(D\subseteq E\) 为交换含幺整环的整扩张。则 \(D\) 是域当且仅当 \(E\) 是域。
\end{lemma}
\begin{proof}
若 \(D\) 是域，对非零 \(e\in E\) 取最低次数的首一方程。其常数项非零，否则可在整环中约去 \(e\) 而降次。把方程除以 \(e\)，再除以非零常数项，可把 \(e^{-1}\) 写成 \(e\) 的多项式，故 \(E\) 是域。

若 \(E\) 是域，对 \(0\ne d\in D\)，\(d^{-1}\in E\) 整。把它的次数 \(n\) 的首一方程乘以 \(d^{n-1}\)，得 \(d^{-1}\in D\)。所以 \(D\) 是域。
\end{proof}

设 \(A\subseteq B\) 为交换含幺环的整扩张，\(\mathfrak q\in\operatorname{Spec}B\)，\(\mathfrak p=\mathfrak q\cap A\)。收缩保证 \(\mathfrak p\) 是素理想，\cref{b3:prop:integral-quotient-localization}给出整包含
\[
A/\mathfrak p\hookrightarrow B/\mathfrak q.
\]
这两个商环都是整环，因而域判据给出
\[
\mathfrak q\;\text{极大}
\ \Longleftrightarrow\ B/\mathfrak q\;\text{是域}
\ \Longleftrightarrow\ A/\mathfrak p\;\text{是域}
\ \Longleftrightarrow\ \mathfrak p\;\text{极大}.
\]

\begin{theorem}[Lying-over]\label{b3:thm:lying-over}
设 \(A\subseteq B\) 为交换含幺环的整扩张。则对每个 \(\mathfrak p\in\operatorname{Spec}A\)，存在 \(\mathfrak q\in\operatorname{Spec}B\) 满足 \(\mathfrak q\cap A=\mathfrak p\)。这一结论称为\term[lying-over]{Lying-over 定理}[lying-over theorem]。
\end{theorem}
\begin{proof}
在 \(\mathfrak p\) 处局部化后，底环的唯一极大理想是 \(\mathfrak pA_{\mathfrak p}\)。这样，上环中任取一个极大理想，只要其收缩仍然极大，就必收缩到这个指定的理想。取 \(S=A\setminus\mathfrak p\)，\cref{b3:prop:integral-quotient-localization}给出整包含 \(A_{\mathfrak p}\subseteq S^{-1}B\)。因 \(A_{\mathfrak p}\ne0\) 且包含仍是单射，\(S^{-1}B\ne0\)，故可取其极大理想 \(\mathfrak n\)。令 \(\mathfrak m=\mathfrak n\cap A_{\mathfrak p}\)，则 \(\mathfrak m\) 为素理想；再次取商得到整包含
\[
A_{\mathfrak p}/\mathfrak m\hookrightarrow (S^{-1}B)/\mathfrak n.
\]
此时两边都是整环，右边还是域。由\cref{b3:lem:integral-field-criterion}，左边也是域，所以 \(\mathfrak m\) 极大，只能是 \(\mathfrak pA_{\mathfrak p}\)。局部化素理想对应\cref{b3:thm:spec-localization}将 \(\mathfrak n\) 收缩为 \(B\) 中不交 \(S\) 的素理想 \(\mathfrak q\)，而
\[
\begin{aligned}
\mathfrak q\cap A
&=\{\,a\in A\mid a/1\in\mathfrak n\cap A_{\mathfrak p}\,\}\\
&=\{\,a\in A\mid a/1\in\mathfrak pA_{\mathfrak p}\,\}
=\mathfrak p.
\end{aligned}
\]
\end{proof}

包含是必要条件。若只给一个不单射的整同态 \(A\to B\)，像是 \(\operatorname{Spec}A\) 中包含同态核的闭集；应对 \(A/\ker\to B\) 应用定理。

\subsection{Incomparability 定理}\label{b3:subsec:0502:02}

\begin{theorem}[Incomparability]\label{b3:thm:incomparability}
设 \(A\subseteq B\) 为交换含幺环的整扩张，\(\mathfrak q\subseteq\mathfrak q'\) 是 \(B\) 的素理想。若二者在 \(A\) 中的收缩都为 \(\mathfrak p\)，则 \(\mathfrak q=\mathfrak q'\)。这称为\term[incomparability]{Incomparability 定理}[incomparability theorem]。
\end{theorem}
\begin{proof}
先模掉 \(\mathfrak q\)，把所要证明的等式变成 \(\mathfrak q'/\mathfrak q=0\)。此时 \(A/\mathfrak p\) 嵌入整环 \(B/\mathfrak q\)；再倒置底环 \(A/\mathfrak p\) 的全部非零元素，把底环变为域 \(\operatorname{Frac}(A/\mathfrak p)\)，并得到在它上面整的整环
\[
C=(A\setminus\mathfrak p)^{-1}(B/\mathfrak q).
\]
由\cref{b3:lem:integral-field-criterion}，\(C\) 也成为域。\(\mathfrak q'/\mathfrak q\) 不交这些分母，故对应于 \(C\) 的一个素理想，而域中只有零素理想。所倒置的分母在整环 \(B/\mathfrak q\) 中都非零，局部化映射因此单射，\(\mathfrak q'/\mathfrak q\) 不可能有非零元素在局部化后消失，故它已经为零。
\end{proof}

同一底层素理想仍可有多个原像。例如整扩张 \(\mathbb Z\subseteq\mathbb Z[i]\) 中，\((2+i)\) 与 \((2-i)\) 的商环都同构于 \(\mathbb F_5\)，其中 \(i\) 分别映为 \(-2\) 与 \(2\)。两理想因此都收缩到 \((5)\)，而且都是极大理想；\(2+i\) 在第二个商环中的像为 \(4\ne0\)，所以两理想不同，便互不包含。

\subsection{Going-up 定理}\label{b3:subsec:0502:03}

\begin{theorem}[Going-up]\label{b3:thm:going-up}
设 \(A\subseteq B\) 为交换含幺环的整扩张，\(\mathfrak p\subseteq\mathfrak p'\) 是 \(A\) 的素理想，\(\mathfrak q\) 是 \(B\) 中满足 \(\mathfrak q\cap A=\mathfrak p\) 的素理想。存在 \(\mathfrak q'\supseteq\mathfrak q\) 满足 \(\mathfrak q'\cap A=\mathfrak p'\)。这称为\term[going-up]{Going-up 定理}[going-up theorem]。
\end{theorem}
\begin{proof}
要保证新素理想包含已指定的 \(\mathfrak q\)，可以把 \(\mathfrak q\) 商成零：\(B/\mathfrak q\) 中每个素理想的原像都会包含它。由\cref{b3:prop:integral-quotient-localization}，商包含 \(A/\mathfrak p\subseteq B/\mathfrak q\) 整。在此包含上对素理想 \(\mathfrak p'/\mathfrak p\) 应用\cref{b3:thm:lying-over}，得到 \(B/\mathfrak q\) 的素理想 \(\mathfrak r\)，其收缩为 \(\mathfrak p'/\mathfrak p\)。取 \(\mathfrak r\) 在 \(B\) 中的原像 \(\mathfrak q'\)，其在 \(A\) 中的收缩就是 \(\mathfrak p'/\mathfrak p\) 在 \(A\) 中的原像 \(\mathfrak p'\)，因而满足要求。
\end{proof}

从任意指定的底层原像出发，重复使用此定理，就能提升一个有限的上升素理想链。反向收缩一个严格链时，\cref{b3:thm:incomparability}保证严格性不会丢失。因此，整扩张两端的素理想链长可以相互比较。

\begin{proposition}\label{b3:prop:integral-spec-closed}
设 \(A\subseteq B\) 为交换含幺环的整扩张。它诱导的映射 \(\pi:\operatorname{Spec}B\to\operatorname{Spec}A\) 满射且为闭映射。更精确地，对理想 \(J\subseteq B\)，
\[
\pi\bigl(V_B(J)\bigr)=V_A(J\cap A).
\]
\end{proposition}
\begin{proof}
左侧的每个素理想显然包含 \(J\cap A\)。反向，对包含 \(J\cap A\) 的素理想，向整包含 \(A/(J\cap A)\subseteq B/J\) 应用\cref{b3:thm:lying-over}，便得到包含 \(J\) 的原像。取 \(J=0\) 给出满射性。
\end{proof}

% !TeX root = ../../Book3.tex
% 纲要标题：### 5.3 Going-down 定理
% 纲要位置：工作记录/计划纲要.md，第 1947 行。
% 纲要细目（写作范围）：
% * 整闭基整环等假设的明确版本
% * 证明及素理想链的应用
% * 不能把 Going-down 当作任意整扩张的性质

\section{Going-down 定理}\label{b3:sec:0503}

Going-up 从链的下端出发，向上寻找原像。若已经指定上端，能否反过来向下寻找？整性本身不足以保证这一点。本节在底环整闭、两端为整环的条件下证明肯定结论，并给出底环不整闭时失败的例子。

\subsection{整闭基整环上的 Going-down 定理}\label{b3:subsec:0503:01}

\begin{theorem}[Going-down]\label{b3:thm:going-down}
设 \(A\subseteq B\) 为交换含幺整环的整扩张，且 \(A\) 在 \(\operatorname{Frac}(A)\) 中整闭。若
\[
\mathfrak p_1\subseteq\mathfrak p_2,\qquad
\mathfrak q_2\cap A=\mathfrak p_2,
\]
其中 \(\mathfrak p_i\in\operatorname{Spec}A\)、\(\mathfrak q_2\in\operatorname{Spec}B\)，则存在素理想 \(\mathfrak q_1\subseteq\mathfrak q_2\)，使 \(\mathfrak q_1\cap A=\mathfrak p_1\)。这称为\term[going-down]{Going-down 定理}[going-down theorem]。
\end{theorem}

证明将使用整闭性控制最小多项式的系数。设 \(K=\operatorname{Frac}(A)\)，则 \(B\) 的分式域是 \(K\) 的代数扩张：\(B\) 中每个元素代数，两个代数元素的商仍代数。这一观察保证下列最小多项式确实存在。

\begin{lemma}\label{b3:lem:minimal-polynomial-integral-coefficients}
设 \(A\) 是交换含幺整闭整环，\(K=\operatorname{Frac}(A)\)，\(L/K\) 为域扩张，\(b\in L\) 在 \(A\) 上整。则 \(b\) 在 \(K\) 上的首一最小多项式属于 \(A[T]\)。
\end{lemma}
\begin{proof}
取 \(b\) 在 \(A\) 上的首一方程 \(F(b)=0\)。其在 \(K\) 上的最小多项式 \(f\) 整除 \(F\)。取 \(f\) 在 \(K\) 上的分裂域 \(\Omega\)，则其中每个根也满足 \(F\)，因而在 \(A\) 上整。把根按重数计入，\(f\) 的系数是这些根的初等对称多项式，由\cref{b3:prop:integral-closure-ring}也在 \(A\) 上整。系数原来属于 \(K\)，整闭性便给出它们属于 \(A\)。
\end{proof}

\begin{lemma}\label{b3:lem:minimal-polynomial-prime-coefficients}
设 \(A\subseteq B\) 为交换含幺整环的整扩张，\(A\) 在 \(K=\operatorname{Frac}(A)\) 中整闭。若 \(\mathfrak p\in\operatorname{Spec}A\)、\(b\in\mathfrak pB\)，则 \(b\) 在 \(K\) 上的首一最小多项式可写为
\[
f(T)=T^n+c_1T^{n-1}+\cdots+c_n,
\qquad c_i\in\mathfrak p.
\]
\end{lemma}
\begin{proof}
由\cref{b3:lem:minimal-polynomial-integral-coefficients}，\(f\in A[T]\)。写 \(b=\sum_{j=1}^r a_ju_j\)，其中 \(a_j\in\mathfrak p\)、\(u_j\in B\)。每个 \(u_j\) 都整，逐次用\cref{b3:prop:integral-finite-criterion}便知 \(M=A[u_1,\ldots,u_r]\) 为有限 \(A\) 代数。它包含 \(b\)，且 \(bM\subseteq\mathfrak pM\)。对有限组模生成元，乘 \(b\) 的表示矩阵各项均可取在 \(\mathfrak p\) 中。\cref{b3:thm:determinant-trick}因而给出首一多项式
\[
h(T)=T^m+d_1T^{m-1}+\cdots+d_m,
\qquad d_i\in\mathfrak p^i,
\qquad h(b)=0.
\]
这里 \(M\) 含有一，故算子等式确实推出环内等式。由于 \(f\) 是 \(b\) 在 \(K\) 上的最小多项式，\(f\mid h\) 于 \(K[T]\)，所以 \(f\) 的每个根也满足 \(h\)。

取 \(f\) 的一个分裂域 \(\Omega/K\)，令 \(C\) 为 \(A\) 在 \(\Omega\) 中的整闭包。由\cref{b3:thm:lying-over}取 \(\mathfrak P\in\operatorname{Spec}C\) 位于 \(\mathfrak p\) 上。\(f\) 的每个根 \(\beta\) 都在 \(A\) 上整，因而属于 \(C\)，且满足 \(h(\beta)=0\)。模掉 \(\mathfrak P\) 后有 \(\bar\beta^m=0\)，而 \(C/\mathfrak P\) 为整环，所以 \(\beta\in\mathfrak P\)。\(f\) 的首项以外系数作为这些根的对称多项式均属于 \(\mathfrak P\)，又属于 \(A\)，从而属于 \(\mathfrak P\cap A=\mathfrak p\)。把根按重数计入即可，故不要求 \(f\) 可分。
\end{proof}

\subsection{定理的证明与素理想链的应用}\label{b3:subsec:0503:02}

\begin{proof}[\cref{b3:thm:going-down}的证明]
要使所求素理想位于指定的 \(\mathfrak q_2\) 以下，先进入 \(B_{\mathfrak q_2}\)：其中的素理想收缩到 \(B\) 后都包含于 \(\mathfrak q_2\)。还要使它在 \(A\) 中的收缩恰为 \(\mathfrak p_1\)，因此须包含 \(\mathfrak p_1\)，同时避开 \(A\setminus\mathfrak p_1\)。后一次局部化能保留真理想，正好要求 \(\mathfrak p_1B_{\mathfrak q_2}\) 不碰到这个乘法集。为此先证
\begin{equation}\label{b3:eq:going-down-contraction}
\mathfrak p_1B_{\mathfrak q_2}\cap A=\mathfrak p_1.
\end{equation}
若 \(a\) 属于左侧而不在 \(\mathfrak p_1\)，则在 \(B_{\mathfrak q_2}\) 中可写成有限和
\[
\dfrac a1=\sum_{j=1}^r\dfrac{u_jb_j}{s_j},
\qquad u_j\in\mathfrak p_1,\quad b_j\in B,\quad
s_j\in B\setminus\mathfrak q_2.
\]
令 \(s=\prod_js_j\)。统一清除这些分母后，在局部环中有 \(sa=\sum_j u_jb_j\prod_{i\ne j}s_i\)。由于 \(B\) 是整环，\(B\to B_{\mathfrak q_2}\) 单射，这个等式也在 \(B\) 中成立。因此
\[
sa\in\mathfrak p_1B,\qquad s\in B\setminus\mathfrak q_2.
\]
在 \(A_{\mathfrak p_1}\subseteq(A\setminus\mathfrak p_1)^{-1}B\) 中，\(a\) 可逆，故
\(s\in\mathfrak p_1(A\setminus\mathfrak p_1)^{-1}B\)。
由\cref{b3:cor:integrally-closed-localization}，\(A_{\mathfrak p_1}\) 整闭；其分式域仍是 \(K=\operatorname{Frac}(A)\)，而 \((A\setminus\mathfrak p_1)^{-1}B\) 在 \(A_{\mathfrak p_1}\) 上整。将\cref{b3:lem:minimal-polynomial-prime-coefficients}用于此整包含，得 \(s\) 在 \(K\) 上的首一最小多项式的首项以外系数属于 \(\mathfrak p_1A_{\mathfrak p_1}\)。\cref{b3:lem:minimal-polynomial-integral-coefficients}用于原包含 \(A\subseteq B\)，说明这些系数也属于 \(A\)。因此它们属于
\[
A\cap\mathfrak p_1A_{\mathfrak p_1}=\mathfrak p_1\subseteq\mathfrak p_2.
\]
把该方程模掉 \(\mathfrak q_2\)，得到 \(s^n=0\)，于是 \(s\in\mathfrak q_2\)，矛盾。这证明了\eqref{b3:eq:going-down-contraction}。

式\eqref{b3:eq:going-down-contraction}说明理想 \(\mathfrak p_1B_{\mathfrak q_2}\) 不交 \(A\setminus\mathfrak p_1\)。把这些元素再倒置，所得理想仍为真理想，故包含在某个极大理想中。将这个极大理想依次收缩到 \(B_{\mathfrak q_2}\) 和 \(B\)，得到素理想 \(\mathfrak q_1\subseteq\mathfrak q_2\)。它包含 \(\mathfrak p_1\)，且不交 \(A\setminus\mathfrak p_1\)，故收缩恰为 \(\mathfrak p_1\)。
\end{proof}

在\cref{b3:thm:going-down}的假设下，给定 \(A\) 的素理想链
\(\mathfrak p_0\subsetneq\cdots\subsetneq\mathfrak p_r\)
及收缩为 \(\mathfrak p_r\) 的 \(B\) 中素理想 \(\mathfrak q_r\)，重复向下应用定理，就能把这条链提升为顶端为 \(\mathfrak q_r\) 的严格链。严格性来自收缩理想不同。相反，\(\mathfrak q_r\) 以下的严格链收缩后仍严格，由\cref{b3:thm:incomparability}保证。因此 \(\mathfrak q_r\) 与其收缩以下的有限素理想链，具有相同的可能最大长度；若长度无界，两边同时无界。

\begin{table}[H]
\centering
\caption[整扩张中四个素理想定理]{整扩张中四个素理想定理。设 \(A\subseteq B\) 为交换含幺环的整扩张，\(\pi:\operatorname{Spec}B\rightarrow\operatorname{Spec}A\) 为收缩映射；各 \(\mathfrak p_i\) 为 \(A\) 的素理想，各 \(\mathfrak q_i\) 为 \(B\) 的素理想。}
\label{b3:tab:integral-prime-theorems}
\begin{tabularx}{\textwidth}{@{}l>{\raggedright\arraybackslash}X>{\raggedright\arraybackslash}X@{}}
\toprule
定理 & 已指定的数据 & 结论及附加条件 \\
\midrule
Lying-over & \(\mathfrak p\in\operatorname{Spec}A\) & 存在 \(\mathfrak q\) 使 \(\pi(\mathfrak q)=\mathfrak p\) \\
Incomparability & \(\mathfrak q_1\subseteq\mathfrak q_2\)，且收缩相同 & \(\mathfrak q_1=\mathfrak q_2\) \\
Going-up & \(\mathfrak p_1\subseteq\mathfrak p_2\)，\(\pi(\mathfrak q_1)=\mathfrak p_1\) & 可找 \(\mathfrak q_2\supseteq\mathfrak q_1\)，收缩为 \(\mathfrak p_2\) \\
Going-down & \(\mathfrak p_1\subseteq\mathfrak p_2\)，\(\pi(\mathfrak q_2)=\mathfrak p_2\) & 可找 \(\mathfrak q_1\subseteq\mathfrak q_2\)，收缩为 \(\mathfrak p_1\)；本书版本还假设 \(A,B\) 为整环且 \(A\) 整闭 \\
\bottomrule
\end{tabularx}
\end{table}

\cref{b3:tab:integral-prime-theorems}中，Going-up 固定链的下端，Going-down 固定上端；前三项只用整包含，而最后一项需要额外假设。

\subsection{Going-down 定理的适用条件}\label{b3:subsec:0503:03}

取特征不为 \(2\) 的域 \(k\)，在 \(B=k[t,z]\) 内取
\[
u=t^2-1,\qquad v=t(t^2-1),\qquad A=k[u,v,z].
\]
这里 \(u,v\) 不是代数独立的变量，而有 \(v^2=u^2(u+1)\)；更确切地，
\[
A\cong k[U,V,Z]/\bigl(V^2-U^2(U+1)\bigr),
\qquad (U,V,Z)\longmapsto(u,v,z).
\]
核的计算可用关于 \(V\) 的首一关系作除法：每个多项式的余式形如 \(P(U,Z)+VQ(U,Z)\)。若代入后为零，分别比较 \(t\) 的偶次项与奇次项，得到 \(P(t^2-1,z)=0\) 与 \((t^2-1)Q(t^2-1,z)=0\)。代换 \(U=t^2-1\) 保持 \(k[z][U]\) 中的非零多项式为非零，因为其最高 \(t\) 次项不能抵消；故 \(P=Q=0\)，所写关系生成全部核。

\ProseParagraphBreak
因为 \(t\) 满足 \(T^2-(u+1)=0\)，且 \(B=A[t]\)，故 \(B\) 为有限 \(A\) 代数，且在 \(A\) 上整。令
\[
\mathfrak q_+=(t-1,z),\qquad
\mathfrak q_-=(t+1,z),\qquad
\mathfrak r=(z-t-1).
\]
\(\mathfrak q_\pm\) 收缩为同一个极大理想
\(\mathfrak p_2=(u,v,z)\)：代入 \((t,z)=(1,0)\) 或 \((-1,0)\)，\(A\) 上的取值都是把 \(u,v,z\) 送到零。令 \(\mathfrak p_1=\mathfrak r\cap A\)。由于 \(\mathfrak r\subseteq\mathfrak q_-\)，有 \(\mathfrak p_1\subseteq\mathfrak p_2\)；此包含严格，因为 \(z\notin\mathfrak r\)，从而 \(z\notin\mathfrak p_1\)。

\ProseParagraphBreak
注意 \(u\notin\mathfrak p_1\)，因为模掉 \(\mathfrak r\) 后 \(u=t^2-1\) 非零。倒置 \(u\) 后，\(t=v/u\) 已属于 \(A_u\)，所以 \(A_u=B_u\)。局部化素理想对应因而说明 \(\mathfrak p_1\) 在 \(B\) 中只有一个原像 \(\mathfrak r\)：任何原像都不含 \(u\)，局部化后的原像唯一，再收缩也唯一。但
\(z-t-1\notin\mathfrak q_+\)，因为它在 \((1,0)\) 取值 \(-2\)。
所以不存在位于 \(\mathfrak q_+\) 以下、收缩为 \(\mathfrak p_1\) 的素理想。

\ProseParagraphBreak
这组理想的选择可以用下面的参数位置与收缩关系理解。左图只示意 \((t,z)\) 平面中的位置；右图中的箭头表示素理想的包含或收缩。
\begin{center}
\begin{tikzpicture}[scale=0.9,>=stealth]
\draw[->] (-1.8,0)--(1.9,0) node[right] {\(t\)};
\draw[->] (0,-0.5)--(0,2.4) node[above] {\(z\)};
\draw[thick] (-1.4,-0.4)--(1.2,2.2) node[right] {\(z=t+1\)};
\fill (-1,0) circle (2pt) node[below left] {\((-1,0)\)};
\fill (1,0) circle (2pt) node[below] {\((1,0)\)};
\node at (0,-1.1) {两点在下方被识别};
\end{tikzpicture}\qquad
\begin{tikzcd}[column sep=2.8em,row sep=2em]
\mathfrak r \arrow[r,hook] \arrow[d,mapsto]
 & \mathfrak q_- \arrow[d,mapsto]
 & \mathfrak q_+ \arrow[dl,mapsto] \\
\mathfrak p_1 \arrow[r,hook] & \mathfrak p_2 &
\end{tikzcd}
\end{center}
直线经过 \((-1,0)\)，不经过 \((1,0)\)，对应 \(\mathfrak r\subset\mathfrak q_-\) 而 \(\mathfrak r\not\subset\mathfrak q_+\)。下方的特殊化 \(\mathfrak p_1\subsetneq\mathfrak p_2\) 虽然存在，一旦把上端指定为 \(\mathfrak q_+\)，唯一可能的下端 \(\mathfrak r\) 就无法接到它。

\ProseParagraphBreak
这里底环确实不整闭。元素 \(t=v/u\) 属于其分式域且整，却不属于 \(A\)：\(A\) 中的元素在 \((1,0)\) 与 \((-1,0)\) 取相同值，而 \(t\) 的两个值不同。因此底环的整闭性不能从这个 Going-down 定理的假设中删去。

% !TeX root = ../../Book3.tex
% 纲要标题：### 5.4 Noether 正规化
% 纲要位置：工作记录/计划纲要.md，第 2054 行。
% 纲要细目（写作范围）：
% * 域上有限生成代数的整子多项式环；与分式域内整闭包的区别
% * 统一的加权证明与无限域上的线性代换；唯一最高幂与有限域的取值障碍
% * 超越次数、整环情形及有限态射的纤维；分式域识别与有限维 Artin 纤维
% ---

\section{Noether 正规化}\label{b3:sec:0504}

有限型代数通常由带关系的生成元给出。Noether 正规化要从中找出一组没有代数关系的元素，使原代数在所选多项式子环上成为有限生成模。这一构造既为维数提供坐标，也为下一章的正规化有限性提供一个可控制的底环。

\subsection{有限生成代数的整子多项式环}\label{b3:subsec:0504:01}

\begin{theorem}[Noether 正规化]\label{b3:thm:noether-normalization}
设 \(k\) 为任意域，\(A\ne0\) 为交换含幺有限型 \(k\) 代数。存在代数独立元素 \(y_1,\ldots,y_d\in A\)，使 \(A\) 是多项式子环
\[
R=k[y_1,\ldots,y_d]\subseteq A
\]
上的有限 \(R\) 代数；这里允许 \(d=0\)，此时底环为 \(k\)。这称为\term[noether-zhengguihua]{Noether 正规化定理}[Noether normalization theorem]。
\end{theorem}

本定理不要求 \(A\) 是整环或约化环。例如
\(k[x,\varepsilon]/(\varepsilon^2)\) 作为 \(k[x]\) 模由 \(1,\varepsilon\) 生成，\(x\) 本身仍然代数独立。对整环而言，选取多项式子环与在分式域内取整闭包是两个不同的构造，\cref{b3:tab:two-normalizations}列出了它们的区别。

\begin{table}[htbp]
\centering
\caption[Noether 正规化与环的正规化]{Noether 正规化与环的正规化。设 \(k\) 为域，\(A\) 为交换含幺有限型 \(k\) 整环，\(K=\operatorname{Frac}(A)\)，\(\widetilde A\) 为 \(A\) 在 \(K\) 中的整闭包。}
\label{b3:tab:two-normalizations}
\begin{tabularx}{\textwidth}{@{}>{\raggedright\arraybackslash}X>{\raggedright\arraybackslash}X@{}}
\toprule
Noether 正规化 & 环的正规化\\
\midrule
选出多项式子环 \(R=k[y_1,\ldots,y_d]\subseteq A\) & 添入 \(K\) 中所有在 \(A\) 上整的元素，得到 \(\widetilde A\)\\
\(A\) 是有限 \(R\) 代数 & \(\widetilde A\) 在 \(A\) 上整\\
\(R\) 是正规整环，\(A\) 未必正规 & \(\widetilde A\) 是分式域为 \(K\) 的正规整环\\
保留 \(A\) 本身，选择其子环 & \(A\subseteq\widetilde A\)，且两者相等当且仅当 \(A\) 正规\\
\bottomrule
\end{tabularx}
\end{table}

\subsection{统一的加权证明与无限域上的线性代换}\label{b3:subsec:0504:02}

\begin{proof}[\cref{b3:thm:noether-normalization}的证明]
写 \(A=k[a_1,\ldots,a_n]\)，对生成元个数 \(n\) 归纳。\(n=0\) 时 \(A=k\)。若 \(a_1,\ldots,a_n\) 代数独立，则已完成构造。否则存在非零多项式 \(f\in k[X_1,\ldots,X_n]\)，使 \(f(a_1,\ldots,a_n)=0\)。若能改选前 \(n-1\) 个生成元，使这个关系成为关于 \(a_n\) 的首一方程，\(a_n\) 就整于前面新生成元组成的子代数，而 \(A\) 在该子代数上有限。这样便能把归纳应用于一个少用一个生成元的子代数。下面的代换正是为使最高次项的系数成为 \(k\) 中的非零常数。

选整数 \(N\) 大于 \(f\) 中出现的全部单个变量指数，对 \(i<n\) 令 \(m_i=N^i\)，并作代换
\[
X_i=Y_i+Y_n^{m_i}\quad(i<n),\qquad X_n=Y_n.
\]
原来指数为 \(\alpha=(\alpha_1,\ldots,\alpha_n)\) 的单项式，在新的 \(Y_n\) 中产生最高次数
\[
w(\alpha)=\alpha_n+\alpha_1N+\cdots+\alpha_{n-1}N^{n-1}.
\]
由于每个指数小于 \(N\)，不同的 \(\alpha\) 有不同的 \(N\) 进制表示，从而这些次数互异。取其中最大的一个，其最高次 \(Y_n\) 项的系数恰为原单项式的非零系数；其他项不能抵消它。例如，对二变量多项式
\[
X_1^2+X_1X_2^2+X_2,
\]
取 \(N=3\)，三项的权重分别为 \(6,5,1\)。代入 \(X_1=Y_1+Y_2^3\)、\(X_2=Y_2\) 后，这三项产生的最高 \(Y_2\) 次数分别就是 \(6,5,1\)，所以 \(Y_2^6\) 只有一个来源，不会被抵消。

回到所取的关系，将 \(f\) 除以变换后最高次项的非零系数，并仍记为 \(f\)。令
\[
g(Y_1,\ldots,Y_n)
=f(Y_1+Y_n^{m_1},\ldots,Y_{n-1}+Y_n^{m_{n-1}},Y_n).
\]
于是 \(g\) 作为 \(Y_n\) 的多项式首一，系数在 \(k[Y_1,\ldots,Y_{n-1}]\) 中。令 \(b_i=a_i-a_n^{m_i}\)，实际代入给出
\[
g(b_1,\ldots,b_{n-1},a_n)=f(a_1,\ldots,a_n)=0.
\]
这是一条 \(C=k[b_1,\ldots,b_{n-1}]\) 系数首一方程，所以 \(a_n\) 在 \(C\) 上整；又 \(a_i=b_i+a_n^{m_i}\)，故 \(A=C[a_n]\) 是有限 \(C\) 代数。\(C\) 包含 \(k\)，因而非零，且至多由 \(n-1\) 个元素作为 \(k\) 代数生成；所减少的是选定的代数生成元个数。归纳假设给出一个多项式子环 \(R=k[y_1,\ldots,y_d]\)，使 \(C\) 为有限 \(R\) 代数。将两层有限组模生成元相乘便知 \(A\) 也是有限 \(R\) 代数，证明完成。
\end{proof}

加权代换对有限域和无限域同样有效。无限域上还有次数通常较低的做法：设 \(f\) 的最高齐次部分为 \(f_m\)。非零多项式在无限域上不可能处处为零，因此可取
\[
f_m(c_1,\ldots,c_{n-1},1)\ne0.
\]
这里代入最后变量为一不会把非零齐次式变成零多项式：其余各变量的指数确定后，最后一个指数由总次数唯一确定。作代换 \(X_i=Y_i+c_iY_n\) 后，\(Y_n^m\) 的系数为上述非零值；除以这一系数，就使变换后的关系对 \(Y_n\) 首一。所用的非零多项式取值事实可对变量个数归纳：固定前面的变量使某个非零系数不消失，再避开最后一个变量的有限个根。

\ProseParagraphBreak
有限域上这一步不能照搬。例如
\(X^qY-XY^q\) 是 \(\mathbb F_q[X,Y]\) 中的非零齐次多项式，却在所有 \((a,b)\in\mathbb F_q^2\) 处取零，因为 \(a^q=a\)、\(b^q=b\)。因此不能保证选到一个域内取值使最高齐次部分非零。上面的加权证明只比较多项式的指数，不需要寻找非零取值，因而同样适用于有限域。

\subsection{超越次数与整环情形}\label{b3:subsec:0504:03}

设 \(k\) 为域，\(A\) 为交换含幺有限型 \(k\) 整环。Noether 正规化得到的多项式子环 \(R=k[y_1,\ldots,y_d]\) 满足
\[
k(y_1,\ldots,y_d)\subseteq\operatorname{Frac}(A)
\]
为有限域扩张。事实上，记 \(K=\operatorname{Frac}(R)\)、\(S=R\setminus\{0\}\)。局部化 \(S^{-1}A\cong A\otimes_RK\) 是有限维 \(K\) 代数，因而在 \(K\) 上整；它又是整环，由\cref{b3:lem:integral-field-criterion}可知它是域。由于 \(R\subseteq A\) 都是整环，\(S^{-1}A\) 自然包含于 \(\operatorname{Frac}(A)\)，且包含 \(A\)。既然它是域，\(A\) 中每个非零元素的逆也在其中，故反过来 \(\operatorname{Frac}(A)\subseteq S^{-1}A\)，两者相等。所以 \(d\) 等于此分式域在 \(k\) 上的超越次数：\(y_i\) 代数独立，整个扩域在它们生成的有理函数域上代数。这里使用第一卷定义14.3.1的超越基判据及定理14.3.6的超越次数不变性；本卷后面的维数理论将把这个数与素理想链长度联系起来。

\begin{example}\label{b3:exm:noether-normalization-curves}
设 \(k\) 为域，\(A=k[x,y]/(xy-1)\)，并以 \(x,y\) 仍记其剩余类。按 \(x\) 坐标将双曲线 \(xy=1\) 投影到 \(x\) 轴，对应的环包含为 \(k[x]\subseteq k[x,x^{-1}]=A\)。\(A\) 不是有限 \(k[x]\) 代数，因为 \(y=x^{-1}\) 不是在 \(k[x]\) 上的整元素。改用投影坐标 \(u=x+y\) 则不同：
\[
x^2-ux+1=0,\qquad y=u-x.
\]
因此 \(A=k[u][x]\) 由 \(1,x\) 作为 \(k[u]\) 模生成。给定一个 \(u\) 值，\(x\) 必须满足上面的首一二次方程，\(y\) 再由 \(y=u-x\) 确定；在任意扩域中，这至多给出两个点。换投影方向后，整性用这一条首一关系控制了剩余坐标，使坐标环在 \(k[u]\) 上有限。\(u\) 代数独立，因为在同构 \(A=k[x,x^{-1}]\) 下，若次数 \(r\) 的非零多项式在 \(x+x^{-1}\) 处取零，其最高 Laurent 次数 \(x^r\) 的系数却非零，矛盾。对尖点代数 \(k[t^2,t^3]\)，直接取 \(u=t^2\) 即得一个由 \(1,t^3\) 作为 \(k[u]\) 模生成的有限 \(k[u]\) 代数。
\end{example}

\subsection{有限态射与纤维}\label{b3:subsec:0504:04}

\begin{definition}\label{b3:def:finite-affine-map}
设 \(R,A\) 为交换含幺环，\(\varphi:R\to A\) 为保单位环同态。若 \(A\) 经 \(\varphi\) 成为有限生成 \(R\) 模，则称诱导的素谱映射 \(\operatorname{Spec}A\to\operatorname{Spec}R\) 为\term[youxian-taishe]{有限态射}[finite morphism]。
\end{definition}

这一条件要求选定有限组元素 \(h_1,\ldots,h_r\in A\)，使 \(A\) 中的每个元素都能写成 \(\sum_i\varphi(c_i)h_i\)，其中 \(c_i\in R\)。因此它控制整个坐标环的模结构。上例中双曲线到 \(x\) 轴的投影对应 \(k[x]\subset k[x,x^{-1}]\)，其素谱映射是单射，每个纤维至多一个点，坐标环却不是有限 \(k[x]\) 模。所以只数点集纤维不能判定有限态射。

\ProseParagraphBreak
Noether 正规化中 \(R=k[y_1,\ldots,y_d]\subseteq A\) 的包含诱导有限态射
\[
\operatorname{Spec}A\longrightarrow\operatorname{Spec}R.
\]
由于这个环同态单射且 \(A\) 在 \(R\) 上整，\cref{b3:thm:lying-over}保证上述态射满射。一般的有限态射也有有限的点集纤维：对定义中的任意 \(\mathfrak p\in\operatorname{Spec}R\)，纤维对应于
\[
\operatorname{Spec}\bigl(A\otimes_R\kappa(\mathfrak p)\bigr).
\]
该纤维代数在剩余域上有限维，所以理想降链中每次严格包含都会使向量空间维数下降，降链必在有限步后稳定；因此它是 Artin 环。若它为零代数，纤维为空；否则由\cref{b3:lem:artin-primes-radical}，其素理想全为极大理想，且只有有限多个。这就证明了点集纤维的有限性，所用的纤维对应正是\cref{b3:prop:fibre-spectrum}。有限态射仍可能把不同点合在一起，也可能有非约化纤维；有限性没有消除这些现象。


\begin{chaptersummary}
首一方程使一个元素的高次幂落在有限个低次幂的线性生成空间中；行列式技巧给出反向推导。因此有限代数生成与各生成元的整性合起来，才得到有限生成模。整性可传递，所有整元素组成整闭包，但整闭性总要相对于指定的扩环或分式域来理解。

\ProseParagraphBreak
整扩张的素谱映射满射且闭，同一纤维中的素理想互不包含。Going-up 不要求整闭性；从指定顶端向下提升则需要本章 Going-down 版本中的整环及整闭假设。证明中整闭性的实际用途，是把最小多项式的系数留在底环内。

\ProseParagraphBreak
Noether 正规化把任意非零有限型域代数写成多项式子环上的有限扩张。加权代换在有限域上仍能找出唯一的最高幂，因而不依赖域中可供选择的点数。这个多项式子环将在正规化有限性、超越次数和维数的比较中反复使用。
\end{chaptersummary}

\BookExercises

\begin{exercise}\label{b3:ex:05-sum-radicals}
给 \(\alpha=\sqrt2+\sqrt3\) 找一个整数系数的首一方程，并给 \(\mathbb Z[\alpha]\) 找一组至多四个元素的模生成元。
\end{exercise}
\begin{solution}
由 \(\alpha^2=5+2\sqrt6\)，得到 \((\alpha^2-5)^2=24\)，即
\(\alpha^4-10\alpha^2+1=0\)。依此递降幂次，\(1,\alpha,\alpha^2,\alpha^3\) 生成 \(\mathbb Z[\alpha]\)。这里不需要先证明这个四次多项式不可约。
\end{solution}

\begin{exercise}\label{b3:ex:05-inverting-integers}
设整数 \(n>1\)。证明 \(\mathbb Z[1/n]\) 是有限型 \(\mathbb Z\) 代数，却不是有限生成的 \(\mathbb Z\) 模。
\end{exercise}
\begin{solution}
它由 \(1/n\) 作为代数生成。若为有限生成模，\cref{b3:prop:integral-finite-criterion}将使 \(1/n\) 整，设其首一方程次数为 \(r\)。乘以 \(n^r\) 后有 \(1+a_1n+\cdots+a_rn^r=0\)，从而 \(n\mid1\)，矛盾。
\end{solution}

\begin{exercise}\label{b3:ex:05-nilpotent-perturbation}
设 \(A\subseteq B\) 为交换含幺环的包含，\(b,u\in B\)，且 \(u\) 幂零。证明 \(b\) 在 \(A\) 上整，当且仅当 \(b+u\) 在 \(A\) 上整。
\end{exercise}
\begin{solution}
若首一 \(f\in A[T]\) 满足 \(f(b)=0\)，则逐项展开得 \(f(b+u)\in uB\)。取 \(u^N=0\)，便有 \(f(b+u)^N=0\)，而 \(f(T)^N\) 仍首一，所以 \(b+u\) 整。把 \(u\) 换成 \(-u\) 得到反向结论。
\end{solution}

\begin{exercise}\label{b3:ex:05-integral-union}
设 \(A\subseteq B\) 为交换含幺环的整扩张。证明 \(B\) 是其中有限 \(A\) 子代数的有向并，并用素数 \(p\) 对应的代数扩域 \(\overline{\mathbb F}_p/\mathbb F_p\) 说明整扩张不必是有限扩张。
\end{exercise}
\begin{solution}
任取有限多个元素，它们生成的子代数由\cref{b3:prop:integral-finite-criterion}有限。两组元素的并仍有限，所以这些子代数有向，且每个元素都属于其中一个。\(\overline{\mathbb F}_p\) 中每个元素代数，因而整。第一卷定理17.1.1用 \(T^{p^r}-T\) 的根集构造 \(p^r\) 元域；在固定的代数闭包 \(\overline{\mathbb F}_p\) 内作这一构造，就得到其中的子域 \(\mathbb F_{p^r}\)。若 \([\overline{\mathbb F}_p:\mathbb F_p]=d<\infty\)，基的坐标计数将给出它只有 \(p^d\) 个元素，却又含有 \(p^{d+1}\) 元子域，矛盾。因此这个整扩张不是有限扩张。
\end{solution}

\begin{exercise}\label{b3:ex:05-localize-closure}
设 \(A\) 为交换含幺整环、\(L\) 为其分式域的扩域、\(C\) 为 \(A\) 在 \(L\) 中的整闭包，且 \(S\subseteq A\setminus\{0\}\) 为含一的乘法集。证明 \(S^{-1}C\) 是 \(S^{-1}A\) 在 \(L\) 中的整闭包。
\end{exercise}
\begin{solution}
局部化保持整性，所以 \(S^{-1}C\) 中的元素整。反向，若 \(z\) 满足 \(S^{-1}A\) 上首一方程，选公共分母 \(s\in S\)，把各系数写成 \(a_i/s\)。与\cref{b3:cor:integrally-closed-localization}的证明一样，\(sz\) 满足
\(T^n+a_1T^{n-1}+a_2sT^{n-2}+\cdots+a_ns^{n-1}\)，故 \(sz\in C\)，于是 \(z\in S^{-1}C\)。
\end{solution}

\begin{exercise}\label{b3:ex:05-quadratic-order}
证明 \(\mathbb Z[\sqrt5]\) 不整闭，并证明 \((1+\sqrt5)/2\) 所生成的环是有限生成的 \(\mathbb Z[\sqrt5]\) 模。
\end{exercise}
\begin{solution}
\(\omega=(1+\sqrt5)/2\) 满足 \(\omega^2-\omega-1=0\)，属于原环的分式域，却不写成 \(a+b\sqrt5\)、\(a,b\in\mathbb Z\)，因为 \(1,\sqrt5\) 在 \(\mathbb Q\) 上线性无关。又 \(\sqrt5=2\omega-1\)，所以 \(\mathbb Z[\omega]\) 包含原环，并由 \(1,\omega\) 在原环上生成。这只证明一个有限的整扩环，不在此宣称已算出全部整闭包。
\end{solution}

\begin{exercise}\label{b3:ex:05-primes-square-map}
设 \(k\) 为代数闭域，\(k[u]\subseteq k[t]\) 由 \(u=t^2\) 给定。对每个 \(a\in k\)，求位于 \((u-a)\) 上的全部素理想，并区分 \(\operatorname{char}k=2\) 的情形。
\end{exercise}
\begin{solution}
原像由 \(k[t]/(t^2-a)\) 的素理想给出，即 \((t-b)\)，其中 \(b^2=a\)。特征不为二且 \(a\ne0\) 时有两个不同根；\(a=0\) 时只有 \((t)\)。特征二时，代数闭性给出一个平方根 \(b\)，且 \(t^2-a=(t-b)^2\)，所以也只有一个素理想。二重因子不产生第二个点，却使商环含有幂零元。
\end{solution}

\begin{exercise}\label{b3:ex:05-integral-homeomorphism}
设 \(A\subseteq B\) 为交换含幺环的整扩张。证明若素谱收缩映射 \(\operatorname{Spec}B\to\operatorname{Spec}A\) 为双射，则它是同胚。说明双射之外还需要用到什么性质。
\end{exercise}
\begin{solution}
环同态诱导的素谱映射连续，因为闭集 \(V_A(I)\) 的原像为 \(V_B(IB)\)。\cref{b3:prop:integral-spec-closed}进一步说明此映射闭。一个连续闭双射的逆映射连续：逆像为原映射的像，而原映射把闭集送到闭集。因此是同胚。单独的连续双射并不提供这一闭性。
\end{solution}

\begin{exercise}\label{b3:ex:05-nonintegral-incomparability}
设 \(k\) 为域。在 \(k[x]\subseteq k[x,y]\) 中给出严格包含而收缩相同的两个素理想，并证明这个扩张不整。
\end{exercise}
\begin{solution}
\((0)\subsetneq(y)\) 都收缩到零，两个商环分别为 \(k[x,y]\) 和 \(k[x]\)，所以确为素理想。若 \(y\) 在 \(k[x]\) 上整，就会满足关于独立变量 \(y\) 的首一方程，和多项式表达唯一性矛盾。因此不能在这里应用\cref{b3:thm:incomparability}。
\end{solution}

\begin{exercise}\label{b3:ex:05-finite-prime-fibres}
设 \(A\subseteq B\) 为交换含幺环的包含，\(B\) 是有限 \(A\) 代数，并可由 \(r\ge1\) 个元素作为 \(A\) 模生成。证明每个 \(\mathfrak p\in\operatorname{Spec}A\) 上至多有 \(r\) 个素理想。
\end{exercise}
\begin{solution}
\cref{b3:prop:fibre-spectrum}把所求纤维识别为 \(\operatorname{Spec}C\)，其中 \(C=B\otimes_A\kappa(\mathfrak p)\)。原来的 \(r\) 个模生成元取像后仍生成 \(C\)，故 \(\dim_{\kappa(\mathfrak p)}C\le r\)。其任意素商是在剩余域上有限的整环，由\cref{b3:lem:integral-field-criterion}为域，所以每个素理想都极大。任取 \(s\) 个不同素理想 \(\mathfrak n_1,\ldots,\mathfrak n_s\)，它们两两互素，中国剩余定理给出满映射
\[
C\twoheadrightarrow\prod_{i=1}^s C/\mathfrak n_i.
\]
每个域商都是非零的 \(\kappa(\mathfrak p)\) 向量空间，至少贡献一个维数，因此
\[
s\le\sum_{i=1}^s\dim_{\kappa(\mathfrak p)}(C/\mathfrak n_i)
\le\dim_{\kappa(\mathfrak p)}C\le r.
\]
若纤维中有多于 \(r\) 个素理想，取其中 \(r+1\) 个便与此式矛盾。
\end{solution}

\begin{exercise}\label{b3:ex:05-going-down-explicit-ideal}
设 \(k\) 为特征不为二的域，在 \(B=k[t,z]\) 中置 \(u=t^2-1\)、\(v=t(t^2-1)\)，并令 \(A=k[u,v,z]\)、\(\mathfrak p_1=(z-t-1)B\cap A\)。这是\subsecref{b3:subsec:0503:03}的反例。证明
\[
\mathfrak p_1=\bigl(u-z^2+2z,\ v-(z-1)u\bigr)_A.
\]
据此再次核验 \(\mathfrak p_1\subsetneq(u,v,z)\)。
\end{exercise}
\begin{solution}
这里要把 \(B\) 中的直线方程改写成 \(A\) 内的理想生成元，才能在底环中直接比较 \(\mathfrak p_1\) 与 \((u,v,z)\)。代入 \(t=z-1\) 后，\(u=z^2-2z\)、\(v=(z-1)u\)，故所列两项都为零。记它们生成的理想为 \(J\)，则 \(J\subseteq\mathfrak p_1\)。在 \(A/J\) 中，\(u,v\) 都成为 \(z\) 的多项式，所以自然映射 \(k[z]\to A/J\) 满射；代入又给出 \(A/J\to k[z]\)，二者的复合在 \(z\) 上恒等，故这两个映射互逆。于是代入映射的核恰为 \(J\)，即 \(\mathfrak p_1=J\)。两生成元都属于 \((u,v,z)\)，而 \(z\in(u,v,z)\) 在商环 \(A/\mathfrak p_1\cong k[z]\) 中非零，因此包含严格。
\end{solution}

\begin{exercise}\label{b3:ex:05-prescribed-top}
设 \(k\) 为特征不为二的代数闭域。在 \(A=k[x,y]\subseteq B=k[s,t]\)、\(x=s^2,y=t^2\) 中，把链
\[
(0)\subsetneq(x-1)\subsetneq(x-1,y-1)
\]
提升为顶端指定为 \((s+1,t-1)\) 的链，并解释 Going-down 的假设为何满足。
\end{exercise}
\begin{solution}
可取 \((0)\subsetneq(s+1)\subsetneq(s+1,t-1)\)。中间理想的商环为 \(k[t]\)，\(A\) 上的核恰为 \((x-1)\)，因为 \(y\mapsto t^2\) 仍超越；顶端的收缩直接取值得到。\(A\) 是唯一分解整环，按\cref{b3:prop:ufd-integrally-closed}整闭，且 \(s,t\) 各满足一个首一二次方程，故 \(B/A\) 整，满足\cref{b3:thm:going-down}。
\end{solution}

\begin{exercise}\label{b3:ex:05-monic-quotient-free}
对任意非零交换含幺环 \(R\) 及次数 \(n\ge1\) 的首一多项式 \(f\in R[T]\)，证明 \(R[T]/(f)\) 是秩 \(n\) 的自由 \(R\) 模，即使 \(R\) 有零因子也成立。
\end{exercise}
\begin{solution}
首一除法使每个商类有次数小于 \(n\) 的代表，给出 \(1,\overline T,\ldots,\overline T^{n-1}\) 的生成性。若低次多项式 \(g\) 属于 \((f)\)，写 \(g=fh\)。只要 \(h\ne0\)，因 \(f\) 首一，最高项不会消失，\(\deg(fh)=n+\deg h\ge n\)，矛盾。因此代表唯一，这组元素线性无关。
\end{solution}

\begin{exercise}\label{b3:ex:05-normal-local-global}
设 \(A\) 为交换含幺整环。证明在 \(\operatorname{Frac}(A)\) 内，
\[
A=\bigcap_{\mathfrak m\in\operatorname{Max}A}A_{\mathfrak m},
\]
并据此证明 \(A\) 整闭当且仅当每个 \(A_{\mathfrak m}\) 整闭。
\end{exercise}
\begin{solution}
若 \(z\notin A\)，则 \(J=\{a\in A\mid az\in A\}\) 是真理想，取极大理想 \(\mathfrak m\supseteq J\)。若 \(z=c/s\in A_{\mathfrak m}\)，其中 \(s\notin\mathfrak m\)，则 \(sz=c\in A\)，即 \(s\in J\)，矛盾。这证明交集等式。整闭性向局部化传递已由\cref{b3:cor:integrally-closed-localization}证明；反向，若 \(z\) 在 \(A\) 上整，它在每个 \(A_{\mathfrak m}\) 上整，从而属于所有局部环，交集等式给出 \(z\in A\)。
\end{solution}

\begin{exercise}[两个 Noether 正规化]\label{b3:ex:05-two-normalizations}
设 \(k\) 为域。在 \(A=k[x,y]/(x^2+y^3-1)\) 中，分别给出 \(A\) 作为 \(k[x]\) 模和 \(k[y]\) 模的一组基，并说明两个变量各自代数独立。
\end{exercise}
\begin{solution}
把关系看成关于 \(y\) 的首一三次式，\cref{b3:ex:05-monic-quotient-free}给出自由 \(k[x]\) 模基 \(1,y,y^2\)，同时说明 \(k[x]\to A\) 单射。把关系看成关于 \(x\) 的首一二次式，得到自由 \(k[y]\) 模基 \(1,x\)，并说明 \(k[y]\to A\) 单射。单射正说明对应变量在 \(k\) 上超越。这不需要预先判定关系多项式是否不可约。
\end{solution}

\begin{exercise}\label{b3:ex:05-finite-field-weights}
令 \(A=\mathbb F_2[x,y]/(x^2y-xy^2)\)。作 \(u=x-y^3\)、\(v=y\)，写出关于 \(v\) 的首一关系，并证明 \(u\) 超越，从而给出一个 Noether 正规化。
\end{exercise}
\begin{solution}
在特征二中代入 \(x=u+v^3,y=v\)，得到
\[
v^7+v^5+u^2v+uv^2=0.
\]
原变量替换是多项式环同构，其逆为 \(u=x-y^3,v=y\)。因此 \(A\) 同构于以所列首一七次式为关系的 \(\mathbb F_2[u,v]\) 商。由\cref{b3:ex:05-monic-quotient-free}的解答，商是自由 \(\mathbb F_2[u]\) 模，基为 \(1,v,\ldots,v^6\)，故 \(\mathbb F_2[u]\) 单射且扩张有限。
\end{solution}

\begin{exercise}\label{b3:ex:05-nonnormalization-projection}
设 \(k\) 为域，\(x\) 为不定元。证明 \(k[x,x^{-1}]\) 不在 \(k[x]\) 上有限，却在 \(k[x+x^{-1}]\) 上自由且秩为二。
\end{exercise}
\begin{solution}
第一项由 \(k[x]\) 整闭而 \(x^{-1}\notin k[x]\)，结合有限扩张必整得到。令 \(u=x+x^{-1}\)。若 \(P\in k[T]\) 是次数 \(r\ge0\) 的非零多项式，则 Laurent 多项式 \(P(x+x^{-1})\) 中 \(x^r\) 的系数恰为 \(P\) 的首项系数，不可能为零。因此 \(u\) 在 \(k\) 上超越。关系 \(x^2-ux+1=0\) 及 \(x^{-1}=u-x\) 给出模生成元 \(1,x\)。若 \(P(u)+x\cdot Q(u)=0\)，其中 \(P,Q\in k[T]\)，在有理函数域中用自同构 \(x\mapsto x^{-1}\) 再写一次，相减得 \((x-x^{-1})Q(u)=0\)。\(x-x^{-1}\ne0\)，而 \(u\) 超越，所以 \(Q=0\)，继而 \(P=0\)。故所列生成元为基，此论证也适用于特征二。
\end{solution}

\begin{exercise}\label{b3:ex:05-finite-nonreduced-fibres}
设 \(k\) 为任意特征的域，\(A=k[x,y]/(y^2-x^3)\)，以 \(k[x]\) 为底环。计算 \(x=0\) 与 \(x=1\) 的纤维代数，比较其点数与约化性。
\end{exercise}
\begin{solution}
两纤维分别为 \(k[y]/(y^2)\) 与 \(k[y]/(y^2-1)\)。第一纤维总有一个素理想且非约化。第二纤维在特征不为二时由中国剩余定理同构于 \(k\times k\)，有两个点且约化；特征二时为 \(k[y]/\bigl((y-1)^2\bigr)\)，只有一个点且非约化。两者作为 \(k\) 向量空间的维数始终为二。
\end{solution}

% !TeX root = ../../Book3.tex
% 纲要标题：## 第6章　Hilbert 零点定理、正规化与导子理想
% 纲要位置：工作记录/计划纲要.md，第 2062 行。

\chapter{Hilbert 零点定理、正规化与导子理想}\label{b3:ch:06}
\BookChapterNumber{b3:ch:06}

\begin{chapterabstract}
零点定理把代数闭域上的多项式零点集与根理想联系起来，正规化则补入分式域或全分式环中的整元素。本章先由 Zariski 引理证明弱零点定理，再用 Rabinowitsch 技巧得到强形式，随后建立正规化的对象，完整证明任意域上有限型代数的正规化有限性。尖点、节点及导子理想的计算，用来区分有限性、双有理性、点集纤维与局部环的变化。
\end{chapterabstract}

\begin{learninggoals}
\begin{itemize}
\item 区分 \(L=k(\alpha_1,\ldots,\alpha_r)\) 的域生成与 \(L=k[\alpha_1,\ldots,\alpha_r]\) 的代数生成，判断极大理想的剩余域。
\item 从零点条件得到根理想成员关系，并写清 Rabinowitsch 技巧中的分母消去。
\item 对整环、约化环和非约化环采用明确的正规化约定。
\item 分开处理可分迹论证与正特征不可分扩张，指出正规化有限性证明中需要添加的有限组系数。
\item 计算曲线正规化、导子及商模，避免把点数、环同构和平坦性混为一谈。
\end{itemize}
\end{learninggoals}

在代数闭域 \(k\) 上，方程 \(x^2=0\) 与 \(x=0\) 有完全相同的零点；然而 \(k[x]/(x^2)\) 中的 \(\bar x\) 仍非零，且平方为零。只看取值点会忘记这个幂零元素。零点定理所对应的不是原理想本身，而是它的根；要把这一点推广到多个变量，先须知道极大理想是否都来自实际的取值点。

\ProseParagraphBreak

另一种失真发生在一个环已没有幂零元时：尖点的坐标环 \(k[t^2,t^3]\) 仍缺少分式域中的元素 \(t\)，虽然它满足首一方程。把所有这样的整元素补入，会得到正规化。它既保持分式域，又可能改变局部环和一点之上的纤维。因此零点集、坐标环与正规环是三个需要分别比较的对象。

\ProseParagraphBreak

本章前半使用 Noether 正规化、极大理想与根理想；后半还用到有限生成模的 Noether 性，以及第一卷的可分扩张、迹配对和纯不可分扩张。零点定理已足以支持\secref{b3:sec:0701}与\secref{b3:sec:0702}的坐标环、正则函数和态射；\secref{b3:sec:0604}的两条曲线与\secref{b3:sec:0606}的导子计算还使用具体的正规化。两条曲线的有限性可由显式模生成式直接核验。\cref{b3:prop:reduced-total-quotient-product}的分支分解则用于\secref{b3:sec:0704}的交叉直线与\chapref{b3:ch:20}的正规性判据。

\ProseParagraphBreak
Matsumura 的《Commutative Ring Theory》\cite[§5, Theorems 5.3--5.4, pp. 33--34]{Matsumura1989CommutativeRingTheory}可用于对照零点定理和 Rabinowitsch 技巧；正规化有限性的论证见该书\cite[§33, Integral closure of a Noetherian integral domain, pp. 261--263]{Matsumura1989CommutativeRingTheory}。本章把其中的不可分扩张处理和有限分支推广单独展开，逐项说明所需条件。

% !TeX root = ../../Book3.tex
% 纲要标题：### 6.1 Zariski 引理
% 纲要位置：工作记录/计划纲要.md，第 2064 行。
% 纲要细目（写作范围）：
% * 作为域的有限生成代数；域生成与代数生成、有限分母的限制
% * 极大理想的剩余域
% * 与代数闭域假设的联系；闭点扩域后的分裂

\section{Zariski 引理}\label{b3:sec:0601}

把一个有限型代数模掉极大理想，得到的剩余域依然由有限个元素作为代数生成。域的倒数运算与代数生成元的有限性相互约束，结果是这种剩余域只能是基域的有限代数扩张。Noether 正规化使这一事实有一个简洁证明。

\subsection{作为域的有限生成代数}\label{b3:subsec:0601:01}

设 \(k\) 为任意域，\(t\) 在 \(k\) 上超越。\(k(t)\) 虽由一个元素 \(t\) 作为域生成，却必须包含 \(k[t]\) 中每个非零多项式的倒数。有限个有理函数的分母乘积若为 \(D\ne0\)，则它们生成的 \(k\) 代数包含于 \(k[t,1/D]\)，其元素的分母都可写成 \(D\) 的幂。取非恒定多项式 \(1+tD\) 的一个不可约因子 \(h\)，则 \(h\nmid D\)，所以 \(h^{-1}\) 不在这个子环中：否则存在 \(P\in k[t]\) 和整数 \(m\ge0\)，使 \(D^m=hP\)，与 \(h\nmid D\) 矛盾。这说明有限个代数生成元不能提供这个超越方向上所需的全部倒数；下面的引理把这种限制推广到任意有限型域代数。

\begin{lemma}[Zariski]\label{b3:lem:zariski}
设 \(k\) 为域，\(L/k\) 为域扩张，且 \(L\) 作为 \(k\) 代数有限生成，则 \([L:k]<\infty\)。这一结论称为\term[zariski-yinli]{Zariski 引理}[Zariski's lemma]。
\end{lemma}
\begin{proof}
由\cref{b3:thm:noether-normalization}，存在多项式子环 \(R=k[y_1,\ldots,y_d]\subseteq L\)，使 \(L\) 为有限 \(R\) 代数，因而在 \(R\) 上整。因为 \(L\) 是域，\cref{b3:lem:integral-field-criterion}迫使 \(R\) 也是域。若 \(d>0\)，变量 \(y_1\) 在多项式环中非零且不可逆，矛盾。因此 \(d=0\)，\(R=k\)，而正规化结论正说明 \(L\) 是有限生成的 \(k\) 模，即有限维 \(k\) 向量空间。
\end{proof}

反过来，有限代数扩张选一个有限的向量空间基即可作为代数生成集，所以引理刻画了哪些域是有限型 \(k\) 代数。

\subsection{极大理想的剩余域}\label{b3:subsec:0601:02}

\begin{corollary}\label{b3:cor:finite-type-residue-fields}
设 \(k\) 为域，\(A\) 为交换含幺有限型 \(k\) 代数。对 \(A\) 的任一极大理想 \(\mathfrak m\)，剩余域 \(A/\mathfrak m\) 是 \(k\) 的有限扩张。对 \(A\) 的一般素理想 \(\mathfrak p\)，剩余域 \(\kappa(\mathfrak p)=\operatorname{Frac}(A/\mathfrak p)\) 则只必然是有限生成的域扩张。
\end{corollary}
\begin{proof}
\(A/\mathfrak m\) 是域，同时是 \(A\) 的有限型代数商，应用\cref{b3:lem:zariski}即可。对一般素理想，\(A/\mathfrak p\) 的有限组代数生成元也作为域生成元生成其分式域，但此时已经允许把任意非零多项式倒置，不能继续应用 Zariski 引理。
\end{proof}

例如 \(k[x,y]\) 的极大理想若取 \((x,y)\)，剩余域为 \(k\)；素理想 \((x)\) 的剩余域却是 \(k(y)\)。闭点和一般点的差别，在这里反映为剩余域的代数性与超越性。

\subsection{代数闭域假设的作用}\label{b3:subsec:0601:03}

若 \(k\) 代数闭，它没有非平凡有限代数扩张：有限扩张中的每个元素有 \(k\) 上的不可约最小多项式，而代数闭性迫使该多项式为一次。因此\cref{b3:cor:finite-type-residue-fields}给出
\[
A/\mathfrak m=k
\]
的典范 \(k\) 代数同构。这里的等号意为结构映射 \(k\to A/\mathfrak m\) 为同构，而不是选定某个任意的抽象域同构。

\ProseParagraphBreak
基域不代数闭时，有限剩余域扩张必须保留。例如 \(\mathbb R[x]/(x^2+1)\cong\mathbb C\)，对应一个实多项式环的极大理想，却不对应实数轴上的点。把系数域扩大到 \(\mathbb C\) 后，\(x^2+1=(x-i)(x+i)\)，两因子互素，中国剩余定理给出
\[
\mathbb C[x]/(x^2+1)\cong\mathbb C\times\mathbb C.
\]
原来的一个闭点由此分裂成 \(i,-i\) 两个复坐标点。另一方面，\(\mathbb F_q[x]\) 的次数 \(r\) 的首一不可约多项式给出剩余域 \(\mathbb F_{q^r}\)。下一节正是把代数闭性用在排除这些额外的剩余域上。

% !TeX root = ../../Book3.tex
% 纲要标题：### 6.2 弱零点定理
% 纲要位置：工作记录/计划纲要.md，第 2070 行。
% 纲要细目（写作范围）：
% * 代数闭域上的极大理想；多变量求值映射的核
% * 公共零点的存在判据；坐标点与素谱点的区别
% * 非代数闭域的对照

\section{弱零点定理}\label{b3:sec:0602}

一个坐标点可以由所有在该点取零的多项式描述。设 \(k\) 为任意域，\(n\ge0\)，\(a=(a_1,\ldots,a_n)\in k^n\)，考虑取值映射 \(\operatorname{ev}_a:k[x_1,\ldots,x_n]\to k\)。逐次将 \(x_1,\ldots,x_n\) 替换成 \(a_1,\ldots,a_n\)，每次作差都是关于 \(x_i\) 的多项式在 \(x_i\) 与 \(a_i\) 处的取值之差，因而被 \(x_i-a_i\) 整除。把这些差相加，对任意 \(f\) 得到
\[
f(x_1,\ldots,x_n)-f(a_1,\ldots,a_n)
=\sum_{i=1}^n(x_i-a_i)g_i(x_1,\ldots,x_n),
\qquad g_i\in k[x_1,\ldots,x_n].
\]
因此 \(\ker\operatorname{ev}_a=(x_1-a_1,\ldots,x_n-a_n)\)。取值映射在常数上是恒等映射，因而满射到域 \(k\)，所以这个核是极大理想。代数闭域上的弱零点定理将说明，每个极大理想都以这种方式产生。

\subsection{代数闭域上的极大理想}\label{b3:subsec:0602:01}

\begin{theorem}[Hilbert 弱零点定理]\label{b3:thm:weak-nullstellensatz}
设 \(k\) 是代数闭域，\(n\ge0\)。则 \(k[x_1,\ldots,x_n]\) 的每个极大理想唯一写成
\[
\mathfrak m_a=(x_1-a_1,\ldots,x_n-a_n),\qquad a\in k^n.
\]
这称为\term[ruo-lingdian-dingli]{弱零点定理}[weak Nullstellensatz]。
\end{theorem}
\begin{proof}
对极大理想 \(\mathfrak m\)，由\cref{b3:cor:finite-type-residue-fields}及代数闭性，结构映射 \(k\to k[x_1,\ldots,x_n]/\mathfrak m\) 是同构。令 \(a_i\in k\) 为 \(x_i\) 在商域中的像，则 \(x_i-a_i\in\mathfrak m\)，故 \(\mathfrak m_a\subseteq\mathfrak m\)。二者都极大，只能相等。若 \(\mathfrak m_a=\mathfrak m_b\)，将 \(x_i-a_i\) 在 \(b\) 处取值，得到 \(b_i=a_i\)，所以点唯一。
\end{proof}

\subsection{公共零点的存在判据}\label{b3:subsec:0602:02}

对理想 \(I\subseteq k[x_1,\ldots,x_n]\)，记
\[
Z_k(I)=\bigl\{\,a\in k^n\mid f(a)=0\text{ 对每个 }f\in I\,\bigr\}.
\]
\symidx{\(Z_k(I)\)：理想 \(I\) 在 \(k^n\) 中的公共零点集}
这与素谱中的 \(V(I)\) 需要区别：\(Z_k(I)\) 只记录 \(k\) 坐标点，而 \(V(I)\) 包含所有含 \(I\) 的素理想。若 \(k\) 代数闭，弱零点定理把 \(Z_k(I)\) 中的坐标点与 \(V(I)\) 中的闭点对应起来。例如取 \(I=(0)\subset k[x,y]\)，则 \(Z_k(0)=k^2\)，而 \(V(0)=\operatorname{Spec}k[x,y]\) 还包含 \((0),(x),(y-x^2)\) 等非闭点；它们的商环分别为 \(k[x,y],k[y],k[x]\)，都是整环却不是域，所以这些素理想都不极大。这些非闭点即使在 \(k\) 代数闭时也仍然存在。

\begin{corollary}\label{b3:cor:proper-ideal-common-zero}
设 \(k\) 为代数闭域，\(n\ge0\)，\(I\) 为 \(k[x_1,\ldots,x_n]\) 的理想。记 \(Z_k(I)=\{a\in k^n:f(a)=0\text{ 对所有 }f\in I\}\)。则 \(I\) 是真理想，当且仅当 \(Z_k(I)\ne\varnothing\)。特别地，给定整数 \(r\ge1\) 及 \(k[x_1,\ldots,x_n]\) 中的多项式 \(f_1,\ldots,f_r\)，它们没有 \(k^n\) 中的公共零点，当且仅当存在 \(g_1,\ldots,g_r\in k[x_1,\ldots,x_n]\) 满足
\[
g_1f_1+\cdots+g_rf_r=1.
\]
\end{corollary}
\begin{proof}
若 \(I\) 是真理想，可先取极大理想 \(\mathfrak m\supseteq I\)。\cref{b3:thm:weak-nullstellensatz}把这个纯代数的选择转为一个坐标点：\(\mathfrak m=\mathfrak m_a\)，其中 \(a\in k^n\)。于是每个 \(f\in I\) 都属于 \(\ker\operatorname{ev}_a\)，即 \(f(a)=0\)，所以 \(a\) 是公共零点。反过来，若 \(a\) 为公共零点，则 \(I\subseteq\ker\operatorname{ev}_a\)，所以 \(1\notin I\)。最后对 \(I=(f_1,\ldots,f_r)\) 使用理想生成的定义即可。
\end{proof}

例如 \(xy-1\) 和 \(x\) 没有公共零点，等式
\(yx-(xy-1)=1\) 就是可以直接核验的证据。零点定理保证任何无公共零点的有限系统都有这样的代数组合；它本身没有给出这些系数的次数上界。第一卷附录F.3的 Buchberger 算法及定理F.3.3给出一种计算办法：系数域支持精确四则运算和零判定，且选定的全局单项式序能够有效比较时，可从有限生成元计算 Gröbner 基；保存各步的生成元表示，再把 \(1\) 的约化过程回代，就能得到上述理想成员证据。

\subsection{非代数闭域上的对照}\label{b3:subsec:0602:03}

\((x^2+1)\subseteq\mathbb R[x]\) 是真理想，却没有实零点。有限域上的差别甚至会出现在零理想的零点理想中：多项式 \(x^q-x\) 在全部 \(\mathbb F_q\) 元素上取零，本身却不是零多项式。这些例子分别说明公共零点存在性及从点集恢复多项式关系，都依赖所使用的基域。

\ProseParagraphBreak
不过任何域上的真理想，在代数闭包中都有公共零点。对真理想 \(I\subseteq k[x_1,\ldots,x_n]\)，选择极大理想 \(\mathfrak m\supseteq I\)，其剩余域 \(L\) 由\cref{b3:cor:finite-type-residue-fields}是有限代数扩张。把 \(L\) 嵌入一个固定代数闭包 \(\overline k\)，变量的像给出所求点。这里有限扩域嵌入代数闭包，可按有限组代数生成元逐次选择最小多项式的根完成；不要求 \(L/k\) 可分。

\ProseParagraphBreak
若只研究实点或有限域上的有理点，就应当保留相应基域的额外关系，而不能把代数闭域版本原样照用。本卷后续的实代数补充也将据此区分不同的零点定理。

% !TeX root = ../../Book3.tex
% 纲要标题：### 6.3 强零点定理
% 纲要位置：工作记录/计划纲要.md，第 2076 行。
% 纲要细目（写作范围）：
% * 根理想与零点集的双向对应
% * Rabinowitsch 技巧；任意交换含幺环上的增广理想判据
% * 几何对象与代数对象反变关系的起点；普通坐标环的约化性

\section{强零点定理}\label{b3:sec:0603}

弱形式只判断方程组是否有解。强形式进一步回答：已知全部解之后，究竟能恢复哪些多项式关系？一个多项式与它的正整数次幂有相同零点，因此期望恢复原理想通常过强；正确对象是根理想。本节除特别说明外，\(k\) 始终为代数闭域。

\subsection{根理想与零点集的对应}\label{b3:subsec:0603:01}

设 \(k\) 为域，\(n\ge0\)，\(X\subseteq k^n\)。定义 \(X\) 的\term[lingdian-lixiang]{零点理想}[vanishing ideal]
\[
I(X)=\bigl\{\,f\in k[x_1,\ldots,x_n]\mid f(a)=0\text{ 对每个 }a\in X\,\bigr\}.
\]
这是根理想，因为 \(f(a)^r=0\) 在域中推出 \(f(a)=0\)。空集的零点理想按全称命题约定为整个多项式环。
\symidx{\(I(X)\)：点集 \(X\) 上恒为零的多项式组成的理想}

\begin{theorem}[Hilbert 强零点定理]\label{b3:thm:strong-nullstellensatz}
设 \(k\) 为代数闭域，\(n\ge0\)，\(J\) 为 \(k[x_1,\ldots,x_n]\) 的理想。记 \(Z_k(J)=\{a\in k^n:f(a)=0\text{ 对所有 }f\in J\}\)；对 \(X\subseteq k^n\)，记 \(I(X)=\{f\in k[x_1,\ldots,x_n]:f(a)=0\text{ 对所有 }a\in X\}\)。则
\[
I\bigl(Z_k(J)\bigr)=\sqrt J.
\]
这称为\term[qiang-lingdian-dingli]{强零点定理}[strong Nullstellensatz]。因此根理想与形如 \(Z_k(J)\) 的代数集，分别经 \(Z_k\) 与 \(I\) 构成反向包含的一一对应。
\end{theorem}

例如 \((x^2,y)\) 和 \((x,y)\) 在 \(k^2\) 中都只给出原点；从这个点集恢复的是根理想 \((x,y)\)。商环 \(k[x,y]/(x^2,y)\) 中保留的非零幂零元，无法由普通点集上的取值看到。下一小节证明等式，再核对对应关系。

\subsection{Rabinowitsch 技巧}\label{b3:subsec:0603:02}

把“不等于零”的条件写成“存在倒数”，是从弱形式过渡到强形式的关键。对 \(f\in k[x_1,\ldots,x_n]\)，再加变量 \(T\)，方程 \(Tf-1=0\) 正是要求 \(f\) 的值可逆。这个方法称为\term[rabinowitsch-jiqiao]{Rabinowitsch 技巧}[Rabinowitsch trick]。\footnote{原论文为《Zum Hilbertschen Nullstellensatz》，发表于 1930 年的《Mathematische Annalen》，见\cite{Rabinowitsch1930Nullstellensatz}。}

其中从增广理想回到根理想的推理，不依赖零点或代数闭性，而是下面的纯代数事实。

\begin{lemma}\label{b3:lem:rabinowitsch-equivalence}
设 \(R\) 为交换含幺环，\(I\subseteq R\) 为理想，\(f\in R\)，\(T\) 为不定元。则
\[
f\in\sqrt I
\quad\Longleftrightarrow\quad
IR[T]+(1-Tf)=R[T].
\]
\end{lemma}
\begin{proof}
若 \(f^N\in I\)、\(N\ge1\)，有限几何级数给出
\[
1=(1-Tf)\sum_{j=0}^{N-1}(Tf)^j+T^Nf^N,
\]
所以右侧理想含一。

反向，设右侧理想含一。由多项式环的泛性质，局部化映射 \(R\to R_f\) 延拓为环同态
\[
R[T]\longrightarrow R_f,\qquad T\longmapsto f^{-1}.
\]
在这个同态下，\(1-Tf\) 变成零，故单位理想的等式给出 \(1\in IR_f\)。将\cref{b3:prop:ideal-localization-saturation}用于乘法集 \(\{1,f,f^2,\ldots\}\)，得到 \(I\) 中含有某个 \(f^N\)。若 \(N=0\)，则 \(1\in I\)，可改取 \(N=1\)。因此 \(f\in\sqrt I\)。这里通过局部化的成员判据清分母，所以即使 \(R\) 有零因子，推理仍成立。
\end{proof}

\begin{proof}[\cref{b3:thm:strong-nullstellensatz}的证明]
若 \(f^r\in J\)，则在 \(Z_k(J)\) 的每个点都有 \(f(a)^r=0\)，从而 \(f(a)=0\)。因此 \(\sqrt J\subseteq I\bigl(Z_k(J)\bigr)\)。

反向，取在 \(Z_k(J)\) 上恒为零的多项式 \(f\)，在 \(k[x_1,\ldots,x_n,T]\) 中考虑
\[
J' = J\,k[x_1,\ldots,x_n,T]+(1-Tf).
\]
若 \((a,t)\) 为其零点，则 \(a\in Z_k(J)\)，所以 \(f(a)=0\)，但又要求 \(1-t\cdot f(a)=0\)，矛盾。由\cref{b3:cor:proper-ideal-common-zero}，无公共零点使 \(J'\) 成为单位理想；再由\cref{b3:lem:rabinowitsch-equivalence}，得到 \(f\in\sqrt J\)。弱零点定理在这里把几何上的无解转成单位理想，随后根理想的结论完全由纯代数引理给出。

两种运算都反向包含。刚证明对根理想 \(J\) 有 \(I\bigl(Z_k(J)\bigr)=J\)。对代数集 \(X=Z_k(J)\)，一方面 \(J\subseteq I(X)\) 给出 \(Z_k\bigl(I(X)\bigr)\subseteq X\)；另一方面每个 \(a\in X\) 按定义使 \(I(X)\) 中全部多项式为零，得到反向包含。于是两种运算在所述对象上互逆。
\end{proof}

\subsection{几何与代数的反变关系}\label{b3:subsec:0603:03}

\begin{definition}\label{b3:def:affine-coordinate-ring}
设 \(k\) 为代数闭域，\(n\ge0\)，\(X\subseteq k^n\) 为多项式方程组的公共零点集。定义商环
\[
k[X]=k[x_1,\ldots,x_n]/I(X)
\]
并称它为 \(X\) 的\term[zuobiao-huan]{坐标环}[coordinate ring]。其中一个元素就是多项式在 \(X\) 上给出的函数；两个多项式给出同一函数，当且仅当它们之差属于 \(I(X)\)。
\end{definition}
\symidx{\(k[X]\)：仿射代数集 \(X\) 的坐标环}

因为 \(I(X)\) 是根理想，普通代数集的坐标环总是约化环。若从理想 \(J\subseteq R=k[x_1,\ldots,x_n]\) 出发，先取零点集再取坐标环，强零点定理给出
\[
k[Z_k(J)]=R/\sqrt J.
\]
这未必是 \(R/J\)：例如前面的 \(J=(x^2,y)\) 得到原点的坐标环 \(k\)，而 \(R/J\) 中的非零幂零元 \(\bar x\) 已经丢失。

设 \(X,Y\subseteq k^n\) 为代数集。强零点定理说明，代数集的包含与零点理想的包含反向对应：
\[
X\subseteq Y\ \Longleftrightarrow\ I(Y)\subseteq I(X).
\]
当 \(X\subseteq Y\) 时，把 \(Y\) 上的多项式函数限制到 \(X\)，就得到满同态 \(k[Y]\rightarrow k[X]\)，其核是 \(I(X)/I(Y)\)。因此
\[
k[X]\cong k[Y]\big/\bigl(I(X)/I(Y)\bigr).
\]
例如 \(Y=k^2\)、\(X=Z_k(y)\) 时，限制函数把 \(k[x,y]\) 送到 \(k[x]\)，正是令 \(y=0\) 的商同态。代数集的包含由此表现为函数环的商，而不是子环。

\ProseParagraphBreak
推广到一般映射时，需要判断函数与映射复合后是否仍属于相应的函数环；在局部允许分母的情形，还要考察这些表达式在各点是否有定义。\secref{b3:sec:0701}与\secref{b3:sec:0702}以此研究代数集的拓扑、正则函数与态射。\secref{b3:sec:0604}的尖点、节点和\secref{b3:sec:0606}的导子计算，则进一步比较零点集与环本身保留的信息。

% !TeX root = ../../Book3.tex
% 纲要标题：### 6.4 正规化
% 纲要位置：工作记录/计划纲要.md，第 2082 行。
% 纲要细目（写作范围）：
% * 域上有限生成代数的正规化及其有限性问题：先处理整环，再处理约化环，非约化情形约定先取约化；定义与模型在本节建立，一般有限性定理在紧接的第6.5节陈述并证明
% * 正规化不有限的情形；缺失常数所造成的无限剩余域维数
% * 尖点和节点曲线的代数例子
% * 整闭包、有限性和双有理性的不同角色；本节的尖点、节点直接通过显式环包含核验，不依赖尚未证明的一般有限性定理

\section{正规化}\label{b3:sec:0604}

同一个有理函数域可能包含不同的坐标环。尖点的 \(k[t^2,t^3]\) 与直线的 \(k[t]\) 就有相同分式域，但前者漏掉了一个已经满足首一方程的函数 \(t\)。正规化把这样的整元素一并加入。它改变的是环；点集是否改变、扩张是否有限，需要分别核验。

\subsection{整环、约化环与正规化的定义}\label{b3:subsec:0604:01}

\begin{definition}\label{b3:def:normalization}
设 \(A\) 为交换含幺环。若 \(A\) 是整环，则称它在 \(\operatorname{Frac}(A)\) 中的整闭包 \(\widetilde A\) 为 \(A\) 的\term[zhengguihua]{正规化}[normalization]。若 \(A\) 是约化环，令 \(S\subseteq A\) 为其全部非零因子的集合，定义\term[quan-fenshi-huan]{全分式环}[total ring of fractions]
\(Q(A)=S^{-1}A\)，并把 \(A\) 在 \(Q(A)\) 中的整闭包称为其正规化。对于非约化环，本书先取 \(A_{\mathrm{red}}=A/\sqrt0\)，再正规化 \(A_{\mathrm{red}}\)，不将幂零元保留在正规化对象中。
\end{definition}
\symidx{\(Q(A)\)：约化环的全分式环}
\symidx{\(\widetilde A\)：环 \(A\) 的正规化}

非零因子乘法封闭，而且 \(A\to Q(A)\) 单射：\(a/1=0\) 意味着某个非零因子 \(s\) 满足 \(sa=0\)，从而 \(a=0\)。当 \(A\) 是整环时，全分式环就是普通分式域，因此两种定义一致。若 \(A\) 约化，则 \(A\subseteq\widetilde A\) 是整扩张。对一般 \(A\)，自然映射为
\[
A\longrightarrow A_{\mathrm{red}}\longrightarrow
\widetilde{A_{\mathrm{red}}}.
\]
第二个箭头单射，所以复合映射的核为 \(\sqrt0\)。目标环中的每个元素都在 \(A_{\mathrm{red}}\) 上整，也就是在 \(A\) 的像上整；因此这是整的环同态，但当 \(A\) 非约化时不是环包含。无论哪种情形，定义本身都不保证正规化是有限生成的 \(A\) 模。

\ProseParagraphBreak
例如，令 \(C=k[x,y]/(xy)\)。把一个商类送到它在两条坐标轴上的限制，得到单射
\[
C\longrightarrow k[x]\times k[y],\qquad
[F]\longmapsto\bigl(F(x,0),F(0,y)\bigr).
\]
每个商类唯一写成 \(c+xp(x)+yq(y)\)，所以其像恰为常数项相等的多项式对，这也说明 \(C\) 约化。这样的对 \((f,g)\) 是非零因子，当且仅当 \(f,g\) 都非零：两者非零时可在两个整环中分别约去；若 \(f=0\)，则乘以非零元素 \((x,0)\) 得到零，若 \(g=0\)，则乘以 \((0,y)\) 得到零。因此上述映射延拓为单射 \(Q(C)\to k(x)\times k(y)\)。对任意 \(F/G\in k(x)\) 和 \(P/H\in k(y)\)，其中 \(G,H\ne0\)，有
\[
\frac{(xF,0)}{(xG,y)}\longmapsto(F/G,0),\qquad
\frac{(0,yP)}{(x,yH)}\longmapsto(0,P/H).
\]
这里分子、分母都是常数项为零的多项式对，且分母的两个坐标都非零。把这两类分式相加便得到任意有理函数对，故 \(Q(C)\cong k(x)\times k(y)\)。

在这个全分式环中，\(e=(1,0)\) 满足首一方程 \(e^2-e=0\)。而且 \(C[e]=k[x]\times k[y]\)：任意 \((f,g)\) 都等于
\(e(f,f(0))+(1-e)(g(0),g)\)，右侧两对多项式属于 \(C\)。所以 \(k[x]\times k[y]\) 在 \(C\) 上整。若一个有理函数对在 \(k[x]\times k[y]\) 上整，投影到每个坐标便得到该坐标在相应多项式环上的首一方程；\(k[x],k[y]\) 都整闭，故两个坐标均为多项式。任何在 \(C\) 上整的有理函数对也在这个上环上整，因此属于它。这证明 \(k[x]\times k[y]\) 恰为 \(C\) 的正规化。原环要求两个分支上的函数在原点取相同值，正规化后这项限制消失。

对一般约化 Noether 环，极小素理想只有有限多个，并有
\[
Q(A)\cong\prod_{\mathfrak p\in\operatorname{Min}(A)}
\operatorname{Frac}(A/\mathfrak p).
\]
这一分支描述将在\cref{b3:prop:reduced-total-quotient-product}证明。对于整环，全分式环是域，正规化是其中的整环。

\subsection{非有限正规化的例子}\label{b3:subsec:0604:02}

\begin{example}\label{b3:exm:nonfinite-normalization}
取无限次代数扩域 \(K/k\)，例如 \(\overline{\mathbb Q}/\mathbb Q\)。允许正次数项的系数来自 \(K\)，却把常数项限制在 \(k\)，便得到
\[
A=k+xK[x]\subseteq B=K[x].
\]
则 \(\operatorname{Frac}(A)=K(x)\)，正规化为 \(B\)，但 \(B\) 不是有限生成的 \(A\) 模；此处 \(A\) 不是 Noether 环。
\end{example}
\begin{proof}
对每个 \(c\in K\)，有 \(c=(cx)/x\in\operatorname{Frac}(A)\)，故分式域为 \(K(x)\)。\(K\) 中每个元素在 \(k\) 上代数，因而在 \(A\) 上整；而 \(B\) 由 \(A\) 与这些常数生成，所以\cref{b3:prop:integral-closure-ring}说明 \(B/A\) 整。又 \(B=K[x]\) 是唯一分解整环，按\cref{b3:prop:ufd-integrally-closed}在 \(K(x)\) 中整闭。任何在 \(A\) 上整的 \(z\in K(x)\) 也在 \(B\) 上整，所以属于 \(B\)。这就证明了正规化等于 \(B\)。

若 \(B\) 是有限生成的 \(A\) 模，模掉同时为 \(A,B\) 理想的 \(xK[x]\)，就会得到 \(K=B/xK[x]\) 是有限维 \(A/xK[x]=k\) 向量空间，与 \([K:k]=\infty\) 矛盾。最后，理想 \(\mathfrak m=xK[x]\) 满足
\(\mathfrak m/\mathfrak m^2\cong K\)，其上 \(A\) 的作用经 \(k\) 因子化；若 \(\mathfrak m\) 有限生成，此商将是有限维 \(k\) 空间，仍矛盾。因此 \(A\) 确实不是 Noether 环。
\end{proof}

这个例子说明没有额外条件时正规化未必有限；其底环是非 Noether 环。域上有限型代数拥有更强的结构；\cref{b3:thm:normalization-finite}将证明这一范围内的正规化有限，任意基域都允许。

\subsection{尖点与节点曲线的正规化}\label{b3:subsec:0604:03}

\begin{proposition}\label{b3:prop:cusp-normalization}
设 \(k\) 为任意域。映射 \(x\mapsto t^2\)、\(y\mapsto t^3\) 给出同构
\[
k[x,y]/(y^2-x^3)\cong k[t^2,t^3].
\]
此整环的分式域是 \(k(t)\)，正规化为 \(k[t]\)，并且 \(k[t]\) 由 \(1,t\) 作为 \(k[t^2,t^3]\) 模生成。
\end{proposition}
\begin{proof}
按首一多项式 \(y^2-x^3\) 除法，每个商类有代表 \(P(x)+y\cdot Q(x)\)。代入后为 \(P(t^2)+t^3\cdot Q(t^2)\)，前一项只含偶次幂，后一项只含不小于三的奇次幂，故等于零恰好要求 \(P=Q=0\)。因此核就是所列主理想，像按定义为 \(k[t^2,t^3]\)。

元素 \(t=t^3/t^2\) 在分式域中，且满足 \(T^2-t^2=0\)，所以 \(k[t]=A+At\) 在 \(A=k[t^2,t^3]\) 上有限整。\(k[t]\) 在 \(k(t)\) 中整闭，任何在 \(A\) 上整的分式也在 \(k[t]\) 上整，因而属于 \(k[t]\)。于是它正是整闭包。
\end{proof}

这里 \(1,t\) 是模生成元而不是模基，因为 \(t^2\cdot t-t^3\cdot1=0\) 是一个非平凡的 \(A\) 系数关系。两项模生成式证明了有限性，而分式域相等说明双有理性；相应的分式域扩张次数为一，并不是模生成元的个数二。

\begin{proposition}\label{b3:prop:node-normalization}
设 \(k\) 为特征不等于 \(2\) 的域。映射
\[
x\longmapsto t^2-1,\qquad y\longmapsto t(t^2-1)
\]
给出同构
\[
k[x,y]/\bigl(y^2-x^2(x+1)\bigr)
 \cong A=k\bigl[t^2-1,t(t^2-1)\bigr].
\]
其分式域为 \(k(t)\)，正规化为 \(B=k[t]\)，且 \(B=A+At\)。这称为本章的\term[jiedian]{节点}[node]模型。
\end{proposition}
\begin{proof}
参数值 \(t=1\) 和 \(t=-1\) 都给出 \(x=y=0\)，所以这两个不同参数值落在同一个原点。只要 \(x\ne0\)，便可由 \(t=y/x\) 恢复参数。

模掉关于 \(y\) 的首一关系后，每个商类写为 \(P(x)+y\cdot Q(x)\)。像为
\[
P(t^2-1)+t(t^2-1)\cdot Q(t^2-1).
\]
第一项只含偶次幂，第二项只含奇次幂。两者之和为零时分别为零；由于 \(t^2-1\) 在 \(k[t]\) 中超越且非零，得到 \(P=Q=0\)。这证明核的断言。分式 \(y/x=t\) 给出分式域相等，也说明 \(A_x=B_x\)。方程 \(t^2=x+1\) 给出 \(B=A+At\)，最后由 \(k[t]\) 的整闭性判定它就是正规化。
\end{proof}

特征不为二用于保证 \(t=1\) 与 \(t=-1\) 是两个不同点；上面的核及整闭包计算本身不需除以二。这个节点的原环还有函数值描述
\[
A=\bigl\{\,f\in k[t]\mid f(1)=f(-1)\,\bigr\}.
\]
\Cref{b3:prop:node-conductor}将证明这一等式，并由此计算正规化商模和导子理想。

\subsection{整闭包、有限性与双有理性}\label{b3:subsec:0604:04}

设 \(A\subseteq B\) 为交换含幺整环的包含。若诱导的分式域映射 \(\operatorname{Frac}(A)\to\operatorname{Frac}(B)\) 是同构，则称 \(B/A\) 为\term[shuangyouli-kuozhang]{双有理扩张}[birational extension]；相应的仿射映射称为双有理。正规化按定义双有理且整，但有限性是额外结论。例如前面的非 Noether 例仍双有理而不有限。

\ProseParagraphBreak
反过来，有限整扩张不必双有理：\(k[t^2]\subseteq k[t]\) 虽然由 \(1,t\) 生成，却诱导二次分式域扩张 \(k(t)/k(t^2)\)。即使特征为二、底层点集可能不能区分平方根，这个域扩张仍有次数二。双有理也不意味着有限，例如 \(k[t]\subseteq k[t,t^{-1}]\) 分式域相同，却把零点处的分母倒置；\(t^{-1}\) 不在整闭环 \(k[t]\) 上整，所以此扩张不有限。

\ProseParagraphBreak
尖点与节点的有限性在\cref{b3:prop:cusp-normalization}及\cref{b3:prop:node-normalization}中都体现为两项模生成式。一般有限型代数未必有这样显式的参数表达；下一节借 Noether 正规化和有限域扩张的结构证明其正规化仍然有限。

% !TeX root = ../../Book3.tex
% 纲要标题：### 6.5 正规化有限性的证明
% 纲要位置：工作记录/计划纲要.md，第 1989 行。
% 纲要细目（写作范围）：
% * 用 Noether 正规化把有限型整环的整闭包问题归约为多项式环在有限分式域扩张中的整闭包问题
% * 对有限可分扩张建立迹配对和公共分母控制引理，利用多项式环的整闭性与 Noether 性得到有限生成
% * 正特征下单独建立不可分情形所需的有限性引理；通过有限个系数和纯不可分扩张的代数处理完成证明，不把可分迹论证无条件套用，也不默设基域完美
% * 从整环推广到约化有限型代数的有限个极小素分支；明确全分式环中的正规化与各分支整闭包的关系
% * 本节接续第6.4节的正规化对象，准确陈述并完整证明有限性定理；不以附录中的一般 Nagata 环定理代替所需引理
% * Nagata 性、卓越性与正规化有限性的适用范围

\section{正规化有限性的证明}
\label{b3:sec:0605}

上一节定义了正规化，并在曲线中找到了有限个新生成元。一般有限型代数是否也只需添有限个元素，不能由定义直接得出。本节的关键是控制整元素的分母：可分扩张中用迹配对，不可分扩张中先进行一个只涉及有限个系数的纯不可分变基。整个论证允许基域不完美。这一有限性论证可对照 \cite[\S33, Integral closure of a Noetherian integral domain, pp. 261--263]{Matsumura1989CommutativeRingTheory}；下文把其中的不可分变基与分支处理分别展开。

\begin{figure}[htbp]
\centering
\begin{tikzcd}[column sep=6em,row sep=2.5em]
\text{有限型 }k\text{ 整环}
  \arrow[r,"\text{Noether 正规化}"]
  & A=k[x_1,\ldots,x_d]
    \arrow[dl,"\text{有限可分扩张}"']
    \arrow[d,"\text{有限不可分扩张}"] \\
\text{迹配对与公共分母}\arrow[dr]
  & \text{有限纯不可分变基}
    \arrow[l,"\text{化为可分}"'] \\
  & \text{整环情形的整闭包有限}\arrow[d] \\
  & \text{有限型约化环逐极小素分支处理}
\end{tikzcd}
\caption[正规化有限性的证明结构]{正规化有限性的证明结构。基域 \(k\) 任意，记 \(K=\operatorname{Frac}(A)\)，两条分支分别处理有限扩张 \(L/K\) 的可分与不可分情形。}
\label{b3:fig:normalization-finiteness-proof}
\end{figure}

\cref{b3:fig:normalization-finiteness-proof}中，迹配对始终用于有限可分扩张；不可分情形先作有限变基，再通过有限模的传递及 Noether 子模性质回到原环。约化环则需要全分式环的极小素分支分解。\secref{b3:sec:0704}的交叉直线计算及\chapref{b3:ch:20}的 Serre 正规性判据也用到这一分解。\footnote{塞尔（Jean-Pierre Serre，1926-），法国数学家，研究代数拓扑、代数几何与数论。}

\subsection{由 Noether 正规化归约整闭包问题}
\label{b3:subsec:0605:01}

设 \(C\) 为域 \(k\) 上的有限型整环。由\cref{b3:thm:noether-normalization}，存在代数独立的 \(x_1,\ldots,x_d\in C\)，使
\[
A=k[x_1,\ldots,x_d]\subseteq C
\]
为有限扩张。记 \(K=\operatorname{Frac}(A)\)、\(L=\operatorname{Frac}(C)\)。\(C\) 的有限组代数生成元在 \(K\) 上代数，且生成 \(L\)，因此 \([L:K]<\infty\)。

\ProseParagraphBreak
更一般地，若 \(L'/L\) 是任意有限域扩张，则 \(C\) 在 \(L'\) 中的整闭包等于 \(A\) 在 \(L'\) 中的整闭包。一个方向由 \(A\subseteq C\) 直接得到；另一个方向用 \(C/A\) 整和\cref{b3:prop:integral-transitivity}。所以只须证明：多项式环 \(A\) 在任意有限分式域扩张中的整闭包是有限生成的 \(A\) 模。若这一步完成，有限组 \(A\) 模生成元也生成同一个 \(C\) 模，因为 \(A\subseteq C\)。

\subsection{可分扩张中的迹配对与公共分母}
\label{b3:subsec:0605:02}

\begin{lemma}\label{b3:lem:normalization-separable-finite}
设 \(A\) 为交换含幺整闭 Noether 整环，\(K=\operatorname{Frac}(A)\)，\(L/K\) 为有限可分扩张。则 \(A\) 在 \(L\) 中的整闭包 \(B\) 是有限生成的 \(A\) 模。
\end{lemma}
\begin{proof}
先取一组 \(K\) 基 \(v_1,\ldots,v_n\)；乘以适当的非零 \(A\) 元素，可使各基元素对 \(A\) 整。确切地，若
\[
v_i^m+c_1v_i^{m-1}+\cdots+c_m=0,\qquad c_j\in K,
\]
选非零 \(d_i\in A\) 清除全部 \(c_j\) 的分母，则 \(u_i=d_iv_i\) 满足首一方程
\[
u_i^m+d_ic_1u_i^{m-1}+\cdots+d_i^mc_m=0
\]
且其系数均在 \(A\) 内。缩放后 \(u_1,\ldots,u_n\) 仍为 \(K\) 基。

对任一整元素 \(z\in L\)，域迹 \(\operatorname{Tr}_{L/K}(z)\) 属于 \(A\)。事实上，第一卷定理18.2.3把可分扩张的迹写成各 \(K\) 嵌入像之和；每个像都满足 \(z\) 的同一个首一 \(A\) 系数方程，故都对 \(A\) 整。它们的和也整，而迹又属于 \(K\)，由 \(A\) 整闭得此结论。

令
\[
T=\bigl(\operatorname{Tr}_{L/K}(u_i u_j)\bigr)_{1\le i,j\le n},
\qquad \Delta=\det T.
\]
各 \(u_i u_j\) 都整，故 \(T\) 的条目在 \(A\) 中。第一卷定理18.3.3说明有限可分扩张的迹配对非退化，所以 \(\Delta\ne0\)。给定 \(b\in B\)，写 \(b=\sum_i\lambda_i u_i\)、\(\lambda_i\in K\)。各 \(bu_j\) 仍整，于是
\[
T(\lambda_1,\ldots,\lambda_n)^{\mathsf T}
=\bigl(\operatorname{Tr}_{L/K}(bu_1),\ldots,
        \operatorname{Tr}_{L/K}(bu_n)\bigr)^{\mathsf T}\in A^n.
\]
用伴随矩阵乘上式，得每个 \(\Delta\lambda_i\in A\)。因此
\begin{equation}\label{b3:eq:normalization-trace-denominator}
B\subseteq\Delta^{-1}\sum_{i=1}^nAu_i.
\end{equation}
因为 \(u_1,\ldots,u_n\) 是 \(K\) 基，它们也在 \(A\) 上线性无关，故 \(\sum_iAu_i\cong A^n\)。乘以非零标量 \(\Delta^{-1}\) 不改变这个模的同构类型，所以右端是有限自由 \(A\) 模，左端由\cref{b3:prop:integral-closure-ring}是一个 \(A\) 子模。因 \(A\) 为 Noether 环，\cref{b3:prop:noether-finite-modules}保证其子模 \(B\) 仍有限生成。
\end{proof}

这一步真正使用可分性的地方是 \(\Delta\ne0\)。在非平凡纯不可分扩张中，域迹恒为零，所有这样的迹矩阵都退化；不能在此情形下仍写 \(\Delta^{-1}\)。下一小节通过有限变基把问题移到可以使用这个引理的地方。

\subsection{正特征不可分扩张中的有限性引理}
\label{b3:subsec:0605:03}

\begin{lemma}\label{b3:lem:finite-purely-inseparable-base-change}
设 \(k\) 为特征 \(p>0\) 的域，\(d\ge0\)，\(x_1,\ldots,x_d\) 为代数独立元，\(A=k[x_1,\ldots,x_d]\)，\(K=\operatorname{Frac}(A)\)，\(L/K\) 为有限扩张。在同一个代数闭包内，可以选取整数 \(e\ge0\)、\(q=p^e\)、有限纯不可分扩张 \(k'/k\) 及元素 \(y_j\) 满足 \(y_j^q=x_j\)，使
\[
K'=k'(y_1,\ldots,y_d),\qquad L'=LK'
\]
满足 \(L'/K'\) 有限可分。并且 \(A'=k'[y_1,\ldots,y_d]\) 是多项式环，是 \(A=k[x_1,\ldots,x_d]\) 上的有限生成模。
\end{lemma}
\begin{proof}
取有限组域生成元 \(z_1,\ldots,z_s\)，使 \(L=K(z_1,\ldots,z_s)\)。将每个首一不可约最小多项式写成
\[
f_i(T)=g_i(T^{q_i}),\qquad q_i=p^{e_i},\qquad g_i'(T)\ne0.
\]
这可以通过反复除去全部指数的公因子 \(p\) 实现。\(g_i\) 仍不可约，否则代入 \(T^{q_i}\) 会分解 \(f_i\)；因而 \(g_i\) 可分。

取 \(q=p^e\) 为所有 \(q_i\) 的公倍数。将各 \(g_i\) 的有限多个系数写成 \(x_1,\ldots,x_d\) 的有理函数，只涉及有限多个 \(k\) 中的分子、分母多项式系数，设其集合为 \(E\)。令
\[
k'=k\bigl(\,c^{1/q}:c\in E\,\bigr),\qquad y_j=x_j^{1/q}.
\]
各次根在固定代数闭包中唯一，\(k'/k\) 有限纯不可分，且 \((k')^q\subseteq k\)。因此 \(y_1,\ldots,y_d\) 在 \(k'\) 上代数独立：若有非零多项式关系 \(\sum_\alpha b_\alpha y^\alpha=0\)，将其取 \(q\) 次幂，Frobenius 加法性给出 \(\sum_\alpha b_\alpha^q x^\alpha=0\)。这里每个 \(b_\alpha^q\in k\)，而 Frobenius 在域 \(k'\) 上单射，至少一个非零系数的 \(q\) 次幂仍非零；这便是 \(x_1,\ldots,x_d\) 在 \(k\) 上的非零关系，矛盾。

每个 \(g_i\) 的系数 \(c\in K\) 都有 \(q_i\) 次根属于 \(K'\)。确切地，写 \(c=P(x)/Q(x)\)，其中 \(Q\ne0\)，并将 \(P\) 按多重指数写成 \(P(x)=\sum_\alpha a_\alpha x^\alpha\)。各 \(a_\alpha\) 的 \(q\) 次根已添入 \(k'\)，所以
\[
a_\alpha^{1/q_i}
  =\bigl(a_\alpha^{1/q}\bigr)^{q/q_i}\in k',
\qquad
P_i^\sharp(y)
  =\sum_\alpha a_\alpha^{1/q_i}y^{(q/q_i)\alpha}.
\]
Frobenius 加法性给出 \((P_i^\sharp(y))^{q_i}=P(x)\)。对 \(Q\) 同样构造 \(Q_i^\sharp(y)\)，则 \((Q_i^\sharp(y))^{q_i}=Q(x)\ne0\)，所以它可以作分母，且
\[
c^{1/q_i}=\dfrac{P_i^\sharp(y)}{Q_i^\sharp(y)}\in K'.
\]
写 \(g_i(T)=\sum_jc_{ij}T^j\)，令
\[
h_i(T)=\sum_jc_{ij}^{1/q_i}T^j\in K'[T].
\]
则 \(h_i(T)^{q_i}=f_i(T)\)，故 \(h_i(z_i)=0\)。而 \(h_i\) 可分：若 \(w\) 同时零化 \(h_i\) 与 \(h_i'\)，取 \(q_i\) 次幂后即得
\[
g_i(w^{q_i})=0,\qquad g_i'(w^{q_i})=0,
\]
其中导数系数中的整数属于素域，取 \(q_i\) 次幂后不变。这与 \(g_i\) 可分矛盾。\(z_i\) 在 \(K'\) 上的最小多项式整除 \(h_i\)，而可分多项式的不可约因子仍可分，故各 \(z_i\) 在 \(K'\) 上可分。它们共同生成的 \(L'/K'\) 也是有限可分扩张。

最后取 \(k'/k\) 的有限基 \(c_1,\ldots,c_t\)。因为 \(y_j^q=x_j\)，\(A'\) 由有限组元素
\[
c_l y_1^{a_1}\cdots y_d^{a_d},
\qquad 1\le l\le t,\quad 0\le a_j<q,
\]
作为 \(A\) 模生成。\(d=0\) 或所有 \(q_i=1\) 时，上述构造分别采用空变量组或 \(q=1\)，同样成立。
\end{proof}

只取“实际用到的有限多个系数”的次根，是这个引理中的有限性所在。任意不完美域的扩张 \(k^{1/p}/k\) 可能无限，因而直接把全部 \(p\) 次根都添进去，反而会失去这里需要的有限生成模控制。

\begin{theorem}\label{b3:thm:polynomial-normalization-finite}
设 \(k\) 为任意域，\(d\ge0\)，\(A=k[x_1,\ldots,x_d]\)，\(K=\operatorname{Frac}(A)\)。对任意有限扩张 \(L/K\)，\(A\) 在 \(L\) 中的整闭包是有限生成的 \(A\) 模。
\end{theorem}
\begin{proof}
多项式环 \(A\) 是 Noether 唯一分解整环，因而由\cref{b3:prop:ufd-integrally-closed}整闭。若 \(L/K\) 可分，直接应用\cref{b3:lem:normalization-separable-finite}。特征零的情形已全部包括在内。

在正特征中，按\cref{b3:lem:finite-purely-inseparable-base-change}构造 \(A'\)、\(K'\)、\(L'\)。\(A'\) 也是整闭 Noether 多项式环，而 \(L'/K'\) 可分，所以 \(A'\) 在 \(L'\) 中的整闭包 \(B'\) 有限于 \(A'\)。由于 \(A'\) 又有限于 \(A\)，\(B'\) 是有限生成的 \(A\) 模。原来每个对 \(A\) 整的 \(b\in L\) 也对 \(A'\) 整，因此原整闭包 \(B\) 包含于 \(B'\)。\(B\) 是 \(A\) 子模，\(A\) 的 Noether 性于是保证 \(B\) 有限。
\end{proof}

\subsection{约化代数的极小素分支与全分式环}
\label{b3:subsec:0605:04}

整环只有一个分式域；约化环可能有多个极小素分支，故应在全分式环中同时处理各分支。下面先给出此前使用的乘积描述的证明。

\begin{proposition}\label{b3:prop:reduced-total-quotient-product}
设 \(C\ne0\) 为交换含幺约化 Noether 环，极小素理想为 \(\mathfrak p_1,\ldots,\mathfrak p_r\)，\(K_i=\operatorname{Frac}(C/\mathfrak p_i)\)。则自然映射给出同构
\begin{equation}\label{b3:eq:reduced-total-quotient-product}
Q(C)\cong\prod_{i=1}^rK_i.
\end{equation}
若 \(B_i\) 为 \(C/\mathfrak p_i\) 在 \(K_i\) 中的整闭包，则 \(C\) 在 \(Q(C)\) 中的整闭包为 \(\prod_iB_i\)。
\end{proposition}
\begin{proof}
因为 \(C\) 同时约化且 Noether，\cref{b3:prop:reduced-ring-no-embedded-primes}给出 \(\operatorname{Ass}_C(C)=\min\operatorname{Spec}C\)。再由\cref{b3:thm:associated-zero-divisors}，\(C\) 的非零因子组成
\[
S=C\setminus\bigcup_i\mathfrak p_i.
\]
约化性给出 \(\bigcap_i\mathfrak p_i=0\)，故 \(C\rightarrow\prod_iK_i\) 单射。每个 \(s\in S\) 在所有 \(K_i\) 中非零，因此该映射延伸为单射 \(Q(C)=S^{-1}C\rightarrow\prod_iK_i\)。

为证明满射，对每个 \(i\) 选
\[
c_i\in\bigcap_{j\ne i}\mathfrak p_j\setminus\mathfrak p_i.
\]
这样的 \(c_i\) 可由逐个取 \(a_j\in\mathfrak p_j\setminus\mathfrak p_i\) 后相乘得到；\(r=1\) 时取 \(c_1=1\)。给定第 \(i\) 个域中的分式 \(\bar a/\bar b\)，其中 \(a,b\in C\)、\(b\notin\mathfrak p_i\)，令
\[
s=bc_i+\sum_{j\ne i}c_j.
\]
它模 \(\mathfrak p_i\) 等于非零的 \(bc_i\)，模每个 \(\mathfrak p_j\)、\(j\ne i\) 等于非零的 \(c_j\)，所以 \(s\in S\)。分式 \(ac_i/s\) 的第 \(i\) 个坐标就是 \(\bar a/\bar b\)，其余坐标全零。逐坐标相加即可得到任意元组，满射得证。

若 \(b\in Q(C)\) 对 \(C\) 整，投影它的首一方程可见每个坐标 \(b_i\) 属于 \(B_i\)。反之，若各 \(b_i\in B_i\)，为每个 \(i\) 取一个零化 \(b_i\) 的首一 \((C/\mathfrak p_i)\) 系数多项式，并把系数提升为首一 \(f_i(T)\in C[T]\)。对元组 \(b=(b_i)_i\)，多项式 \(\prod_i f_i(T)\) 在每个坐标上都有一个为零的因子 \(f_i(b_i)=0\)，故 \(\prod_i f_i(b)=0\)。这个首一方程证明 \(b\) 对 \(C\) 整。
\end{proof}

特别地，全分式环中的不同分支彼此独立，正规化也逐分支进行。这里的有限个分支来自 Noether 性；约化性则保证环嵌入这些分式域的乘积。若原环非约化，本书按\cref{b3:def:normalization}的约定先取约化商，再谈这一正规化。

\subsection{正规化有限性定理的完整证明}
\label{b3:subsec:0605:05}

\begin{theorem}\label{b3:thm:normalization-finite}
设 \(k\) 为任意域。
\begin{enumerate}
\item 若 \(C\) 为交换含幺有限型 \(k\) 整环，\(L'/\operatorname{Frac}(C)\) 为有限域扩张，则 \(C\) 在 \(L'\) 中的整闭包是有限生成的 \(C\) 模。特别地，\(C\) 的正规化有限。
\item 若 \(C\) 为交换含幺约化有限型 \(k\) 代数，则它在全分式环中的正规化是有限生成的 \(C\) 模。
\end{enumerate}
\end{theorem}
\begin{proof}
第一项取\cref{b3:thm:noether-normalization}给出的有限扩张 \(A=k[x_1,\ldots,x_d]\subseteq C\)。\subsecref{b3:subsec:0605:01}已经核验：\(L'/\operatorname{Frac}(A)\) 有限，且 \(A\) 与 \(C\) 在 \(L'\) 中的整闭包相同。\cref{b3:thm:polynomial-normalization-finite}保证此整闭包有限于 \(A\)，故也有限于 \(C\)。

第二项中，\(C=0\) 时结论直接成立。否则 \(C\) 是 Noether 环，有有限个极小素理想 \(\mathfrak p_i\)。每个 \(C/\mathfrak p_i\) 都是有限型 \(k\) 整环，由第一项其正规化 \(B_i\) 为有限生成的 \(C/\mathfrak p_i\) 模，从而也是有限生成的 \(C\) 模。\cref{b3:prop:reduced-total-quotient-product}把 \(C\) 的正规化识别为有限乘积 \(\prod_iB_i\)，所以它仍是有限生成的 \(C\) 模。
\end{proof}

可分性只是证明中的工具，最终定理允许任意基域及有限分式域扩张。上述有限性也为下一节导子理想中的有限环包含提供依据。

\subsection{Nagata 性、卓越性与正规化的有限性}\label{b3:subsec:0605:06}
\label{b3:subsec:0604:05}

上述证明先利用有限型性选出多项式子环 \(A\)，使原整环在 \(A\) 上有限；分式域扩张的有限性再保证只需处理有限组域生成元及其最小多项式中的有限个系数。一般 Noether 整环未必有这样的多项式子环，其正规化也未必有限。

\ProseParagraphBreak
后续资料中常见的\term[nagata-huan]{Nagata 环}[Nagata ring]，\footnote{永田雅宜（假名：ながた まさよし，罗马音：Masayoshi Nagata，1927-2008），日本数学家，研究交换代数与代数几何，构造 Hilbert 第十四问题的反例。}按一种标准定义，是满足下列条件的 Noether 环 \(A\)：对每个素理想 \(\mathfrak p\)，以及 \(\operatorname{Frac}(A/\mathfrak p)\) 的每个有限域扩张 \(L\)，\(A/\mathfrak p\) 在 \(L\) 中的整闭包都是有限生成的 \(A/\mathfrak p\) 模。\cref{b3:thm:normalization-finite}对这些商整环及有限域扩张均适用，因而域上有限型代数都是 Nagata 环。\term[zhuoyue-huan]{卓越环}[excellent ring]还包含关于形式纤维、维数链及正则性的条件，其定义及基本性质见\secref{b3:sec:D04}。

% !TeX root = ../../Book3.tex
% 纲要标题：### 6.6 导子理想与曲线正规化的比较
% 纲要位置：工作记录/计划纲要.md，第 1997 行。
% 纲要细目（写作范围）：
% * 对有限双有理扩张 \(A\subseteq B\) 定义导子理想 \(\mathfrak c=\{a\in A:aB\subseteq A\}\)，说明其同时为 \(A\) 和 \(B\) 的理想
% * 在 \(A=k[t^2,t^3]\subseteq B=k[t]\) 中计算 \(\mathfrak c=t^2k[t]\)、\(B/A\) 与奇点处的差异
% * 以特征不为2的曲线 \(y^2=x^2(x+1)\) 为节点模型，使用 \(x=t^2-1\)、\(y=t(t^2-1)\) 核验正规化及奇点上的两个原像
% * 比较有限双有理映射、点集上的行为和局部环的改变；平坦性与非分歧性的独立检验留到第23.4节
% ---

\section{导子理想与曲线正规化的比较}\label{b3:sec:0606}

正规化中哪些元素原来就能相乘而不离开旧环？这个问题由导子理想回答。它既是旧环的理想，也是新环的理想，因此可以直接比较两侧的局部化与商模。本节继续使用尖点和节点的显式环，不以一般奇点理论解释尚未计算的现象。

\subsection{有限双有理扩张中的导子理想}\label{b3:subsec:0606:01}

\begin{definition}\label{b3:def:conductor}
设 \(A\subseteq B\) 为交换含幺环的包含。定义其\term[daozi-lixiang]{导子理想}[conductor ideal]
\[
\mathfrak c_{B/A}
 =\{\,a\in A\mid aB\subseteq A\,\}
 =\operatorname{Ann}_A(B/A).
\]
本节主要考虑整环的有限双有理扩张。
\end{definition}
\symidx{\(\mathfrak c_{B/A}\)：扩张 \(A\subseteq B\) 的导子理想}

\begin{proposition}\label{b3:prop:conductor-basic}
设 \(A\subseteq B\) 为交换含幺环的包含，\(\mathfrak c_{B/A}=\{a\in A:aB\subseteq A\}\) 为导子理想。则 \(\mathfrak c_{B/A}\) 是包含在 \(A\) 中的最大 \(B\) 理想，因而同时为 \(A,B\) 的理想。若整环扩张 \(B/A\) 有限且双有理，则存在非零 \(d\in A\) 使 \(dB\subseteq A\)。对每个 \(c\in\mathfrak c_{B/A}\)，都有 \(A_c=B_c\)。
\end{proposition}
\begin{proof}
若 \(a\in\mathfrak c_{B/A}\)、\(b\in B\)，则 \(ab\in A\)，且 \((ab)B\subseteq aB\subseteq A\)，所以 \(ab\) 也在导子理想中。任意 \(B\) 理想 \(J\subseteq A\) 满足 \(jB\subseteq J\subseteq A\)，故 \(J\subseteq\mathfrak c_{B/A}\)。

有限双有理情形选 \(B\) 的有限组 \(A\) 模生成元 \(b_i=u_i/v_i\)，其中 \(u_i,v_i\in A\)、\(v_i\ne0\)。令 \(d=\prod_i v_i\)，便有 \(db_i\in A\)，于是 \(dB\subseteq A\) 且 \(d\ne0\)。最后对 \(b\in B\)，有 \(b=(cb)/c\in A_c\)，这证明 \(B_c\subseteq A_c\)，反向包含由原包含得到。
\end{proof}

对有限扩张，差异还可以用支集精确表述：
\[
\operatorname{Supp}_A(B/A)=V_A(\mathfrak c_{B/A}).
\]
这里 \(B/A\) 是有限生成模，故\cref{b3:thm:support-finite}把其支集识别为零化子的零点集。因而一个素理想 \(\mathfrak p\) 不含导子理想，当且仅当 \(A_{\mathfrak p}\to B\otimes_AA_{\mathfrak p}\) 是同构；这只是局部环比较，不是光滑性判据。

\subsection{尖点的导子理想与正规化商模}\label{b3:subsec:0606:02}

\begin{proposition}\label{b3:prop:cusp-conductor}
设 \(k\) 为域，\(A=k[t^2,t^3]\subseteq B=k[t]\)。则
\[
A=k\oplus t^2k[t],\qquad
B/A\cong k\cdot[t],\qquad
\mathfrak c_{B/A}=t^2k[t]=(t^2,t^3)_A.
\]
其中 \(B/A\) 上 \(A\) 的作用经 \(A/(t^2,t^3)\cong k\) 因子化。
\end{proposition}
\begin{proof}
所有不小于二的整数都由二和三作非负整数线性组合，故 \(A\) 恰由没有一次项的多项式组成。于是商空间由 \([t]\) 生成，且 \(t^2,t^3\) 都零化它。若 \(a=c+t^2\cdot h(t)\in A\)，则
\(at=ct+t^3\cdot h(t)\) 属于 \(A\) 当且仅当 \(c=0\)。由于 \(B=A+At\)，条件 \(aB\subseteq A\) 等价于 \(at\in A\)，故导子为 \(t^2k[t]\)。最后任意次数至少四的单项式可除以 \(t^2\) 而仍属于 \(A\)，再加上 \(t^2,t^3\) 两项，得到其作为 \(A\) 理想的生成式。
\end{proof}

导子只在尖点极大理想 \(\mathfrak m=(t^2,t^3)\) 处消失，离开该点两环局部化相同。该点上只有一个原像 \((t)\subseteq k[t]\)，但纤维代数是
\[
B/\mathfrak mB=k[t]/(t^2),
\]
并不是 \(k\)。所以“只有一个原像点”仍会遗漏非约化纤维。局部化后的差商
\(B_{(t)}/A_{\mathfrak m}\) 仍为一维 \(k\) 空间；这里两个局部化嵌在 \(k(t)\) 中，并且 \(B\otimes_AA_{\mathfrak m}=B_{(t)}\)。后一个等式可由\cref{b3:lem:integral-field-criterion}与局部化素理想对应看出：左环的极大理想收缩为 \(A_{\mathfrak m}\) 的唯一极大理想，因此左环也只有极大理想 \((t)\)，因而所有不属于 \((t)\) 的多项式都成为单位。

\subsection{节点正规化与奇点上的两个原像}\label{b3:subsec:0606:03}

\begin{proposition}\label{b3:prop:node-conductor}
设 \(k\) 为特征不等于 \(2\) 的域，令
\[
A=k\bigl[t^2-1,t(t^2-1)\bigr]\subseteq B=k[t].
\]
则
\[
A=k+(t^2-1)k[t]
 =\bigl\{\,f\in k[t]\mid f(1)=f(-1)\,\bigr\},
\]
并有
\[
B/A\cong k,\qquad
\mathfrak c_{B/A}=(t^2-1)k[t]=(x,y)_A,
\]
其中 \(x=t^2-1\)、\(y=t(t^2-1)\)。原点的两个原像为 \((t-1)\) 和 \((t+1)\)。
\end{proposition}
\begin{proof}
记 \(h=t^2-1\)。若 \(j\) 为偶数，则 \(ht^j=h(h+1)^{j/2}\in k[h]\)；若 \(j\) 为奇数，则
\(ht^j=ht(h+1)^{(j-1)/2}\in k[h,ht]\)。
故 \(hB\subseteq A\)，而 \(A\) 的生成元都属于 \(k+hB\)，得到第一项等号。

任意 \(f\in k[t]\) 除以 \(h\) 的余式为 \(a+bt\)。条件 \(f(1)=f(-1)\) 等价于 \(2b=0\)，由特征假设等价于 \(b=0\)，这证明第二项等号。线性映射 \(f\mapsto f(1)-f(-1)\) 满射到 \(k\)，核为 \(A\)，所以 \(B/A\) 一维；\(A\) 的作用经公共值 \(a(1)=a(-1)\) 给出。

因为 \(hB\subseteq A\)，导子包含 \(hB\)。若 \(a\) 在导子中，则 \(a\in A\) 且 \(at\in A\)，从而
\(a(1)=a(-1)\) 与 \(a(1)=-a(-1)\) 同时成立。因二可逆，二者均为零，所以 \(h\mid a\)。这证明导子等于 \(hB\)。前面对 \(ht^j\) 的表达又说明其作为 \(A\) 理想由 \(h,ht\) 生成，即 \((x,y)\)。

最后 \((x,y)B=hB\)，而 \(h=(t-1)(t+1)\) 有两个不同线性因子。商环及中国剩余定理给出
\[
B/(x,y)B\cong k[t]/(t^2-1)\cong k\times k.
\]
因此其两个素理想正给出所述两个原像。
\end{proof}

\subsection{有限双有理映射的点集与局部环}\label{b3:subsec:0606:04}

尖点和节点的正规化都有限且双有理，差商的 \(k\) 维数也都为一；它们仍以不同方式改变特殊点。尖点只增加同一点附近缺失的函数，节点则把同一个原点上的两种参数值分开。前者纤维为 \(k[t]/(t^2)\)，后者纤维为 \(k\times k\)，所以点的个数与纤维的幂零结构必须同时保留。

\ProseParagraphBreak
对节点令 \(\mathfrak m=(x,y)\)、\(S=A\setminus\mathfrak m\)。局部正规化
\[
B\otimes_AA_{\mathfrak m}=S^{-1}k[t]
\]
有两个极大理想，分别来自 \((t-1)\) 与 \((t+1)\)，是一个半局部整环。它并不是两个局部环的直积：左侧仍嵌在域 \(k(t)\) 中，而非零环的二项直积有零因子。再分别在这两个极大理想处局部化，才得到
\[
k[t]_{(t-1)},\qquad k[t]_{(t+1)}.
\]
中国剩余分解出现在模掉 \(\mathfrak m\) 的纤维中，不能提前搬到整个局部正规化环上。

\ProseParagraphBreak
在两例中，导子理想之外的局部化都是同构；在特殊点，有限性与双有理性则不足以决定平坦性和非分歧性。它们需要对模或微分再作独立检验，留到\secref{b3:sec:2304}进行。本章没有把正规化称为这些更强性质的实例。


\begin{chaptersummary}
有限型 \(k\) 代数的极大理想 \(\mathfrak m\) 具有有限次剩余域扩张，即 \([A/\mathfrak m:k]<\infty\)。当 \(k\) 代数闭时，结构映射 \(k\to A/\mathfrak m\) 就是同构，于是闭点由坐标值表示；强零点定理进一步说明，普通零点集能恢复根理想，不能恢复被幂零元记录的加厚信息。

\ProseParagraphBreak
正规化是在全分式环中取整闭包，非约化对象先取约化。域上有限型代数的有限性证明先借 Noether 正规化得到多项式底环：可分部分由非退化迹配对控制公共分母，不可分部分则只添有限多个系数的根及有限个参数的根。约化代数再按有限个极小素分支合并；没有假设基域完美。

\ProseParagraphBreak
对有限双有理整环扩张 \(A\subseteq B\)，导子是包含于 \(A\) 的最大 \(B\) 理想，也是 \(A\) 模 \(B/A\) 的零化子。尖点与节点的差商都只有一维，前者在特殊点给出非约化的单点纤维，后者给出两个点。节点的局部正规化仍是半局部整环，不能把纤维的直积分解当成整个局部环的直积分解。
\end{chaptersummary}

\BookExercises

\begin{exercise}\label{b3:ex:06-rational-field-algebra}
设 \(k\) 为域，\(t\) 为变量，\(0\ne D\in k[t]\)，令 \(R=k[t,1/D]\)。把 \(D\) 写成 \(c\prod_{i=1}^s p_i^{e_i}\)，其中 \(c\in k^\times\)、\(e_i\ge1\)，\(p_i\) 是互不相同的首一不可约多项式；\(D\) 为常数时取 \(s=0\)。求 \(R\) 的全部单位，并证明对任意不可约多项式 \(h\)，有 \(1/h\in R\) 当且仅当 \(h\mid D\)。说明为何 \(R\) 是有限型 \(k\) 代数且 \(\operatorname{Frac}(R)=k(t)\)，但 \(R\ne k(t)\)。
\end{exercise}
\begin{solution}
每个 \(p_i\) 在 \(R\) 中可逆，因为 \(1/p_i=(D/p_i)/D\)。若 \(F/D^n\) 为单位，写其逆为 \(G/D^m\)，便有 \(FG=D^{n+m}\)。由 \(k[t]\) 的唯一分解，\(F\) 的不可约因子只能是各个 \(p_i\)。因此
\[
R^\times=\left\{a\prod_{i=1}^s p_i^{n_i}:a\in k^\times,\ n_i\in\mathbb Z\right\}.
\]
同样，\(1/h=F/D^n\) 给出 \(hF=D^n\)，故 \(h\mid D\)；反向已由 \(1/h=(D/h)/D\) 得到。

环 \(R\) 由 \(t,1/D\) 代数生成。又 \(k[t]\subseteq R\subseteq k(t)\)，取分式域得到 \(\operatorname{Frac}(R)=k(t)\)。然而 \(1+tD\) 是非常数多项式，且与 \(D\) 互素；取它的一个不可约因子 \(h\)，便有 \(1/h\notin R\)。所以 \(R\) 不是域。代数生成只允许对所选生成元作多项式运算，取分式域则允许对每个非零元素取逆，二者因此给出不同的对象。
\end{solution}

\begin{exercise}\label{b3:ex:06-finite-field-closed-points}
列出 \(\mathbb F_2[x]\) 中剩余域次数为一、二、三的全部极大理想。
\end{exercise}
\begin{solution}
一次不可约式为 \(x,x+1\)；二次为 \(x^2+x+1\)；三次为
\(x^3+x+1,x^3+x^2+1\)。二次或三次多项式可约时必有一次因子，所以只需排除在零或一取零的首一多项式，便得到完整清单。对应剩余域分别有 \(2,4,8\) 个元素。后两类不是 \(\mathbb F_2\) 有理点，但都是素谱中的闭点。
\end{solution}

\begin{exercise}\label{b3:ex:06-real-residue-field}
证明 \(\mathfrak m=(x^2+1,y-x)\) 是 \(\mathbb R[x,y]\) 的极大理想，并求其在 \(\mathbb C^2\) 中的零点。
\end{exercise}
\begin{solution}
先用 \(y=x\) 消元，商环为 \(\mathbb R[x]/(x^2+1)\cong\mathbb C\)，所以极大。它没有实零点；复数零点恰为 \((i,i)\) 与 \((-i,-i)\)。这说明一个基域上的闭点在扩域后可以分裂为多个坐标点。
\end{solution}

\begin{exercise}\label{b3:ex:06-unit-ideal-certificate}
证明 \(xy-1\) 与 \(y^2\) 在任意域上都没有公共零点，并给出单位理想的显式证据。
\end{exercise}
\begin{solution}
等式
\[
x^2y^2-(xy+1)(xy-1)=1
\]
成立，因此两者生成单位理想。若有公共零点，在该点取值将给出 \(0=1\)。此处已经给出证据，所以不需代数闭域假设。
\end{solution}

\begin{exercise}\label{b3:ex:06-empty-real-zero-set}
证明 \((x^2+y^2+1)\subseteq\mathbb R[x,y]\) 是真理想，却没有实零点，并给一个复零点。
\end{exercise}
\begin{solution}
非零的非常数多项式不是 \(\mathbb R[x,y]\) 的单位，因为非零乘积的全次数相加，所以所列主理想是真理想。实数上 \(x^2+y^2+1>0\)，无零点；复点 \((i,0)\) 满足方程。不能由实零点集为空推出单位理想。
\end{solution}

\begin{exercise}\label{b3:ex:06-fat-point-radical}
设 \(k\) 为域，令 \(I=(x^2,xy,y^3)\subseteq k[x,y]\)。计算 \(\sqrt I\) 和 \(k[x,y]/I\) 的一组 \(k\) 基；若 \(k\) 代数闭，说明点集遗漏了什么信息。
\end{exercise}
\begin{solution}
因为 \(x^2,y^3\in I\)，有 \((x,y)\subseteq\sqrt I\)。又 \(I\subseteq(x,y)\)，后者极大且为根理想，所以 \(\sqrt I=(x,y)\)。模掉 \(I\) 后只剩单项式 \(1,x,y,y^2\)；理想由单项式生成，这四项中没有任何非零线性组合属于 \(I\)，故为基。若 \(k\) 代数闭，普通零点集是 \(\{(0,0)\}\)，其坐标环为 \(k\)；原商环 \(k[x,y]/I\) 则有四维，且含有非零幂零元，这些信息不能由普通点集恢复。
\end{solution}

\begin{exercise}\label{b3:ex:06-hyperbola-vanishing}
设 \(k\) 代数闭，\(X=Z_k(xy-1)\)。证明 \(I(X)=(xy-1)\)，并计算 \(k[X]\)。
\end{exercise}
\begin{solution}
映射 \(k[x,y]\to k[t,t^{-1}]\)、\(x\mapsto t,y\mapsto t^{-1}\) 诱导商环同构：反向映射把 \(t\) 送到 \([x]\)，其逆元为 \([y]\)。商环为整环，所以 \((xy-1)\) 是素理想，因而为根理想。\cref{b3:thm:strong-nullstellensatz}于是给出 \(I(X)=(xy-1)\)，坐标环即 \(k[t,t^{-1}]\)。
\end{solution}

\begin{exercise}\label{b3:ex:06-unions-intersections}
对代数闭域上的代数集 \(X,Y\subseteq k^n\)，证明
\[
I(X\cup Y)=I(X)\cap I(Y),\qquad
I(X\cap Y)=\sqrt{I(X)+I(Y)}.
\]
\end{exercise}
\begin{solution}
多项式在并集上恒为零当且仅当在两集合上分别为零，给出第一式。第二式中，和理想的公共零点恰为 \(X\cap Y\)，因为 \(Z_k\bigl(I(X)\bigr)=X\)、\(Z_k\bigl(I(Y)\bigr)=Y\)；再应用\cref{b3:thm:strong-nullstellensatz}得到其零点理想为该和的根。
\end{solution}

\begin{exercise}\label{b3:ex:06-squaring-pullback}
设 \(k\) 为代数闭域。考察点集映射 \(k\to k\)、\(a\mapsto a^2\)，以及它诱导的坐标环同态 \(k[x]\to k[t]\)、\(x\mapsto t^2\)。证明点集映射满射，坐标环同态单射但不是满射；在特征二中，点集映射甚至双射。
\end{exercise}
\begin{solution}
代数闭性保证每个数有平方根，故点集映射满射。环同态为 \(P(x)\mapsto P(t^2)\)，不同次数产生不同的偶次幂，故单射；像只含偶次幂，所以不含 \(t\)，不是满射。特征二时 \(a^2=b^2\) 推出 \((a-b)^2=0\)，域中于是 \(a=b\)。因此点集双射也不能替代坐标环同构。
\end{solution}

\begin{exercise}\label{b3:ex:06-rabinowitsch-equivalence}
设 \(k\) 为任意域，\(R=k[x,y]\)、\(I=(x^2,xy,y^3)\)。在 \(R[T]\) 中令
\[
J_x=IR[T]+(1-Tx),\qquad J_{1+x}=IR[T]+(1-T(1+x)).
\]
给出 \(1\in J_x\) 的显式多项式组合；再证明 \(J_{1+x}\) 是真理想，并说明在 \(R/I\) 中倒置 \(x\) 与倒置 \(1+x\) 分别发生了什么。
\end{exercise}
\begin{solution}
因为 \(x^2\in I\)，有
\[
1=(1-Tx)(1+Tx)+T^2x^2,
\]
这就是第一项所需的单位理想证据。另一方面，取值 \((x,y,T)=(0,0,1)\) 使 \(J_{1+x}\) 的所有生成元为零，因此它是真理想，这一取值在任意域中都可作出。

在 \(C=R/I\) 中，\(\bar x^2=0\)，故将 \(\bar x\) 变成单位会迫使 \(C_{\bar x}=0\)。而 \(1+\bar x\) 已是单位，其逆为 \(1-\bar x\)，所以 \(C_{1+\bar x}=C\ne0\)。两个增广理想分别记录了这两种倒置的结果；\cref{b3:lem:rabinowitsch-equivalence}把它们与 \(x\in\sqrt I\)、\(1+x\notin\sqrt I\) 对应起来。
\end{solution}

\begin{exercise}\label{b3:ex:06-finite-field-vanishing}
设 \(q\) 为素数幂，\(n\ge1\) 为整数。证明在 \(\mathbb F_q[x_1,\ldots,x_n]\) 中，
\[
I(\mathbb F_q^n)=(x_1^q-x_1,\ldots,x_n^q-x_n).
\]
\end{exercise}
\begin{solution}
右侧的生成元在每个点上为零，故包含于左侧。反向，依次用这些关于不同变量的首一多项式除法，得到余式 \(r\)，其中每个变量的次数都小于 \(q\)，且 \(r\) 仍在全部点上为零。对变量数归纳：\(n=1\) 时非零次数小于 \(q\) 的多项式不可能有 \(q\) 个根；一般情形固定前 \(n-1\) 个变量，将 \(r\) 视为最后一个变量的多项式，所有系数在这些固定值上为零，再由归纳假设得系数多项式为零。因此 \(r=0\)，原多项式属于右侧理想。
\end{solution}

\begin{exercise}\label{b3:ex:06-nonreduced-normalization}
设 \(k\) 为域。按本书约定求 \(A=k[t,\varepsilon]/(\varepsilon^2)\) 的正规化，并说明为何 \(A\) 到正规化的自然映射不是单射。
\end{exercise}
\begin{solution}
写元素为 \(P(t)+\varepsilon\cdot Q(t)\)。幂零元恰好为 \((\varepsilon)\)：若一个幂为零，模掉 \(\varepsilon\) 后 \(P(t)\) 在整环 \(k[t]\) 中幂零，只能为零；反向该理想平方为零。因此 \(A_{\mathrm{red}}=k[t]\)，后者已整闭，正规化还是 \(k[t]\)。自然映射的核为非零理想 \((\varepsilon)\)。
\end{solution}



\begin{exercise}\label{b3:ex:06-three-axes}
设 \(k\) 为域，\(A=k[x,y,z]/(xy,xz,yz)\)。求全分式环、正规化、正规化商模的 \(k\) 维数及导子理想。
\end{exercise}
\begin{solution}
把一个元素送到三条坐标轴上的限制，得到单射
\[
A\longrightarrow B=k[x]\times k[y]\times k[z].
\]
确切地说，\(A\) 的元素唯一写成
\(c+x\cdot f(x)+y\cdot g(y)+z\cdot h(z)\)，其像恰为三个常数项相等的多项式组。因此环约化，极小素理想是三条轴对应的 \((y,z),(x,z),(x,y)\)。\cref{b3:prop:reduced-total-quotient-product}给出全分式环
\(k(x)\times k(y)\times k(z)\)。

在 \(Q(A)\) 中添入三个坐标幂等元 \(e_1,e_2,e_3\)，它们满足 \(e_i^2-e_i=0\)，且生成 \(B\)，所以 \(B/A\) 整。\(B\) 的三个分量各自整闭，任何在 \(A\) 上整的分数组逐分量属于这些多项式环。因此正规化为 \(B\)。

商模由三常数项模去对角常数组成，维数为二。一个 \(a\in A\) 若要满足 \(aB\subseteq A\)，其公共常数必须为零，因为分别乘 \(e_i\) 后三个常数仍须相等；反过来常数全零时乘任意 \(B\) 元素仍满足要求。因此导子为 \((x,y,z)_A\)，在 \(B\) 中即 \(xk[x]\times yk[y]\times zk[z]\)。
\end{solution}

\begin{exercise}\label{b3:ex:06-normal-base-birational}
设 \(A\subseteq B\) 为整环的整双有理扩张，且 \(A\) 整闭。证明 \(A=B\)。说明此处不需有限性。
\end{exercise}
\begin{solution}
双有理使 \(B\) 嵌在 \(\operatorname{Frac}(A)\) 中，每个 \(b\in B\) 又在 \(A\) 上整。按整闭的定义，\(b\in A\)。因此原包含为等号，仅用了逐元素整性。
\end{solution}

\begin{exercise}\label{b3:ex:06-infinite-conductor}
设 \(K/k\) 为无限次代数扩域，\(A=k+xK[x]\subseteq B=K[x]\)，如\cref{b3:exm:nonfinite-normalization}所示。计算导子与 \(B/A\)，并说明导子非零不保证扩张有限。
\end{exercise}
\begin{solution}
正次数部分都在 \(A\)，所以 \(B/A\cong K/k\) 作为 \(k\) 向量空间，且 \(xK[x]\) 零化此商。若导子中的元素常数项 \(c\in k\) 非零，取 \(\lambda\in K\setminus k\)，则与 \(\lambda\) 相乘的常数项 \(c\lambda\notin k\)，不再属于 \(A\)，矛盾。因此导子恰为 \(xK[x]\)，非零；但正文已由无限剩余域维数证明扩张不有限。
\end{solution}

\begin{exercise}\label{b3:ex:06-three-four-cusp}
设 \(k\) 为域，令 \(A=k[t^3,t^4]\)。证明其正规化为 \(B=k[t]\)，计算 \(B/A\) 的维数与导子。
\end{exercise}
\begin{solution}
\(t=t^4/t^3\) 在分式域中，且 \(t^3\in A\)，所以 \(B=A+At+At^2\) 有限；\(B\) 整闭使其成为正规化。由三、四生成的非负整数包括 \(0,3,4\) 及全部不小于六的整数，遗漏 \(1,2,5\)，因此 \(B/A\) 的基为 \([t],[t^2],[t^5]\)，维数三。

显然 \(t^6B\subseteq A\)，所以它包含于导子。若 \(a\in A\) 有非零常数项，则 \(at\) 的一次项非零；若常数项为零而三次项非零，则 \(at^2\) 的五次项非零；若前两项为零而四次项非零，则 \(at\) 的五次项非零。这些情况都不在导子中。因此导子中元素的这三个系数都为零，而 \(A\) 中其余可能出现的幂次都至少为六，故这些元素都属于 \(t^6B\)。于是导子恰为 \(t^6B=(t^6,t^7,t^8)_A\)。
\end{solution}

\begin{exercise}\label{b3:ex:06-two-five-cusp}
设 \(k\) 为域。对 \(A=k[t^2,t^5]\subseteq B=k[t]\)，求正规化商模与导子，并比较正文尖点例子。
\end{exercise}
\begin{solution}
\(t=t^5/(t^2)^2\)，又满足 \(T^2-t^2=0\)，故正规化为 \(k[t]\)。可出现的幂次是零、全部正偶数和不小于五的奇数，遗漏一、三，因此 \(B/A\) 有基 \([t],[t^3]\)，维数二。全部不小于四的单项式都在 \(A\)，所以 \(t^4B\) 在导子中；若 \(a\) 的常数项或二次项非零，乘 \(t\) 后会出现一次项或三次项，故不在导子中。因此导子为 \(t^4B=(t^4,t^5)_A\)。与 \(k[t^2,t^3]\) 相比，缺失的函数和导子的消失阶数都增加了。
\end{solution}

\begin{exercise}\label{b3:ex:06-trace-bound-quadratic}
在 \(\mathbb Q(\sqrt5)/\mathbb Q\) 中取由整元素组成的域基 \(u_1=1,u_2=\sqrt5\)。计算迹矩阵 \(T=\bigl(\operatorname{Tr}(u_iu_j)\bigr)\)，并用\cref{b3:lem:normalization-separable-finite}说明所有代数整数的公共分母受到怎样的控制。
\end{exercise}
\begin{solution}
迹把 \(a+b\sqrt5\) 送到 \(2a\)，故
\(T=\operatorname{diag}(2,10)\)、\(\det T=20\)。
若 \(z=a+b\sqrt5\) 在 \(\mathbb Z\) 上整，则 \(\operatorname{Tr}(z)=2a\) 与 \(\operatorname{Tr}(z\sqrt5)=10b\) 都是整数，故
\[
z\in \tfrac12\mathbb Z+\tfrac1{10}\mathbb Z\sqrt5,
\qquad
10z\in\mathbb Z+\mathbb Z\sqrt5.
\]
因此 \(10\) 已是统一公共分母，当然 \(20=\det T\) 也可以。伴随矩阵方法只需给出一个统一非零分母，并不保证所得分母最优。这给出有限格内的控制，不声称所有该格中的元素都整。例如 \((1+\sqrt5)/2\) 整，而 \(1/2\) 不整。
\end{solution}

\begin{exercise}\label{b3:ex:06-inseparable-finite-coefficients}
设 \(k\) 为特征 \(p>0\) 的域，\(a\in k\setminus k^p\)，\(x\) 为变量。令 \(K=k(x)\)，在 \(K\) 的一个代数闭包中取 \(z\) 满足 \(z^p=ax^p+x\)。在添入 \(\alpha=a^{1/p}\)、\(y=x^{1/p}\) 后，说明 \(z\) 如何消除，并给出 \(k[\alpha,y]\) 作为 \(k[x]\) 模的一组有限生成元。
\end{exercise}
\begin{solution}
在共同代数闭包中，
\[
(\alpha y^p+y)^p=ay^{p^2}+y^p=ax^p+x.
\]
特征 \(p\) 的域中取 \(p\) 次幂单射，所以 \(z=\alpha y^p+y\)。因而加入有限个系数根与参数根后，这个纯不可分生成元已属于新的有理函数域。环 \(k[\alpha,y]\) 在 \(k[x]\) 上由
\(\alpha^i y^j\)、\(0\le i,j<p\) 生成，因为 \(\alpha^p=a\)、\(y^p=x\)。本题不要求确定原环在 \(K(z)\) 中的整闭包；它展示的是有限基变换如何处理不可分部分。
\end{solution}

\begin{exercise}\label{b3:ex:06-inseparable-trace-zero}
设 \(k\) 为特征 \(p>0\) 的域，\(t\) 为变量，\(L=k(t)\)、\(K=k(t^p)\)。证明 \(L/K\) 次数为 \(p\)，且迹映射恒为零。为什么不能把可分情形的迹分母论证原样用于此例？
\end{exercise}
\begin{solution}
\(1,t,\ldots,t^{p-1}\) 在 \(K\) 上线性无关：清分母后，不同幂次模 \(p\) 的余数不同，非零项不能抵消；又 \(t^p\in K\)，故它们是基。对 \(1\le i<p\)，乘 \(t^i\) 在这组基上的矩阵没有对角项，因为指数的余数都被非零的 \(i\) 改变；乘常数 \(c\in K\) 的迹为 \(pc=0\)。线性性给出全部元素的迹为零。因此对任一 \(K\) 基 \(u_1,\ldots,u_p\)，迹矩阵 \(T=(\operatorname{Tr}_{L/K}(u_i u_j))\) 为零矩阵，\(\Delta=\det T=0\)。可分证明中用来限制分母的 \(\Delta^{-1}\) 在这里不存在。
\end{solution}

\begin{exercise}\label{b3:ex:06-conductor-local-support}
设 \(k\) 为域，\(A=k[t]\subseteq B=k[t,t^{-1}]\)，令 \(M=B/A\)。求导子 \(\mathfrak c=\operatorname{Ann}_A(M)\) 与 \(\operatorname{Supp}_A(M)\)，并证明 \(B\) 不是有限生成的 \(A\) 模。比较 \(\operatorname{Supp}_A(M)\) 与 \(V_A(\mathfrak c)\)，说明有限扩张下的导子局部同构判据为什么不能无条件沿用。
\end{exercise}
\begin{solution}
若 \(a\in\mathfrak c\)，则对每个 \(n\ge1\) 都有 \(at^{-n}\in k[t]\)，所以 \(a\) 被任意高次的 \(t\) 整除，只能为零。因此 \(\mathfrak c=0\)。有限多个 Laurent 多项式有共同的最低幂次下界，它们的 \(A\) 线性组合也有这个下界，不可能生成全部 \(t^{-n}\)，故 \(B\) 不是有限生成的 \(A\) 模。

当 \(t\notin\mathfrak p\) 时，\(t\) 在 \(A_{\mathfrak p}\) 中已经可逆，故 \(B\otimes_AA_{\mathfrak p}=A_{\mathfrak p}\)，即 \(M_{\mathfrak p}=0\)。在 \(\mathfrak p=(t)\) 处，类 \(t^{-1}+A\) 的局部化却非零：若它被某个 \(s\in A\setminus(t)\) 零化，就有 \(st^{-1}\in A\)，迫使 \(t\mid s\)，矛盾。因此
\[
\operatorname{Supp}_A(M)=\{(t)\}=V_A(t)
\quad\subsetneq\quad V_A(\mathfrak c)=\operatorname{Spec}A.
\]
例如在 \(\mathfrak p=(0)\) 处两环局部化相同，但 \(\mathfrak c\setminus\mathfrak p\) 为空。每个 \(M\) 中的元素都能被某个 \(t\) 的幂零化，却没有一个共同的非零元素零化整个 \(M\)；有限生成假设正是将逐元素的零化条件化为有限个条件、再取共同乘积的依据。
\end{solution}

\begin{exercise}\label{b3:ex:06-node-local-idempotents}
设 \(k\) 为特征不等于二的域，\(A=k[t^2-1,t(t^2-1)]\subseteq B=k[t]\)，\(\mathfrak m=(t^2-1,t(t^2-1))_A\)，\(S=A\setminus\mathfrak m\)。解释为什么节点的局部正规化 \(S^{-1}B\) 没有非平凡幂等元，而它模掉 \(\mathfrak mA_{\mathfrak m}\) 后的纤维有非平凡幂等元，并写出后者。
\end{exercise}
\begin{solution}
局部正规化是整环，若 \(e^2=e\)，则 \(e(e-1)=0\)，所以 \(e=0\) 或一。纤维为 \(k[t]/(t^2-1)\)，特征不为二时，
\[
e_+=\frac{t+1}{2},\qquad e_-=\frac{1-t}{2}
\]
满足 \(e_\pm^2=e_\pm\)、\(e_++e_-=1\)、\(e_+e_-=0\)，分别在 \(t=1,-1\) 取值 \((1,0)\)、\((0,1)\)。这两个幂等元是商环中的类，不在整环中满足同一恒等式。
\end{solution}

\begin{exercise}\label{b3:ex:06-node-characteristic-two}
设 \(k\) 为特征二的域，\(A=k\bigl[t^2-1,t(t^2-1)\bigr]\)。用新变量 \(u=t-1\) 把 \(A\) 化成尖点形式，并计算其导子和原点纤维。
\end{exercise}
\begin{solution}
此时 \(t^2-1=u^2\)，且 \(t(t^2-1)=u^3+u^2\)，所以
\(A=k[u^2,u^3]\)。正规化是 \(k[u]\)，由\cref{b3:prop:cusp-conductor}导子为 \(u^2k[u]\)。原点理想扩张为 \((u^2)\)，纤维为 \(k[u]/(u^2)\)，只有一个点且非约化。于是正文节点的“两点”结论确实需要特征不为二。
\end{solution}

\begin{exercise}\label{b3:ex:06-finite-type-normalization-algebra}
设 \(k\) 为域，\(A\) 是约化有限型 \(k\) 代数。证明其正规化 \(\widetilde A\) 仍为有限型 \(k\) 代数。再说明若 \(A\) 为整环，正规化不改变分式域的超越次数。
\end{exercise}
\begin{solution}
由\cref{b3:thm:normalization-finite}，取有限个 \(A\) 模生成元 \(b_1,\ldots,b_s\) 生成 \(\widetilde A\)。若 \(a_1,\ldots,a_r\) 作为 \(k\) 代数生成 \(A\)，则
\(a_1,\ldots,a_r,b_1,\ldots,b_s\) 作为 \(k\) 代数生成 \(\widetilde A\)：每个元素都是 \(b_j\) 的 \(A\) 线性组合。整环情形
\(A\subseteq\widetilde A\subseteq\operatorname{Frac}(A)\)，且两侧取分式域得到同一域，因此超越次数不变。
\end{solution}

% !TeX root = ../../Book3.tex
% 纲要标题：## 第7章　仿射代数集与坐标环
% 纲要位置：工作记录/计划纲要.md，第 2006 行。


\chapter{仿射代数集与坐标环}
\label{b3:ch:07}
\BookChapterNumber{b3:ch:07}

\begin{chapterabstract}
零点定理使仿射代数集与约化有限生成代数相互对应。本章说明这项对应怎样作用于不可约分支、连通性、正则函数和态射，证明仿射代数集上的全局正则函数恰为坐标环元素，并把局部环解释为正则函数的芽环。随后区分集合意义的纤维与保留幂零元的纤维代数，讨论基域扩张和几何点，最后用尖点、节点、交叉直线及加厚点检验代数与几何之间的对应。
\end{chapterabstract}

\begin{learninggoals}
\begin{itemize}
\item 从坐标环辨认不可约分支和开闭分解，区分代数集与不可约的仿射簇。
\item 证明全局正则函数定理及仿射反等价，解释局部环所记录的函数信息。
\item 计算纤维代数和基域扩张，辨认约化性、闭点与几何点的不同含义。
\item 在曲线例子中计算正规化、导子与局部不变量，说明双射为何不足以给出同构。
\end{itemize}
\end{learninggoals}

平面抛物线 \(y=x^2\) 的坐标环是 \(k[x,y]/(y-x^2)\)，把 \(y\) 换成 \(x^2\) 后它同构于 \(k[x]\)。从方程所在的平面看，抛物线与直线不同；从抽象仿射代数集看，它们由同一个多项式环描述。这个例子提醒我们，嵌入方式和对象自身要分开，而多项式函数正是连接两种看法的桥梁。

\ProseParagraphBreak

若在 \(x\ne0\) 的部分研究函数，\(1/x\) 可以使用；它却不能作为整条直线上的多项式函数。于是要在各点附近允许不同分母，并检查函数在重叠区域是否相合。另一方面，把一个点映到另一个空间会使目标上的函数反向拉回。这些局部与反向的运算，最终由坐标环及其局部化清楚表达。

\ProseParagraphBreak

前六章已经给出了本章所需的代数工具：素谱把理想组织成拓扑空间，局部化描述一点附近的运算，有限性控制分支与生成元，整扩张和零点定理则把环与多项式的零点联系起来。这些代数结论在几何中有具体含义。态射把点向前送，函数则通过复合向后拉回。

\secref{b3:sec:0701}与\secref{b3:sec:0702}使用坐标环、局部化与强零点定理。\secref{b3:sec:0704}的曲线例子则以\secref{b3:sec:0604}的尖点、节点及\secref{b3:sec:0606}的导子计算为基础；交叉直线的全分式环还用到\cref{b3:prop:reduced-total-quotient-product}。

\ProseParagraphBreak
前两节以及最后一节以代数闭域为基础；第三节讨论基域扩张时会明确放宽这一假设。全局正则函数与局部环的经典叙述可参阅 Hartshorne 的《Algebraic Geometry》\cite[Chapter I, \S 3, pp. 14--18]{Hartshorne1977}。其中定理 3.2 以不可约簇为对象；本章给出的拼接论证同时适用于可约仿射代数集，因而不必把每个例子预先拆成分支。

% !TeX root = ../../Book3.tex
% 纲要标题：### 7.1 代数集
% 纲要位置：工作记录/计划纲要.md，第 2109 行。
% 纲要细目（写作范围）：
% * Zariski 拓扑与闭子集；经典点、极大谱与一般点的区别
% * 不可约性与素理想；严格反包含及极小素与最大不可约分量
% * 连通性和幂等元；开闭分解的特征函数

\section{代数集}
\label{b3:sec:0701}

前六章主要从环和模出发；本章把其中一些对象重新解释为多项式方程的解集。为使点集确实能够恢复全部约化关系，本节和下一节固定代数闭域 \(k\)。\secref{b3:sec:0703}才明确放宽基域，并用域值点及纤维代数保留普通点集看不到的信息。

\subsection{Zariski 拓扑与闭子集}
\label{b3:subsec:0701:01}

\begin{definition}\label{b3:def:affine-algebraic-set}
设 \(k\) 为代数闭域，\(n\ge0\)，记 \(\mathbb A_k^n=k^n\)。给定 \(k[x_1,\ldots,x_n]\) 的理想 \(I\)，若 \(X=Z_k(I)=\{a\in k^n:f(a)=0\text{ 对每个 }f\in I\}\)，则称 \(X\) 为\term[fangshe-daishu-ji]{仿射代数集}[affine algebraic set]。以这些子集为闭集，得到 \(\mathbb A_k^n\) 的 Zariski 拓扑；每个代数集赋予子空间拓扑。本书允许代数集为空或可约；若代数集非空且不可约，则称它为\term[fangshe-daishu-cu]{仿射代数簇}[affine variety]。
\end{definition}

此处采用的“簇”约定比“代数集”多一个不可约条件。拓扑公理可与\eqref{b3:eq:zariski-closed-identities}相同地验证：任意交对应理想之和，有限并对应理想之积。对闭集 \(X\)，\cref{b3:thm:strong-nullstellensatz}给出 \(I(X)\) 为根理想，且 \(X=Z_k\bigl(I(X)\bigr)\)。

\ProseParagraphBreak
沿用\cref{b3:def:affine-coordinate-ring}记 \(A=k[X]\)。对 \(f\in A\)，记
\[
D_X(f)=\bigl\{\,a\in X:f(a)\ne0\,\bigr\},\qquad
Z_X(J)=\bigl\{\,a\in X:g(a)=0\text{ 对全部 }g\in J\,\bigr\}.
\]
这里 \(J\) 是 \(A\) 的理想；这些定义不依赖在多项式环中选择的代表。\(D_X(f)\) 构成 \(X\) 的开集基。由弱零点定理，映射
\begin{equation}\label{b3:eq:closed-points-affine}
X\longrightarrow\operatorname{MaxSpec}A,
\qquad a\longmapsto\mathfrak m_a=\bigl\{f\in A:f(a)=0\bigr\}
\end{equation}
是双射，且把上述拓扑对应为素谱中极大理想子集的子空间拓扑。因为 \(a\in Z_X(J)\) 恰好等价于 \(J\subseteq\mathfrak m_a\)，它同时识别了闭集。
\symidx{\(D_X(f)\)：代数集上的主开集}
\symidx{\(\operatorname{MaxSpec}A\)：环的极大理想集合，赋予素谱的子空间拓扑}

因此 \(X\) 的每个点都闭，而 \(\operatorname{Spec}k[X]\) 还可能有非闭的一般点。两种空间记录的信息不同。把\eqref{b3:eq:closed-points-affine}与极大理想集合的包含放在一起，可写成
\[
X\xrightarrow{\sim}\operatorname{MaxSpec}k[X]
\lhook\joinrel\longrightarrow\operatorname{Spec}k[X].
\]
例如仿射直线的坐标点 \(a\in k\) 对应闭点 \((t-a)\)，素谱还含零理想 \((0)\)。后者的闭包是整个 \(\operatorname{Spec}k[t]\)，所以它是非闭的一般点；在经典坐标点集合 \(k\) 中没有与它对应的点。在 \(\operatorname{Spec}k[x,y]\) 中，素理想 \((x)\) 的闭包则为 \(V(x)\)，其中的闭点恰为直线 \(x=0\) 上的坐标点。\((x)\) 是整条直线的一般点，并不是这条直线上另加的一个坐标点。因此极大谱把经典点编码为理想，而素谱还为各个不可约闭子集保留了它的一般点。

\subsection{不可约性与素理想}
\label{b3:subsec:0701:02}

\begin{proposition}\label{b3:prop:affine-irreducible}
设 \(k\) 为代数闭域，\(X\subseteq k^n\)（\(n\ge0\)）为仿射代数集。若 \(X\ne\varnothing\)，则 \(X\) 不可约，当且仅当它的零点理想 \(I(X)\) 为素理想，也当且仅当坐标环 \(k[X]\) 为整环。每个代数集有有限个不可约分量，它们对应 \(k[X]\) 的极小素理想。
\end{proposition}
\begin{proof}
若 \(fg\in I(X)\)，则每个点上至少一个因子为零，故
\(X=Z_X(f)\cup Z_X(g)\)。不可约性使其中一个等于 \(X\)，即 \(f\) 或 \(g\) 属于 \(I(X)\)，所以这个真理想素。

反过来，若 \(I(X)\) 素，而 \(X=Y\cup Z\) 是两个闭子集之并，假设两者都真。闭子集 \(Y\subsetneq X\) 给出 \(I(X)\subseteq I(Y)\)；若理想相等，由强零点定理便有 \(Y=Z_k(I(Y))=Z_k(I(X))=X\)，矛盾。所以理想包含也严格，\(Z\) 的情形相同，因而可取 \(f\in I(Y)\setminus I(X)\)、\(g\in I(Z)\setminus I(X)\)。乘积在全部 \(X\) 上为零，违反素性。商整环的等价性由\cref{b3:prop:prime-maximal-quotient}给出。

最后 \(A=k[X]\) 为 Noether 环。由\cref{b3:prop:noether-spectrum-components}，它有有限个极小素理想 \(\mathfrak p_i\)，且因 \(A\) 约化而 \(\bigcap_i\mathfrak p_i=0\)。零点运算把这个有限交变成 \(X=\bigcup_i Z_X(\mathfrak p_i)\)。前面的判据说明各成员不可约。素理想包含越小，相应零点集越大：\(\mathfrak p\subseteq\mathfrak q\) 等价于 \(Z_X(\mathfrak q)\subseteq Z_X(\mathfrak p)\)。所以极小素理想对应按包含关系极大的不可约闭子集，也就是全部不可约分量。
\end{proof}

在 \(X=Z_k(xy)\) 中，两个不可约分量为两条坐标轴，它们都不可约，却在原点相交。相比之下，\(Z_k\bigl(x(x-1)\bigr)\subseteq\mathbb A_k^1\) 是两个分开的点。不可约分量的数目和连通分支的数目因而是不同的数据。

\subsection{连通性与幂等元}
\label{b3:subsec:0701:03}

\begin{theorem}\label{b3:thm:affine-connected-idempotents}
设 \(k\) 为代数闭域，\(X\subseteq k^n\)（\(n\ge0\)）为仿射代数集。则 \(X\) 连通，当且仅当坐标环 \(k[X]\) 没有除零与一以外的幂等元。
\end{theorem}
\begin{proof}
设 \(e^2=e\)。在每个点上，\(e(a)\) 为零或一，所以
\[
X=Z_X(e)\sqcup Z_X(1-e).
\]
两部分闭且互为补集，故也开。函数 \(e\) 在 \(Z_X(e)\) 上取零，在 \(Z_X(1-e)\) 上取一，因此正是后一开闭子集的特征函数。如果 \(e\ne0,1\)，因 \(k[X]\) 的元素确为函数，两部分都非空，得到不连通性。

反过来，若 \(X=Y\sqcup Z\)，其中 \(Y,Z\) 非空且均开闭，取它们在 \(A=k[X]\) 中的零点理想 \(J,K\)。因为并为 \(X\)，有 \(J\cap K=0\)；因为交为空，由\cref{b3:cor:proper-ideal-common-zero}，\(J+K=A\)。第一卷定理10.4.3的中国剩余定理给出 \(A\cong A/J\times A/K\)。两因子非零，元素 \((1,0)\) 的原像就是在 \(Y\) 上取一、在 \(Z\) 上取零的非平凡幂等元。中国剩余定理在这里保证这项开闭分解的特征函数确实属于坐标环。空集的坐标环为零环，零与一相同，也符合所述判据。
\end{proof}

例如两条交叉直线的坐标环可嵌入 \(k[x]\times k[y]\)，其像为两坐标常数项相同的多项式对。一个幂等元的两个坐标都只能为零或一，再要求常数项相同，只剩 \((0,0)\)、\((1,1)\)。因此这个可约代数集仍连通。相比之下，两点代数的坐标环为 \(k\times k\)，其两个非平凡幂等元把两个点分开。

% !TeX root = ../../Book3.tex
% 纲要标题：### 7.2 正则函数与态射
% 纲要位置：工作记录/计划纲要.md，第 2014 行。
% 纲要细目（写作范围）：
% * 坐标函数与多项式映射
% * 仿射代数集的坐标环
% * 态射和环同态的反向对应

\section{正则函数与态射}
\label{b3:sec:0702}

坐标多项式显然给出函数，但在一个点附近允许除以不为零的函数，也应得到可用的局部表达。正则性先按这个局部条件定义。本节证明，在整个仿射代数集上处处正则的函数恰好来自坐标环，并由此确定代数集之间的映射应当如何与环同态对应。

\subsection{坐标函数与多项式映射}
\label{b3:subsec:0702:01}

\begin{definition}\label{b3:def:regular-function}
设 \(k\) 为代数闭域，\(n\ge0\)，\(X\subseteq k^n\) 为仿射代数集，\(U\subseteq X\) 为开子集，\(f:U\to k\) 为函数。若对每个 \(a\in U\)，存在满足 \(a\in W\subseteq U\) 的相对开邻域 \(W\) 及 \(g,h\in k[X]\)，使 \(h\) 在 \(W\) 上处处非零且 \(f|_W=g|_W/h|_W\)，则称 \(f\) 为 \(U\) 上的\term[zhengze-hanshu]{正则函数}[regular function]。全部这样的函数构成环 \(\mathcal O_X(U)\)。
\end{definition}

两个局部表达式相加或相乘时，分母取乘积；在共同邻域内它仍非零，所以正则函数对环运算封闭。它们也连续：在局部表达 \(g/h\) 上，方程 \(f=c\) 等价于 \(g-ch=0\)，是相对闭条件；仿射直线的真闭集为有限点集，由此得到 Zariski 连续性。

\ProseParagraphBreak
设 \(k\) 为代数闭域，\(X\subseteq k^n\)、\(Y\subseteq k^m\) 为仿射代数集，\(\varphi:X\to Y\) 为映射。若 \(\varphi\) 的每个坐标函数都在 \(X\) 上正则，则称 \(\varphi\) 为代数集的\term[daishu-ji-taishe]{态射}[morphism of algebraic sets]。多项式关系的代入仍为正则函数，故这种映射连续，并且正则函数的复合仍正则：在目标的局部表达式中代入各坐标，其分母在相应原像邻域不为零。因而态射对复合封闭。对代数集的开子集，也用同一个局部条件定义其间的态射。

\subsection{仿射代数集的坐标环}
\label{b3:subsec:0702:02}

\begin{theorem}\label{b3:thm:affine-global-regular}
设 \(k\) 为代数闭域，\(X\subseteq k^n\) 为仿射代数集。自然映射 \(k[X]\to\mathcal O_X(X)\) 是同构。故整个 \(X\) 上的正则函数都能由一个多项式表示。
\end{theorem}
\begin{proof}
按坐标环的定义，不同商类给不同的多项式函数，所以映射单射。设 \(f\) 处处正则。在每个点附近，将一个局部表达式 \(a/b\) 的邻域缩小为主开集 \(D_X(h)\)，使它仍包含该点且 \(b\) 在其中非零。令 \(a'=ah,b'=bh\)。在 \(D_X(b')\) 上仍有 \(f=a'/b'\)，这些主开集覆盖 \(X\)。以下将这组分子分母记为 \(a_i,b_i\)。

全部 \(b_i\) 没有公共零点，\(b_i^2\) 也没有。由\cref{b3:cor:proper-ideal-common-zero}在坐标环中的形式，它们生成单位理想；因此存在有限多个指标及 \(c_i\in k[X]\)，使 \(\sum_i c_i b_i^2=1\)。对每个 \(i\)，函数等式
\[
b_i^2f=a_i b_i
\]
在 \(D_X(b_i)\) 上来自局部表达式，在其补集上则两边都为零，所以在全部 \(X\) 上成立。于是
\[
f=\sum_i c_i b_i^2f=\sum_i c_i a_i b_i,
\]
右侧是坐标环中的一个函数，证明满射。这一步没有将 \(k[X]\) 假定为整环，所以可约代数集也包括在内。
\end{proof}

该结论可与 Hartshorne 的《Algebraic Geometry》\cite[Chapter I, Theorem 3.2(a), p. 17]{Hartshorne1977}对照。原书此处以不可约仿射簇为对象；上面的证明直接处理本书允许的可约代数集。它说的是全局正则函数；在真开子集上，分母通常不能全部消去。例如 \(1/t\) 在 \(\mathbb A_k^1\setminus\{0\}\) 上正则，却不是全直线上的多项式函数。

\begin{proposition}\label{b3:prop:regular-germ-local-ring}
设 \(k\) 为代数闭域，\(X\subseteq k^n\) 为仿射代数集，\(A=k[X]\)，\(a\in X\)。记 \(\mathfrak m_a=\{f\in A:f(a)=0\}\) 为该点的极大理想。正则函数在 \(a\) 处的芽组成的环 \(\mathcal O_{X,a}\) 自然同构于 \(A_{\mathfrak m_a}\)。
\end{proposition}
\begin{proof}
一个芽由包含 \(a\) 的邻域上的正则函数表示；两代表若在某个更小的含 \(a\) 邻域相等，就表示同一芽。分式 \(g/s\)、\(s(a)\ne0\)，在 \(D_X(s)\) 上定义正则函数，局部化的等价关系保证这给出同态 \(A_{\mathfrak m_a}\to\mathcal O_{X,a}\)。每个芽按定义局部都是这样的分式，故同态满射。

若这个分式的芽为零，缩小到 \(a\in D_X(h)\subseteq D_X(s)\) 后有 \(g=0\)。因此函数 \(hg\) 在 \(D_X(h)\) 上为零，在补集上也为零，故 \(hg=0\) 于 \(A\)。又 \(h(a)\ne0\)，分式零判据使 \(g/s=0\) 于 \(A_{\mathfrak m_a}\)，得到单射。
\end{proof}
\symidx{\(\mathcal O_{X,a}\)：代数集在点 \(a\) 处的正则函数芽环}

\subsection{态射与环同态的反向对应}
\label{b3:subsec:0702:03}

\begin{theorem}\label{b3:thm:affine-antiequivalence}
设 \(k\) 为代数闭域，\(X,Y\) 为 \(k\) 上的仿射代数集。取态射的函数拉回给出双射
\[
\operatorname{Hom}(X,Y)\cong
\operatorname{Hom}_{k\text{-}\mathrm{alg}}\bigl(k[Y],k[X]\bigr).
\]
它与恒等映射及反向的复合相容。每个约化有限型 \(k\) 代数都同构于某个代数集的坐标环，因此代数集与约化有限型 \(k\) 代数之间得到反变等价。
\end{theorem}
\begin{proof}
态射的坐标正则函数由\cref{b3:thm:affine-global-regular}可用多项式表示。若 \(\varphi:X\to Y\)，令 \(\varphi^*(g)=g\circ\varphi\)，则这是 \(k\) 代数同态。

反过来，给定 \(\psi:k[Y]\to k[X]\)，选 \(Y\subseteq k^m\) 的坐标函数 \(y_1,\ldots,y_m\)，令 \(f_i=\psi(\bar y_i)\)。定义 \(\varphi(a)=\bigl(f_1(a),\ldots,f_m(a)\bigr)\)。对任意 \(h\in I(Y)\)，有 \(h(f_1,\ldots,f_m)=\psi(\bar h)=0\) 于 \(k[X]\)，故该坐标点确实落在 \(Y\)。其坐标为多项式函数，所以是态射。所得拉回与 \(\psi\) 在代数生成元上相同，故完全相同；而从原态射开始也恢复相同坐标，所以两个过程互逆。复合次序由函数复合直接得到。

最后，约化有限型代数写成 \(A=k[x_1,\ldots,x_n]/I\)，其中约化性使 \(I\) 为根理想。对 \(X=Z_k(I)\)，强零点定理给出 \(I(X)=I\)，故 \(A\cong k[X]\)。空集对应零环，也符合这组同态与态射的对应。
\end{proof}

一个双射态射未必是同构。例如尖点的参数式 \(t\mapsto(t^2,t^3)\) 在域值点上双射：\(x\ne0\) 时逆参数为 \(y/x\)，原点则对应零。但是坐标环包含 \(k[t^2,t^3]\subsetneq k[t]\) 不满射，依\cref{b3:thm:affine-antiequivalence}，它不能来自同构。形式上的逆函数存在，不说明它在原点正则。

\ProseParagraphBreak
对 \(f\in k[X]\)，主开集 \(D_X(f)\) 可识别为闭代数集
\[
Y=\bigl\{\,(a,t)\in X\times k:t\cdot f(a)=1\,\bigr\}
\]
的第一坐标投影；反向为 \(a\mapsto\bigl(a,1/f(a)\bigr)\)，两者都是态射。其坐标环为 \(k[X][t]/(tf-1)\cong k[X]_f\)，这个商是约化环，因为约化环的局部化约化。故主开集也有自然的仿射代数集结构，且 \(\mathcal O_X\bigl(D_X(f)\bigr)\cong k[X]_f\)。这里使用的是具体的正反映射，不要求把一般开子集都当成仿射代数集。

% !TeX root = ../../Book3.tex
% 纲要标题：### 7.3 纤维与基域扩张
% 纲要位置：工作记录/计划纲要.md，第 2020 行。
% 纲要细目（写作范围）：
% * 纤维的方程
% * 张量积和坐标环
% * 非代数闭域上的点与几何点

\section{纤维与基域扩张}
\label{b3:sec:0703}

把参数固定为一个值，就得到映射的一条纤维。这个操作在方程中是代入，在代数中是取商或张量积。取商可能产生幂零元，即使原来的源和目标都约化；因此纤维的点集与记录全部关系的代数必须分别保留。

\subsection{纤维的方程}
\label{b3:subsec:0703:01}

设 \(k\) 仍代数闭，\(\varphi:X\to Y\) 为态射，\(A=k[X],B=k[Y]\)。对 \(b\in Y\)，记其极大理想为 \(\mathfrak m_b\)，并令 \(J=\varphi^*(\mathfrak m_b)A\)。若坐标表达为 \(\varphi=(f_1,\ldots,f_m)\)、\(b=(b_1,\ldots,b_m)\)，则
\[
\varphi^{-1}(b)=Z_X(f_1-b_1,\ldots,f_m-b_m)=Z_X(J).
\]
这个闭集称为 \(b\) 上的\term[xianwei]{纤维}[fiber]。用强零点定理得到它作为代数集的坐标环
\begin{equation}\label{b3:eq:reduced-fiber-ring}
k\bigl[\varphi^{-1}(b)\bigr]\cong A/\sqrt J.
\end{equation}

\begin{proposition}\label{b3:prop:fiber-algebra}
设 \(k\) 为代数闭域，\(\varphi:X\to Y\) 为 \(k\) 上仿射代数集之间的态射，\(A=k[X]\)、\(B=k[Y]\)，\(b\in Y\)。令 \(\mathfrak m_b=\{g\in B:g(b)=0\}\)，\(J=\varphi^*(\mathfrak m_b)A\)。则有自然同构
\[
A\otimes_B\kappa(b)\cong A/J,
\qquad \kappa(b)=B/\mathfrak m_b\cong k.
\]
因此纤维的坐标环是该张量代数的约化商；只有当 \(J\) 已为根理想时，才能直接把坐标环写成 \(A/J\)。
\end{proposition}
\begin{proof}
将 \(a\otimes(c+\mathfrak m_b)\) 送到 \(a\varphi^*(c)+J\)。改变 \(c\) 的代表所产生的差在 \(J\) 中，且该公式对 \(B\) 平衡、保持代数乘法，故得到同态。反向由 \(a+J\mapsto a\otimes1\) 给出；\(J\) 的每个生成元都在这个映射下为零，所以良定义。两个复合在纯张量及 \(a+J\) 上分别为恒等，得到同构。再结合\eqref{b3:eq:reduced-fiber-ring}即得最后断言。
\end{proof}

\begin{table}[htbp]
\centering
\caption[纤维代数与集合纤维的坐标环]{纤维代数与集合纤维的坐标环。设 \(k\) 代数闭，\(\varphi:X\to Y\) 为仿射代数集的态射，\(A=k[X]\)、\(B=k[Y]\)、\(b\in Y\)，且 \(J=\varphi^*(\mathfrak m_b)A\)。}
\label{b3:tab:fiber-algebra-and-reduction}
\begin{tabularx}{\textwidth}{@{}l>{\raggedright\arraybackslash}X>{\raggedright\arraybackslash}p{0.3\textwidth}@{}}
\toprule
对象 & 代数表达 & 平方映射 \(t\mapsto t^2\) 的零纤维 \\
\midrule
纤维代数 & \(A\otimes_B\kappa(b)\cong A/J\) & \(k[t]/(t^2)\) \\
集合纤维的坐标环 & \(A/\sqrt J\) & \(k[t]/(t)\cong k\) \\
\bottomrule
\end{tabularx}
\end{table}

在\cref{b3:tab:fiber-algebra-and-reduction}的零纤维代数中，\(\bar t\) 非零却平方为零；普通点集只有零，因而其坐标环不保留这个类。对同一个平方映射，若 \(\operatorname{char}k\ne2\)，在一处的纤维代数为 \(k[t]/(t^2-1)\cong k\times k\)，有两个不同点；在特征二时它变成 \(k[t]/\bigl((t-1)^2\bigr)\)，仍只有一个点，却保留非零平方零元。相同的代入过程可以产生不同的根数与幂零结构。

\subsection{张量积与坐标环}
\label{b3:subsec:0703:02}

经典仿射代数集的坐标环把在整个点集上为零的多项式看作零，因此总是约化的。纤维代数则要保留代入后的全部关系，已经可以出现 \(k[t]/(t^2)\) 这样的非约化环。为了保留这些幂零信息及其扩域后的变化，研究对象便从普通点集转为有限型代数，相应的点空间是代数的素谱。设 \(k\) 为任意域，\(A\) 为交换含幺有限型 \(k\) 代数，不要求约化。对域扩张 \(K/k\)，记 \(A_K=K\otimes_kA\)，称为将标量扩张到 \(K\) 的代数。

\begin{proposition}\label{b3:prop:base-change-presentation}
设 \(k\) 为域，\(K/k\) 为域扩张，\(I\) 为 \(k[x_1,\ldots,x_n]\) 的理想，\(A=k[x_1,\ldots,x_n]/I\)，其中 \(n\ge0\)。则
\[
K\otimes_kA\cong K[x_1,\ldots,x_n]/I K[x_1,\ldots,x_n].
\]
进一步，设 \(m\ge0\)，\(J\) 为 \(k[y_1,\ldots,y_m]\) 的理想，\(B=k[y_1,\ldots,y_m]/J\)。令
\[
R=k[x_1,\ldots,x_n,y_1,\ldots,y_m],
\]
并记 \(a_i=x_i+I\in A\)、\(b_j=y_j+J\in B\)。
则两组关系在 \(R\) 中的扩张理想 \(IR,JR\) 给出自然同构
\[
\begin{aligned}
A\otimes_kB&\cong R/(IR+JR),\\
a_i\otimes1&\longleftrightarrow x_i+(IR+JR),\\
1\otimes b_j&\longleftrightarrow y_j+(IR+JR).
\end{aligned}
\]
\end{proposition}
\begin{proof}
第一式中，将 \(c\otimes\bar f\) 送到 \(cf\) 的商类，得到良定义的代数同态。反向将 \(K\) 中的标量送到 \(c\otimes1\)，变量送到 \(1\otimes\bar x_i\)；原关系都为零，所以该赋值可下降到商。两复合固定所有标量和代数生成元，因此互逆。

第二项把 \(x_i\) 送到 \(a_i\otimes1\)，把 \(y_j\) 送到 \(1\otimes b_j\)。两组像元彼此交换，各自满足 \(I\) 与 \(J\) 中的关系，故得到 \(R/(IR+JR)\to A\otimes_kB\)。反向将 \((f+I)\otimes(g+J)\) 送到 \(fg+(IR+JR)\)，其中 \(f\in k[x_1,\ldots,x_n]\)、\(g\in k[y_1,\ldots,y_m]\)。这个公式对 \(k\) 平衡，改变两个代表元产生的差分别属于 \(IR\) 与 \(JR\)，因而良定义。两个复合固定两组生成元及 \(k\) 中的标量，得到所述同构。
\end{proof}

这个结论说明张量积怎样保存方程，却不自动保证所得环约化。若把一组方程解释为代数闭域中的普通零点集，强零点定理仍要求取其关系理想的根。在不代数闭的域上，仅看本域中的点还可能损失更多信息。

\ProseParagraphBreak
例如 \(A=\mathbb R[u]/(u^2+1)\cong\mathbb C\) 没有到 \(\mathbb R\) 的实代数同态，但
\[
\mathbb C\otimes_{\mathbb R}A
\cong\mathbb C[u]/(u^2+1)\cong\mathbb C\times\mathbb C,
\]
两个因子对应 \(u=i,-i\)。从环的角度，原对象并不为空；只是它没有实坐标点。不能把实零点集为空与原代数为零环混为一谈。

\subsection{非代数闭域上的点与几何点}
\label{b3:subsec:0703:03}

\begin{definition}\label{b3:def:field-valued-point}
设 \(k\) 为域，\(A\) 为交换含幺有限型 \(k\) 代数，\(K/k\) 为域扩张。若 \(\varphi:A\to K\) 为 \(k\) 代数同态，则称 \(\varphi\) 为一个 \(K\) 值点。若 \(K=k\)，则称它为 \(k\) 有理点；若 \(K\) 为代数闭域，则称它为以 \(K\) 为值的\term[jihe-dian]{几何点}[geometric point]。
\end{definition}

若 \(A=k[x_1,\ldots,x_n]/I\)，一个这样的同态恰由使全部 \(I\) 关系为零的一组 \(K\) 坐标确定。它的核在 \(\operatorname{Spec}A\) 中确定一个点，但未必是极大理想：例如自然包含 \(k[t]\to\overline{k(t)}\)、\(t\mapsto t\) 给出一个几何点，其在 \(\operatorname{Spec}k[t]\) 中的像为 \((0)\)，是全谱的一般点而非闭点。这里 \(t\) 的像是超越元；经典仿射直线的 \(k\) 坐标点 \(a\) 则对应极大理想 \((t-a)\)。如果只取固定代数闭包 \(\bar k\) 中的值，则由有限型性，坐标的像生成有限代数扩域，核为极大理想。

\ProseParagraphBreak
对任意 \(\mathfrak p\in\operatorname{Spec}A\)，点处的剩余域 \(\kappa(\mathfrak p)\) 都可嵌入某个代数闭域，从而给出位于它上面的几何点。这与只枚举 \(k\) 有理点不同。更一般地，给定交换含幺有限型 \(k\) 代数之间的同态 \(B\to A\) 及底点 \(\mathfrak q\in\operatorname{Spec}B\)，纤维代数 \(A\otimes_B\kappa(\mathfrak q)\) 就是在固定这个底点后，继续研究剩余域扩张中坐标的对象。

\begin{example}[变基后的非约化性]\label{b3:exm:inseparable-geometric-thickening}
设 \(p\) 为素数，\(s\) 为 \(\mathbb F_p\) 上的超越元，\(k=\mathbb F_p(s)\)。在 \(k\) 的一个代数闭包中取 \(\alpha\) 满足 \(\alpha^p=s\)，并令 \(L=k(\alpha)\)。则 \(L\) 为域，但
\[
L\otimes_kL\cong L[\varepsilon]/(\varepsilon^p).
\]
\end{example}
\begin{proof}
\(s\) 不是 \(k\) 中的 \(p\) 次幂：若 \((f/g)^p=s\)，把 \(f,g\in\mathbb F_p[s]\) 取互素，比较不可约因子 \(s\) 的指数便矛盾。因为 \(\alpha^p=s\)，其首一最小多项式 \(m_\alpha(T)\) 在 \(k[T]\) 中整除 \(T^p-s\)。在代数闭包中，\(T^p-s=(T-\alpha)^p\)，所以 \(m_\alpha\) 的全部根都为 \(\alpha\)，即 \(m_\alpha(T)=(T-\alpha)^r\)，其中 \(1\le r\le p\)。若 \(r<p\)，其 \(T^{r-1}\) 系数 \(-r\alpha\in k\) 会使 \(\alpha\in k\)，与上述结论矛盾。故 \(L\cong k[T]/(T^p-s)\)。应用\cref{b3:prop:base-change-presentation}后，在 \(L\) 中写 \(T^p-s=(T-\alpha)^p\)，令 \(\varepsilon=T-\alpha\) 即得所述同构。
\end{proof}

因此“约化”与“经过任意扩域仍约化”是不同要求。上例变基后仍只出现一个点，却有 \(p\) 次幂零结构。本章保留这个具体区别；一般几何约化性、可分性和平展性的关系将在后面的微分与平展章节中继续讨论。

% !TeX root = ../../Book3.tex
% 纲要标题：### 7.4 曲线分支与非约化结构
% 纲要位置：工作记录/计划纲要.md，第 2127 行。
% 纲要细目（写作范围）：
% * 尖点、节点、交叉直线与加厚点；经典约化模型与非约化测试代数的区别
% * 约化、不可约与正规性的区别；局部整闭性及素谱与经典参数映射的同胚
% * 不用图像替代局部环计算；导子零点集之外的局部同构、余切空间与一阶切空间
% ---

\section{曲线分支与非约化结构}
\label{b3:sec:0704}

尖点、节点与交叉直线都属于经典仿射代数集，具有约化坐标环；加厚点 \(k[\varepsilon]/(\varepsilon^2)\) 则是非约化有限型代数，不能作为普通代数集的坐标环。它用来检验域值点会遗漏什么信息。固定代数闭域 \(k\)，节点例额外要求 \(\operatorname{char}k\ne2\)。本节使用第六章已经计算的正规化与导子，比较这四个模型的纤维代数和局部环。

尖点与节点的计算使用\cref{b3:prop:cusp-normalization,b3:prop:node-normalization}的显式环包含及\secref{b3:sec:0606}的导子计算；交叉直线的全分式环与正规化已在\secref{b3:sec:0604}通过常数项相等的多项式对具体求出。

\subsection{尖点、节点、交叉直线与加厚点}
\label{b3:subsec:0704:01}

尖点 \(C=Z_k(y^2-x^3)\) 的坐标环由\cref{b3:prop:cusp-normalization}识别为 \(A=k[t^2,t^3]\)，正规化为 \(B=k[t]\)。参数映射 \(t\mapsto(t^2,t^3)\) 双射：原点来自 \(t=0\)，其余点 \((a,b)\) 满足 \(a\ne0\)，由 \(t=b/a\) 唯一恢复参数，其平方、立方恰为 \(a,b\)。环包含却严格。\cref{b3:prop:cusp-conductor}给出导子和商模；原点的极大理想 \(\mathfrak m=(t^2,t^3)_A\) 在 \(B\) 中生成 \((t^2)\)，因为 \(t^3=t\cdot t^2\)。于是
\[
\mathfrak c=t^2k[t],\qquad B/A\cong k\cdot[t],
\qquad B\otimes_Ak=B/\mathfrak mB\cong k[t]/(t^2).
\]
其中 \([t]=t+A\) 张成一维商空间 \(B/A\)。所以即使点的个数没有改变，原点的纤维代数仍有非零幂零元；集合纤维只有一个点，其约化坐标环为 \(k\)。原坐标环也确实少了一个函数 \(t\)。

\ProseParagraphBreak
节点 \(N=Z_k\bigl(y^2-x^2(x+1)\bigr)\) 使用参数 \(x=t^2-1,y=t(t^2-1)\)。由\cref{b3:prop:node-normalization}和\cref{b3:prop:node-conductor}，坐标环与正规化为
\[
A=k+(t^2-1)k[t]\subseteq B=k[t],\qquad
\mathfrak c=(t^2-1)k[t].
\]
原点理想 \(\mathfrak m=(x,y)_A\) 在 \(B\) 中生成 \((t^2-1)\)，所以纤维代数为 \(B/\mathfrak mB=k[t]/(t^2-1)\cong k\times k\)，对应 \(t=1,-1\) 两个不同点。其余点满足 \(x\ne0\)，可以用 \(t=y/x\) 恢复参数。

\begin{proposition}\label{b3:prop:crossed-lines-normalization}
设 \(k\) 为代数闭域，\(A=k[x,y]/(xy)\) 为交叉直线的坐标环。限制到两轴给出嵌入
\[
A\hookrightarrow k[x]\times k[y],\qquad
\bar f\longmapsto\bigl(f(x,0),f(0,y)\bigr).
\]
其像为常数项相等的多项式对，正规化为 \(B=k[x]\times k[y]\)，导子为 \((\bar x,\bar y)_A\)，且 \(B/A\cong k\)。
\end{proposition}
\begin{proof}
\secref{b3:sec:0604}已证明限制映射的单射性与常数项相等的像，并在 \(Q(A)=k(x)\times k(y)\) 中添入 \(e=(1,0)\)，将正规化识别为 \(A[e]=B\)。原环要求两轴上的函数在交点取相同值；\(e\) 在第一条轴上取一、第二条轴上取零，正是原环所不允许的函数。正规化取消了这项值相等的限制。

若 \((f,g)\in A\) 属于导子，分别乘以 \(e\)、\(1-e\)，要求 \((f,0),(0,g)\) 的常数项仍相等，所以 \(f(0)=g(0)=0\)。反向这个条件保证与任意多项式对相乘仍在 \(A\) 内。因此导子就是两坐标常数项都为零的理想，由 \(\bar x,\bar y\) 生成。最后，差值映射 \((f,g)\mapsto f(0)-g(0)\) 的核为 \(A\)，像为 \(k\)，得到商空间的描述；\(A\) 在此商上的作用经两坐标的公共常数值给出。
\end{proof}

交叉直线的原点理想在 \(B\) 中扩张为 \(xk[x]\times yk[y]\)，因而原点的正规化纤维也为 \(k\times k\)。这里两个点分别来自两条全局不可约分量，与不可约节点曲线的两个上方点有所不同。

加厚点则用代数 \(D=k[\varepsilon]/(\varepsilon^2)\) 及其素谱描述。它只有一个素理想，却含非零的 \(\varepsilon\)。任何到域的同态都把 \(\varepsilon\) 送到零，所以所有域值点仍只能看到普通一点。若测试 \(k\) 代数自身含有非零平方零元，就可把 \(\varepsilon\) 送到这样的元，从而辨认域值点遗漏的函数变化。作为普通代数集，这一点的坐标环是 \(k\)；要研究加厚结构，必须保留 \(D\) 本身。按\cref{b3:def:normalization}的约定，本书先取 \(D_{\mathrm{red}}=k\) 再正规化，所得正规化为 \(k\)，自然映射 \(D\to k\) 的核为 \((\varepsilon)\)。这不是把 \(D\) 本身作为约化环在全分式环中取整闭包。

\subsection{约化性、不可约性与正规性}
\label{b3:subsec:0704:02}

坐标环按定义约化，所以普通代数集本身不带被点集遗漏的幂零元。不可约性进一步要求坐标环为整环。设 \(k\) 为代数闭域，\(X\subseteq k^n\) 为不可约仿射代数集。若坐标整环 \(k[X]\) 在其分式域中整闭，则称 \(X\) 为\term[zhenggui-fangshe-cu]{正规仿射簇}[normal affine variety]。正规化正是针对最后这一项条件的构造。

\begin{proposition}\label{b3:prop:normality-local-detection}
设 \(A\) 为交换含幺整环。则 \(A\) 整闭，当且仅当对 \(A\) 的每个极大理想 \(\mathfrak m\)，局部环 \(A_{\mathfrak m}\) 整闭。
\end{proposition}
\begin{proof}
正向由\cref{b3:cor:integrally-closed-localization}得到。反向设 \(z\in K=\operatorname{Frac}(A)\) 对 \(A\) 整。同一个首一方程在每个 \(A_{\mathfrak m}\) 上也成立，局部环整闭便给出 \(z\in A_{\mathfrak m}\)。\secref{b3:sec:0204}已把\cref{b3:cor:submodule-local-detection}应用于 \(A\subseteq K\)，证明共同分式域中的交式 \(A=\bigcap_{\mathfrak m\;\text{极大}}A_{\mathfrak m}\)。因此 \(z\in A\)，即 \(A\) 整闭。
\end{proof}

尖点和节点都约化且不可约，因为它们的坐标环是 \(k[t]\) 的子环。但 \(t\) 属于其分式域、满足首一方程且不在原环中，因此都不正规。交叉直线约化而可约，其正规化将两个不可约分量分开；加厚点的代数模型则首先就不约化。不能只用“有几个点”或“有几条曲线”代替这些环条件。

\ProseParagraphBreak
尖点的参数映射还给出同胚而非同构的例子。记 \(A=k[t^2,t^3]\)、\(B=k[t]\)，分别考虑谱映射 \(\pi_{\mathrm{Spec}}:\operatorname{Spec}B\to\operatorname{Spec}A\) 和经典参数映射 \(\pi(k):k\to C\)、\(t\mapsto(t^2,t^3)\)。对任一素理想 \(\mathfrak q\subseteq B\)，将\cref{b3:lem:integral-field-criterion}用于整包含 \(A/(\mathfrak q\cap A)\subseteq B/\mathfrak q\)，便知 \(\mathfrak q\) 极大当且仅当其收缩极大。在弱零点定理的识别下，\(\pi_{\mathrm{Spec}}\) 对极大谱的限制就是前面已经核验双射性的 \(\pi(k)\)。\(B=k[t]\) 的唯一非极大素理想为 \((0)\)；由整扩张谱映射的满射性，\(A\) 的非极大素理想只能是它的收缩 \((0)\)。所以两个谱的非闭点也唯一对应，\(\pi_{\mathrm{Spec}}\) 确为双射。谱映射本来连续，\cref{b3:prop:integral-spec-closed}又说明它是闭映射，因而是同胚。

经典点集上的同胚也可由闭映射直接核验。对任一理想 \(J\subseteq B\)，极大性与收缩相容，故同一闭映射公式限制为
\[
\pi_{\mathrm{Spec}}\bigl(V_B(J)\cap\operatorname{MaxSpec}B\bigr)
=V_A(J\cap A)\cap\operatorname{MaxSpec}A.
\]
因此连续的经典参数映射 \(\pi(k)\) 也是闭双射，从而为同胚；但其坐标环同态 \(A\hookrightarrow B\) 不满，因为 \(t\notin A\)，所以它不是代数集的同构。拓扑空间相同并不意味着正则函数相同。

\subsection{几何图像与局部环计算}
\label{b3:subsec:0704:03}

对尖点取 \(A=k[t^2,t^3]\)、原点极大理想 \(\mathfrak m=(t^2,t^3)_A\)。此时 \(t\notin A_{\mathfrak m}\)。否则存在 \(a,b\in A\)、\(b\notin\mathfrak m\)，使 \(tb=a\)。\(b\) 的常数项非零，所以 \(tb\) 的一次项非零；\(A\) 中元素的一次项却全为零，矛盾。这证明不整闭性确实保留在特殊点的局部环中。在导子零点集 \(V_A(\mathfrak c)\) 之外，每个素理想都避开某个导子元素，由\cref{b3:prop:conductor-basic}两侧局部化相同，因而没有这项差别。

\ProseParagraphBreak
对节点则取 \(A=k+(t^2-1)k[t]\)、原点极大理想 \(\mathfrak m=\bigl(t^2-1,t(t^2-1)\bigr)_A\)。这里同样不能把 \(t\) 写成原局部环的元素。若 \(tb=a\)、\(a,b\in A\) 且 \(b\) 在节点原点处非零，即 \(b(1)=b(-1)\ne0\)，在 \(t=1,-1\) 处分别取值得 \(a(1)=b(1)\)、\(a(-1)=-b(-1)\)。但 \(a\in A\) 要求 \(a(1)=a(-1)\)，与特征不为二矛盾。原局部环仍是整环，正规化上方却有两个点；\secref{b3:sec:0606}已具体核验其半局部正规化，并说明这两个上方点与形式方程的两条分支分别属于什么环境。

\ProseParagraphBreak
设 \(R\) 为交换含幺局部环，\(\mathfrak m\) 为其唯一极大理想，\(\kappa=R/\mathfrak m\) 为剩余域。则 \(\mathfrak m/\mathfrak m^2\) 是 \(\kappa\) 向量空间，称为该局部环的\term[yuqie-kongjian]{余切空间}[cotangent space]；其 \(\kappa\) 对偶 \(\operatorname{Hom}_{\kappa}(\mathfrak m/\mathfrak m^2,\kappa)\) 是 Zariski 切空间。在以上平面尖点、节点和交叉直线的原点，关系多项式都没有一次项，故 \(\bar x,\bar y\) 给出二维余切空间：若其线性组合落在 \(\mathfrak m^2\)，比较多项式的一次项只能使两个系数都为零。局部化后这个计算不变，因为不在原点极大理想中的元素在该商上按其非零常数值可逆作用。

这一计算只保留一阶条件。例如在 \(k[\varepsilon]/(\varepsilon^2)\) 中代入 \(x=\alpha\varepsilon\)、\(y=\beta\varepsilon\)，尖点关系 \(y^2-x^3\) 对任意 \(\alpha,\beta\in k\) 都变成零，所以 Zariski 切空间是整个 \(k^2\)。若再看方程最低次数的非零项，则得到 \(y^2\)；它所确定的切锥约化后只有方向 \(y=0\)。因此这一条方向与二维的一阶切空间并不矛盾，二者保留的阶数信息不同。

\ProseParagraphBreak
直线 \(k[t]_{(t)}\) 的余切空间却只有基 \(t+(t^2)\)，维数为一；加厚点 \(k[\varepsilon]/(\varepsilon^2)\) 的余切空间也为一维，而普通点 \(k\) 的余切空间为零。\cref{b3:tab:curve-thick-point-data}汇集了上述四个模型的数据：节点只有一个全局不可约分量，正规化在原点上方却有两个点；加厚点与尖点的集合纤维都只有一个点，纤维代数却不同。仅凭余切空间维数尚不能判断光滑性；正则局部环、Jacobian 判据及光滑性之间的联系见后面的章节。

\begin{table}[htbp]
\centering
\caption[曲线与加厚点的局部数据]{曲线与加厚点的局部数据。设 \(k\) 为代数闭域，节点例要求 \(\operatorname{char}k\ne2\)。各模型的代数记为 \(A\)，原点极大理想记为 \(\mathfrak m\)，\(A/\mathfrak m=k\)。不可约分量数指 \(\operatorname{Spec}A\) 的分量数，纤维代数为 \(\widetilde A\otimes_Ak\)，余切维数为 \(\dim_k\mathfrak m/\mathfrak m^2\)。加厚点按本书约定先约化，其正规化为 \(k\)。}
\label{b3:tab:curve-thick-point-data}
\begin{tabularx}{\textwidth}{@{}l c>{\raggedright\arraybackslash}X c c@{}}
\toprule
模型 & 不可约分量数 & 正规化纤维代数 & 纤维点数 & 余切维数\\
\midrule
尖点 & \(1\) & \(k[t]/(t^2)\) & \(1\) & \(2\)\\
节点 & \(1\) & \(k\times k\) & \(2\) & \(2\)\\
交叉直线 & \(2\) & \(k\times k\) & \(2\) & \(2\)\\
加厚点 & \(1\) & \(k\) & \(1\) & \(1\)\\
\bottomrule
\end{tabularx}
\end{table}


\begin{chaptersummary}
在代数闭域上，仿射代数集由根理想刻画，其坐标环是约化有限生成代数。不可约性对应整性，开闭分解对应幂等元。坐标环不仅记录多项式函数，也包含全部全局正则函数；局部化则给出一点处的正则函数芽。由此，正则态射与坐标环同态构成方向相反的对应。

\ProseParagraphBreak
点集并不包含全部代数信息。纤维代数可能有幂零元，取根理想才得到普通代数集的坐标环；基域扩张既可能使一个闭点分裂，也可能在不可分情形下产生幂零元。几何点还可能落在非闭点上，不能一概等同于原底域上的坐标点。

\ProseParagraphBreak
正规化把整闭性落实到具体函数环中。尖点、节点与交叉直线有不同的分支和导子，双射也未必保持正则函数。这些区别都能在坐标环及其局部化中检验。下一编转向 Gröbner 基、消元及若干有效计算，使理想之间的关系能够通过有限步骤求出。
\end{chaptersummary}

\BookExercises
以下各题除另有说明外，\(k\) 为代数闭域，代数集及其态射采用本章约定。

\begin{exercise}\label{b3:ex:07-parabola}
证明抛物线 \(X=Z_k(y-x^2)\) 与仿射直线同构。分别写出两个方向的态射及坐标环同构，并说明嵌入空间维数不同不妨碍代数集同构。
\end{exercise}
\begin{solution}
参数映射 \(t\mapsto(t,t^2)\) 与投影 \((a,b)\mapsto a\) 都由多项式给出，复合分别为恒等映射。相应的环同构为 \(k[x,y]/(y-x^2)\to k[t]\)，将 \(\bar x,\bar y\) 送到 \(t,t^2\)，逆映射将 \(t\) 送到 \(\bar x\)。这里同构要求互逆的正则态射，并不要求所选嵌入具有相同的坐标数。
\end{solution}

\begin{exercise}\label{b3:ex:07-three-lines}
求 \(Z_k\bigl(xy(x-1)\bigr)\) 的不可约分支，证明它连通但不可约性不成立。
\end{exercise}
\begin{solution}
由乘积为零，零点集是 \(Z_k(x)\)、\(Z_k(y)\) 与 \(Z_k(x-1)\) 三条直线的并。每条直线的坐标环为一元多项式环，故不可约，且三者互不包含，正是不可约分支。水平直线分别在 \((0,0)\)、\((1,0)\) 与两条竖直线相交。两个相交的连通子空间之并连通：若其并分成不交开闭部分，每个连通子空间必须全落在其中一部分，交点迫使二者落在同一部分。连续使用这项观察便得整个代数集连通。它有三个分支，故不不可约。
\end{solution}

\begin{exercise}\label{b3:ex:07-disjoint-lines}
求 \(X=Z_k\bigl(x(x-1)\bigr)\) 的坐标环与全部幂等元，并用它们描述 \(X\) 的连通分支。
\end{exercise}
\begin{solution}
理想 \((x)\) 与 \((x-1)\) 互素，其交为积。中国剩余定理给出
\[
k[X]\simeq k[y]\times k[y],\qquad
\bar x\longmapsto(0,1).
\]
每个因子为整环，幂等元只有零与一。因此乘积环的幂等元为 \(0,1,\bar x,1-\bar x\)。后两者的取值分别选出两条竖直线；每条直线不可约，因而连通，这就是全部连通分支。
\end{solution}

\begin{exercise}\label{b3:ex:07-line-automorphisms}
证明仿射直线的全部 \(k\)-自同构为 \(t\mapsto at+b\)，其中 \(a\in k^\times,b\in k\)。
\end{exercise}
\begin{solution}
由\cref{b3:thm:affine-global-regular}，一个态射由多项式 \(f(t)\) 给出。若逆态射由 \(g(t)\) 给出，则 \(g\circ f=t\)，二者都非常数，次数公式给出 \(\deg g\deg f=1\)。因此 \(f=at+b\) 且 \(a\ne0\)。反之，\(t\mapsto a^{-1}(t-b)\) 是这种映射的正则逆。
\end{solution}

\begin{exercise}\label{b3:ex:07-frobenius-bijection}
设 \(\operatorname{char}k=p>0\)。证明 \(F:\mathbb A_k^1\to\mathbb A_k^1\)、\(a\mapsto a^p\) 是有限的双射，并诱导拓扑同胚，但不是代数集同构。
\end{exercise}
\begin{solution}
代数闭性保证每个元素有 \(p\) 次根，而 \(a^p=b^p\) 蕴含 \((a-b)^p=0\)，故该根唯一。环映射 \(k[s]\to k[t]\)、\(s\mapsto t^p\) 使后者由 \(1,t,\ldots,t^{p-1}\) 生成为模，所以态射有限。仿射直线的真闭集都是有限点集，任意双射都保持这类闭集，因而此处是同胚。环映射的像却只有 \(k[t^p]\)，不含 \(t\)；由\cref{b3:thm:affine-antiequivalence}，态射不是同构。两边坐标环都是整闭整环也不能改变这个结论。
\end{solution}

\begin{exercise}\label{b3:ex:07-closed-embedding}
设 \(f:X\to Y\) 为仿射代数集的态射。证明 \(f^*:k[Y]\to k[X]\) 满射，当且仅当 \(f\) 给出 \(X\) 与 \(Y\) 的一个闭代数子集之间的同构。
\end{exercise}
\begin{solution}
设拉回满射，记核为 \(J\)。因商环 \(k[Y]/J\simeq k[X]\) 约化，\(J\) 为根理想。零点定理给出 \(Z_Y(J)\) 的坐标环恰为这个商。拉回经商环分解得到一个同构，再由\cref{b3:thm:affine-antiequivalence}还原为 \(X\to Z_Y(J)\) 的同构，其与包含映射的复合就是 \(f\)。反之，闭代数子集的坐标环由取商得到，限制映射满射；与一个环同构复合后仍满射。
\end{solution}

\begin{exercise}\label{b3:ex:07-dense-image}
证明 \(f^*:k[Y]\to k[X]\) 单射当且仅当 \(f(X)\) 在 \(Y\) 中稠密。举例说明单射不能保证 \(f\) 满射。
\end{exercise}
\begin{solution}
核恰是 \(Y\) 上在 \(f(X)\) 处处为零的函数所成的理想，也就是在其闭包上为零的函数理想。闭包为 \(Y\) 时，这个理想为零；反之，若闭包是真闭代数子集，零点定理表明它的零点理想非零，否则零理想的零点集将是整个 \(Y\)。例如双曲线 \(xy=1\) 到横轴的投影在环上给出 \(k[t]\hookrightarrow k[t,t^{-1}]\)，是单射；其像为 \(k^\times\)，稠密但遗漏原点。
\end{solution}

\begin{exercise}\label{b3:ex:07-cubic-fibers}
计算态射 \(t\mapsto t^3\) 在零、一上方的纤维代数及集合意义的纤维。讨论底域特征为三时发生的变化。
\end{exercise}
\begin{solution}
由\cref{b3:prop:fiber-algebra}，在 \(a\) 上方的代数为 \(k[t]/(t^3-a)\)。零纤维代数为 \(k[t]/(t^3)\)，有非零幂零元 \(\bar t\)，而点集只有原点，点集的坐标环为 \(k\)。若特征不为三，\(t^3-1\) 有三个不同的根，因为它与导数 \(3t^2\) 无公共根；中国剩余定理给出一纤维代数 \(k^3\)，点集有三个点。特征为三时，\(t^3-1=(t-1)^3\)，一纤维代数为 \(k[t]/\bigl((t-1)^3\bigr)\)，约化后为 \(k\)，点集只有一个。
\end{solution}

\begin{exercise}\label{b3:ex:07-product-fibers}
对 \(f:\mathbb A_k^2\to\mathbb A_k^1\)、\((x,y)\mapsto xy\)，计算各纤维的坐标环，判断它们是否约化、是否不可约。
\end{exercise}
\begin{solution}
纤维代数为 \(k[x,y]/(xy-a)\)。当 \(a=0\) 时，\((xy)=(x)\cap(y)\) 为根理想，故它已经约化；两个轴给出两个不可约分支。当 \(a\ne0\) 时，映射 \(x\mapsto t,y\mapsto at^{-1}\) 给出该环与 \(k[t,t^{-1}]\) 的同构，逆映射取 \(t\mapsto\bar x,t^{-1}\mapsto a^{-1}\bar y\)。后者为整环，故相应纤维约化且不可约。
\end{solution}

\begin{exercise}\label{b3:ex:07-tensor-square-map}
通过 \(s\mapsto u^2\) 和 \(s\mapsto v^2\) 把两个多项式环视为 \(k[s]\)-代数。计算 \(k[u]\otimes_{k[s]}k[v]\)，并比较特征为二与不为二的情形。
\end{exercise}
\begin{solution}
张量积的生成元为 \(u\otimes1,1\otimes v\)，平衡关系使二者的平方相等。反向将 \(u^i\otimes v^j\) 送到商环中 \(u^iv^j\)，得到互逆同构
\[
k[u]\otimes_{k[s]}k[v]\simeq k[u,v]/(u^2-v^2).
\]
若特征不为二，令 \(r=u-v,q=u+v\) 是可逆线性换元，右边成为 \(k[r,q]/(rq)\)，是两条交叉直线的约化坐标环。若特征为二，令 \(r=u-v\)，右边成为 \(k[v,r]/(r^2)\)，其中 \(\bar r\ne0\) 而平方为零。集合条件 \(u^2=v^2\) 只给出对角线，其坐标环是再取约化得到的 \(k[v]\)。
\end{solution}

\begin{exercise}\label{b3:ex:07-rational-closed-point}
令 \(A=\mathbb Q[t]/(t^2-2)\)。求其 \(\mathbb Q\)-值点、\(\mathbb R\)-值点和 \(\mathbb C\)-值点，并计算 \(\mathbb C\otimes_{\mathbb Q}A\)。为什么 \(\operatorname{Spec}A\) 只有一个点？
\end{exercise}
\begin{solution}
\(\mathbb Q\)-代数同态由 \(\bar t\) 的像决定，这个像必须满足平方等于二。因此没有 \(\mathbb Q\)-值点，而实值点和复值点各有两个，分别取 \(\pm\sqrt2\)。基域扩张的展示给出 \(\mathbb C[t]/(t^2-2)\simeq\mathbb C\times\mathbb C\)。另一方面，\(t^2-2\) 在 \(\mathbb Q\) 上不可约，故 \(A\) 本身是域，谱只有零理想这一个闭点。扩张后的两个点来自这个剩余域到 \(\mathbb C\) 的两种嵌入，并非原谱中已有两个素理想。
\end{solution}

\begin{exercise}\label{b3:ex:07-dual-number-valued-points}
设 \(D=k[\varepsilon]/(\varepsilon^2)\)。证明对每个扩域 \(K/k\)，\(D\) 只有一个 \(K\)-值点；但到 \(E=k[\delta]/(\delta^2)\) 的 \(k\)-代数同态由 \(k\) 中的元素参数化。
\end{exercise}
\begin{solution}
域无非零幂零元，故任何同态 \(D\to K\) 必须把 \(\bar\varepsilon\) 送到零，唯一性随之成立。\(E\) 中平方为零的元素恰为 \(c\bar\delta\)：若 \((a+c\bar\delta)^2=0\)，其常数项给出 \(a^2=0\)，故 \(a=0\)；反之显然成立。因此每个 \(c\in k\) 给出且仅给出一个同态 \(\bar\varepsilon\mapsto c\bar\delta\)。域值点看不见的幂零方向，可以被含幂零元的测试代数辨认。
\end{solution}

\begin{exercise}\label{b3:ex:07-crossed-lines-localization}
设 \(A=k[x,y]/(xy)\)。计算主开集 \(D_X(\bar x)\) 的正则函数环及点 \((1,0)\) 的局部环，再与原点的局部环比较。
\end{exercise}
\begin{solution}
将 \(\bar x\) 变成单位后，关系 \(\bar x\bar y=0\) 迫使 \(\bar y=0\)，所以 \(A_{\bar x}\simeq k[x,x^{-1}]\)。由\cref{b3:prop:regular-germ-local-ring}，点 \((1,0)\) 的局部环为 \(A_{(\bar x-1,\bar y)}\simeq k[x]_{(x-1)}\)，是整环。原点处 \(\bar x,\bar y\) 都仍非零：若分母在原点不为零，其与 \(x\) 的乘积属于 \((xy)\)，在整环 \(k[x,y]\) 中约去 \(x\) 将迫使分母属于 \((y)\)，与常数项非零矛盾；对 \(y\) 同理。但二者乘积为零，所以原点局部环不是整环。
\end{solution}

\begin{exercise}\label{b3:ex:07-multiplicative-line-automorphisms}
把 \(\mathbb G_m\) 视为双曲线 \(xy=1\)。证明其全部自同构为 \(t\mapsto at\) 或 \(t\mapsto at^{-1}\)，其中 \(a\in k^\times\)。
\end{exercise}
\begin{solution}
坐标环为 \(k[t,t^{-1}]\)。非零 Laurent 多项式的最高次数与最低次数之差在相乘时相加，故可逆元的差必须为零；所有单位恰为 \(at^n\)，其中 \(a\ne0,n\in\mathbb Z\)。自同构将 \(t\) 送到一个单位 \(at^n\)，逆同构将 \(t\) 送到 \(bt^m\)，复合后指数为 \(mn=1\)，所以 \(n=\pm1\)。这两种映射均有同类的逆，因而穷尽全部自同构。
\end{solution}

\begin{exercise}\label{b3:ex:07-punctured-plane}
设 \(U=\mathbb A_k^2\setminus\bigl\{(0,0)\bigr\}\)，正则函数按开子集上的定义理解。证明 \(\mathcal O(U)=k[x,y]\)，但 \(U\) 不与任何仿射代数集同构。
\end{exercise}
\begin{solution}
写 \(R=k[x,y]\)。主开覆盖 \(U=D(x)\cup D(y)\) 把正则函数确定为 \(R_x\) 和 \(R_y\) 中在交集上相等的两个元素。交集的函数环是 \(R_{xy}\)，并嵌入 \(k(x,y)\)，故这个条件给出 \(\mathcal O(U)=R_x\cap R_y\)。若 \(a/x^m=b/y^n\)，则 \(ay^n=bx^m\)。第一卷定理12.2.6保证 \(R\) 是唯一分解整环；因 \(x,y\) 互素，得到 \(x^m\mid a\)，从而该分式在 \(R\) 中。这证明交等于 \(R\)。

若存在仿射代数集 \(Y\) 和同构 \(h:Y\to U\)，与包含映射复合得到 \(j:Y\to\mathbb A_k^2\)。刚才的计算和\cref{b3:thm:affine-global-regular}说明 \(j^*:R\to k[Y]\) 为同构。仿射反等价便迫使 \(j\) 为同构，特别必须满射；但其像遗漏原点，矛盾。这里反等价只对假定存在的仿射代数集 \(Y\) 与 \(\mathbb A_k^2\) 使用，没有预先把 \(U\) 当成仿射代数集。
\end{solution}

\begin{exercise}\label{b3:ex:07-finite-fibers-not-finite}
证明包含态射 \(\mathbb G_m\to\mathbb A_k^1\) 的每个集合纤维都有限，但该态射不是有限态射。
\end{exercise}
\begin{solution}
非零点上方恰有一个点，零上方为空。若态射有限，\(k[t,t^{-1}]\) 将是有限生成的 \(k[t]\) 模，从而 \(t^{-1}\) 在 \(k[t]\) 上整。由于 \(k[t]\) 是整闭整环，而 \(t^{-1}\) 属于其分式域但不属于它，这不可能。也可将一个假定存在的首一整方程乘以足够次 \(t\)，得到常数项为一却能被 \(t\) 整除的多项式矛盾。
\end{solution}

\begin{exercise}\label{b3:ex:07-monomial-curve-conductor}
设 \(C=Z_k(y^3-x^4)\)。结合\cref{b3:ex:06-three-four-cusp}，证明参数映射 \(t\mapsto(t^3,t^4)\) 是有限双射而非同构。计算原点的纤维代数，并比较 \(C\) 原点与参数直线原点的余切空间。
\end{exercise}
\begin{solution}
先核实坐标环为 \(A=k[t^3,t^4]\)。将多项式对首一多项式 \(y^3-x^4\) 除法，余式为 \(\sum_{j=0}^2 a_j(x)\cdot y^j\)。代入参数后，各组单项式的次数 \(3i+4j\) 在模三意义下由 \(j\) 区分，同组的次数也互异，所以余式只有为零才可能消失。这证明参数代入的核恰为 \((y^3-x^4)\)。商环嵌入 \(k[t]\)，故该理想素，零点定理给出所需坐标环。

由\cref{b3:ex:06-three-four-cusp}，正规化 \(B=k[t]\) 是有限生成的 \(A\) 模，且包含严格，所以参数态射有限而非同构。若 \((a,b)\in C\) 且 \(a\ne0\)，取 \(t=b/a\)，则由 \(b^3=a^4\) 得 \(t^3=a\)、\(t^4=b\)，且这样的参数唯一；若 \(a=0\)，必有 \(b=0\)，参数只能为零。因此点集上双射。原点的纤维代数为
\[
B/(t^3,t^4)B\simeq k[t]/(t^3),
\]
含非零幂零元，而普通点集的坐标环只有 \(k\)。

原点极大理想 \(\mathfrak m=(\bar x,\bar y)\) 的余切空间以 \(\bar x,\bar y\) 的类为基：关系 \(y^3-x^4\) 没有一次项，故二者模 \(\mathfrak m^2\) 线性无关。维数为二，参数直线原点的余切空间则为 \((t)/(t^2)\)，维数为一。这进一步从局部环说明双射不能成为同构。
\end{solution}

\begin{exercise}\label{b3:ex:07-node-semilocal-units}
设 \(\operatorname{char}k\ne2\)，\(A=k+(t^2-1)k[t]\)，\(\mathfrak m=\bigl(t^2-1,t(t^2-1)\bigr)_A\)，\(S=A\setminus\mathfrak m\)。证明 \(S^{-1}k[t]\) 正是把所有在 \(1,-1\) 处都不为零的多项式变成单位得到的环。它为何没有非平凡幂等元？
\end{exercise}
\begin{solution}
\(A\) 的元素在 \(1,-1\) 的值相同，属于 \(S\) 当且仅当这个公共值非零。因此先把所有在这两点均非零的多项式变成单位，必使 \(S\) 可逆。反过来，若 \(f(1)f(-1)\ne0\)，则 \(f(t)f(-t)\) 是偶多项式，属于 \(k[t^2]\subseteq A\)，且在两点的值均为非零的 \(f(1)f(-1)\)，所以属于 \(S\)。恒等式 \(f(t)^{-1}=f(-t)/\bigl(f(t)f(-t)\bigr)\) 的分母由这个元素给出，故 \(f\) 已在 \(S^{-1}k[t]\) 中可逆。两次局部化因此相同。该环是域 \(k(t)\) 的子环，为整环，幂等元只能为零或一；虽有两个极大理想，也不能写成两个非零环的直积。
\end{solution}

\begin{exercise}\label{b3:ex:07-crossing-versus-node}
设 \(A=k[x,y]/(xy)\)，\(B=k[x]\times k[y]\)，\(\mathfrak m=(\bar x,\bar y)\)，\(S=A\setminus\mathfrak m\)。证明
\[
S^{-1}B\simeq k[x]_{(x)}\times k[y]_{(y)},\qquad B/\mathfrak mB\simeq k\times k.
\]
将前一个结果与\cref{b3:ex:07-node-semilocal-units}比较。
\end{exercise}
\begin{solution}
把 \(A\) 认作常数项相同的多项式对，\(S\) 就是公共常数项非零的那些对。局部化后两因子的分母都在零处非零。反之，给定两个分母 \(g(x),h(y)\) 满足 \(g(0)h(0)\ne0\)，用 \(h(0)g(x),g(0)h(y)\) 组成公共常数项相同的分母对，就能表示乘积局部环中的任意分式对，得到满射；单射由两个因子的分式表示立即成立。

理想 \(\mathfrak mB\) 为 \((x)\times(y)\)，故商为 \(k\times k\)。此处两个分支的正规化局部环确实成直积；节点原环却是整环，其分式域内的正规化仍为整环。两个例子在闭点上方都有两个点，甚至纤维代数都为 \(k\times k\)，但这并不能判定正规化的半局部环是否为直积。
\end{solution}

\begin{exercise}\label{b3:ex:07-fat-point-cotangent}
设 \(r\ge2\)，\(A=k[x,y]/(x,y)^r\)，\(\mathfrak m=(\bar x,\bar y)\)。计算 \(\dim_k A\)、\(\dim_k\mathfrak m/\mathfrak m^2\) 与 \(\operatorname{Spec}A\) 的点数。
\end{exercise}
\begin{solution}
全次数小于 \(r\) 的单项式的像组成 \(A\) 的基，各次依次有 \(1,2,\ldots,r\) 个，所以 \(\dim_k A=r(r+1)/2\)。模去 \(\mathfrak m^2\) 后，\(\mathfrak m\) 中只有 \(\bar x,\bar y\) 的一次项留下，且 \(r\ge2\) 保证它们线性无关，故余切空间维数为二。每个素理想都包含幂零元 \(\bar x,\bar y\)，而商 \(A/\mathfrak m=k\)，所以谱只有 \(\mathfrak m\) 一个点。点数、代数的向量空间维数和余切空间维数分别记录不同信息。
\end{solution}

\begin{exercise}\label{b3:ex:07-regular-denominators}
设 \(X=Z_k(xy-1)\)。证明 \(1/x\) 是全局正则函数，而 \(1/(x+y)\) 只在一个真开子集上正则，不能延拓成 \(X\) 上的全局正则函数。
\end{exercise}
\begin{solution}
在坐标环 \(k[x,x^{-1}]\) 中，\(x^{-1}=y\)，所以第一项本来就是多项式函数的限制。第二项的分母对应 \(x+x^{-1}\)。代数闭域中 \(x^2+1\) 有非零根，故分母确实有零点；其非零集是非空真主开集。若逆函数延拓为全局函数 \(g\)，则 \((x+x^{-1})g=1\) 在这个稠密主开集成立，因而作为整环 \(k[x,x^{-1}]\) 中的等式处处成立。在分母的零点取值便得到 \(0=1\)，矛盾。
\end{solution}

\begin{exercise}\label{b3:ex:07-connected-to-finite}
设 \(X\) 为非空连通仿射代数集，\(Y\) 为非空有限仿射代数集。证明每个态射 \(X\to Y\) 都是常值，并用坐标环解释这个结论。
\end{exercise}
\begin{solution}
有限代数集的每个单点闭，其补集也是有限闭集，故拓扑离散。连续像连通，只能包含一个点。环上，设 \(Y\) 有 \(n\) 个点；零点定理及中国剩余定理给出 \(k[Y]\simeq k^n\)。其标准幂等元 \(e_i\) 的像在 \(k[X]\) 中仍为两两正交、和为一的幂等元。由\cref{b3:thm:affine-connected-idempotents}，每个像只能为零或一；两两正交保证至多一个为一，和为一保证恰有一个为一。于是拉回先取 \(k^n\) 的某个坐标，再嵌入 \(k[X]\) 的常数，正是常值态射。
\end{solution}

